% \iffalse meta-comment
% arara: lualatex: { branch: developer, draft: true }
% arara: lualatex: { branch: developer, interaction: batchmode, shell: true }
% arara: clean: { extensions: [ aux, log, out, glo, gls, hd, ind, ilg, idx, toc, atfi ] }
% arara: clean: { files: [ userdoc.idx, userdoc.ilg, userdoc.ind, mylhmc.lua, mylhmc.sty ] }
% arara: halt
%
% File: enumext.dtx
%
% Copyright (C) 2024-2026 by Pablo González L <pablgonz@educarchile.cl>
%
% This work may be distributed and/or modified under the conditions of the
% LaTeX Project Public License, either version 1.3c of this license or (at
% your option) any later version. The latest version of this license is in
%
%    https://www.latex-project.org/lppl.txt
%
% and version 1.3c or later is part of all distributions of LaTeX version
% 2005/12/01 or later.
%
% This work is "maintained" (as per the LPPL maintenance status)
% by Pablo González L.
%
% This work consists of the files enumext.dtx and enumext.sty.
%
% \fi
%
% \iffalse
%<*driver>
\documentclass[full]{l3doc}
\usepackage[svgnames]{xcolor}
\usepackage[sf,bf,compact,medium,pagestyles]{titlesec}
\usepackage{lastpage,microtype,attachfile2,multicol,listings,titletoc,fontawesome5}
\usepackage[english]{babel}
\usepackage[top=0.5in,bottom=0.3in,left=1.70in,right=0.7in,footskip=0.2in,headheight=1cm,headsep=0.27cm]{geometry}
% Set \parindent and \parskip
\setlength{\parindent}{0pt}
\setlength{\parskip}{\smallskipamount}
% Extra stretch before overflowing the line
\setlength{\emergencystretch}{2em}
% Only for arara... I like arara :)
\usepackage[scale=0.85]{comfortaa}
\definecolor{araracolor}{rgb}{0, 0.72, 0.28}
\newcommand{\araratext}[1]{{\small\normalfont\comfortaa\color{araracolor}\bfseries#1}}
\newcommand*\arara{\araratext{ar\kern-.03emar\kern-.03ema}}
% Fonts
\usepackage[osf,nomath,mono=false,ScaleSF=0.95,ScaleRM=0.95]{libertinus-otf}
% Set Monospace font, fix 'quotes' in verbatim
\usepackage{sourcecodepro}
\defaultfontfeatures[\ttfamily]
  {
    Numbers   = OldStyle,
    Scale     = 0.84249,
    Extension = .otf,
  }
\setmonofont[
    UprightFont       = *-Regular,
    ItalicFont        = *-RegularIt,
    BoldFont          = *-Medium,
    BoldItalicFont    = *-MediumIt,
    RawFeature        = {+zero,+ss01}]{SourceCodePro}
\newfontfamily\mysourcecodeprolight{SourceCodePro-Regular.otf}[
    Scale = 0.80,FakeStretch =0.75,
    FontFace = {m}{it}{Font =SourceCodePro-RegularIt.otf,Color=FF0000}]
\RenewDocumentCommand{\MacroLongFont}{}
  {
    \mysourcecodeprolight\small
  }
% Internal
\newfontfamily\mylmrotf{lmroman10-regular.otf}[NFSSFamily=lmrotf,Scale = 0.95]
\newfontfamily\mylmmotf{lmmono10-regular.otf}[NFSSFamily=lmmotf,Scale = 0.95]
% The character of visible space is now taken from Latin Modern Mono
% to prevent fonts in T1. The original definition for xetex/luatex is
% \def\verbvisiblespace{\usefont{OT1}{cmtt}{m}{n}\asciispace}
\def\verbvisiblespace{{\fontfamily{lmmotf}\fontshape{n}\selectfont\char"2423}}

% Math font
\usepackage{unicode-math}
\setmathfont[Scale = 0.95]{latinmodern-math.otf}
\newfontfamily\fetamontotf{ffmw10.otf}[
   Scale             = 0.95,%
   RawFeature        = {+latn,+rand,+kern,+size},%
   FontFace          = {b}{n}{ffmw10.otf},% fix replace font
   ]
% Remove `"` and `|` for short verbatim defined in l3doc, then set short |...|
\AtBeginDocument
  {
    \DeleteShortVerb \"
    \DeleteShortVerb \|
    \lstMakeShortInline[language=enumext-doc,basicstyle=\ttfamily]\|
  }
% Index
\usepackage{imakeidx}
\indexsetup{level=\section,firstpagestyle=myheader,othercode=\pagestyle{myheader}}
% imakeidx also opens a write stream for the .idx file, and that conflicts
% with the one opened by l3doc. Here we close that write stream:
\ExplSyntaxOn
\iow_close:c { @indexfile }
% and later we will copy the write stream opened by imakeidx into \@indexfile
% so that entries written to both streams end up in the same file.
\makeindex[name=userdoc,options={-q ~ -s ~ gind.ist},columnsep=15pt,title={Index ~ of ~ Documentation}]
\makeindex[options={-q ~ -s ~ gind.ist},columnsep=15pt,title={Index ~ of ~ Implementation}]
\EnableCrossrefs
% \PageIndex, \CodelineIndex undoes what \PageIndex does
% \CodelineIndex tries to open another write stream for the index file. We don't
% want that, so we temporarily make \makeindex a no-op:
\cs_set_eq:NN \l__mydoc_tmp_makeindex_cs \makeindex
\cs_set_nopar:Npn \makeindex { }
\CodelineIndex
\cs_set_eq:NN \makeindex \l__mydoc_tmp_makeindex_cs
\DoNotIndex{\ , \1,\^}
\exp_args:Ne \DoNotIndex { \tl_to_str:n { \{ } }
\exp_args:Ne \DoNotIndex { \tl_to_str:n { \} } }
\exp_args:Ne \DoNotIndex { \tl_to_str:n { \begin } }
% \StartImplementation
\cs_new_protected:Npn \StartImplementation
  {
    \bool_set_true:N \l__codedoc_in_implementation_bool
  }
\bool_set_false:N \l__codedoc_in_implementation_bool
% Links and index
\newcommand{\HP}[1]{\emph{\hyperpage{#1}}\normalsize}
\cs_new_protected:Nn \__mydoc_sort_index:nn
  {
    \bool_if:NTF \l__codedoc_in_implementation_bool
      {
        \index{#1\actualchar#2|HP}
      }
      {
        \index[userdoc]{#1\actualchar#2|HP}
      }
  }
% The text of the entry in the first language (#2) for the user part and in
% the second one (#3) for the implementation, chosen when the command is called.
\cs_new_protected:Nn \__mydoc_sort_index:nnn
  {
    \bool_if:NTF \l__codedoc_in_implementation_bool
      { \__mydoc_sort_index:nn {#1} {#3} }
      { \__mydoc_sort_index:nn {#1} {#2} }
  }
% Now, after imakeidx opens the write stream for the index file, we copy
% the reference to \@indexfile:
\cs_set_eq:cc { @indexfile } { \c_sys_jobname_str @idxfile }
\ExplSyntaxOff
% Package documented and packages for the examples
\usepackage{tikz,enumext} % main
\usepackage{scontents}
\usepackage{shellesc}
% LuaLaTeX syntax highlighter for l3doc (it also defines the color palette)
\usepackage{mylhmc}
\setmylhmc
  {
    module  = enumext,
    provide =
      { setenumext, setenumextmeta, resetenumext, printkeyans, foreachkeyans,
        getkeyans, anskey, anspic, miniright, keyans, keyanspic }
  }
% Colors
\definecolor{newlinkcolor}{HTML}{0A602C}
\definecolor{rulecolor}{rgb}{0.96,0.96,0.96}
% Get file info
\GetFileInfo{enumext.sty}
% Config hyperref
\hypersetup
  {
    allcolors          = newlinkcolor,
    colorlinks         = true,
    linktoc            = all,
    pdftitle           = {.:: The enumext package --- \fileinfo{} ::.},
    pdfauthor          = {Pablo González L},
    pdfsubject         = {Documentation for \fileversion{} [\filedate] },
    pdfkeywords        = {worksheets, multiple choice, expl3, l3prop, l3seq, list, enumerate, store answers, tagged},
    bookmarksopenlevel = 2,
  }
% User level commands
\ExplSyntaxOn
% Logo
\NewDocumentCommand \pkglogo { }
  {
    \group_begin:
      \LibertinusSansOsF
        { \textcolor { module_name } { enum } \textcolor { OrangeRed } { ext } }
    \group_end:
  }
% Custom \meta[...]{...}, \marg[...]{...} and \oarg[...]{...} with color
\NewDocumentCommand \mymeta { O{} m }
  {
    \__mydoc_meta_generic:Nnn \__mydoc_meta:n { #1 } { #2 }
  }
\NewDocumentCommand \mymarg { O{} m }
  {
    \__mydoc_meta_generic:Nnn \__mydoc_marg:n { #1 } { #2 }
  }
\NewDocumentCommand \myoarg { O{} m }
  {
    \__mydoc_meta_generic:Nnn \__mydoc_oarg:n { #1 } { #2 }
  }
\NewDocumentCommand \myparg { O{} m }
  {
    \__mydoc_meta_generic:Nnn \__mydoc_parg:n { #1 } { #2 }
  }
% Variables and keys
\tl_new:N \l_userdoc_meta_font_tl
\keys_define:nn { userdoc / mymeta }
  {
    type .choice:,
    type / tt .code:n = \tl_set:Nn \l_userdoc_meta_font_tl { \ttfamily },
    type / rm .code:n = \tl_set:Nn \l_userdoc_meta_font_tl { \rmfamily },
    type .initial:n = tt,
    cf   .tl_set:N = \l_userdoc_meta_color_tl,
    cf   .initial:n = keyname,
    ac   .tl_set:N = \l_userdoc_meta_anglecolor_tl,
    ac   .initial:n = lightgray,
    sbc  .tl_set:N = \l_userdoc_meta_brackcolor_tl,
    sbc  .initial:n = gray,
    cbc  .tl_set:N = \l_userdoc_meta_bracecolor_tl,
    cbc  .initial:n = gray,
    pbc  .tl_set:N = \l_userdoc_meta_parencolor_tl,
    pbc  .initial:n = gray,
  }
% Internal commands
\cs_new_protected:Npn \__mydoc_meta_generic:Nnn #1 #2 #3
  {
    \group_begin:
      \keys_set:nn { userdoc / mymeta } { #2 }
      \color{ \l_userdoc_meta_color_tl }
      \l_userdoc_meta_font_tl
      #1 { #3 } % #1 is \__mydoc_meta:n, \__mydoc_marg:n, \__mydoc_oarg:n or \__mydoc_parg:n
    \group_end:
  }
\cs_new_protected:Npn \__mydoc_meta:n #1
  {
    \__mydoc_meta_angle:n { \textlangle }
    \__mydoc_meta_meta:n { #1 }
    \__mydoc_meta_angle:n { \textrangle }
  }
\cs_new_protected:Npn \__mydoc_marg:n #1
  {
    \__mydoc_meta_brace:n { \textbraceleft }
    \__mydoc_meta:n { #1 }
    \__mydoc_meta_brace:n { \textbraceright }
  }
\cs_new_protected:Npn \__mydoc_oarg:n #1
  {
    \__mydoc_meta_brack:n { [ }
    \__mydoc_meta:n { #1 }
    \__mydoc_meta_brack:n { ] }
  }
\cs_new_protected:Npn \__mydoc_parg:n #1
  {
    \__mydoc_meta_brace:n { ( }
    \__mydoc_meta:n { #1 }
    \__mydoc_meta_brace:n { ) }
  }
\cs_new_protected:Npn \__mydoc_meta_meta:n #1
  {
    \textnormal{\textit{#1}}
  }
\cs_new_protected:Npn \__mydoc_meta_angle:n #1
  {
    \group_begin:
      \fontfamily{lmrotf}\fontshape{n}\selectfont
      \textcolor{\l_userdoc_meta_anglecolor_tl}{#1}
    \group_end:
  }
\cs_new_protected:Npn \__mydoc_meta_brace:n #1
  {
    \group_begin:
      \color{\l_userdoc_meta_bracecolor_tl}
        #1
    \group_end:
  }
\cs_new_protected:Npn \__mydoc_meta_brack:n #1
  {
    \textcolor{\l_userdoc_meta_brackcolor_tl}{#1}
  }
% Color the characters matched by the regex #2 in the token list #1 with the
% palette color #3 (only the first match, or every match)
\cs_new_protected:Npn \__mydoc_color_once:Nnn #1 #2 #3
  {
    \regex_replace_once:nnN { ( #2 ) } { \c{textcolor}\cB\{#3\cE\}\cB\{\1\cE\} } #1
  }
\cs_new_protected:Npn \__mydoc_color_all:Nnn #1 #2 #3
  {
    \regex_replace_all:nnN { ( #2 ) } { \c{textcolor}\cB\{#3\cE\}\cB\{\1\cE\} } #1
  }
% Rewrite \textborn | textasteriskcentered
\DeclareTextSymbol{\textbornstd}{TS1}{98}
\UndeclareTextCommand{\textborn}{TS1}
\makeatletter
\DeclareTextCommandDefault{\textbornstd}{\tc@check@symbol2\textbornstd}
\makeatother
\DeclareRobustCommand{\textborn}{{\fontfamily{lmr}\selectfont\textbornstd}}
\newcommand{\oldtextasteriskcentered}{{\fontfamily{lmr}\selectfont\textasteriskcentered}}
% \envexamp{m}
\newsavebox{\boxexaenv}
\tl_new:N \l_tempc_tl
\NewDocumentCommand \envexamp { m }
  {
    \tl_set:Nn \l_tempc_tl { #1 }
    \__mydoc_color_once:Nnn \l_tempc_tl { \* } { star }
    \begin{lrbox}{\boxexaenv}%
      \begin{minipage}[t]{3\marginparwidth}%
        \raggedright\ttfamily\small
        \textcolor{gray}{\textbackslash begin\{\textcolor{module_name}{{\tl_use:N \l_tempc_tl}}\}\myoarg
        [type=tt,sbc=gray,ac=lightgray,cf=keyname]{keyval~list}}\par%
        \hspace{0.25cm}\ics*{item}\space\mymeta[ac=lightgray,cf=gray]{item~content}\par%
        \hspace{0.25cm}\ics*{item}\space\myoarg[sbc=gray,ac=lightgray,cf=keyname]{custom}\space\mymeta[ac=lightgray,cf=gray]{item~content}\par
        \hspace{0.25cm}\ics*{item*}\myoarg[sbc=gray,ac=lightgray,cf=keyname]{symbol}\myoarg[sbc=gray,ac=lightgray,cf=keyname]{offset}\space\mymeta[ac=lightgray,cf=gray]{item~content}\par
        \textcolor{gray}{\textbackslash end\{{\textcolor{module_name}{{\tl_use:N \l_tempc_tl}}}\}}\par
      \end{minipage}%
    \end{lrbox}%
    \usebox{\boxexaenv}
  }
% \cmdexamp{o m o m o}
\tl_new:N \l_tempd_tl
\NewDocumentCommand \cmdexamp { o m o m o }
  {
    \group_begin:
      \small\ttfamily
      \textcolor{module_name}{\textbackslash#2}
      \IfValueT{#1}{ \textcolor{star}{#1} }
      \IfValueT{#3}
        {
          \tl_set:Nn \l_tempd_tl { #3 }
          \regex_replace_once:nnN { (\*) } { \c{texttt}\cB\{\c{textcolor}\cB\{star\cE\}\cB\{\1\cE\} \cE\} } \l_tempd_tl
          \myoarg[type=tt,sbc=gray,ac=lightgray,cf=keyname]{\tl_use:N \l_tempd_tl}
        }
      \IfValueTF{#5}
        {
          \mymeta[ac=gray,type=tt,cf=star]{#5}%
          \mymeta[type=tt,ac=gray,cf=keyname]{#4}
          \mymeta[ac=gray,type=tt,cf=star]{#5}%
        }
        { \mymarg[type=tt,cbc=gray,ac=lightgray,cf=keyname]{#4} }
      \par
    \group_end:
  }
% Index entries for the keys. #1 is the target: 'env' (environments in
% general), a name ending in '*' (an environment such as anskey*) or the
% name of a command (printed with a backslash). #2 is the key.
\tl_new:N \l__mydoc_key_target_tl
\cs_new_protected:Npn \__mydoc_key_index:nn #1 #2
  {
    \str_case:nnF {#1}
      {
        { env }
          {
            \__mydoc_sort_index:nn { Keys }
              { Keys~for~environments~provided~by~\pkglogo :> \texttt {#2} }
          }
      }
      {
        \tl_set:Nn \l__mydoc_key_target_tl {#1}
        \tl_if_in:NnTF \l__mydoc_key_target_tl { * }
          {
            \__mydoc_color_once:Nnn \l__mydoc_key_target_tl { \* } { star }
            \__mydoc_sort_index:nn { Keys }
              {
                Keys~for~\texttt { \tl_use:N \l__mydoc_key_target_tl }~
                provided~by~\pkglogo :> \texttt {#2}
              }
          }
          {
            \__mydoc_sort_index:nn { Keys }
              { Keys~for~\texttt { \textbackslash #1 }~provided~by~\pkglogo :> \texttt {#2} }
          }
      }
  }
\cs_new_protected:Npn \__mydoc_key_index_list:nn #1 #2
  {
    \clist_map_inline:nn {#1}
      { \__mydoc_key_index:nn {##1} {#2} }
  }
% Key line: name hanging in the margin, value, and default
% \__mydoc_key_line:Nnnn {boolean} {key} {value} {default}
\cs_new_protected:Npn \__mydoc_key_line:Nnnn #1 #2 #3 #4
  {
    \group_begin:
      \noindent
      \mode_leave_vertical:
      \hbox_to_wd:nn { -\marginparsep }
        {
          \tex_hss:D
          \textcolor { keyname } { \small \ttfamily {#2} }
        }
      \IfBooleanTF {#1}
        {
          \hphantom { \textcolor { white } { \, \bfseries \texttt{=} } \, {} }
          \mymeta [ type=tt, cbc=gray, ac=lightgray, cf=lightgray ] { \small {#3} }
        }
        {
          \textcolor { gray } { \, \bfseries \texttt{=} } \, {}
          \mymarg [ type=tt, cbc=gray, ac=lightgray, cf=keyname ] { \small {#3} }
        }
      \hfill
      \textcolor { gray } { \small \textsf { default } :~ \emph{#4} }
      \par
    \group_end:
  }
% \keyexampcmd{s m m m o}, the optional #5 is a clist of targets for the
% index (env, anskey*, a command name); without it there is no index entry
\NewDocumentCommand \keyexampcmd { s m m m o }
  {
    \__mydoc_key_line:Nnnn #1 {#2} {#3} {#4}
    \IfValueT {#5}
      { \__mydoc_key_index_list:nn {#5} {#2} }
  }
% \mykey{o m}, the optional #1 is a clist of targets for the index
\NewDocumentCommand \mykey { o m }
  {
    \group_begin:
      \tl_set:Nn \l_tmpa_tl { #2 }
      \__mydoc_color_once:Nnn \l_tmpa_tl { \* } { star }
      \textcolor { keyname } { \ttfamily \tl_use:N \l_tmpa_tl }
      \IfValueT {#1}
        { \__mydoc_key_index_list:nn {#1} { \tl_use:N \l_tmpa_tl } }
    \group_end:
  }
% \mypkg{s O{} m}, use the optional argument [no-index] to skip the index
\NewDocumentCommand \mypkg { s O{} m }
  {
    \group_begin:
      \IfBooleanTF {#1}
        { \pkglogo }
        { \textcolor { pkgname } { \textbf { \texttt {#3} } } }
      \str_if_eq:nnF {#2} { no-index }
        { \__mydoc_sort_index:nn { packages } { Packages~:>~\texttt {#3} } }
    \group_end:
  }
% \myclass{m}
\NewDocumentCommand \myclass { m }
  {
    \mbox { \textcolor { clsname } { \textbf { \texttt {#1} } } }
    \__mydoc_sort_index:nn { class } { Document~class~:>~\texttt {#1} }
  }
% \myenv{s O{} m}, use the optional argument [no-index] to skip the index
\NewDocumentCommand \myenv { s O{} m }
  {
    \IfBooleanTF {#1}
      {
        \tl_set:Nn \l_tmpa_tl { #3 }
        \__mydoc_color_once:Nnn \l_tmpa_tl { \* } { star }
        \textcolor { module_name } { \ttfamily { \tl_use:N \l_tmpa_tl } }
        \str_if_eq:nnF {#2} { no-index }
          {
            \__mydoc_sort_index:nn { environment }
              { Environments~provided~by~\pkglogo :> \texttt { \tl_use:N \l_tmpa_tl } }
          }
      }
      {
        \textcolor { pkgname } { \ttfamily {#3} }
        \str_if_eq:nnF {#2} { no-index }
          { \__mydoc_sort_index:nn { environment } { Environments~:>~\texttt {#3} } }
      }
  }
% \ics{s m}
\NewDocumentCommand \ics { s m }
  {
    \IfBooleanTF {#1}
      {
        \tl_set:Nn \l_tmpa_tl { #2 }
        \__mydoc_color_once:Nnn \l_tmpa_tl { \* } { star }
        \__mydoc_color_once:Nnn \l_tmpa_tl { ! } { star }
        \textcolor { module_name } { \ttfamily \textbackslash { \tl_use:N \l_tmpa_tl } }
        \__mydoc_sort_index:nn { Commands }
          { Commands~provided~by~\pkglogo :> \texttt { \textbackslash #2 } }
      }
      {
        \textcolor { latex_cmd } { \ttfamily \textbackslash {#2} }
        \__mydoc_sort_index:nn {#2} { \texttt { \textbackslash #2 } }
      }
  }
% \icounter{O{} m}, use the optional argument [no-index] to skip the index
\NewDocumentCommand \icounter { O{} m }
  {
    \textcolor { keyname } { \ttfamily #2 }
    \str_if_eq:nnF {#1} { no-index }
      {
        \__mydoc_sort_index:nn { Counters }
          { Counters~defined~by~\pkglogo :> \texttt {#2} }
      }
  }
% \idimen{O{} m}, use the optional argument [no-index] to skip the index
\NewDocumentCommand \idimen { O{} m }
  {
    \textcolor { module_name } { \ttfamily \textbackslash #2 }
    \str_if_eq:nnF {#1} { no-index }
      {
        \__mydoc_sort_index:nn {#2} { \texttt { \textbackslash #2 } }
        \__mydoc_sort_index:nn { Dimension }
          { Dimension~defined~by~\pkglogo :> \texttt { \textbackslash #2 } }
      }
  }
% \ilabel{s O{} m}, use the optional argument [no-index] to skip the index
\NewDocumentCommand \ilabel { s O{} m }
  {
    \tl_set:Nn \l_tmpb_tl { #3 }
    \tl_remove_once:Nn \l_tmpb_tl { . }
    \tl_remove_once:Nn \l_tmpb_tl { ) }
    \tl_remove_once:Nn \l_tmpb_tl { ( }
    \str_if_eq:nnF {#2} { no-index }
      {
        \__mydoc_sort_index:nn { Labels }
          { Labels~provided~by~\pkglogo :> \texttt { \textbackslash \tl_use:N \l_tmpb_tl } }
      }
    \tl_set:Nn \l_tmpa_tl { #3 }
    \__mydoc_color_once:Nnn \l_tmpa_tl { \* } { star }
    \__mydoc_color_once:Nnn \l_tmpa_tl { \. } { star }
    \__mydoc_color_once:Nnn \l_tmpa_tl { \( } { star }
    \__mydoc_color_once:Nnn \l_tmpa_tl { \) } { star }
    \tl_put_left:Nn \l_tmpa_tl { \textbackslash }
    \IfBooleanT {#1}
      { \tl_put_left:Nn \l_tmpa_tl { \textcolor { star } { ( } } }
    \textcolor { keyname } { \ttfamily \tl_use:N \l_tmpa_tl }
  }
% \mydim{m}
\NewDocumentCommand \mydim { m }
  {
    \group_begin:
      \tl_set:Nn \l_tmpa_tl { #1 }
      \__mydoc_color_all:Nnn \l_tmpa_tl { \d } { number }
      \__mydoc_color_all:Nnn \l_tmpa_tl { \. } { number }
      \textcolor { unit } { \ttfamily \tl_use:N \l_tmpa_tl }
    \group_end:
  }
% \myfile{m}
\NewDocumentCommand \myfile { m }
  {
    \group_begin:
      \tl_set:Nn \l_tmpa_tl { #1 }
      \__mydoc_color_all:Nnn \l_tmpa_tl { \d } { number }
      \__mydoc_color_all:Nnn \l_tmpa_tl { \. } { star }
      \textcolor { filename } { \ttfamily \tl_use:N \l_tmpa_tl }
    \group_end:
  }
% \myattachfile{s m}, the star version attaches only the \texttt{.tex} file
\NewDocumentCommand \myattachfile { s m }
  {
    \group_begin:
      \IfBooleanTF { #1 }
        {
          \IfFileExists { #2.tex }
            {
              Source~file~
              \textattachfile[color=newlinkcolor,print=false]{#2.tex}{\faFile*[regular]}.
            }
            { }
        }
        {
          \IfFileExists { #2.pdf }
            {
              Source~file~
              \textattachfile[color=newlinkcolor,print=false]{#2.tex}{\faFile*[regular]}~
              and~\emph{tagged}~\texttt{PDF}~
              \textattachfile[color=newlinkcolor,print=false]{#2.pdf}{\faFilePdf[regular]}.
            }
            { }
        }
    \group_end:
  }
% Icons in the margin for the environment important (internal use)
\cs_new_protected:Npn \__mydoc_margin_icon:n #1
  {
    \par
    \mode_leave_vertical:
    \skip_horizontal:n { -0.5 \marginparsep }
    \hbox_overlap_left:n
      {
        \hbox_to_wd:nn { \marginparsep }
          {
            \tex_hss:D
            \textcolor { module_at } { #1 } % privatefun
          }
      }
    \skip_horizontal:n { 0.5 \marginparsep }
  }
% Space before the function environment
\hook_gput_code:nnn { env/function/before } { mydoc }
  {
    \vspace* { -\baselineskip }
  }
% Environment for notes
\NewDocumentEnvironment { important } { s }
  {
    \par\addvspace { \smallskipamount }
    \small\sffamily
    \IfBooleanTF { #1 }
      {
        \__mydoc_margin_icon:n { \footnotesize \faIcon { bomb } }
      }
      {
        \__mydoc_margin_icon:n { \scriptsize \reflectbox { \faIcon { tag } } }
      }
    \ignorespaces
  }
  { }
% Parse and format an e-mail address
\seq_new:N \l__mydoc_email_seq
\cs_new_protected:Npn \__mydoc_format_email:n #1
  {
    \seq_set_split:Nnn \l__mydoc_email_seq { @ } { #1 }
    \href { mailto:#1 }
      {
        \ttfamily
        \seq_item:Nn \l__mydoc_email_seq { 1 }
        \textcolor { arroba } { \footnotesize @ }
        \seq_item:Nn \l__mydoc_email_seq { 2 }
      }
  }
\NewDocumentCommand \email { m }
  {
    \__mydoc_format_email:n { #1 }
  }
% Don't copy numbers in code example
\newcommand*{\noaccsupp}[1]
  {
    \__mylhmc_luafun_actual_text_begin: #1 \__mylhmc_luafun_actual_text_end:
  }
\ExplSyntaxOff
% Create a language for documentation
\lstdefinelanguage{enumext-doc}{
    texcsstyle    = *,%
    escapechar    =`,%
    extendedchars = true,%
    keepspaces    = true,%
    stringstyle   = {\color{red}},%
    basicstyle    = \small\ttfamily,% \small
    alsoletter    = {-,.},% @,
% comments
    morecomment   = [l]{\%},%
    commentstyle  = \ttfamily\small\color{comment},%
% Important \documentclass \begin \end \makeatletter \NewDocumentCommand
    keywordstyle  = [1]{\bfseries\color{structure}},%
    keywords      = [1]{begin, end, document, documentclass, makeatletter, makeatother, NewDocumentCommand,%
                    RenewDocumentCommand, NewCommandCopy},%
% Other words \usepackage
    keywordstyle  = [2]{\color{pkgstruct}},%
    keywords      = [2]{usepackage},%
% Other words 3 tagpdf
    keywordstyle  = [3]{\color{mathtag}},%
    keywords      = [3]{DocumentMetadata, IfDocumentMetadataTF, IfDocumentMetadataF, SuspendTagging, ResumeTagging},%
% Other words 4 (LaTeX2e)
    keywordstyle  = [4]{\color{latex_cmd}},%
    keywords      = [4]{section, input, listparindent, linewidth, emph, textbf, makesavenoteenv, makelabel,%
                    columnseprule, footnote, footnotemark, footnotetext, raggedcolumns, hspace, parindent,%
                    removelastskip, subsection, list, endlist, leftmargin, itemindent, bfseries, endminipage,%
                    minipage, parskip, raisebox, strut, smash, labelwidth, labelsep, labelindent, label,%
                    topsep, parsep, includegraphics, itemize, description, multicolsep, leftmargini,%
                    leftmarginii, hfill, small, vspace, par, parbox, value, noindent, makebox, textcolor,%
                    leftmarginiii, leftmarginiv, itemsep, strutbox, raggedright, cs_new_eq, fbox, addvspace,%
                    Alph, alph, arabic, Roman, roman, partopsep, baselineskip, mode_leave_vertical, textborn,%
                    textbullet, textendash, textasteriskcentered, textperiodcentered, raggedleft,%
                    peek_remove_spaces, skip_horizontal, triangleright, diamond, circ, parleft, star,%
                    centering, ast, hyperlink, hypertarget, newtheorem, counterwithin, aux_i, aux_ii, aux_iii,%
                    exp_args, if@noitemarg, if, if@hyper@item, hyper, noitemarg, regex_match, regex_if_match,%
                    skip_vertical, obeyedline, tl_replace_all, mbox, tl_rescan, tl_set_rescan, tex_scantokens,%
                    setcounter, .code, .meta, marginpar, c_zero_skip, tl_retokenize, keys_define, exp_after,%
                    settowidth, dimeval, dimexpr, IfBooleanTF, IfHyperBoolean, newlength, setlength, pagestyle,%
                    definecolor, chapter},%
% Other words 5 (package name)
    keywordstyle  = [5]{\bfseries\color{pkgname}},%
    keywords      = [5]{footnotehyper},%
% Reserved words 6 (document class)
    keywordstyle  = [6]{\bfseries\color{clsname}},%
    keywords      = [6]{article},%
% Reserved words 7 (enumext cmd)
    keywordstyle  = [7]{\color{module_name}},%
    keywords      = [7]{enumext, keyans, keyanspic, getkeyans, anskey, anspic, item, anskey*, itemwidth,%
                    enumextmeta, setenumext, printkeyans, miniright, labelbx, mywrap, descitemwd, theenumXi,%
                    theenumXii, theenumXiii, setenumextmeta, foreachkeyans, theenumX, theenumXiv, theenumXv,%
                    theenumXvii, theenumXviii, resetenumext},%
% Reserved words 8 (keyname) KEYS and counters
    keywordstyle  = [8]{\color{keyname}},%
    keywords      = [8]{save-ans, wrap-ans, save-ref, show-pos, show-ans, columns, list-offset, enumXi, enumX,%
                    nosep, before, after, start, right, item-sym, item-sep, hyperfootnotes, item-star, scale,%
                    wrap-label, font, first, widest, mini-env, list-indent, no-store, save-key, save-sep, ref,%
                    store-env, print-env, write-out, force-eol, write-env, overwrite, unknown, base-fix,%
                    midsep, wrapper, label-pos, label-sep, layout-sty, layout-sep, layout-top, lang, tagging,%
                    pdfversion, pdfstandard, testphase, alt, mark-sep, wrap-ans*, mark-ans, beginpenalty,%
                    midpenalty, endpenalty, reset, resume, resume*, reset*},%
% Reserved words 9 (values)
    keywordstyle  = [9]{\color{ltcmd_arg}},%
    keywords      = [9]{m, s, M, VIII, viii, v, E, t, l, A, XXIII, true, false, test, sample, above, below, ua-2,%
                    phase-III, en-US, math, title, table, firstaid, on, Verbatim, env, myseq, the, red, blue,%
                    gray, pt, em},%
% Reserved words 10 (signature)
    keywordstyle  = [10]{\color{sigargs}},%
    keywords      = [10]{J, D, N, NN, Nee, Nnn, nn, Ne, w, n, wN},%
% Reserved words 12 (plain TeX or internal)
    keywordstyle  = [12]{\color{danger}},%
    keywords      = [12]{let, verb, def, prevdepth, glue, if@minipage, splittopskip, showboxdepth, topskip,%
                    showoutput, nointerlineskip, ignorespaces, unskip, unkern, lastnodetype, rightskip,%
                    @rightskip, flushglue, @flushglue, hbadness, hss, scantokens, beginparpenalty, itempenalty,%
                    endparpenalty, lowpenalty},%
% literate
    literate      =*{\{}{{{\color{brace}{\char`\{}}}}{1}
                    {\}}{{{\color{brace}{\char`\}}}}}{1}
                    {\|}{{{\color{gray}{\char`\|}}}}{1}
                    {\$}{{{\color{math}{\char`\$}}}}{1}
                    {\#}{{{\color{argument}{\char`\#}}}}{1}
                    {(}{{\textcolor{gray}{(}}}{1}
                    {)}{{\textcolor{gray}{)}}}{1}
                    {[}{{\textcolor{bracket}{[}}}{1}
                    {]}{{\textcolor{bracket}{]}}}{1}
                    {,}{{\textcolor{gray}{,}}}{1}
                    {;}{{\textcolor{gray}{;}}}{1}
                    {.}{{\textcolor{signature}{.}}}{1}
                    {:}{{\textcolor{signature}{:}}}{1}
                    {`}{{\textcolor{backtick}{\textasciigrave}}}{1}
                    {\&}{{\textcolor{gray}{\&}}}{1}
                    {>}{{\textcolor{quote}{\guillemotright}}}{1}
                    {<}{{\textcolor{quote}{\guillemotleft}}}{1}
                    {*}{{\textcolor{star}{*}}}{1}
                    {"}{{\textcolor{quote}{"}}}{1}
                    {!}{{\textcolor{star}{!}}}{1}
                    {@}{{\textcolor{arroba}{@}}}{1}%
                    {\^}{{\textcolor{keyname}{\textasciicircum}}}{1}
                    {0}{{\textcolor{argnum}{0}}}{1}
                    {1}{{\textcolor{argnum}{1}}}{1}
                    {2}{{\textcolor{argnum}{2}}}{1}
                    {3}{{\textcolor{argnum}{3}}}{1}
                    {4}{{\textcolor{argnum}{4}}}{1}
                    {5}{{\textcolor{argnum}{5}}}{1}
                    {6}{{\textcolor{argnum}{6}}}{1}
                    {7}{{\textcolor{argnum}{7}}}{1}
                    {8}{{\textcolor{argnum}{8}}}{1}
                    {9}{{\textcolor{argnum}{9}}}{1}
                    {.0}{{\textcolor{argnum}{.0}}}{2}% Following is to ensure that only periods
                    {.1}{{\textcolor{argnum}{.1}}}{2}% followed by a digit are changed.
                    {.2}{{\textcolor{argnum}{.2}}}{2}
                    {.3}{{\textcolor{argnum}{.3}}}{2}
                    {.4}{{\textcolor{argnum}{.4}}}{2}
                    {.5}{{\textcolor{argnum}{.5}}}{2}
                    {.6}{{\textcolor{argnum}{.6}}}{2}
                    {.7}{{\textcolor{argnum}{.7}}}{2}
                    {.8}{{\textcolor{argnum}{.8}}}{2}
                    {.9}{{\textcolor{argnum}{.9}}}{2}
                    {=}{{\textcolor{compare}{=}}}{1},%
}[keywords,tex,comments,strings]% end language

% \begin{examplecode}[key=val]...\end{examplecode}
\lstnewenvironment{examplecode}[1][]{%
\lstset{
    language    = enumext-doc,%
    stringstyle = {\color{red}},%
    basicstyle  = \ttfamily\small,%
    columns     = fullflexible,%
    fontadjust  = true,%
    aboveskip   = 8pt,%
    belowskip   = 5pt,%
    rulecolor   = \color{gray!50},%
    framesep    = \fboxsep,%
    framerule   = \fboxrule,%
    xleftmargin = \dimexpr\fboxsep+\fboxrule\relax,%
    xrightmargin= \dimexpr\fboxsep+\fboxrule\relax,%
           #1,%
    }% close lstset
}{\vspace{-\parskip}}% close examplecode
% \lstinline[style=inline]|...|
\lstdefinestyle{inline}
  {
    language    = enumext-doc,%
    upquote     = true,%
    numbersep   = 1em,%
    escapechar  =`,%
    basicstyle  =\ttfamily,%\small
    numberstyle = \tiny\color{lightgray}\noaccsupp,%
    literate    =*{\%}{{\textcolor{percent}{\%}}}{1}
                  {i}{{\textcolor{keyname}{i}}}{1}
                  {ii}{{\textcolor{keyname}{ii}}}{2}
                  {iii}{{\textcolor{keyname}{iii}}}{3}
                  {iv}{{\textcolor{keyname}{iv}}}{2}
                  {v}{{\textcolor{keyname}{v}}}{1}
                  {vi}{{\textcolor{keyname}{vi}}}{1}
  }
% Set default \lstinline style
\lstset{style=inline}
% Configuration titleps and titlesec
\settitlemarks{section}
\renewpagestyle{plain}[\color{keyname}\small\sffamily]{%
\setfoot{\rlap{\hskip\dimexpr-\oddsidemargin-1in\relax%
         \parbox{1.93\paperwidth}{\hfil\thepage\,/\,\pageref{LastPage}}}}%
        {}%
        {}%
}
\newpagestyle{myheader}[\color{keyname}\small\sffamily]{%
\renewcommand\makeheadrule{%
\rlap{\hskip\dimexpr-\oddsidemargin-1in\relax
      \color{rulecolor}\rule[0.3cm]{\paperwidth}{0.7cm}}\hss
}
\setfoot{\rlap{\hskip\dimexpr-\oddsidemargin-1in\relax%
         \parbox{1.93\paperwidth}{\hfil\thepage\,/\,\pageref{LastPage}}}}%
        {\parbox{\textwidth}{\raggedright \textcolor{gray}{\raisebox{-1pt}{\textcopyright}{}2024--2026 by Pablo González L}}}%
        {}%
\sethead{\llap{\raisebox{0.55cm}{\parbox{\dimexpr\oddsidemargin+1in\relax}{\makebox[0pt][l]{\hspace{15pt}\pkglogo\space\fileversion}}}}}
        {\raisebox{0.55cm}{\parbox{\textwidth}{\hspace*{-\oddsidemargin}\centering\small\S.\thesection\space\sectiontitle}}}%
        {}%
}
\titlecontents{section}[0mm]{}%
    {\bfseries\contentspush{\makebox[5mm][l]{\thecontentslabel\hfill}}}%
    {\hspace*{-5mm}}% numberless
    {\hspace{0.25em}\titlerule*[6pt]{.}\contentspage}%
\titlecontents{subsection}[5mm]{}%
    {\contentspush{\makebox[6mm][l]{\thecontentslabel\hfill}}}
    {\hspace*{-11mm}}% numberless
    {\hspace{0.25em}\titlerule*[6pt]{.}\contentspage}%
\titlecontents{subsubsection}[11mm]{}%
    {\contentspush{\makebox[8mm][l]{\thecontentslabel\hfill}}}
    {\hspace*{-19mm}}% numberless
    {\hspace{0.25em}\titlerule*[6pt]{.}\contentspage}%
% Table of contents
\ExplSyntaxOn
\cs_new_eq:Nc \__mydoc_mkboth:nn { @mkboth }
\cs_new_eq:Nc \__mydoc_starttoc:n { @starttoc }
\RenewDocumentCommand \tableofcontents { }
  {
    \group_begin:
    % Restore the l3doc \parskip locally (only for the TOC)
    \skip_set:Nn \parskip { 0pt~plus~3pt~minus~0pt }
    \section*
      {
        \contentsname \quad
          {
            \color { gray }
            \tex_leaders:D \tex_hrule:D height 5pt depth -4.4pt \tex_hfill:D
          }
      }
    \__mydoc_mkboth:nn { \text_uppercase:n { \contentsname } } { \text_uppercase:n { \contentsname } }
    \vspace* { -14pt }
    \dim_set:Nn \columnsep { 10pt }
    \begin{multicols}{2}
      \__mydoc_starttoc:n { toc }
    \end{multicols}
    \vspace* { -3pt }
      {
        \color { gray }
        \tex_hrule:D height 0.6pt
      }
    \vspace* { 5pt }
    \group_end:
  }
% \thanks
\cs_new:Npn \__mydoc_fnsymbol:n #1
  {
    \ensuremath
      {
        \int_case:nnF { #1 }
          {
            { 0 } { }
            { 1 } { \textcolor { module_name } { * } }
            { 2 } { \textcolor { module_name } { \dagger } }
            { 3 } { \ddagger }
            { 4 } { \mathsection }
            { 5 } { \mathparagraph }
            { 6 } { \| }
            { 7 } { ** }
            { 8 } { \dagger \dagger }
            { 9 } { \ddagger \ddagger }
          }
          { \use:c { @ctrerr } }
      }
  }
\hook_gput_code:nnn { begindocument } { enumext }
  {
    \cs_set_eq:cN { @fnsymbol } \__mydoc_fnsymbol:n
  }
\ExplSyntaxOff
% \usepackage{lua-visual-debug}
\usepackage[font=small,skip=2pt]{caption}
% For input samples
\usepackage{tikzducks,siunitx,amssymb}
\sisetup{output-decimal-marker={,}}
\newcommand*{\myCircled}[1]{%
  \begin{tikzpicture}[baseline=(number.base)]
    \node[circle, draw=green!10!gray, inner sep=1pt, fill=blue!30!green!10] (number) {#1};
  \end{tikzpicture}%
}
\usetikzlibrary{shapes.geometric}
\newcommand*{\mySquared}[1]{%
  \begin{tikzpicture}[baseline=(number.base),square/.style={regular polygon,regular polygon sides=4}]
    \node[square, rounded corners=1pt, inner sep=1pt, draw=none, fill=gray!40] (number) {\makebox[\height]{#1}};
  \end{tikzpicture}%
}
% For SAT sample
\newsavebox{\satbox}
\sbox\satbox{\small\textbf{67}}%
\NewDocumentCommand \stylesat { m }
  {%
    \setlength{\fboxsep}{2pt}%
    \raisebox{2\fboxsep}{%
    \noindent\colorbox{black}{\textcolor{white}{\makebox[\dimeval{\labelwidth-2\fboxsep}][c]{\small\textbf{\vphantom{67}#1}}}}%
    \colorbox{lightgray}{\phantom{\rule[-\dp\satbox]{\dimeval{\linewidth-2\fboxsep}}{\dimeval{\dp\satbox+\ht\satbox}}}}}%
  }
% For non-numbered list samples, need ´%´ at end (prevent wrong unwanted spaces)
\NewDocumentCommand \labelbx { s +m }
  {%
    \IfBooleanTF{#1}
      {\strut\smash{\parbox[t]{\labelwidth}{\raggedright{#2}}}}%
      {\strut\smash{\parbox[t]{\labelwidth}{\raggedleft{#2}}}}%
  }
% For azetina sample
\newlength{\itemwd}
\settowidth{\itemwd}{\bfseries 10.} % max width
\NewDocumentCommand \rating { }{8.5a, KSC}
\ExplSyntaxOn
\bool_new:N \l_azetina_rating_bool
\NewDocumentCommand \showrating { m }
  {
    \str_case:nn {#1}
      {
        { true  }
                {
                  \bool_set_true:N \l_azetina_rating_bool
                  \setenumext{show-ans=true,save-sep={ ~ }}
                }
        { false }
                {
                  \bool_set_false:N \l_azetina_rating_bool
                  \setenumext{show-ans=false,save-sep={ ~ }}
                }
      }
  }
\NewDocumentCommand \myans { m +m }
  {
    \bool_if:NTF \l_azetina_rating_bool
      {
        \item*[(#1)] \textcolor{red}{#2}
      }
      {
        \item*[(#1)] \textcolor{black}{#2}
      }
  }
\NewDocumentCommand \azwrapp { m }
  {
    \bool_if:NTF \l_azetina_rating_bool
      {
        \makebox[\dim_eval:n{\marginparwidth}][r]
          {
            \textcolor{red}{\bfseries\footnotesize\rating}
          }
        \hspace{\marginparsep}\textbf{#1}
      }
      {
        \makebox[\dim_eval:n{\marginparwidth+\marginparsep}][l]{ }\textbf{#1}
      }
  }
\ExplSyntaxOff
\setenumextmeta[1]{azetinasty}{label={\arabic*.},labelwidth={\marginparwidth+\marginparsep+\itemwd},list-offset={-\marginparwidth-\marginparsep},}
\begin{document}
  \DocInput{\jobname.dtx}
\end{document}
%</driver>
% \fi
%
% \title{
%    \scalebox{2.67}{\pkglogo}\\[-10pt]
%    {\large\textcolor{gray}{\fetamontotf{enumerate exercise sheets}}}\\[0pt]
%    \fileversion \quad \filedate\thanks{
%    This file describes a documentation for \fileversion, last revised
%    \filedate.}\\[25pt]
%    \author{
%    \large
%    \raisebox{-1pt}{\textcopyright{}2024--2026 by Pablo González L}\thanks{
%    E-mail: \textcolor{OrangeRed}{\textsf{\guillemotleft}}\email{pablgonz@educarchile.cl}\textcolor{OrangeRed}{\textsf{\guillemotright}}.
%       }%
%    }
% \small
% \textsc{ctan}: \url{https://www.ctan.org/pkg/enumext}\\
% \textcolor{gray}{\scriptsize\faIcon[regular]{github}}\,\,\,\url{https://github.com/pablgonz/enumext}
% \vspace*{-2cm}
% }
% \date{}
% \maketitle
%
% \begin{abstract}
% This package provides enumerated list environments compatible
% with \emph{tagging} \texttt{PDF} for creating \emph{\enquote{simple
% exercise sheets}} along with \emph{\enquote{multiple choice questions}},
% storing the \emph{\enquote{answers}} to these in memory using \mypkg{multicol}
% package.
% \end{abstract}
%
% \tableofcontents
%
% \section*{Motivation and acknowledgments}
%
% Usually, it is enough to use the classic \myenv[no-index]{enumerate} environment to
% generate \emph{\enquote{simple exercise sheets}} or
% \emph{\enquote{multiple choice questions}}, the basic idea behind
% \mypkg*[no-index]{enumext} is to cover three points:
%
% \begin{enumerate}[left=0pt, nosep,labelsep=5pt, label=\small\arabic*.,topsep=2pt]
%   \item To have a simple interface to be able to write \emph{\enquote{lists of
%         exercises}} with \emph{\enquote{answers}}.
%   \item To have a simple interface for writing \emph{\enquote{multiple choice
%         questions}}.
%   \item To have a simple interface for placing \emph{\enquote{columns}}
%         and \emph{\enquote{drawings}} or \emph{\enquote{tables}}.
% \end{enumerate}
%
% This package would not be possible without Phelype Oleinik who has
% collaborated and adapted a large part of the code and all
% \hologo{LaTeX} team for their great work and to the different members
% of the \href{https://tex.stackexchange.com}{TeX-SX} community who have
% provided great answers and ideas. Here are notes on the main ones:
%
% \begin{enumerate}[left=0pt, labelsep=5pt, nosep, label=\small\arabic*.,topsep=2pt]
% \item Answer given by Alan Munn in \href{https://tex.stackexchange.com/a/300418}{\textbackslash topsep,  \textbackslash itemsep, \textbackslash partopsep, \textbackslash parsep - what do they each mean (and what about the bottom)?}
%
% \item Answer given by Enrico Gregorio in \href{https://tex.stackexchange.com/a/34172}{Understanding minipages - aligning at top}
%
% \item Answer given by Ulrich Diez in \href{https://tex.stackexchange.com/a/553839/7832}{Different mechanics of hyperlink vs. hyperref}
%
% \item Answer given by Enrico Gregorio in \href{https://tex.stackexchange.com/a/294460/7832}{Minipage and multicols, vertical alignment}
%
% \end{enumerate}
%
% \thispagestyle{plain}
%
% \section*{License and Requirements}
% \label{sec:licence}
%
% Permission is granted to copy, distribute and/or modify this software under
% the terms of the LaTeX Project Public License (lppl), version 1.3 or later
% (\url{https://www.latex-project.org/lppl.txt}).
% The software has the status \enquote{maintained}.
%
% The \mypkg*[no-index]{enumext} package loads and requires \mypkg[no-index]{multicol}\cite{multicol}
% package, need to have a modern \hologo{TeX} distribution such as
% \hologo{TeX} Live or \hologo{MiKTeX}. It has been tested with the standard
% classes provided by \hologo{LaTeX}: \myclass{book}, \myclass{report},
% \myclass{article} and \myclass{letter} on \mydim{10pt}, \mydim{11pt} and \mydim{12pt}.
%
% \begin{important}*
%  The minimum requirement is \hologo{LaTeX} release 2026-11-01.
% \end{important}
%
% \pagestyle{myheader}
%
% \section{Introduction}
%
% In the \hologo{LaTeX} world there are many useful packages and
% classes for creating \emph{\enquote{lists of exercises}}, \emph{\enquote{worksheets}}  or
% \emph{\enquote{multiple choice questions}}, classes like \myclass{exam}\cite{exam}
% and packages like \mypkg{xsim}\cite{xsim} do the job perfectly, but
% they don't always fit the basic day to day needs.
%
% In the work of many teachers it is common to use
% \emph{\enquote{simple exercise sheets}} also known as
% \emph{\enquote{informal lists of exercises}}, for example:
%
% \begin{enumext}[save-ans=temp2,columns=2,nosep,save-ref=true]%
%  \item Factor $x^{2}-2x+1$  \anskey{$\left(x-1\right)^{2}$}
%  \item Factor $3x+3y+3z$    \anskey{$3(x+y+z)$}
%  \item True False
%    \begin{enumext}[nosep]
%      \item $\alpha > \delta$ \anskey{False}
%      \item \LaTeX2e\ is cool? \anskey{Very True!}
%    \end{enumext}
%  \item Related to Linux
%    \begin{enumext}[nosep]
%      \item You use Linux? \anskey{Yes}
%      \item Usually uses the package manager? \anskey{Yes, \texttt{dnf}}
%      \item Rate the following package and class
%        \begin{enumext}[nosep]
%          \item \texttt{xsim-exam} \anskey{doesn't exist for now :(}
%          \item \texttt{xsim} \anskey{very good}
%          \item \texttt{exsheets} \anskey{obsolete}
%        \end{enumext}
%    \end{enumext}
% \end{enumext}
%
% Sometimes it is also useful to show the \emph{\enquote{answers}}
% along with the questions:
%
% \begin{enumext}[save-ans=worksheet,mark-ans={\oldtextasteriskcentered},columns=2,show-ans=true,nosep,save-ref=true]
%  \item Factor $x^{2}-2x+1$ \anskey{$\left(x-1\right)^{2}$}
%  \item Factor $3x+3y+3z$ \anskey{$3(x+y+z)$}
%  \item True False
%    \begin{enumext}[nosep]
%      \item $\alpha > \delta$ \anskey{False}
%      \item \LaTeX2e\ is cool? \anskey{Very True!}
%    \end{enumext}
%  \item Related to Linux
%    \begin{enumext}[nosep]
%      \item You use Linux? \anskey{Yes}
%      \item Usually uses the package manager? \anskey{Yes, \texttt{dnf}}
%      \item Rate the following package and class
%        \begin{enumext}[nosep]
%          \item \texttt{xsim-exam} \anskey{doesn't exist for now :(}
%          \item \texttt{xsim} \anskey{very good}
%          \item \texttt{exsheets} \anskey{obsolete}
%        \end{enumext}
%    \end{enumext}
% \end{enumext}
%
% Or to refer to a specific question and its
% \emph{\enquote{answer}}, for example:
%
% The answer to \ref{worksheet:4} is \enquote{\getkeyans{worksheet:4}}
% and the answer to \ref{worksheet:8} is \enquote{\getkeyans{worksheet:8}}.
%
% Or to print all the \emph{\enquote{answers}}:
%
% \printkeyans{worksheet}
%
% Another very common thing to use is \emph{\enquote{multiple choice
% questions}}, for example:
%
% \begin{enumext}[save-ans=mytest,columns=2,nosep,save-ref=true,save-sep={\space}]
%   \item First type of questions
%      \begin{keyans}[columns=2,noitemsep]
%        \item value
%        \item*[$x=5$] correct
%        \item value
%        \item value
%      \end{keyans}
%   \item Second type of questions
%      \begin{enumext}[no-store,label=\Roman*.,noitemsep]
%        \item $2\alpha+2\delta=90^{\circ}$
%        \item $\alpha=\delta$
%        \item $\angle EDF=45^{\circ}$
%      \end{enumext}
%      \begin{keyans}[columns=2,noitemsep]
%        \item I only
%        \item II only
%        \item I and II only
%        \item* I and III only
%        \item I, II, and III
%      \end{keyans}
% \item* Third type of questions
%        \begin{enumext}[no-store,list-offset=0.5cm,label=(\arabic*),noitemsep]
%            \item $2\alpha+2\delta=90^{\circ}$
%            \item $\angle EDF=45^{\circ}$
%        \end{enumext}
%    \begin{keyans}[columns=2,noitemsep]
%        \item value
%        \item value
%        \item*[some note] value
%        \item value
%        \item value
%    \end{keyans}
%
% \item Question with image and label below:
%
%  \begin{keyanspic}[label-pos=below, layout-sty={3,2}]
%    \anspic{\includegraphics[scale=0.15]{example-image-a}}
%    \anspic{\includegraphics[scale=0.15]{example-image-b}}
%    \anspic{\includegraphics[scale=0.15]{example-image-a}}
%    \anspic{\includegraphics[scale=0.15]{example-image-a}}
%    \anspic*[A duck]{\includegraphics[scale=0.30]{example-image-duck}}
%  \end{keyanspic}
%
% \item Question with image on right side:
%
%    \begin{keyans}[mini-env=0.5\linewidth, columns=1,noitemsep]
%        \item value
%        \item value
%        \item value
%        \item*[\enquote{other note}] correct
%        \item value
%        \miniright
%        \includegraphics[scale=0.15]{example-image-b}
%    \end{keyans}
% \end{enumext}
%
% Here the \mymeta{label} and a
% \emph{\enquote{short note}} left as an explanation are stored, and then
% printed:
%
% \printkeyans{mytest}
%
% The \mypkg*[no-index]{enumext} package was created and designed to meet these
% small requirements in the creation of \emph{\enquote{simple
% worksheets}} and \emph{\enquote{multiple choice questions}}.
%
% \begin{important}*
% These \emph{\enquote{simple worksheets}} or \emph{\enquote{multiple
% choice questions}} appear to be easy to obtain using a combination of
% the \myenv[no-index]{enumerate}, \myenv[no-index]{minipage} and \myenv[no-index]{multicols}
% environments, but like many things, what \emph{\enquote{looks
% simple}} is not so simple.
% \end{important}
%
% \subsection{Description and usage}\label{subsec:description}
%
% The \mypkg*[no-index]{enumext} package defines enumerated environments using the
% \myenv[no-index]{list} environment provided by \hologo{LaTeX}\cite{latexref}, but
% \emph{\enquote{does not redefine}} any internal commands associated with
% it such as |\list|, |\endlist| or |\item| outside of the
% \emph{\enquote{scope}} in which they are defined.
%
% \begin{important}*
% This package is NOT intended to replace the \myenv[no-index]{enumerate} environment
% nor replace the powerful \mypkg{enumitem}\cite{enumitem}, the approach
% is intended to work without hindering either of them.
% \end{important}
%
% This package can be used with |xelatex|, |lualatex|, |pdflatex| and
% the classical |latex>dvips>ps2pdf| and is present in \hologo{TeX} Live and \hologo{MiKTeX},
% use the package manager to install. For manual installation, download
% \href{http://mirrors.ctan.org/macros/latex/contrib/enumext.zip}{enumext.zip}
% and unzip it, run \lstinline[language=enumext-doc,basicstyle=\ttfamily]+luatex enumext.ins+
% and move all files to appropriate locations, then run |mktexlsr|. To produce the documentation run
% \lstinline[language=enumext-doc,basicstyle=\ttfamily]+arara enumext.dtx+.
%
% \iffalse
%<*example>
% \fi
\begin{examplecode}[frame=single,language={}]
  enumext.sty  >  TDS:tex/latex/enumext/
  enumext.pdf  >  TDS:doc/latex/enumext/
  README.md    >  TDS:doc/latex/enumext/
  enumext.dtx  >  TDS:source/latex/enumext/
  enumext.ins  >  TDS:source/latex/enumext/
\end{examplecode}
% \iffalse
%</example>
% \fi
%
% The package is loaded in the usual way:
%
% \iffalse
%<*example>
% \fi
\begin{examplecode}[frame=single]
\usepackage{enumext}
\end{examplecode}
% \iffalse
%</example>
% \fi
%
% \subsection{The concept of left margin}\label{subsec:left-margin}
%
% There is a direct relationship between the parameters |\leftmargin|,
% |\itemindent|, |\labelwidth| and |\labelsep| plus an
% \emph{\enquote{extra space}}, which makes it difficult to obtain the
% desired \emph{horizontal spaces} in a \myenv[no-index]{list} environment.
% Usually the \myenv[no-index]{list} should not go beyond the left margin of the
% page, but since these four values are related, that causes a problem.
%
% The \mypkg[no-index]{enumitem}\cite{enumitem} package adds the |\labelindent| parameter to
% solve some of these problems. A simplified representation of this in
% the figure~\ref{fig:one}.
% \setlength\intextsep{3pt}
% \begin{figure}[ht]
%   \centering
%   \begin{tikzpicture}[line cap=round,line join=round,x=0.8cm,y=0.8cm,every node/.style={font=\footnotesize}]
%     \draw[<->,color=gray] (-0.5,6.175) --node[below,,text=lightgray,font=\footnotesize\ttfamily]{margin page} (2,6.175);
%     \draw[line width=0.5pt,dash pattern=on 2pt off 2pt,color=gray] (2,4.85) -- (2,7.5);
%     \draw[<->,color=gray] (2.0,7.0) --node[above,text=blue!75]{\texttt{\textbackslash{}labelindent}} (4.5,7.0);
%     \draw[<->,color=gray] (4.5,7.0) --node[above,text=blue!75]{\texttt{\textbackslash{}labelwidth}} (7.5,7.0);
%     \draw[<->,color=gray] (7.5,7.0) --node[above,text=blue!75]{\texttt{\textbackslash{}labelsep}} (9.5,7.0);
%     \draw[color=lightgray,fill=lightgray,fill opacity=0.15] (4.5,6.5) rectangle (7.5,6.85);% labelwidth
%     \draw[color=lightgray,fill=lightgray,fill opacity=0.15] (7.5,6.5) rectangle (9.5,6.85);% labelsep
%     \draw[color=lightgray,fill=lightgray,fill opacity=0.15] (9.5,6.5) rectangle (15,6.85);% item width
%     \draw[color=lightgray,fill=lightgray,fill opacity=0.15] (9.5,6.0) rectangle (15,6.35);
%     \draw[color=lightgray,fill=lightgray,fill opacity=0.15] (9.5,5.5) rectangle (15,5.85);
%     \draw[<->,color=gray] (2.0,5.35) --node[below,text=blue!75]{\texttt{\textbackslash{}leftmargin}} (7.5,5.35);
%     \draw[<->,color=gray] (7.5,5.35) --node[below,text=blue!75]{\texttt{\textbackslash{}itemindent}} (9.5,5.35);
%   \end{tikzpicture}
%   \caption{Representation of horizontal lengths in \mypkg[no-index]{enumitem}.}
%   \label{fig:one}
% \end{figure}
%
% The \mypkg*[no-index]{enumext} package does NOT provide a user interface to set the values for
% |\leftmargin| and |\itemindent|, instead it provides the keys
% \mykey{list-offset} and \mykey{list-indent} which internally set the
% values for |\leftmargin| and |\itemindent|. The concepts of |\leftmargin| and |\itemindent|
% are different in \mypkg*[no-index]{enumext}. The figure~\ref{fig:two} shows the visual representation of idea.
%
% \begin{figure}[ht]
%   \centering
%   \begin{tikzpicture}[line cap=round,line join=round,x=0.8cm,y=0.8cm,every node/.style={font=\footnotesize}]
%     \draw[<->,color=gray] (-0.5,6.175) --node[below,text=lightgray,font=\footnotesize\ttfamily]{margin page} (2,6.175);
%     \draw[line width=0.5pt,dash pattern=on 2pt off 2pt,color=gray] (2,4.85) -- (2,7.5);
%     \draw[<->,color=gray] (2.0,7.0) --node[above,text=keyname]{\vphantom{p}\texttt{list-offset}} (4.5,7.0);
%     \draw[<->,color=gray] (4.5,7.0) --node[above,text=blue!75]{\vphantom{p}\texttt{\textbackslash{}labelwidth}} (7.5,7.0);
%     \draw[<->,color=gray] (7.5,7.0) --node[above,text=blue!75]{\texttt{\textbackslash{}labelsep}} (9.5,7.0);
%     \draw[color=lightgray,fill=lightgray,fill opacity=0.15] (4.5,6.5) rectangle (7.5,6.85);% labelwidth
%     \draw[color=lightgray,fill=lightgray,fill opacity=0.15] (7.5,6.5) rectangle (9.5,6.85);% labelsep
%     \draw[color=lightgray,fill=lightgray,fill opacity=0.15] (9.5,6.5) rectangle (15,6.85);% item width
%     \draw[color=lightgray,fill=lightgray,fill opacity=0.15] (9.5,6.0) rectangle (15,6.35);
%     \draw[color=lightgray,fill=lightgray,fill opacity=0.15] (9.5,5.5) rectangle (15,5.85);
%     \draw[<->,color=gray] (4.5,5.35) --node[below,text=keyname]{\texttt{list-indent}} (9.5,5.35);
%   \end{tikzpicture}
%   \caption{Representation of horizontal lengths concept in \mypkg*[no-index]{enumext}.}
%   \label{fig:two}
% \end{figure}
%
% In this way the number of parameters to
% pass is reduced a \emph{little}. With the default values of keys \mykey{list-offset},
% \mykey{list-indent}, \mykey{labelwidth} and \mykey{labelsep} the lists
% will have the (usually) expected output for \emph{\enquote{simple
% worksheets}}. The figure~\ref{fig:three} shows the visual representation.
%
% \begin{figure}[ht]
%   \centering
%   \begin{tikzpicture}[line cap=round,line join=round,x=0.8cm,y=0.8cm,every node/.style={font=\footnotesize}]
%     \draw[<->,color=gray] (-0.5,6.175) --node[below,text=lightgray,font=\footnotesize\ttfamily]{margin page} (2,6.175);
%     \draw [line width=0.5pt,dash pattern=on 2pt off 2pt,color=gray] (2,4.85) -- (2,7.5);
%     \draw[<->,color=gray] (2.0,7.0) --node[above,text=keyname]{\vphantom{p}\texttt{labelwidth}} (5.0,7.0);
%     \draw[<->,color=gray] (5.0,7.0) --node[above,text=keyname]{\texttt{labelsep}} (7.0,7.0);
%     \draw[color=lightgray,fill=lightgray,fill opacity=0.15] (2.0,6.5) rectangle (5.0,6.85);% labelwidth
%     \draw[color=lightgray,fill=lightgray,fill opacity=0.15] (5.0,6.5) rectangle (7.0,6.85);% labelsep
%     \draw[color=lightgray,fill=lightgray,fill opacity=0.15] (7.0,6.5) rectangle (15,6.85);% item width
%     \draw[color=lightgray,fill=lightgray,fill opacity=0.15] (7.0,6.0) rectangle (15,6.35);
%     \draw[color=lightgray,fill=lightgray,fill opacity=0.15] (7.0,5.5) rectangle (15,5.85);
%     \draw[<->,color=gray] (2.0,5.35) --node[below,text=keyname]{\texttt{list-indent}} (7.0,5.35);
%   \end{tikzpicture}
%   \caption{Default horizontal lengths \mykey{list-offset}=\mydim{0pt}, \mykey{list-indent}=\ics{labelwidth}+\ics{labelsep} in \mypkg*[no-index]{enumext}.}
%   \label{fig:three}
% \end{figure}
%
% \subsection{User interface}\label{subsec:user-interface}
%
% The user interface consists of two main list environments |enumext|
% (vertical) and |enumext*| (horizontal), the environment |anskey*| and
% the command |\anskey| to \emph{\enquote{store content}} and the
% environments |keyans|, |keyans*| and |keyanspic| for multiple choice.
% It also provides the commands |\getkeyans| to print individual
% \emph{stored content}, |\printkeyans| to print all \emph{stored content},
% |\foreachkeyans| to print a range of \emph{stored content}, |\miniright|
% for \myenv[no-index]{minipage}, |\setenumext| to config \myoarg{key \textnormal{\textcolor{gray}{=}} val},
% |\setenumextmeta| to add a \emph{\enquote{meta-key}} and |\resetenumext|
% to reset counters.
%
% \subsubsection{Public counters}
%
% The package \mypkg*[no-index]{enumext} uses the \icounter{enumXi},
% \icounter{enumXii}, \icounter{enumXiii}, \icounter{enumXiv} counters
% for the \emph{four} nesting levels of the \myenv*{enumext} environment,
% the \icounter{enumXv} counter for the \myenv*[no-index]{keyans} environment, the
% \icounter{enumXvi} counter for the \myenv*[no-index]{keyanspic} environment,
% the counter \icounter{enumXvii} for \myenv*{enumext*} environment and
% the counter \icounter{enumXviii} for \myenv*[no-index]{keyans*} environment.
%
% \begin{important}*
% If any package defines these \emph{counters} or they are user-defined
% in the document, the package will return a \emph{\enquote{fatal error}}
% and abort the load.
% \end{important}
%
% \subsubsection{Public dimension}
%
% The package \mypkg*[no-index]{enumext} only provides a \emph{single public dimension}
% \idimen{itemwidth} and is intended for user convenience only and is NOT
% for internal use as such. The dimension |\itemwidth| is a \emph{rigid length}
% and contains the \emph{\enquote{width of the content}} of each |\item|
% regardless of \mykey{labelwidth} and \mykey{labelsep}.
%
% \begin{important}*
% If any package defines |\itemwidth| or the user defines |\itemwidth|
% in the document, the package will overwrite it without warning.
% \end{important}
%
% \subsubsection{Support for \pkg{multicol}}
%
% The package provides direct support for using the \mypkg[no-index]{multicol}\cite{multicol}
% package. This allows for a two-column output as shown
% in the figure~\ref{fig:fourth}.
%
% \begin{figure}[ht]
%   \centering
%   \begin{tikzpicture}[line cap=round,line join=round,x=0.8cm,y=0.8cm,every node/.style={font=\footnotesize}]
%     \draw[<->,color=gray] (2.0,7.0) --node[above,text=keyname]{\vphantom{p}\texttt{labelwidth}} (5.0,7.0);
%     \draw[<->,color=gray] (5.0,7.0) --node[above,text=keyname]{\texttt{labelsep}} (7.0,7.0);
%     \draw[color=lightgray,fill=lightgray,fill opacity=0.15] (2.0,6.5) rectangle (5.0,6.85);% labelwidth
%     \draw[color=lightgray,fill=lightgray,fill opacity=0.15] (5.0,6.5) rectangle (7.0,6.85);% labelsep
%     \draw[color=lightgray,fill=lightgray,fill opacity=0.15] (7.0,6.5) rectangle (17.5,6.85);% item width
%     \draw[color=lightgray,fill=lightgray,fill opacity=0.15] (7.0,5.5) rectangle (11.5,6.35) node[pos=.5,text=black]%
%        {\begin{tabular}{c}
%          \texttt{column one} \\
%          \textcolor{blue}{\emph{nested items}}
%        \end{tabular}};
%     \draw[color=lightgray,fill=lightgray,fill opacity=0.15] (13.5,5.5) rectangle (17.5,6.35) node[pos=.5,text=black]%
%        {\begin{tabular}{c}
%          \texttt{column two} \\
%          \textcolor{blue}{\emph{nested items}}
%        \end{tabular}};
%     \draw[<->,color=gray] (2.0,5.35) --node[below,text=keyname]{\texttt{list-indent}} (7.0,5.35);
%     \draw[<->,color=gray] (11.5,5.35) --node[below,text=keyname]%
%       {\begin{tabular}{c}
%         \texttt{column-sep} \\
%           \textcolor{blue}{\texttt{\textbackslash{}columnbreak}}
%         \end{tabular}}
%       (13.5,5.35);
%   \end{tikzpicture}
%   \caption{Representation of the two column output for a nested level in \myenv*{enumext} environment.}
%   \label{fig:fourth}
% \end{figure}
%
% The \emph{\enquote{non starred}} version of the \myenv[no-index]{multicols}
% environment is always used together with the |\raggedcolumns|
% command and is controlled by \mykey{columns} and \mykey{columns-sep}
% keys. It can be used in all nesting levels of the
% environment |enumext| and the environment |keyans| and can
% together with the \mykey{mini-env} key. If you need to force a
% start a new column |\columnbreak| must be used (see~\S\ref{subsec:keys-multicols}).
%
% \begin{important}*
% The |\columnseprule| command is not available as a key and is set to
% \emph{\enquote{zero}} for the inner levels and the \myenv*[no-index]{keyans} environment. If the value of
% this is set inside the document, it will affect \emph{\enquote{all environments}} that
% use the \mykey{columns} key.
% \end{important}
%
% \subsubsection{Support for \env{minipage}}
%
% The package provides direct support for \myenv[no-index]{minipage}
% environment, this allows you to obtain an output like the
% one shown in figure~\ref{fig:five}.
%
% \begin{figure}[!htbp]
%   \centering
%   \begin{tikzpicture}[line cap=round,line join=round,x=0.8cm,y=0.8cm,every node/.style={font=\footnotesize}]
%     \draw[<->,color=gray] (2.0,7.0) --node[above,text=keyname]{\vphantom{p}\texttt{labelwidth}} (5.0,7.0);
%     \draw[<->,color=gray] (5.0,7.0) --node[above,text=keyname]{\texttt{labelsep}} (7.0,7.0);
%     \draw[color=lightgray,fill=lightgray,fill opacity=0.15] (2.0,6.5) rectangle (5.0,6.85);% labelwidth
%     \draw[color=lightgray,fill=lightgray,fill opacity=0.15] (5.0,6.5) rectangle (7.0,6.85);% labelsep
%     \draw[color=lightgray,fill=lightgray,fill opacity=0.15] (7.0,6.5) rectangle (17.5,6.85);% item width
%     \draw[color=lightgray,fill=lightgray,fill opacity=0.15] (7.0,5.5) rectangle (12.5,6.35) node[pos=.5,text=black]%
%       {\begin{tabular}{c}
%          \texttt{minipage left} \\
%          \textcolor{blue}{\emph{nested list}}
%        \end{tabular}};
%     \draw[color=lightgray,fill=lightgray,fill opacity=0.15] (14.5,5.5) rectangle (17.5,6.35) node[pos=.5,text=black]%
%       {\begin{tabular}{c}
%          \texttt{minipage right} \\
%          \textcolor{blue}{\emph{drawing or tabular}}
%        \end{tabular}};
%     \draw[<->,color=gray] (2.0,5.35) --node[below,text=keyname]{\texttt{list-indent}} (7.0,5.35);
%     \draw[<->,color=gray] (7.0,5.35) --node[below,text=keyname]{\texttt{mini-env}} (12.5,5.35);
%     \draw[<->,color=gray] (12.5,5.35) --node[below,text=keyname]%
%       {\begin{tabular}{c}
%         \texttt{mini-sep} \\
%           \textcolor{module_name}{\texttt{\textbackslash{}miniright}}
%         \end{tabular}} (14.5,5.35);
%   \end{tikzpicture}
%   \caption{Representation of the \mykey{mini-env} output for a nested level \myenv*{enumext} environment.}
%   \label{fig:five}
% \end{figure}
%
% The \myenv[no-index]{minipage} environments on the \emph{\enquote{left side}} and
% \emph{\enquote{right side}} are always used with \emph{\enquote{aligned
% on top}} |[t]|. It can be used in all nesting levels of the
% environment |enumext| and the environment |keyans| and is controlled
% by \mykey{mini-env} and \mykey{mini-sep} keys. In order to switch from
% the \emph{\enquote{left}} side \myenv[no-index]{minipage} environment to the
% \emph{\enquote{right}} side one must use the command |\miniright|
% (see~\S\ref{subsec:keys-minipage}).
%
% \subsubsection{The \cs[no-index]{label} and \cs[no-index]{ref} system}
%
% This package provides a user interface like the
% \mypkg[no-index]{enumitem}\cite{enumitem} package to customize the references
% which is activated by the \mykey{ref} key (\S\ref{subsec:label-and-ref}), the standard \hologo{LaTeX}
% |\label| and |\ref| commands work as usual.
% It also provides an \emph{\enquote{internal reference}} system for the
% \emph{\enquote{stored content}} by means of the key \mykey{save-ref}
% (\S\ref{subsubsec:label-ref}) when the key \mykey{save-ans} is active (\S\ref{subsec:storage-keys}).
%
% \subsubsection{Support for \cs[no-index]{footnote}}\label{subsubsec:footnote}
%
% The \myenv*[no-index]{enumext*} and \myenv*[no-index]{keyans*} environments and the
% \mykey{mini-env} key use the \myenv[no-index]{minipage} environment in their
% implementation but in a transparent way for the user, it is only
% used for typesetting. The implementation supports |\footnote| and is
% compatible with \mypkg{hyperref} and works the same way as if used
% anywhere in the document.
%
% Unfortunately, if \mypkg[no-index]{hyperref} is loaded and \emph{tagged} \texttt{PDF} is not
% active, it will not produce the \enquote{links}, the internal implementation
% uses |\footnotetext|\myoarg{number} and |\footnotemark|\myoarg{number}\mymarg{text}
% and \enquote{links} for this are not supported by the \mypkg[no-index]{hyperref} package.
%
% The best way to solve this if \emph{tagged} \texttt{PDF} is NOT active
% is to use Jean-François Burnol's \mypkg{footnotehyper}\cite{footnotehyper}
% package. It will support keeping the \enquote{links} if
% \mypkg[no-index]{hyperref} is loaded with the |hyperfootnotes=true| option (default).
% Load it is as follows:
%
% \iffalse
%<*example>
% \fi
\begin{examplecode}[frame=single]
\IfDocumentMetadataF
  {
    \usepackage{`\mypkg{footnotehyper}`}
    \makesavenoteenv{enumext}
    \makesavenoteenv{enumext*}
  }
\end{examplecode}
% \iffalse
%</example>
% \fi
%
% \begin{important}
%  At the moment the \mypkg[no-index]{footnotehyper} package v1.1f (2025/11/15) is not compatible
%  with \emph{tagged} \texttt{PDF}.
% \end{important}
%
% \section{The environments provided}
%
% The package \mypkg*[no-index]{enumext} provides two main list environments, the
% \emph{vertical} environment \myenv*[no-index]{enumext} and the \emph{horizontal}
% environment \myenv*[no-index]{enumext*}.
%
% \begin{function}{enumext,enumext*}
%   \begin{minipage}[t]{0.45\linewidth}
%       \envexamp{enumext}
%   \end{minipage}\hfill
%   \begin{minipage}[t]{0.45\linewidth}
%       \envexamp{enumext*}
%   \end{minipage}
% \end{function}
%
% \subsection{The environment \texttt{enumext}}\label{subsec:envenumext}
%
% The \myenv*[no-index]{enumext} is an environment that works in the same way as
% the standard \myenv[no-index]{enumerate} environment provided by \hologo{LaTeX},
% |\item| and |\item|\myoarg{custom} commands work in
% the usual way. The environment can be nested with at most \emph{\enquote{four levels}}
% and the options can be configured globally using |\setenumext|
% command and locally using \myoarg{key \textnormal{\textcolor{gray}{=}} val}
% in the environment.
%
% \subsubsection*{Example with \texttt{columns=2}}
%
% \begin{enumext}[columns=2, topsep=0pt]
%   \item This text is in the first level.
%     \begin{enumext}
%        \item This text is in the second level.
%           \begin{enumext}
%              \item This text is in the third level.
%                \begin{enumext}
%                   \item This text is in the fourth level.
%                \end{enumext}
%            \end{enumext}
%     \end{enumext}
%   \item[X] This text is in the first level.
%   \item* This text is in the first level.
% \end{enumext}
%
% \subsection{The environment \texttt{enumext*}}\label{subsec:enumext-starred}
%
% The \myenv*[no-index]{enumext*} is a \emph{horizontal list environment}
% similar to the \myenv[no-index]{shortenumerate} or \myenv[no-index]{tasks} environments
% provided by the \mypkg{shortlst}\cite{shortlst} and \mypkg{tasks}\cite{tasks}
% packages, |\item| and |\item|\myoarg{custom} work as usual.
% The options can be configured globally using |\setenumext| command
% and locally using \myoarg{key \textnormal{\textcolor{gray}{=}} val}
% in the environment.
%
% Some considerations to take into account for this environment:
%
% \begin{itemize}[nosep,parsep=1pt,topsep=2pt]
% \item The environment cannot be nested within itself or in the
%   environment \myenv*[no-index]{keyans*}, but it can be nested within
%   \myenv*[no-index]{enumext} and vice versa.
%
% \item Each \emph{\enquote{item content}} in the environment is placed
%   within a \myenv[no-index]{minipage} environment whose \emph{width} is stored
%   in the dimension \idimen[no-index]{itemwidth} does NOT include \mykey{labelwidth} and
%   \mykey{labelsep}, only the \emph{width of the content}.
%
% \item You cannot have floating environments like \myenv[no-index]{figure} or
%   \myenv[no-index]{table} or |\marginpar| but |\footnote| with \mypkg[no-index]{hyperref}
%   is supported. If you want to activate links \emph{\enquote{without}}
%   \emph{tagged} \texttt{PDF} active you must load \mypkg[no-index]{footnotehyper}
%   and \mypkg[no-index]{hyperref} packages (see~\S\ref{subsubsec:footnote}).
%
% \item You cannot have any standard list environments like \myenv[no-index]{itemize},
%   \myenv[no-index]{enumerate}, \myenv[no-index]{description}, \myenv[no-index]{quote}, \myenv[no-index]{quotation},
%   \myenv[no-index]{verse}, \myenv[no-index]{center}, \myenv[no-index]{flushleft}, \myenv[no-index]{flushright},
%   \myenv[no-index]{verbatim}, \myenv[no-index]{tabbing}, \myenv[no-index]{trivlist}, \myenv[no-index]{list}
%   and all environments created with |\newtheorem|.
% \end{itemize}
%
% \subsubsection*{Example with \texttt{columns=2}}
%
% \begin{enumext*}[columns=2]
%   \item This text is in the first level.
%   \item This text is in the first level.
%   \item[X] This text is in the first level.
%   \item* This text is in the first level.
% \end{enumext*}
%
% \subsection{The command \cs[no-index]{item*}}\label{subsec:item-star}
%
% \begin{function}{\item*}
%   \begin{syntax}
%     |\item*|\myoarg{symbol}\myoarg{offset}
%   \end{syntax}
%    The \ics*{item*}, |\item*|\myoarg{symbol} and |\item*|\myoarg{symbol}\myoarg{offset}
%    work like the numbered |\item|, but places a \mymeta{symbol} to
%    the \emph{\enquote{left}} of the \mymeta{label} separated from it by
%    the \mymeta{offset} set by the \emph{second optional argument}.
%
%    The \emph{starred argument} `|*|' cannot be separated by spaces
%    `\verb*| |' from the command, i.e. |\item*| and the
%    \emph{first optional argument} does NOT support \emph{verbatim content}.
%    It can be configured with the keys  \mykey{item-sym*} and \mykey{item-pos*}
%    locally in the environment or globally using |\setenumext| command (\S\ref{sec:setenumext}).
% \end{function}
%
% \begin{important}*
% The behavior of |\item*| in the \myenv*[no-index]{enumext} and \myenv*[no-index]{enumext*} environments is NOT
% the same as in the \myenv*[no-index]{keyans} and \myenv*[no-index]{keyans*} environments.
% \end{important}
%
% \subsubsection{Keys for \cs[no-index]{item*}}
%
% \keyexampcmd{item-sym\textcolor{star}{*}}{symbol}{\textbackslash{}textborn}[env]
% Sets the \emph{symbol} to be displayed to the \emph{\enquote{left}} of the
% box containing the current \mymeta{label} set by \mykey{labelwidth} key for |\item*| in
% \myenv*[no-index]{enumext} and \myenv*[no-index]{enumext*}. The \emph{symbol} can be in
% \emph{text} or \emph{math} mode, for example \mykey{item-sym*}|={$\star$}|.
%
% \keyexampcmd{item-pos\textcolor{star}{*}}{rigid length}{by levels}[env]
% Sets the \emph{offset} between the box containing the
% current \mymeta{label} defined by \mykey{labelwidth} key and the
% \mymeta{symbol} set by \mykey{item-sym*} key. The default values are
% set by \mykey{labelsep} key at each level. If positive values are
% passed, it will \emph{offset to the left}; if negative values are
% passed it will \emph{offset to the right}.
%
% \subsection{The command \cs[no-index]{item} in \texttt{enumext*}}\label{subsec:item-in-starred}
%
% \begin{function}{\item}
%   \begin{syntax}
%     |\item|\myparg{columns}
%   \end{syntax}
%   The \ics*{item} command for the \myenv*[no-index]{enumext*} environment provides
%   a \emph{\enquote{first optional argument}} |\item|\myparg{columns} which \emph{\enquote{joins items}}
%   between columns.
% \end{function}
%
% Let's consider the following examples adapted directly from
% the \mypkg[no-index]{tasks} package:
%
% \iffalse
%<*example>
% \fi
\begin{examplecode}[frame=single]
\begin{enumext*}[widest=10,columns=4]
  \item The first
  \item* The second
  \item The third
  \item The fourth
  \item(3)* The fifth item is way too long for this and needs three columns
  \item The sixth
  \item The seventh
  \item(2)[X] The eighth item is way too long for this and needs two columns
    (\the\itemwidth)
  \item The ninth
  \item[Z] The tenth (\the\itemwidth)
\end{enumext*}
\end{examplecode}
% \iffalse
%</example>
% \fi
%
% \begin{enumext*}[widest=10,columns=4]
%  \item The first
%  \item* The second
%  \item The third
%  \item The fourth
%  \item(3)* The fifth item is way too long for this and needs three columns
%  \item The sixth
%  \item The seventh
%  \item(2)[X] The eighth item is way too long for this and needs two columns (\the\itemwidth)
%  \item The ninth
%  \item[Z] The tenth (\the\itemwidth)
% \end{enumext*}
%
% \section{The command \cs{setenumext}}\label{sec:setenumext}
%
% \begin{function}{\setenumext}
%   \begin{minipage}[t]{0.45\linewidth}
%     \cmdexamp{setenumext}{key \textnormal{\textcolor{gray}{=}} val}
%     \cmdexamp{setenumext}[enumext\textnormal{\textcolor{gray}{,}} level]{key \textnormal{\textcolor{gray}{=}} val}
%     \cmdexamp{setenumext}[enumext*]{key \textnormal{\textcolor{gray}{=}} val}
%     \cmdexamp{setenumext}[keyans]{key \textnormal{\textcolor{gray}{=}} val}
%   \end{minipage}\hfill
%   \begin{minipage}[t]{0.45\linewidth}
%     \cmdexamp{setenumext}[keyans*]{key \textnormal{\textcolor{gray}{=}} val}
%     \cmdexamp{setenumext}[print\textnormal{\textcolor{gray}{,}} level]{key \textnormal{\textcolor{gray}{=}} val}
%     \cmdexamp{setenumext}[print\textnormal{\textcolor{gray}{,}} \textnormal{\textcolor{red}{*}}]{key \textnormal{\textcolor{gray}{=}} val}
%     \cmdexamp{setenumext}[print*]{key \textnormal{\textcolor{gray}{=}} val}
%   \end{minipage}
%
% The command \ics*{setenumext} sets the \mymeta{keys} on a global basis
% for environments \myenv*[no-index]{enumext}, \myenv*[no-index]{enumext*}, \myenv*[no-index]{keyans},
% \myenv*[no-index]{keyans*} and the |\printkeyans| command.
% It can be used both in the preamble and in the body of the document
% as many times as desired.
% \end{function}
%
% The \mymeta{keys} set in the \emph{optional argument} of environments and
% commands have the \emph{highest precedence}, overriding both
% options passed by |\setenumext|. If the \emph{optional argument} is not
% passed, the \emph{first level} of the environment \myenv*[no-index]{enumext} will be
% taken by default.
%
% \begin{important}*
% For security reasons the keys \mykey{resume} \emph{\enquote{with value}},
% \mykey{resume*}, \mykey{reset}, \mykey{reset*}, \mykey{series} and
% \mykey{save-ans} they can NOT be set by this command and are ignored.
% The key \mykey{save-ans} that activates the \emph{\enquote{storage system}}
% must be passed directly in the \emph{optional argument} of the
% \emph{\enquote{first level}} of the environment in which they are executed.
% \end{important}
%
% \section{The command \cs[no-index]{setenumextmeta}}
%
% \begin{function}{\setenumextmeta}
%   \begin{syntax}
%     |\setenumextmeta|\myoarg{\textnormal{\textcolor{argnum}{\texttt{1}}}}\mymarg{key name}\mymarg{key-one \textnormal{\textcolor{gray}{=}} val\textnormal{\textcolor{gray}{,}} key-two \textnormal{\textcolor{gray}{=}} val\textnormal{\textcolor{gray}{,}} \dots}
%     |\setenumextmeta|\myoarg{\textnormal{\textcolor{argnum}{\texttt{2}}}}\mymarg{key name}\mymarg{key-one \textnormal{\textcolor{gray}{=}} val\textnormal{\textcolor{gray}{,}} key-two \textnormal{\textcolor{gray}{=}} val\textnormal{\textcolor{gray}{,}} \dots}
%     |\setenumextmeta|\myoarg{\textnormal{\textcolor{argnum}{\texttt{3}}}}\mymarg{key name}\mymarg{key-one \textnormal{\textcolor{gray}{=}} val\textnormal{\textcolor{gray}{,}} key-two \textnormal{\textcolor{gray}{=}} val\textnormal{\textcolor{gray}{,}} \dots}
%     |\setenumextmeta|\myoarg{\textnormal{\textcolor{argnum}{\texttt{4}}}}\mymarg{key name}\mymarg{key-one \textnormal{\textcolor{gray}{=}} val\textnormal{\textcolor{gray}{,}} key-two \textnormal{\textcolor{gray}{=}} val\textnormal{\textcolor{gray}{,}} \dots}
%     |\setenumextmeta|\myoarg{\textnormal{\textcolor{star}{\texttt{*}}}}\mymarg{key name}\mymarg{key-one \textnormal{\textcolor{gray}{=}} val\textnormal{\textcolor{gray}{,}} key-two \textnormal{\textcolor{gray}{=}} val\textnormal{\textcolor{gray}{,}} \dots}
%     |\setenumextmeta*|\mymarg{key name}\mymarg{key-one \textnormal{\textcolor{gray}{=}} val\textnormal{\textcolor{gray}{,}} key-two \textnormal{\textcolor{gray}{=}} val\textnormal{\textcolor{gray}{,}} \dots}
%   \end{syntax}
%   The command \ics*{setenumextmeta} adds a new \emph{\enquote{meta-key}}
%   for the environments \myenv*[no-index]{enumext} and \myenv*[no-index]{enumext*}, the
%   \mymarg{key name} must be different from those defined by the package.
%   The \emph{optional argument} of the form |[1]|, |[2]|, |[3]|, |[4]|
%   adds a new \emph{\enquote{meta-key}} for levels |1|, |2|, |3| and |4|
%   of the \myenv*[no-index]{enumext} environment, the form |[*]| adds a new
%   \emph{\enquote{meta-key}} for the \myenv*[no-index]{enumext*} environment. If it is
%   run \emph{without} the \emph{optional argument}, it will adds a new
%   \emph{\enquote{meta-key}} for the \emph{\enquote{first level}} of
%   the \myenv*[no-index]{enumext} environment.
% \end{function}
%
% The \emph{starred argument} `|*|' will create the new \emph{\enquote{meta-key}}
% for the environment \myenv*[no-index]{enumext*} and for \emph{\enquote{all levels}} of the
% environment \myenv*[no-index]{enumext}. For example: |\setenumextmeta*{midsep}{topsep=3pt, partopsep=0pt}|
% will create a new key |midsep| available for all levels of the
% \myenv*[no-index]{enumext} environment and the \myenv*[no-index]{enumext*} environment
% and it can be used like any other key so |\begin{enumext}[midsep]|
% and |\begin{enumext*}[midsep]| will be valid.
%
% \section{The \texttt{keyval} system}
%
% The \mymeta{key = val} system used by the \mypkg*[no-index]{enumext} package is
% implemented using \mypkg{l3keys} so it must be noted that \mymeta{keys}
% marked as \emph{\enquote{value forbidden}}, that is \mymeta{key} is
% different from \mymeta{key=}.
%
% All \mymeta{keys} described in this section are available for the
% \myenv*[no-index]{enumext}, \myenv*[no-index]{enumext*}, \myenv*[no-index]{keyans} and \myenv*[no-index]{keyans*}
% environments with the exception of the keys \mykey{series}, \mykey{resume},
% \mykey{resume*}, \mykey{reset} and \mykey{reset*}  which are only available
% for the \myenv*[no-index]{enumext} and \myenv*[no-index]{enumext*} environments; and the keys
% \mykey{mini-right}, \mykey{mini-right*} which are only available for
% the \myenv*[no-index]{enumext*} and \myenv*[no-index]{keyans*} environments.
%
% All \mymeta{keys} related to vertical or horizontal spacing accept a
% \emph{\enquote{skip}} or \emph{\enquote{dim}} expression if passed
% between braces, i.e. you do not need to use |\dimeval| or |\dimexpr|
% to perform calculations.
%
% \begin{important}*
%  It should be kept in mind that using any \mymeta{key} that sets
%  \emph{rubber lengths} or \emph{rigid lengths} for vertical or
%  horizontal space on a level will influence the vertical and horizontal
%  space for \emph{inner levels} and \myenv*[no-index]{keyans}, \myenv*[no-index]{keyans*} and
%  \myenv*[no-index]{keyanspic} environments.
% \end{important}
%
% \subsection{Keys for \texttt{label} and \texttt{ref}}\label{subsec:label-and-ref}
%
% \keyexampcmd*{mode-box}{value forbidden}{not used}[env]
% This is a \emph{\enquote{switch-key}} that does not receive an
% argument and is \enquote{only} available for the \emph{\enquote{first level}}
% of the \myenv*[no-index]{enumext} environment and the \myenv*[no-index]{enumext*} environment.
% When this is set the \mykey{label}, \mykey{font}, \mykey{wrap-label}
% and \mykey{wrap-label*} keys are executed within |\makebox| for the
% \myenv*[no-index]{enumext} and \myenv*[no-index]{keyans} environments.
%
% \begin{important}
%   This key is intended for compatibility with \emph{tagged} \texttt{PDF}
%   and is forcibly \emph{\enquote{enabled}} when |\DocumentMetadata|\cite{userguide} is
%   present. If you want to get the same document output whether |\DocumentMetadata|
%   is active or not, you must enable this key.
% \end{important}
%
% \begin{important}*
% In the \myenv*[no-index]{enumext*} and \myenv*[no-index]{keyans*} environments |\makelabel|
% is redefined using |\makebox| by default. If \myenv*[no-index]{enumext} or
% \myenv*[no-index]{keyans} is used in the \myenv*[no-index]{enumext*} environment the key
% must be activated manually.
% \end{important}
%
% \keyexampcmd{label}
%   {\textnormal
%       {\ilabel{alph*} \textcolor{lightgray}{\textbar} \ilabel{Alph*}
%         \textcolor{lightgray}{\textbar} \ilabel{arabic*}
%         \textcolor{lightgray}{\textbar} \ilabel{roman*}
%         \textcolor{lightgray}{\textbar} \ilabel{Roman*}}
%   }{by levels}[env]
%
% Sets the \mymeta{label} that will be printed at the \emph{current level}
% and default value for \mykey{labelwidth} key.
% The default value for the first level of the environments
% \myenv*[no-index]{enumext} and \myenv*[no-index]{enumext*} is \ilabel{arabic*.}, for
% second level are \ilabel*{alph*)}, for third level are \ilabel{roman*.}
% and for fourth level is \ilabel{Alph*.}. For \myenv*[no-index]{keyans}
% and \myenv*[no-index]{keyans*} environments the default value is \ilabel{Alph*)}.
%
% \begin{important}*
% This key is intended to give the basic structure with which the
% \mymeta{label} will be displayed, and the form in which it is
% used by standard \emph{\enquote{label and ref}} and the
% \emph{\enquote{internal label and ref}} system with the
% \mykey{save-ref} key. You cannot use commands with \mymeta{label} as an
% argument, for example |\emph|\mymarg{\textnormal{\ilabel[no-index]{alph*}}} will return an
% error. For full customization of how \mymeta{label} is displayed use
% the \mykey{font}, \mykey{wrap-label} and/or \mykey{wrap-label*} keys.
% \end{important}
%
% \keyexampcmd{labelsep}{rigid length}{\mydim{0.3333}em}[env]
% Sets the \emph{horizontal space} between the box containing the
% current \mymeta{label} defined by \mykey{label} key and the text of an
% item on the first line. Internally sets the value of |\labelsep| for
% the current level.
%
% \keyexampcmd{labelwidth}{rigid length}{by label}[env]
% Sets the \emph{width} of the label box containing the current \mymeta{label}
% set by the \mykey{label} key.
% Internally sets the value of |\labelwidth| for the current level.
% The default values are calculated by means of the \emph{width} of a box
% by setting a \emph{value} to the current counter set by \mykey{label}
% key using `\mydim{0}' for \ilabel[no-index]{arabic*}, `|M|' for \ilabel[no-index]{Alph*},
% `|m|' for \ilabel[no-index]{alph*}, `|VIII|' for  \ilabel[no-index]{Roman*} and `|viii|'
% for \ilabel[no-index]{roman*}.
%
% \keyexampcmd{widest}{integer \textnormal{\textcolor{lightgray}{\textbar}} string}{empty}[env]
% Sets the \mykey{labelwidth} key passing the \mymeta{integer} or converting
% the \mymeta{string} of the form |\Alph|, |\alph|, |\Roman| or |\roman|
% to a \emph{value} for the current counter defined by \mykey{label} key,
% then calculating the \emph{width} by means of a box. For example |widest={XXIII}|
% or |widest={23}| are equivalent. This key is useful when the default
% values of the \mykey{labelwidth} key are smaller than those actually used.
%
% \keyexampcmd{font}{font commands}{empty}[env]
% Sets the \emph{font style} for the current \mymeta{label} defined by \mykey{label} key.
% You cannot use commands with an argument, for example |font={\small\textbf{#1}}| will
% return an error, but |font={\bfseries\small}| it will work. For full customization
% of how \mymeta{label} is displayed use \mykey{wrap-label} and/or \mykey{wrap-label*} keys.
%
% \keyexampcmd{align}{left \textnormal{\textcolor{lightgray}{\textbar}} right \textnormal{\textcolor{lightgray}{\textbar}} center}{left}[env]
% Sets the \emph{aligned} of \mymeta{label} defined by \mykey{label}
% key on the current level in the label box of width set by \mykey{labelwidth}.
%
% \newsavebox{\hashbox}
% \begin{lrbox}{\hashbox}
% \lstinline[language=enumext-doc,basicstyle=\ttfamily]+{#1}+
% \end{lrbox}
%
% \keyexampcmd{wrap-label}{code{} \space\usebox{\hashbox} more code}{empty}[env]
%
% Wraps the \emph{current} \mymeta{label} defined by \mykey{label} key
% referenced by |{#1}| after executing the \mykey{align} and \mykey{font}
% keys. The \mymarg{code} must be passed between braces and it does not modify
% the value set by the \mykey{labelwidth} key and is applied \emph{only}
% on |\item| and |\item*|.
% When using it in the |\setenumext| command it is necessary to use the
% \emph{double} `|{##1}|'. For example |wrap-label={\fbox{#1}}|
% or you can create a command:
%
% \iffalse
%<*example>
% \fi
\begin{examplecode}[frame=none]
\NewDocumentCommand \mywrap { s m }
  {
    \IfBooleanTF{#1}
      {\textcolor{red}{\textbf{Q}}\textcolor{blue}{\textbf{.}}\textcolor{gray}{#2}}
      {\textcolor{blue}{\textbf{Q}}\textcolor{red}{\textbf{.}}\textcolor{gray}{#2}}
  }
\end{examplecode}
% \iffalse
%</example>
% \fi
%
% and then pass it through the key |wrap-label={\mywrap{#1}}| or |wrap-label={\mywrap*{#1}}|.
%
% \keyexampcmd{wrap-label\textcolor{star}{*}}{code{} \space\usebox{\hashbox} more code}{empty}[env]
%
% The same as the \mykey{wrap-label} key but also applies on |\item|\myoarg{custom}.
%
%
% \newsavebox{\refboxsample}
% \begin{lrbox}{\refboxsample}
% \small\textcolor{gray}{\ttfamily\{}\textnormal{\small\ilabel[no-index]{alph*}\textcolor{lightgray}{\textbar} \ilabel[no-index]{Alph*}\textcolor{lightgray}{\textbar} \ilabel[no-index]{arabic*}\textcolor{lightgray}{\textbar} \ilabel[no-index]{roman*}\textcolor{lightgray}{\textbar} \ilabel[no-index]{Roman*}}\textcolor{gray}{\ttfamily\}}
% \end{lrbox}
% \keyexampcmd{ref}{code{} \space\usebox{\refboxsample} more code}{empty}[env]
%
% Modifies the way \emph{cross references} are displayed. The
% \mykey{label} key sets the default form of the \emph{cross references},
% by using this key you can define a different format, for example:
% \mykey{ref}|=\emph|\mymarg{\textnormal{\ilabel[no-index]{alph*}}} is valid.
%
% Internally it renews the command associated with each counter when it
% is executed, i.e., in the environment \myenv*[no-index]{enumext} the command |\theenumXi|
% is modified when the key is executed at the first level, |\theenumXii|
% when it is executed at the second level and |\theenumXiii| together
% with |\theenumXiv| when it is executed at the
% third and fourth levels.
%
% \begin{important}*
% This must be kept in mind, since the values set by the \mykey{label} and
% \mykey{ref} keys are not cumulative by levels, so if you have used the \mykey{ref}
% key in the first level and then want to associate the counter with
% \mykey{label} or \mykey{ref} in the second level you must use the direct commands,
% i.e. |\arabic{enumXi}| to indicate the count of the first level
% instead of using |\theenumXi|.
% \end{important}
%
% \subsection{Keys for penalties}\label{subsec:penalties}
%
% Page breaks in the provided environments are controlled by the
% following three parameters, which work together to ensure they look
% good, avoiding unsightly page breaks that could distort the output.
%
% \keyexampcmd{beginpenalty}{integer}{\mydim{-51}}[env]
% Set the \emph{page breaking} penalty for breaking at the beginning of
% the \myenv*[no-index]{enumext}, \myenv*[no-index]{enumext*}, \myenv*[no-index]{keyans} and
% \myenv*[no-index]{keyans*} environments. Internally sets the value of \textcolor{red}{\texttt{\textbackslash}}|@beginparpenalty|.
%
% \keyexampcmd{midpenalty}{integer}{\mydim{-51}}[env]
% Set the \emph{page breaking} penalty for breaking between items of
% the \myenv*[no-index]{enumext}, \myenv*[no-index]{enumext*}, \myenv*[no-index]{keyans} and
% \myenv*[no-index]{keyans*} environments. Internally sets the value of \textcolor{red}{\texttt{\textbackslash}}|@itempenalty|.
%
% \keyexampcmd{endpenalty}{integer}{\mydim{-51}}[env]
% Set the \emph{page breaking} penalty for breaking at the end of
% the \myenv*[no-index]{enumext}, \myenv*[no-index]{enumext*}, \myenv*[no-index]{keyans} and
% \myenv*[no-index]{keyans*} environments. Internally sets the value of \textcolor{red}{\texttt{\textbackslash}}|@endparpenalty|.
%
% \begin{important}*
% The values passed to these \mymeta{keys} affect the nested
% environments in which they were set and cannot be reset.
% \hologo{LaTeX} default is \textcolor{red}{\texttt{-}}\textcolor{red}{\texttt{\textbackslash}}|@lowpenalty|,
% that is, \mydim{-51}. Because it is negative, it somewhat encourages a
% page break at each spot. Change it with, e.g., \textcolor{red}{\texttt{\textbackslash}}|@beginparpenalty=|\mydim{9999}; a value
% of \mydim{10000} prohibits a page break. Please, refer to your
% \hologo{LaTeX} or \hologo{TeX} manual about how penalties control page
% breaks.
% \end{important}
%
% \subsection{Keys for spaces}\label{subsec:keys-spaces}
%
% \keyexampcmd{show-length}{true \textnormal{\textcolor{lightgray}{\textbar}} false}{false}[env]
% Displays on the terminal the values for \emph{all list parameters} at the current level.
% For \emph{vertical spaces} show the values of |\topsep|, |\itemsep|, |\parsep| and |\partopsep|.
% For \emph{horizontal spaces} show the values of |\labelwidth|, |\labelsep|, |\itemindent|, |\listparindent| and |\leftmargin|.
%
% \subsubsection{Vertical spaces}
%
% \keyexampcmd{topsep}{rubber length \textnormal{\textcolor{lightgray}{\textbar}} rigid length}{by levels}[env]
%
% Set the \emph{vertical space} added to both the top and bottom of the list.
% Internally sets the value of |\topsep| for the current level.
% The default value for the first level of the environments
% \myenv*[no-index]{enumext} and \myenv*[no-index]{enumext*} are \mydim{8.0pt plus 2.0pt minus 4.0pt},
% for second level are \mydim{4.0pt plus 2.0pt minus 1.0pt}, for third and
% fourth level are \mydim{2.0pt plus 1.0pt minus 1.0pt}. For \myenv*[no-index]{keyans}
% and \myenv*[no-index]{keyans*} environments the default value is \mydim{4.0pt plus 2.0pt minus 1.0pt}.
%
% \keyexampcmd{parsep}{rubber length \textnormal{\textcolor{lightgray}{\textbar}} rigid length}{by levels}[env]
%
% Set the \emph{vertical space} between paragraphs within an item.
% Internally sets the value of |\parsep| for the current level.
% The default value for the first level of the environments \myenv*[no-index]{enumext}
% and \myenv*[no-index]{enumext*} are \mydim{4.0pt plus 2.0pt minus 1.0pt},
% for second level are \mydim{2.0pt plus 1.0pt minus 1.0pt}, for third and
% fourth level are \mydim{0pt}. For \myenv*[no-index]{keyans} and \myenv*[no-index]{keyans*}
% environments the default value is \mydim{2.0pt plus 1.0pt minus 1.0pt}.
%
% \begin{important}*
%  In the \myenv*[no-index]{enumext*} and \myenv*[no-index]{keyans*} environments this value
%  is passed to |\parskip| within the \myenv[no-index]{minipage} environment where
%  \emph{\enquote{item content}} is placed.
% \end{important}
%
% \keyexampcmd{partopsep}{rubber length \textnormal{\textcolor{lightgray}{\textbar}} rigid length}{by levels}[env]
%
% Set the \emph{vertical space} added, beyond \mykey{topsep}, to the \enquote{top} and \enquote{bottom}
% of the entire environment if the environment instance is preceded by a \emph{\enquote{blank
% line}} or |\par| command. Internally sets the value of |\partopsep| for the current level.
% The default values for first and second level in environment
% \myenv*[no-index]{enumext} are \mydim{2.0pt plus 1.0pt minus 1.0pt},
% for third and fourth level are \mydim{1.0pt minus 1.0pt}.
% For the \myenv*[no-index]{keyans} environment the default value is \mydim{2.0pt
% plus 1.0pt minus 1.0pt}, and for the \myenv*[no-index]{keyans*} and
% \myenv*[no-index]{enumext*} environments it is available but \emph{without}
% effect.
%
% \begin{important}*
% The value of this parameter also affects the \emph{inner levels} and the
% environments \myenv*[no-index]{keyans}, \myenv*[no-index]{keyanspic} and \myenv*[no-index]{keyans*}.
% Caution should be taken with \emph{\enquote{blank lines}} or |\par|
% command \emph{\enquote{before}} each environment or nested level when
% formatting the source code of document. \hologo{TeX} will enter \mymeta{vertical mode}
% and apply this value to the \enquote{top} and \enquote{bottom} the
% environment or nested level.
% \end{important}
%
% \keyexampcmd{itemsep}{rubber length \textnormal{\textcolor{lightgray}{\textbar}} rigid length}{by levels}[env]
% Set the \emph{vertical space} between items, beyond the \mykey{parsep}.
% Internally sets the value of |\itemsep| for the current level.
% The default value for the first level of the environments \myenv*[no-index]{enumext}
% and \myenv*[no-index]{enumext*} are \mydim{4.0pt plus 2.0pt minus 1.0pt},
% for the rest of the levels are \mydim{2.0pt plus 1.0pt minus 1.0pt}.
% For \myenv*[no-index]{keyans} and \myenv*[no-index]{keyans*} environments the default
% value is \mydim{4.0pt plus 2.0pt minus 1.0pt}.
%
% \begin{important}*
%   In the \myenv*[no-index]{enumext*} and \myenv*[no-index]{keyans*} environments this value
%   corresponds to the separation between rows.
% \end{important}
%
% \keyexampcmd*{noitemsep}{value forbidden}{not used}[env]
% This is a \emph{\enquote{meta-key}} that does not receive an argument. Set
% \mykey{itemsep} and \mykey{parsep} equal to \mydim{0pt} the entire level of environment.
%
% \keyexampcmd*{nosep}{value forbidden}{not used}[env]
% This is a \emph{\enquote{meta-key}} that does not receive an argument. Sets all keys
% for vertical spacing equal to \mydim{0pt} the entire level of environment.
%
% \keyexampcmd*{base-fix}{value forbidden}{not used}[env]
%
% This is a \emph{\enquote{switch-key}} that does not receive an
% argument available \emph{only} for the \emph{\enquote{first level}} of environment
% \myenv*[no-index]{enumext}. Fix the \emph{baseline} when an environment \myenv*[no-index]{enumext}
% is nested  in \myenv*[no-index]{enumext*} and there is no material between the
% |\item| and the start of the environment for example |\item \begin{enumext}|
% within the environment \myenv*[no-index]{enumext*}. Internally sets the keys \mykey{topsep},
% \mykey{above} and \mykey{above}|*| at \mydim{0pt}.
%
% \begin{important}*
%  This key is provided as a way to work around this minor issue, but you
%  should be aware that if for some reason you have the \mykey{itemindent}
%  key set in the \myenv*[no-index]{enumext*} environment it will be lost and you
%  will need to adjust it using the \mykey{list-offset} key in the
%  \myenv*[no-index]{enumext} environment.
% \end{important}
%
% \subsubsection*{Extra vertical spaces}
%
% \begin{important}*
%  The following \mymeta{keys} should be used with \emph{\enquote{caution}}, they
%  are intended to be used at the \enquote{top} and \enquote{bottom} of
%  the environment when the \mykey{columns} or \mykey{mini-env} keys do
%  not provide adequate \emph{vertical spaces}.
%  The values passed can be \emph{rubber} or \emph{rigid} lengths, the
%  way they are applied is the way you differ, using the \emph{star} `|*|'
%  \mymeta{keys} applies |\vspace*| so that \hologo{LaTeX} does
%  \emph{not discard} this space at page break.
% \end{important}
%
% \keyexampcmd{above}{rubber length \textnormal{\textcolor{lightgray}{\textbar}} rigid length}{not used}[env]
%
% Set the \emph{extra vertical space} added, beyond \mykey{topsep}, to
% the top of the entire level of environment. This key is intended to
% give a \emph{\enquote{fine adjustment}} of the vertical space
% \emph{\enquote{above}} the environment without hindering the value of
% the \mykey{topsep} key. The space is added with |\vspace| so is
% \emph{\enquote{discardable}}.
%
% \keyexampcmd{above\textcolor{star}{*}}{rubber length \textnormal{\textcolor{lightgray}{\textbar}} rigid length}{not used}[env]
%
% Set the \emph{extra vertical space} added, beyond \mykey{topsep}, to
% the top of the entire level of environment. This key is intended to
% give a \emph{\enquote{fine adjustment}} of the vertical space
% \emph{\enquote{above}} the environment without hindering the value of
% the \mykey{topsep} key. The space is added with |\vspace*| so is
% \emph{\enquote{not discardable}}.
%
% \keyexampcmd{below}{rubber length \textnormal{\textcolor{lightgray}{\textbar}} rigid length}{not used}[env]
%
% Set the \emph{extra vertical space} space added, beyond \mykey{topsep}, to
% the bottom of the entire level of environment. This key is intended to
% give a \emph{\enquote{fine adjustment}} of the vertical space on the
% \emph{\enquote{below}} the environment without hindering the value of
% the \mykey{topsep} key. The space is added with |\vspace| so is
% \emph{\enquote{discardable}}.
%
% \keyexampcmd{below\textcolor{star}{*}}{rubber length \textnormal{\textcolor{lightgray}{\textbar}} rigid length}{not used}[env]
%
% Set the \emph{extra vertical space} space added, beyond \mykey{topsep}, to
% the bottom of the entire level of environment. This key is intended to
% give a \emph{\enquote{fine adjustment}} of the vertical space on the
% \emph{\enquote{below}} the environment without hindering the value of
% the \mykey{topsep} key. The space is added with |\vspace*| so is
% \emph{\enquote{not discardable}}.
%
% \subsubsection{Horizontal spaces}
%
% \keyexampcmd{list-offset}{rigid length}{\mydim{0}pt}[env]
% Sets the \emph{horizontal translation} of the entire environment level from the left edge
% of the box defined by the \mykey{labelwidth} key. Internally sets the values of
% |\leftmargin| and |\itemindent| for the current level.
%
% \keyexampcmd{list-indent}{rigid length}{labelwidth + labelsep}[env]
% Sets the \emph{indentation} of the whole environment under the box defined by
% \mykey{labelwidth} and \mykey{labelsep} keys. Internally sets the value of
% |\leftmargin| and |\itemindent| for the current level.
% If |list-indent=|\mydim{0pt} is set in the environments \myenv*[no-index]{enumext}
% and \myenv*[no-index]{keyans} the \mymeta{label} will be part of the text, separated
% by the value of the \mykey{labelsep} key and the \emph{first word}, in
% simple terms it will look like a \emph{\enquote{common paragraph}}.
%
% \begin{important}*
%  The \myenv*[no-index]{enumext*} and \myenv*[no-index]{keyans*} environments are
%  implemented using |\makebox| and \myenv[no-index]{minipage} which causes
%  \emph{\enquote{list indent}} to always be equal to the value passed to
%  \mykey{labelwidth} plus \mykey{labelsep}. Passing a value to this key is
%  equivalent to setting the value for the \mykey{list-offset} key.
% \end{important}
%
% \keyexampcmd{itemindent}{rigid length}{\mydim{0}pt}[env]
%
% Sets the extra \emph{horizontal indentation}, beyond \mykey{labelsep},
% of the \emph{\enquote{first line}} of each |\item| that is not followed
% by a \emph{\enquote{blank line}} or the |\par| command . This value must
% be greater than or equal to \mydim{0pt} and is applied internally using
% |\hspace| \emph{without} modifying the value of |\itemindent|.
%
% \begin{important}*
%  This key is intended for the \myenv*[no-index]{enumext*} and \myenv*[no-index]{keyans*}
%  environments where, by their implementation, it is not possible to
%  adjust \mykey{labelwidth} and \mykey{list-indent} \emph{without}
%  modifying the output. If you use \myenv*[no-index]{enumext} or \myenv*[no-index]{keyans}
%  and want to get around the \emph{blank line} limitation or the |\par|
%  command followed by |\item| you can modify \mykey{labelwidth} and
%  \mykey{list-indent} and get the same effect.
% \end{important}
%
% \keyexampcmd{rightmargin}{rigid length}{\mydim{0}pt}[env]
%
% Set the \emph{horizontal space} between the right margin of the environment and the
% right margin of the enclosing environment, the value it takes must be
% greater than or equal to \mydim{0pt}. Internally sets the value of
% |\rightmargin| for the current level.
%
% \keyexampcmd{listparindent}{rigid length}{\mydim{0}pt}[env]
% Sets the \emph{horizontal space} indentation, beyond \mykey{list-indent},
% for second and subsequent paragraphs within a list item. Internally
% sets the value of |\listparindent| for the current level.
%
% \begin{important}*
%  In the \myenv*[no-index]{enumext*} and \myenv*[no-index]{keyans*} environments this value
%  is passed to |\parindent| within the \myenv[no-index]{minipage} environment where
%  \emph{\enquote{item content}} is placed.
% \end{important}
%
% \subsection{Keys for add \texttt{code}}\label{subsec:keys-code}
%
% The following \mymeta{keys} should be used with \emph{\enquote{caution}},
% they are intended to inject \mymarg{code} into different parts of the
% defined environments. It must be kept in mind that the defined
% environments are based on the \myenv[no-index]{list} base environment provided by
% \hologo{LaTeX} which is defined (simplified) as plain form |\list|\mymarg{arg one}\mymarg{arg two}.
% Using the \mykey{before*} key does not allow access to the \myenv[no-index]{list}
% parameters defined by \myoarg{key
% \textnormal{\textcolor{gray}{=}} val}.
%
% \keyexampcmd{before}{code}{not used}[env]
% Execute \mymarg{code} \emph{\enquote{before}} the environment
% starts. The \mymarg{code} must be passed between braces, is executed
% \emph{\enquote{after}} all calculations related to the \emph{list
% parameters} in the environment and the \mymeta{keys} sets by
% \myoarg{key \textnormal{\textcolor{gray}{=}} val} have been performed,
% that is, in the \emph{second argument} of the list: |\begin{list}|\mymarg{arg one}|{|\mymeta{arg two}\mymarg[cbc=module_name,ac=araracolor]{code}|}|.
%
% \keyexampcmd{before\textcolor{star}{*}}{code}{not used}[env]
% Execute \mymarg{code} \emph{\enquote{before}} the
% environment starts. The \mymarg{code} must be passed between braces, is executed
% \emph{\enquote{before}} performing all calculations related to the \emph{list
% parameters} and the \mymeta{keys} sets in \myoarg{key \textnormal{\textcolor{gray}{=}}
% val} of the environment that is, \emph{\enquote{before}} the arguments
% defining the list environment are executed:
% \mymarg[type=tt,cbc=module_name,ac=araracolor]{code}|\begin{list}|\mymarg{arg
% one}\mymarg{arg two}.
%
% \keyexampcmd{first}{code}{not used}[env]
% Executes \mymarg{code} when \emph{\enquote{starting}} the
% environment. The \mymarg{code} must be passed between
% braces, is executed right \emph{\enquote{after}} all \emph{list
% parameters} are done, after the second argument of list, just before the first occurrence of |\item|:
% |\begin{list}|\mymarg{arg one}\mymarg{arg two}\mymarg[cbc=module_name,ac=araracolor]{code}|\item|.
%
% \begin{important}*
%  Keep in mind that the \mymarg{code} set in this \mymeta{key} will affect
%  the entire \emph{\enquote{body}} of the environment and therefore the
%  inner levels of the list and the \myenv*[no-index]{keyans}, \myenv*[no-index]{keyans*} and
%  \myenv*[no-index]{keyanspic} environments. It is recommended to set this \mymeta{key}
%  per level. In the \myenv*[no-index]{enumext*} and \myenv*[no-index]{keyans*} environments
%  this \mymeta{key} is executed \emph{\enquote{after}} the \mykey{listparindent},
%  \mykey{parsep} and \mykey{itemindent} \mymeta{keys} within the \myenv[no-index]{minipage}
%  environment in which the \emph{\enquote{item content}} is placed.
%  Making modifications to |\setcounter| in these environments using this \mymeta{key}
%  will not have the expected effect.
% \end{important}
%
% \keyexampcmd{after}{code}{not used}[env]
% Execute \mymarg{code} \emph{\enquote{after}} finishing the
% environment. The \mymarg{code} must be passed between braces.
%
% \subsection{Keys for \texttt{start}, \texttt{series} and \texttt{resume}}\label{subsec:keys-start-series-resume}
%
% \keyexampcmd{start}{integer \textnormal{\textcolor{lightgray}{\textbar}} integer expression}{\mydim{1}}[env]
% Sets the \emph{start value} of the numbering on the \emph{\enquote{current level}}.
% The \mymarg{integer expression} must be passed between braces, internally is evaluated and pass
% to the \emph{\enquote{counter}} defined by \mykey{label} key on the current level, i.e. it is equivalent to enter
% |start={\dimeval{100*\value{chapter}}| or |start={100*\value{chapter}}|.
%
% \keyexampcmd{start\textcolor{star}{*}}{integer \textnormal{\textcolor{lightgray}{\textbar}} string}{not used}[env]
% Sets the \emph{start value} of the numbering on the \emph{\enquote{current level}}.
% Internally \mymeta{string} is converted and passed as value to the \emph{\enquote{counter}}
% defined by \mykey{label} key on the current level, i.e. it is equivalent to enter
% |start*=5|, |start*=E| or |start*=v|.
%
% \begin{important}
%   The \emph{start value} in |enumext| and |keyans| are set in the \emph{second argument} of the \myenv[no-index]{list}
%   environment and \emph{\enquote{before}} the execution of the keys \mykey{before}
%   and \mykey{first}, but in the |enumext*| and |keyans*| it is set much later: |\begin{list}|\mymarg{arg one}\mymarg{arg two}|\setcounter|.
% \end{important}
%
% \begin{important}*
%   The following \mymeta{keys} are available only for the \myenv*[no-index]{enumext}
%   and \myenv*[no-index]{enumext*} environments.
% \end{important}
%
% \keyexampcmd{series}{series name}{not used}[env]
%
% Stores the \emph{keys} of the \emph{optional argument} of the \emph{\enquote{current level}}
% of the environment in which it is executed in \mymarg{series name} which
% is used as an \emph{argument} in the \mykey{resume} key. The \mymeta{keys}
% stored in \mymarg{series name} are NOT cumulative and are overwritten
% if the same \mymarg{series name} is used again at the \emph{\enquote{same level}}
% at which the key was executed.
%
% \begin{important}*
% For security reasons the \mykey{series} key will never save in
% \mymarg{series name} the \mymeta{keys} \mykey{series}, \mykey{resume},
% \mykey{resume*}, \mykey{reset}, \mykey{reset*}, \mykey{save-ans}, \mykey{save-key},
% \mykey{start*} and \mykey{start}.
% \end{important}
%
% \keyexampcmd{resume}{series name}{not used}[env]
% Sets the \emph{start value} and \emph{options} for the \emph{\enquote{current level}}
% continuing the numbering and \emph{options} of the \emph{\enquote{same level}} as
% the environment in which the \mykey{series}|=|\mymarg{series name}
% key was executed, the \emph{start value} will continue numbering according to
% the last execution of \mykey{resume}|=|\mymarg{series name}. If passed
% \emph{\enquote{without value}} this will only set \emph{start value}
% continue the numbering of the \emph{\enquote{same level}}
% from the last environment and level in which \mykey{series}|=|\mymarg{series name}
% or \mykey{resume}|=|\mymarg{series name} is NOT present and if the
% \mykey{save-ans} key is active (on the left) it will continue the numbering from the
% \enquote{last} environment in which it was executed. The \emph{start value}
% can be overwritten using \mykey{start} or \mykey{start*} keys.
%
% \begin{important}*
% The \mykey{resume} key passed \emph{\enquote{without value}} must be
% exactly \emph{\enquote{without value}}, i.e. \mykey{resume}|=| cannot
% be used and if executed before \mykey{resume*} it will affect the
% \emph{start value}.
% \end{important}
%
% \keyexampcmd*{resume\textcolor{star}{*}}{value forbidden}{not used}[env]
% Sets the \emph{start value} and \emph{options} for the \emph{\enquote{current level}}
% continuing the numbering and options of the \emph{\enquote{same level}}
% as the last environment and level in which the \mykey{series}|=|\mymarg{series name} or
% \mykey{resume}|=|\mymarg{series name} keys are NOT present and if the
% \mykey{save-ans} key is active (on the left) it will continue the numbering and options
% from the \enquote{last} environment in which it was executed. The \emph{start value}
% can be overwritten using \mykey{start} or \mykey{start*} keys.
%
% \begin{important}*
%  When using the key \mykey{resume}|=|\mymarg{series name} or
%  \mykey{resume*} you will have hierarchy in the \mymeta{keys} that are
%  stored in \mymarg{series name} or in an internal version of
%  \mymarg{series name} in the case of \mykey{resume*}. If you want to
%  \emph{reset} the value of a \mymeta{key} that is already stored in
%  \mymarg{series name} or in an internal version of \mymarg{series name}
%  this must be placed to the \emph{right} of the key
%  \mykey{resume}|=|\mymarg{series name} or \mykey{resume*}.
% \end{important}
%
% \begin{important}*
%  When the \mykey{resume}|*| key is executed consecutively, it does not
%  rewrite the \mymeta{keys} stored in the internal version of
%  \mymarg{series name} and if the environment that precedes it does not
%  have the \emph{optional argument}, it will just continue with the
%  numbering.
% \end{important}
%
% \subsection{Keys for \texttt{reset}}\label{subsec:reset}
%
% \keyexampcmd*{reset}{value forbidden}{not used}[env]
% Resets the \emph{start value} of the \emph{\enquote{counters}} in the
% \myenv*[no-index]{enumext} and \myenv*[no-index]{enumext*} environments along with the
% \emph{\enquote{internal counters}} used by the \mykey{resume} \emph{without value}
% and \mykey{resume*} keys at the \emph{\enquote{level}} at which it is
% executed. The \emph{start value} can be overwritten using the \mykey{start}
% or \mykey{start*} keys.
%
% \keyexampcmd*{reset\textcolor{star}{*}}{value forbidden}{not used}[env]
% Resets the \emph{start value} of the \emph{\enquote{counters}} in the
% \myenv*[no-index]{enumext} and \myenv*[no-index]{enumext*} environments along with the
% \emph{\enquote{internal counters}} used by the \mykey{resume} \emph{without value}
% and \mykey{resume*} keys at the \emph{\enquote{level}} at which it is
% executed and in the \emph{\enquote{levels below}} it in the case of
% the \myenv*[no-index]{enumext} environment. The \emph{start value} can be overwritten
% using the \mykey{start} or \mykey{start*} keys.
%
% \begin{important}*
% These keys are intended to be used in cases where the |\resetenumext|
% command does not work, e.g. after an unnumbered chapter. It should
% preferably be set \emph{only} on the \emph{first level}, although it is
% available for all levels.
% \end{important}
%
% \subsubsection{The command \cs{resetenumext}}\label{subsubsec:resetenumext}
%
% \begin{function}{\resetenumext}
%   \begin{minipage}[t]{0.45\linewidth}
%     \cmdexamp{resetenumext}[\textnormal{\textcolor{argnum}{\texttt{1}}}]{some counter}
%     \cmdexamp{resetenumext}[\textnormal{\textcolor{argnum}{\texttt{2}}}]{some counter}
%     \cmdexamp{resetenumext}[\textnormal{\textcolor{argnum}{\texttt{3}}}]{some counter}
%   \end{minipage}\hfill
%   \begin{minipage}[t]{0.45\linewidth}
%     \cmdexamp{resetenumext}[\textnormal{\textcolor{argnum}{\texttt{4}}}]{some counter}
%     \cmdexamp{resetenumext}[\textnormal{\textcolor{star}{\texttt{*}}}]{some counter}
%     \cmdexamp{resetenumext\textcolor{star}{*}}{some counter}
%   \end{minipage}
%
%  The \ics*{resetenumext} command \emph{\enquote{resets}} the \emph{start value}
%  of the \emph{\enquote{counters}} for the \myenv*[no-index]{enumext} and \myenv*[no-index]{enumext*}
%  environments along with the \emph{\enquote{internal counters}} used by
%  the keys \mykey{resume} \emph{without value} and \mykey{resume*}
%  according to the value of \mymarg{some counter}. For example |\resetenumext{chapter}|
%  will \emph{\enquote{reset}} the numbering of \emph{\enquote{all levels}}
%  of the \myenv*[no-index]{enumext} environment for each execution of a \emph{\enquote{numbered}} chapter.
% \end{function}
%
% The \emph{optional argument} of the form |[1]|, |[2]|, |[3]|, |[4]|
% \emph{\enquote{reset}} the values for levels |1|, |2|, |3| and |4|
% of the \myenv*[no-index]{enumext} environment, the form |[*]|
% \emph{\enquote{reset}} the values for the \myenv*[no-index]{enumext*}
% environment. If is run \emph{without} the \emph{optional argument}, it
% will \emph{\enquote{reset}} the values for \emph{\enquote{all
% levels}} of the \myenv*[no-index]{enumext} environment.
%
% The \emph{starred argument} `|*|' will \emph{\enquote{reset}} the
% values for \emph{\enquote{all levels}} of the \myenv*[no-index]{enumext} and
% \myenv*[no-index]{enumext*} environments.
%
% \subsection{Keys for \env{multicols}}\label{subsec:keys-multicols}
%
% \keyexampcmd{columns}{integer}{\mydim{1}}[env]
%
% Set the \emph{number of columns} to be used by the \myenv[no-index]{multicols}
% environment within the environments \myenv*[no-index]{enumext} and \myenv*[no-index]{keyans}. The value must be a positive
% integer less than or equal to \mydim{10}.
% In the \myenv*[no-index]{enumext*} and \myenv*[no-index]{keyans*} environments they
% correspond to the default number of columns (without joining) and
% internally adjust the value of \idimen[no-index]{itemwidth}.
%
% \keyexampcmd{columns-sep}{rigid length}{by level}[env]
%
% Set the \emph{space between} columns used by the \myenv[no-index]{multicols}
% environment within the environments \myenv*[no-index]{enumext} and \myenv*[no-index]{keyans}. Internally sets the value of
% |\columnsep|, by default its value is equal to the sum of the
% values set in the keys \mykey{labelwidth} and \mykey{labelsep} of the
% current level.
% In the \myenv*[no-index]{enumext*} and \myenv*[no-index]{keyans*} environments they
% correspond to the \emph{space between} columns (without joining) and
% internally adjust the value of \idimen[no-index]{itemwidth}.
%
% \subsection{Keys for \env{minipage}}\label{subsec:keys-minipage}
%
% \keyexampcmd{mini-env}{rigid length}{not used}[env]
%
% Sets the \emph{width} of the \myenv[no-index]{minipage} environment on
% the \emph{\enquote{right side}}. This value added to the value set by
% the \mykey{mini-sep} key to determines the \emph{width} of the
% \myenv[no-index]{minipage} environment on the \emph{\enquote{left side}}, taking
% |\linewidth| as the maximum reference value.
%
% \keyexampcmd{mini-sep}{rigid length}{\mydim{0.3333}em}[env]
%
% Sets the \emph{space between} the \myenv[no-index]{minipage} environment on the
% \emph{\enquote{left side}} and the \myenv[no-index]{minipage} environment on
% the \emph{\enquote{right side}}. This separation is applied together
% with |\hfill|.
%
% \subsubsection{The command \cs[no-index]{miniright}}
%
% \begin{function}{\miniright}
%   \begin{syntax}
%     |\begin{enumext}[mini-env=|\mymarg{rigid length}|]| \mymeta[cf=gray]{items before} |\item \miniright| \mymeta[cf=gray]{content} |\end{enumext}|\par
%     |\begin{enumext}[mini-env=|\mymarg{rigid length}|]| \mymeta[cf=gray]{items before} |\item \miniright*|\mymeta[cf=gray]{content} |\end{enumext}|
%   \end{syntax}
% \end{function}
%
% The \ics*{miniright} command close the \myenv[no-index]{minipage} environment on
% the \emph{\enquote{left side}} and opens the \myenv[no-index]{minipage} environment
% on the \emph{\enquote{right side}} by starting it with the |\centering|
% command. It must be placed \emph{\enquote{after}} the last |\item|
% of the current environment and \emph{\enquote{before}} starting the
% material to be placed on the \emph{\enquote{right side}}.
%
% The \emph{starred argument} `|*|' inhibits the use of |\centering|
% command i.e. the usual \hologo{LaTeX} justification is maintained in
% the \myenv[no-index]{minipage} on the \emph{\enquote{right side}}.
%
% \subsubsection{The key \texttt{mini-right}}
%
% In the \emph{horizontal list environments} \myenv*[no-index]{enumext*} and \myenv*[no-index]{keyans*}
% it is not possible to use the |\miniright| command and the \mykey{mini-right}
% key must be used instead.
%
% \keyexampcmd{mini-right}{content}{not used}[env]
%
% Set the \emph{content} for the drawing or tabular to be placed in the
% \myenv[no-index]{minipage} environment on the \emph{\enquote{right side}} by
% starting it with |\centering|. The \mymarg{content} must be passed between braces.
%
% \keyexampcmd{mini-right\textcolor{star}{*}}{content}{not used}[env]
%
% Same as above, but \emph{without} starting with |\centering|.
%
% \section{The storage system}\label{sec:storage-system}
%
% The entire mechanism for \emph{\enquote{storing content}} it is
% activated according to \mykey{save-ans} key on the
% \emph{\enquote{first level}} of \myenv*[no-index]{enumext} or \myenv*[no-index]{enumext*}
% environments and it is ignored if they are established when they are
% nested inside each other. Only when this \mymeta{key} is \emph{\enquote{active}}
% the |\anskey| command and the environments \myenv*[no-index]{anskey*},
% \myenv*[no-index]{keyans}, \myenv*[no-index]{keyans*} and \myenv*[no-index]{keyanspic} are available.
%
% \iffalse
%<*example>
% \fi
\vspace*{-1\baselineskip}
\begin{minipage}[t]{0.45\linewidth}
\centering
\begin{examplecode}[frame=none,mathescape=true]
\begin{enumext}[save-ans=`\mymarg{store name}`]
  \item Text \anskey{answer}
  \item Text
    \begin{keyans}
          $\cdots$
    \end{keyans}
\end{enumext}
\end{examplecode}
\end{minipage}\hfill
\begin{minipage}[t]{0.45\linewidth}
\vspace{0pt}\centering
\begin{examplecode}[frame=none,mathescape=true]
\begin{enumext}[save-ans=`\mymarg{store name}`]
  \item Text \anskey{answer}
  \item Text
    \begin{keyanspic}
          $\cdots$
    \end{keyanspic}
\end{enumext}
\end{examplecode}
\end{minipage}
% \iffalse
%</example>
% \fi
%
% By executing the key \mykey{save-ans}|=|\mymarg{store name} the entire
% \emph{\enquote{structure}} of the environment (excluding the \emph{first level})
% including the \emph{optional argument} passed to the inner levels or
% the environment nested in it, along with the \mymeta{content} passed
% to |\anskey| or \myenv*[no-index]{anskey*}, the current \mymeta{labels} for
% |\item*| and |\anspic*| in the environments \myenv*[no-index]{keyans},
% \myenv*[no-index]{keyans*} and \myenv*[no-index]{keyanspic} will be \emph{\enquote{stored}}
% in a \emph{sequence} \mymarg{store name} and at the same time will be
% \emph{\enquote{stored}} (without the \emph{\enquote{structure}} or
% \emph{optional argument}) in a \emph{prop list} \mymarg{store name}.
%
% For security reasons the \emph{optional argument} of the inner levels
% or the nested environment are \emph{filtered} by excluding all \mymeta{keys}
% related to the \emph{\enquote{storage system}} (\S\ref{subsec:storage-keys})
% along with the keys \mykey{mini-env}, \mykey{mini-sep}, \mykey{mini-right},
% \mykey{mini-right*}, \mykey{series}, \mykey{resume} and \mykey{resume*}
% when storing in \emph{sequence} \mymarg{store name} set by \mykey{save-ans} key.
%
% \subsection{Keys for storage system}\label{subsec:storage-keys}
%
% The only \mymeta{keys} available for all levels of the \myenv*[no-index]{enumext}
% environment and the \myenv*[no-index]{enumext*} environment are \mykey{no-store}
% and \mykey{save-key}, the rest of the \mymeta{keys} described in this
% section must be passed directly in the \emph{optional argument} of the
% \emph{\enquote{first level}} of the environment in which the key
% \mykey{save-ans} is executed. The key \mykey{save-ans} should NOT be passed
% with the command |\setenumext|.
%
% \keyexampcmd{save-ans}{store name}{not set}[env]
% Sets the \emph{name} of the \emph{sequence} and \emph{prop list} in
% which the \mymarg{contents} will be \emph{\enquote{stored}} by |\anskey|
% and \myenv*[no-index]{anskey*} in \myenv*[no-index]{enumext} and \myenv*[no-index]{enumext*}
% environments and the current \mymeta{labels} for |\item*| and
% |\anspic*| in the environments \myenv*[no-index]{keyans}, \myenv*[no-index]{keyans*}
% and \myenv*[no-index]{keyanspic}. If the \emph{sequence} or \emph{prop list}
% \mymarg{store name} does not exist, it will be created globally and
% will not be \emph{overwritten} if the key is used again.
%
% \keyexampcmd{save-key}{key list}{not set}[env]
% This key \emph{overrides} the default \emph{\enquote{stored keys}} of
% the \emph{optional argument} of the inner levels or nested environment
% that will be passed to the \emph{sequence} \mymarg{store name} set by
% \mykey{save-ans} key. The \mymeta{key list} passed to this key ignores
% any \mymeta{keys} in the \emph{\enquote{stored structure}} and must be
% passed between braces. For example, if executed at a second level:
%
% \iffalse
%<*example>
% \fi
\begin{examplecode}[frame=none,mathescape=true]
\begin{enumext}[save-ans=`\mymarg{store name}`]
  \item Text \anskey{answer}
  \item Text
    \begin{enumext}[nosep, columns=2, save-key={columns=3}]
          $\cdots$
    \end{enumext}
\end{enumext}
\end{examplecode}
% \iffalse
%</example>
% \fi
%
% The \emph{\enquote{stored keys}} by default in the \emph{sequence}
% \mymarg{store name} would be |nosep| and |columns=2|, but using the key
% \mykey{save-key}|={columns=3}| will overwrite and the
% \emph{\enquote{stored key}} in the \emph{sequence} \mymarg{store name}
% are only |columns=3| ignoring all the others.
%
% \keyexampcmd{save-sep}{text symbol}{\{, \}}[env]
% Sets the \emph{text symbol} that will separate the current \mymeta{label}
% to the \emph{optional argument} passed to the |\item*| and |\anspic*|
% in the environments \myenv*[no-index]{keyans}, \myenv*[no-index]{keyans*} and \myenv*[no-index]{keyanspic}
% and storing them in the \emph{sequence} and \emph{prop list} \mymarg{store name}
% set by \mykey{save-ans} key.
% The \mymarg{text symbol} must always be passed between
% braces, whitespace `\verb*+ +' is preserved within the braces and only
% affects the \emph{\enquote{stored content}} and not what is displayed
% when using the \mykey{show-ans} or \mykey{show-pos} keys.
%
% \keyexampcmd*{no-store}{value forbidden}{not used}[env]
%
% This is a \emph{\enquote{switch-key}} that does not receive an argument
% and disables the \emph{\enquote{storing content}} in the \emph{sequence}
% and \emph{prop list} \mymarg{store name} set by \mykey{save-ans} key at
% the entire level or a nested environment in which it runs. This key
% is intended for use in internal levels or nested \myenv*[no-index]{enumext} or
% \myenv*[no-index]{enumext*} environments in which you want to use \myenv*[no-index]{enumext}
% or \myenv*[no-index]{enumext*} but \emph{\enquote{without}} using the |\anskey|
% command or use \myenv*[no-index]{anskey*} environment and \emph{\enquote{without}}
% interfering with the \mykey{check-ans} key.
%
% \subsubsection{Keys for \texttt{label} and \texttt{ref}}\label{subsubsec:label-ref}
%
% \keyexampcmd{save-ref}{true \textnormal{\textcolor{lightgray}{\textbar}} false}{false}[env]
% Activates the \emph{\enquote{internal label and ref}} mechanism for
% referencing \emph{\enquote{stored content}} in \emph{prop list}
% \mymarg{store name} set by \mykey{save-ans} key. To reference the
% location of the \emph{\enquote{stored content}} within the environment
% you must use |\ref|\mymarg{store name \textnormal{\textcolor{red}{:}} position},
% where \mymeta{position} corresponds to the position occupied by the
% \emph{\enquote{stored content}} in the \emph{prop list} \mymarg{store name}
% returned by the \mykey{show-pos} key. For example |\ref{test:4}| will
% return |3.(b)| which corresponds to the location of the \emph{\enquote{stored content}}
% at position |4| in \emph{prop list} |test| within the environment in
% which the key \mykey{save-ans}|=test| was set.
%
% \keyexampcmd{mark-ref}{symbol}{\textbackslash{}textreferencemark}[env]
% Sets the \emph{symbol} that will be displayed by the |\printkeyans|
% command only if the \mypkg[no-index]{hyperref} package is detected and the
% \mykey{save-ref} key are active.
% This \emph{\enquote{symbol}} is used as a \emph{\enquote{link}} between the environment
% in which the \mykey{save-ans} key was used and the place where the
% command is executed.
%
% \subsubsection{Keys for \texttt{wrap} and \texttt{marks}}\label{subsubsec:wrap-display}
%
% The \mypkg*[no-index]{enumext} package provides a set of \mymeta{keys} to set and
% manipulate \emph{\enquote{symbol marks}} associated with
% \emph{\enquote{answers}} and how they are displayed and stored in the
% \emph{sequence} and \emph{prop list}.
%
% The \mymeta{keys} available for the |\anskey| command and the
% \myenv*[no-index]{anskey*} environment can be passed \emph{\enquote{only}} in the
% \emph{optional argument} in the \emph{\enquote{first level}} of the
% \myenv*[no-index]{enumext} or \myenv*[no-index]{enumext*} environment.
%
% The \mymeta{keys} available for the \myenv*[no-index]{keyans} and
% \myenv*[no-index]{keyans*} environments can be passed locally in the
% \emph{optional argument}, at the \emph{\enquote{first level}} of the
% \myenv*[no-index]{enumext} or \myenv*[no-index]{enumext*} environment or via the
% |\setenumext| command with one minor difference, when \mymeta{keys} are passed
% through the \emph{\enquote{first level}} of the \myenv*[no-index]{enumext} or
% \myenv*[no-index]{enumext*} environment they are set in \emph{\enquote{both}}
% environments, but when they are passed using the |\setenumext| command
% they are set \emph{\enquote{individually}} in each environment.
%
% \keyexampcmd{show-ans}{true \textnormal{\textcolor{lightgray}{\textbar}} false}{false}[env]
%
% Display the \emph{symbol} set by the \mykey{mark-ans} key to the left of
% the \emph{mandatory argument} \mymeta{content} passed to the |\anskey|
% command and \mymeta{body} for the \myenv*[no-index]{anskey*} environment using
% the \mykey{wrap-ans} key if set.
% For |\item*| and |\anspic*| the \myenv*[no-index]{keyans},
% \myenv*[no-index]{keyans*} and \myenv*[no-index]{keyanspic} environments it will display
% the \emph{symbol} set by the \mykey{mark-ans*} key to the left of the
% current \mymeta{label} and \emph{optional argument}. If the
% \emph{optional argument} is present in |\item*| or |\anspic*|
% it will be shown using \mykey{wrap-opt} key.
%
% \subsubsection*{Keys for \texttt{\textbackslash{}anskey} and \texttt{anskey*}}
%
% \keyexampcmd{mark-ans}{symbol}{\textbackslash{}textasteriskcentered}[env]
% Sets the \emph{symbol} to be displayed in the left margin for
% |\anskey| command and \myenv*[no-index]{anskey*} environment when using the
% key \mykey{show-ans}. The \emph{\enquote{symbol}} is placed in a box
% of width equal to the value of \mykey{labelwidth} at the current level,
% separated by the value of the key \mykey{mark-sep} and aligned by
% the value of the key \mykey{mark-pos}. This key is not affected by the
% keys \mykey{font} or \mykey{wrap-label} so if you want to apply
% \emph{styling} you have to do it directly, for example: |mark-ans={\textcolor{red}{\textbf{\textasteriskcentered}}|.
%
% \keyexampcmd{mark-pos}{left \textnormal{\textcolor{lightgray}{\textbar}} right \textnormal{\textcolor{lightgray}{\textbar}} center}{left}[env]
% Sets the \emph{aligned} of the \emph{\enquote{symbol}} defined by
% \mykey{mark-ans} key for |\anskey| command and \myenv*[no-index]{anskey*}
% environment. The \emph{\enquote{symbol}} is aligned in a box with the
% same dimensions of the label box defined by \mykey{labelwidth} key
% on the current level and separated by the value of the \mykey{mark-sep}
% key.
%
% \keyexampcmd{mark-sep}{rigid length}{labelsep}[env]
% Sets the \emph{horizontal space} between the box containing the
% \emph{\enquote{symbol}} defined by \mykey{mark-ans} key and the
% \emph{mandatory argument} \mymeta{content} passed to the |\anskey|
% command and the \mymeta{body} in \myenv*[no-index]{anskey*} environment.
%
% \keyexampcmd{wrap-ans}{code{} \space\usebox{\hashbox} more code}{\textbackslash{}fbox+\textbackslash{}parbox\{\#1\}}[env]
% Wraps the \emph{mandatory argument} \mymeta{content} passed to the |\anskey|
% and the \mymeta{body} in \myenv*[no-index]{anskey*} environment referenced by |{#1}|
% when using the \mykey{show-ans} or \mykey{show-pos} keys.
% The \mymarg{code} must be passed between  braces and only affects how the
% \mymeta{content} or \mymeta{body} is displayed and NOT the \emph{\enquote{stored
% content}} in the \emph{sequence} and \emph{prop list} \mymarg{store
% name} set by \mykey{save-ans} key. If this key is passed using
% |\setenumext| it is necessary to use double `|{##1}|'.
%
% \subsubsection*{Keys for \texttt{keyans}, \texttt{keyans*} and \texttt{keyanspic}}
%
% \keyexampcmd{mark-ans\textcolor{star}{*}}{symbol}{\textbackslash{}textasteriskcentered}[env]
% Sets the \emph{symbol} to be displayed in the left margin for
% |\item*| and |\anspic*| for the \myenv*[no-index]{keyans}, \myenv*[no-index]{keyans*}
% and \myenv*[no-index]{keyanspic} environments when using the key \mykey{show-ans}.
% The \emph{\enquote{symbol}} is placed in a box of width equal to the value of
% \mykey{labelwidth} of the environment in which it is executed, separated by the value of the
% key \mykey{mark-sep*} and aligned by the value of the key
% \mykey{mark-pos*}. This key is not affected by the keys \mykey{font} or
% \mykey{wrap-label} so if you want to apply \emph{styling} you have to do it
% directly, for example: |mark-ans*={\textcolor{red}{\textbf{\textasteriskcentered}}|.
%
% \keyexampcmd{mark-pos\textcolor{star}{*}}{left \textnormal{\textcolor{lightgray}{\textbar}} right \textnormal{\textcolor{lightgray}{\textbar}} center}{left}[env]
% Sets the \emph{aligned} of the \emph{\enquote{symbol}} defined by
% \mykey{mark-ans*} key for the \myenv*[no-index]{keyans}, \myenv*[no-index]{keyans*}
% and \myenv*[no-index]{keyanspic} environments. The \emph{\enquote{symbol}} is
% aligned in a box with the same dimensions of the label box defined by
% \mykey{labelwidth} key of the environment in which it is executed and
% separated by the value of the \mykey{mark-sep*} key.
%
% \keyexampcmd{mark-sep\textcolor{star}{*}}{rigid length}{labelsep}[env]
% Sets the \emph{horizontal space} between the box containing the
% \emph{\enquote{symbol}} defined by \mykey{mark-ans*} key and the current \mymeta{label} for
% |\item*| and |\anspic*| in the \myenv*[no-index]{keyans}, \myenv*[no-index]{keyans*}
% and \myenv*[no-index]{keyanspic} environments.
%
% \keyexampcmd{wrap-ans\textcolor{star}{*}}{code{} \space\usebox{\hashbox} more code}{not used}[env]
% Wraps the \emph{current} \mymeta{label} when using the \mykey{show-ans}
% key for |\item*| and |\anspic*| referenced by |{#1}| in the
% \myenv*[no-index]{keyans}, \myenv*[no-index]{keyans*} and \myenv*[no-index]{keyanspic} environments
% after executing the \mykey{align} and \mykey{font} keys.
% The \mymarg{code} must be passed between braces and \emph{only} affects
% how the \mymeta{label} is displayed and NOT the \emph{\enquote{stored label}}
% in the \emph{sequence} and \emph{prop list} \mymarg{store name} set by
% \mykey{save-ans} key. This key overwrites the key \mykey{wrap-label}
% and if is passed using |\setenumext| it is necessary to use double `|{##1}|'.
% For example, if you want the \mymeta{label} to be displayed in red when
% using \mykey{show-ans} you just set |wrap-ans*={\textcolor{red}{#1}}|.
%
% \keyexampcmd{wrap-opt}{code{} \space\usebox{\hashbox} more code}{[\{\#1\}]}[env]
% Wraps the \emph{optional argument} passed to the
% |\item*| and |\anspic*| referenced by |{#1}|
% in the \myenv*[no-index]{keyans}, \myenv*[no-index]{keyans*} and \myenv*[no-index]{keyanspic}
% environments when using the \mykey{show-ans} or \mykey{show-pos} keys.
% The \mymarg{code} must be passed between braces
% and only affects the current \emph{optional argument} and NOT the
% \emph{\enquote{stored content}} in the \emph{sequence} and \emph{prop list}
% \mymarg{store name} set by \mykey{save-ans} key.
% If this key is passed using |\setenumext| it is necessary to use double `|{##1}|'.
%
% \subsubsection{Keys for debug and checking}
%
% \keyexampcmd{show-pos}{true \textnormal{\textcolor{lightgray}{\textbar}} false}{false}[env]
%
% Displays the \emph{position} occupied by the \emph{\enquote{stored content}} by
% |\anskey|, \myenv*[no-index]{anskey*}, |\item*| and |\anspic*| in the
% \emph{prop list} \mymarg{store name} set by \mykey{save-ans} key.
% This position is used by the |\getkeyans| command and by the |\ref|
% command if the \mykey{save-ref} key is active.
%
% \keyexampcmd{check-ans}{true \textnormal{\textcolor{lightgray}{\textbar}} false}{false}[env]
%
% Enables the \emph{checking answer} mechanism displaying an
% appropriate message on the terminal. This key works under the logic
% that each |\item| or |\item*| that does not open an
% inner level or nested environment contains \emph{\enquote{only one
% answer}} or \emph{\enquote{only one execution}} of the |\anskey|
% or \myenv*[no-index]{anskey*}. It is intended to be used in conjunction with the
% \mykey{no-store} key.
%
% \subsection{The command \cs[no-index]{anskey}}\label{subsec:anskey-cmd}
%
% \begin{function}{\anskey}
%   \begin{syntax}
%    \cmdexamp{anskey}[keys]{content}
%   \end{syntax}
%  The command \ics*{anskey} takes a mandatory non empty argument \mymarg{content}
%  and \emph{\enquote{stores}} it in the \emph{sequence} and \emph{prop
%  list} \mymarg{store name} set by \mykey{save-ans} key.
%  By design the command cannot be nested or passed \emph{verbatim material} in the
%  argument and it is assumed that each \emph{numbered} |\item| or |\item*| within the
%  environment in which it is active it has a \emph{\enquote{single
%  execution}} of |\anskey| unless |\item| or |\item*| open a nested level
%  or use the \mykey{no-store} key.
%
%  If \mykey{save-ref} key are active and the
%  \mypkg[no-index]{hyperref}\cite{hyperref} package is detected, |\hyperlink| and
%  |\hypertarget| will be used, otherwise the usual \emph{\enquote{label
%  and ref}} system provided by \hologo{LaTeX} will be used.
% \end{function}
%
% The |\anskey| command is available for all levels of the
% \myenv*[no-index]{enumext} environment and the \myenv*[no-index]{enumext*} environment, but
% is disabled for the \myenv*[no-index]{keyans}, \myenv*[no-index]{keyans*} and
% \myenv*[no-index]{keyanspic} environments.
%
% \subsubsection{Keys for \cs[no-index]{anskey}}
%
% By default the \emph{mandatory argument} \mymeta{content} passed to |\anskey|
% when \emph{\enquote{storing}} in the \emph{sequence} \mymarg{store name}
% has the form |\item| \mymeta{content}, the following \mymeta{keys} allow
% modifying the way in which it is \emph{\enquote{stored}} in the \emph{sequence}.
%
% \keyexampcmd*{break-col}{value forbidden}{not used}[anskey,anskey*]
%
% Stores \mymarg{content} in the \emph{sequence} \mymarg{store name} of
% the form |\columnbreak| |\item| \mymeta{content}.
%
% \keyexampcmd{item-join}{columns}{not set}[anskey,anskey*]
%
% Set the \emph{number of columns} to be used for |\item|\myparg{columns} and stores \mymarg{content} in the \emph{sequence}
% \mymarg{store name} of the form |\item|\myparg{columns} \mymeta{content}.
%
% \keyexampcmd*{item-star}{value forbidden}{not used}[anskey,anskey*]
%
% Stores \mymarg{content} in the \emph{sequence} \mymarg{store name} of
% the form |\item*| \mymeta{content}.
%
% \keyexampcmd{item-sym\textcolor{star}{*}}{symbol}{not set}[anskey,anskey*]
%
% Sets the \emph{symbol} for |\item*| when using the key
% \mykey[anskey,anskey*]{item-star} and stores \mymarg{content} in the \emph{sequence}
% \mymarg{store name} of the form |\item*|\myoarg{symbol} \mymeta{content}.
% The \emph{symbol} can be in text or math mode, for example
% \mykey[anskey,anskey*]{item-sym*}|={$\ast$}| stores |\item*[$\ast$]| \mymeta{content}.
%
% \keyexampcmd{item-pos\textcolor{star}{*}}{rigid length}{not set}[anskey,anskey*]
% Sets the \emph{offset} for |\item*| when using the keys
% \mykey[anskey,anskey*]{item-star} and \mykey[anskey,anskey*]{item-sym*} and stores \mymarg{content} in the \emph{sequence}
% \mymarg{store name} of the form |\item*|\myoarg{symbol}\myoarg{offset} \mymeta{content}.
%
% \subsubsection*{Example}
%
% \iffalse
%<*example>
% \fi
\begin{examplecode}[frame=single]
\begin{enumext}[save-ans=test,show-ans=true]
  \item* Text containing our instructions or questions. `\ics*{anskey}\mymarg{first answer}`
  \item Text containing our instructions or questions.
    \begin{enumext}
      \item Question.`\ics*{anskey}\mymarg{second answer}`
    \end{enumext}
  \item Text containing our instructions or questions. `\ics*{anskey}\mymarg{third answer}`
  \item Text containing our instructions or questions. `\ics*{anskey}\mymarg{fourth answer}`
\end{enumext}
\end{examplecode}
% \iffalse
%</example>
% \fi
%
% \begin{enumext}[save-ans=test,mark-ans={\oldtextasteriskcentered},show-ans=true,columns=2,nosep]
%   \item* Text containing our instructions or questions.\anskey{first answer}
%   \item Text containing our instructions or questions.
%     \begin{enumext}
%       \item Question.\anskey{second answer}
%     \end{enumext}
%   \item Text containing our instructions or questions. \anskey{third answer}
%   \item Text containing our instructions or questions. \anskey{fourth answer}
% \end{enumext}
%
% \subsection{The environment \env{anskey*}}\label{subsec:anskey-env}
%
% \begin{function}{anskey*}
%   \begin{syntax}
%     |\begin{anskey*}|\myoarg{key \textnormal{\textcolor{gray}{=}} val} \mymeta{body content} |\end{anskey*}|
%   \end{syntax}
%   The environment \myenv*{anskey*} takes a mandatory \mymarg{body content}
%   and \emph{\enquote{stores it}} in the \emph{sequence} and \emph{prop
%   list} \mymarg{store name} set by \mykey{save-ans} key. If \mykey{save-ref}
%   key are active and the \mypkg[no-index]{hyperref}\cite{hyperref} package is
%   detected |\hyperlink| and |\hypertarget| will be used, otherwise the
%   usual \emph{\enquote{label and ref}} system provided by \hologo{LaTeX} will be used.
% \end{function}
%
% By design the environment cannot be nested but full supports
% \emph{\enquote{verbatim material}} in the \mymeta{body} and it is assumed
% that \emph{\enquote{each numbered}} |\item| or |\item*| within the
% environment in which it is active it has a \emph{\enquote{single
% execution}} unless |\item| or |\item*| open a nested level
% or use the \mykey{no-store} key.
%
% The \myenv*[no-index]{anskey*} environment is implemented using the new
% \emph{\enquote{collect code}} \texttt{c-type} argument
% part of \hologo{LaTeX} release 2025-06-01\cite{ltnews}.
% |\begin{anskey*}| and |\end{anskey*}| must be in different lines and
% should not appear within verbatim environments or commands. All \mymeta{keys}
% must be passed separated by commas and \enquote{without separation}
% of the start of the environment.
%
% Comments `\textcolor{gray}{\texttt{\%}}' or \emph{any character}
% after |\begin{anskey*}| or \myoarg{key \textnormal{\textcolor{gray}{=}} val}
% on the same line are NOT supported, \hologo{LaTeX} will return an
% \enquote{error} message if this happens. In a similar way comments
% `\textcolor{gray}{\texttt{\%}}' or \emph{any character} after
% |\end{anskey*}| on the same line \hologo{LaTeX} will
% return a \enquote{warning} message.
%
% \subsubsection{Keys for \texttt{anskey*}}
%
% The \myenv*[no-index]{anskey*} environment uses the same \mymeta{keys} as the
% |\anskey| command  next to the \mymeta{keys} |write-env|, |overwrite|
% and |force-eol|. The environment is available for all levels of the
% \myenv*[no-index]{enumext} environment and the \myenv*[no-index]{enumext*} environment, but
% it is disabled for the \myenv*[no-index]{keyans}, \myenv*[no-index]{keyans*} and
% \myenv*[no-index]{keyanspic} environments.
%
% \keyexampcmd{write-env}{file.ext}{not used}[anskey*]
%
% Sets the \emph{name} of the \mymeta{external file} in which the \mymeta{contents} of
% the environment will be written. The \mymeta{file.ext} will be created
% in the working directory, relative or absolute paths are not supported. If \mymeta{file.ext}
% does not exist, it will be created or overwritten if the \mykey{overwrite} key is used.
%
% \keyexampcmd{overwrite}{true \textnormal{\textcolor{lightgray}{\textbar}} false}{false}[anskey*]
% Sets whether the \mymeta{file.ext} generated by \mykey{write-env}
% from the \myenv*[no-index]{anskey*} environment will be rewritten.
%
% \keyexampcmd{force-eol}{true \textnormal{\textcolor{lightgray}{\textbar}} false}{false}[anskey*]
%
% Sets if the \emph{last end of line} for the \mymeta{stored content} is hidden or not.
% This key is necessary only if the last line is the closing of some
% environment defined by the \mypkg{fancyvrb} package as |\end{|\myenv[no-index]{Verbatim}|}|
% or another environment that does not support a comments `\textcolor{gray}{\texttt{\%}}'
% after closing |\end{|\myenv[no-index]{Verbatim}|}%|.
%
% \subsubsection*{Example}
%
% \iffalse
%<*example>
% \fi
\begin{examplecode}[frame=single]
\begin{enumext}[save-ans=test,show-pos=true,start=5]
  \item* Text containing our instructions or questions.
    \begin{anskey*}[item-star]
      `\mymeta{first answer}`
    \end{anskey*}

  \item Text containing our instructions or questions.
    \begin{enumext}
      \item Question.
        \begin{anskey*}
          `\mymeta{second answer}`
        \end{anskey*}
    \end{enumext}

  \item Text containing our instructions or questions.
    \begin{anskey*}
      `\mymeta{third answer}`
    \end{anskey*}

  \item Text containing our instructions or questions.
    \begin{anskey*}
      `\mymeta{fourth answer}`
    \end{anskey*}
\end{enumext}
\end{examplecode}
% \iffalse
%</example>
% \fi
%
% \begin{enumext}[save-ans=test,show-pos=true,columns=2,nosep,start=5]
%   \item* Text containing our instructions or questions.
%      \begin{anskey*}
%        First answer with \verb+verbatim+
%      \end{anskey*}
%   \item Text containing our instructions or questions.
%     \begin{enumext}
%       \item Question.\anskey{second answer}
%     \end{enumext}
%   \item Text containing our instructions or questions. \anskey{third answer}
%   \item Text containing our instructions or questions. \anskey{fourth answer}
% \end{enumext}
%
% \subsection[The environment \texttt{keyans}]{The environments \env{keyans} and \env{keyans*}}\label{subsec:keyans}
%
% \begin{function}{keyans,keyans*}
%   \begin{syntax}
%     |\begin{keyans}|\myoarg{key \textnormal{\textcolor{gray}{=}} val} |\item| |\item|\myoarg{custom} |\item*| |\item*|\myoarg{content} |\end{keyans}|\par
%     |\begin{keyans*}|\myoarg{key \textnormal{\textcolor{gray}{=}} val} |\item| |\item|\myoarg{custom} |\item*| |\item*|\myoarg{content} |\end{keyans*}|
%   \end{syntax}
%  The \myenv*{keyans} and \myenv*{keyans*} environments are
%  \emph{\enquote{enumerated list}} environments designed for
%  \emph{\enquote{multiple choice}} questions activated by the \mykey{save-ans} key.
%
%  This environments can NOT be nested and must always be at the
%  \emph{\enquote{first level}} of the \myenv*[no-index]{enumext} environment, the
%  commands |\item| and |\item|\myoarg{custom} work in the usual
%  and the command |\item|\myparg{columns} is available for the \myenv*[no-index]{keyans*} environment.
%
%   \begin{important}*
%     The behavior of |\item*| in \myenv*[no-index]{keyans} and \myenv*[no-index]{keyans*} environments
%     is NOT the same as in the \myenv*[no-index]{enumext} or \myenv*[no-index]{enumext*} environments.
%   \end{important}
% \end{function}
%
% \iffalse
%<*example>
% \fi
\vspace*{-1\baselineskip}
\begin{minipage}[t]{0.45\linewidth}
\begin{examplecode}[frame=none]
\begin{enumext}[save-ans=test]
  \item `\mymeta[type=tt,ac=lightgray,cf=gray]{item content}`
    \begin{keyans}`\myoarg{key \textnormal{\textcolor{gray}{=}} val}`
      \item `\mymeta[type=tt,ac=lightgray,cf=gray]{item content}`
      \item `\myoarg{custom} \mymeta[type=tt,ac=lightgray,cf=gray]{item content}`
      \item* `\mymeta[type=tt,ac=lightgray,cf=gray]{item content}`
      `\ics*{item*}\myoarg{content} \mymeta[type=tt,ac=lightgray,cf=gray]{item content}`
    \end{keyans}
\end{enumext}
\end{examplecode}
\end{minipage}
\hfill
\begin{minipage}[t]{0.45\linewidth}
\begin{examplecode}[frame=none]
\begin{enumext}[save-ans=test]
  \item `\mymeta[type=tt,ac=lightgray,cf=gray]{item content}`
    \begin{keyans*}`\myoarg{key \textnormal{\textcolor{gray}{=}} val}`
      \item `\mymeta[type=tt,ac=lightgray,cf=gray]{item content}`
      \item `\myoarg{custom} \mymeta[type=tt,ac=lightgray,cf=gray]{item content}`
      \item* `\mymeta[type=tt,ac=lightgray,cf=gray]{item content}`
      `\ics*{item*}\myoarg{content} \mymeta[type=tt,ac=lightgray,cf=gray]{item content}`
    \end{keyans*}
\end{enumext}
\end{examplecode}
\end{minipage}
% \iffalse
%</example>
% \fi
%
% The \mymeta{keys} set in the \emph{optional argument} of the environment are the same
% (almost) as those of the \myenv*[no-index]{enumext} and \myenv*[no-index]{enumext*}
% environments and have \emph{higher precedence} than those set by
% |\setenumext|\myoarg{keyans}\mymarg{key \textnormal{\textcolor{gray}{=}} val} or
% |\setenumext|\myoarg{keyans\textnormal{\textcolor{star}{\texttt{*}}}}\mymarg{key \textnormal{\textcolor{gray}{=}} val}.
% If the \emph{optional argument} is not passed or the \mymeta{keys}
% are not set by |\setenumext|, the default values will be the same
% as the \emph{\enquote{second level}} of the \myenv*[no-index]{enumext} environment
% with the difference in the \mymeta{label} which will be set to \mykey{label}=\ilabel[no-index]{Alph*)}.
%
% The keys \mykey{mark-ans*}, \mykey{mark-pos*}, \mykey{mark-sep*}, \mykey{save-sep},
% \mykey{wrap-opt}, \mykey{wrap-ans*}, \mykey{show-ans} and \mykey{show-pos}
% are available for both environments.
%
% \subsubsection[The \texttt{\textbackslash item*} in \texttt{keyans}]{The \cs[no-index]{item*} in \texttt{keyans} and \texttt{keyans*}}
%
% \begin{function}{\item*}
%   \begin{syntax}
%     |\item*|
%     |\item*|\myoarg{content}
%   \end{syntax}
%    The \ics*{item*} and |\item*|\myoarg{content} command \emph{\enquote{store}}
%    the current \mymeta{label} set by \mykey{label} key next to the \emph{optional argument} \mymeta{content}
%    in \emph{sequence} and \emph{prop list} \mymarg{store name} set by
%    \mykey{save-ans} key in the \emph{\enquote{first level}} of the
%    \myenv*[no-index]{enumext} or \myenv*[no-index]{enumext*} environments.
% \end{function}
%
% The \emph{starred argument} `|*|' cannot be separated by spaces `\verb*| |'
% from the command, i.e. |\item*| and the \emph{optional argument}
% does \enquote{NOT} support \emph{verbatim content}. By design it is
% assumed that the |\item*| will only appear \emph{\enquote{once}}
% within the environment.
%
% \subsubsection*{Example}
%
% \iffalse
%<*example>
% \fi
\begin{examplecode}[frame=single]
\begin{enumext}[save-ans=test,columns=2,show-ans=true]
  \item Text containing a question.

    \begin{keyans*}[nosep,columns=2]
      \item Choice
      \item* Correct choice
      \item Choice
      \item Choice
      \item Choice
    \end{keyans*}

  \item Text containing a question and image.

    \begin{keyans}[nosep,mini-env={0.4\linewidth}]
      \item Choice
      \item Choice
      \item Choice
      \item Choice
      `\ics*{item*}\myoarg{note}` Correct choice
      \miniright
      \includegraphics[scale=0.25]{example-image-a}
      Some text
    \end{keyans}
\end{enumext}
\end{examplecode}
% \iffalse
%</example>
% \fi
%
% \begin{enumext}[save-ans=test,mark-ans*={\oldtextasteriskcentered},show-ans=true,columns=2,nosep]
%  \item Text containing a question.
%
%    \begin{keyans*}[nosep,columns=2]
%      \item Choice
%      \item* Correct choice
%      \item Choice
%      \item Choice
%      \item Choice
%    \end{keyans*}
%
%   \item Text containing a question and image.
%
%     \begin{keyans}[nosep,mini-env={0.4\linewidth}]
%       \item Choice
%       \item Choice
%       \item Choice
%       \item Choice
%       \item*[note] Correct choice
%       \miniright
%       \includegraphics[scale=0.15]{example-image-a}
%
%       Some text
%     \end{keyans}
% \end{enumext}
%
% \subsection{The environment \texttt{keyanspic}}\label{subsec:keyanspic}
%
% \begin{function}{keyanspic}
%   \begin{syntax}
%     |\begin{keyanspic}|\myoarg{key \textnormal{\textcolor{gray}{=}} val} |\anspic*|\myoarg{content}\mymarg{drawing or tabular} |\end{keyanspic}|
%   \end{syntax}
%  The \myenv*{keyanspic} environment is an \emph{\enquote{enumerated
%  list}} environment activated by the \mykey{save-ans} key that has the
%  same configuration for \emph{\enquote{spacing}} and \mymeta{label}
%  as the \myenv*[no-index]{keyans} environment that uses the |\anspic| command
%  instead of |\item|. It is intended for placing \emph{drawings or tabular}
%  with \mymeta{label} centered \emph{above} or \emph{below} in a \emph{single line}
%  or \emph{upper and lower} layout style.
%
%  When the \myenv*[no-index]{keyanspic} environment is used \emph{without keys}
%  the \mymeta{labels} are centered \emph{below} the \emph{drawings
%  or tabular} in a \emph{single line} layout style.
% \end{function}
%
% A representation of the output can be seen in the figure~\ref{fig:keyanspicdoc}.
%
% \begin{figure}[ht]
%   \centering
%   \begin{tikzpicture}[line cap=round,line join=round,x=0.8cm,y=0.8cm,every node/.style={font=\footnotesize}]
%     % space above
%     \draw[color=white] (7.0,6.10) rectangle (17.5,6.5) node[pos=.5,text=black]%
%       { \textcolor{lightgray}{\texttt{number of drawings or tabular upper}} };
%     \draw[color=white] (7.0,4.5) rectangle (17.5,4.9) node[pos=.5,text=black]%
%       { \textcolor{keyname}{\texttt{layout-sep}} }; %%%[+itemsep]
%     \draw[color=white] (7.0,2.9) rectangle (17.5,3.3) node[pos=.5,text=black]%
%       { \textcolor{lightgray}{\texttt{number of drawings or tabular lower}} };
%     \draw[<->,color=gray] (2.0,7.0) --node[above,text=keyname]{\vphantom{p}\texttt{labelwidth}} (5.0,7.0);
%     \draw[<->,color=gray] (5.0,7.0) --node[above,text=keyname]{\texttt{labelsep}} (7.0,7.0);
%     \draw[color=lightgray,fill=lightgray,fill opacity=0.15] (2.0,6.5) rectangle (5.0,6.85);% labelwidth
%     \draw[color=lightgray,fill=lightgray,fill opacity=0.15] (5.0,6.5) rectangle (7.0,6.85);% labelsep
%     \draw[color=lightgray,fill=lightgray,fill opacity=0.15] (7.0,6.5) rectangle (17.5,6.85) node[pos=.5,text=black]%
%       {\begin{tabular}{c}
%          \textcolor{blue}{\emph{item text above the keyanspic environment}}
%        \end{tabular}};% item width
%     \draw[color=lightgray,fill=lightgray,fill opacity=0.15] (7.0,5.25) rectangle (10.5,6.10) node[pos=.5,text=black]%
%       {\begin{tabular}{c}
%          \texttt{above one} \\
%          \textcolor{blue}{\emph{drawing or tabular}}
%        \end{tabular}};
%     \draw[color=lightgray,fill=lightgray,fill opacity=0.15] (7.0,4.9) rectangle (10.5,5.25) node[pos=.5,text=black]%
%       { \textcolor{blue}{\texttt{label}} };
%     \draw[color=lightgray,fill=lightgray,fill opacity=0.15] (10.5,5.25) rectangle (14,6.10) node[pos=.5,text=black]%
%       {\begin{tabular}{c}
%          \texttt{above two} \\
%          \textcolor{blue}{\emph{drawing or tabular}}
%        \end{tabular}};
%     \draw[color=lightgray,fill=lightgray,fill opacity=0.15] (10.5,4.9) rectangle (14,5.25) node[pos=.5,text=black]%
%       { \textcolor{blue}{\texttt{label}} };
%     \draw[color=lightgray,fill=lightgray,fill opacity=0.15] (14,5.25) rectangle (17.5,6.10) node[pos=.5,text=black]%
%      {\begin{tabular}{c}
%          \texttt{above three} \\
%          \textcolor{blue}{\emph{drawing or tabular}}
%        \end{tabular}};
%     \draw[color=lightgray,fill=lightgray,fill opacity=0.15] (14,4.9) rectangle (17.5,5.25) node[pos=.5,text=black]%
%       { \textcolor{blue}{\texttt{label}} };
%     \draw[color=lightgray,fill=lightgray,fill opacity=0.15] (7.0,3.65) rectangle (12.25,4.5) node[pos=.5,text=black]%
%      {\begin{tabular}{c}
%          \texttt{below one} \\
%          \textcolor{blue}{\emph{drawing or tabular}}
%        \end{tabular}};
%     \draw[color=lightgray,fill=lightgray,fill opacity=0.15] (7.0,3.3) rectangle (12.25,3.65) node[pos=.5,text=black]%
%       { \textcolor{blue}{\texttt{label}} };
%     \draw[color=lightgray,fill=lightgray,fill opacity=0.15] (12.25,3.65) rectangle (17.5,4.5) node[pos=.5,text=black]%
%      {\begin{tabular}{c}
%          \texttt{below two} \\
%          \textcolor{blue}{\emph{drawing or tabular}}
%        \end{tabular}};
%     \draw[color=lightgray,fill=lightgray,fill opacity=0.15] (12.25,3.3) rectangle (17.5,3.65) node[pos=.5,text=black]%
%       { \textcolor{blue}{\texttt{label}} };
%     \draw[color=lightgray,fill=lightgray,fill opacity=0.15] (2.0,2.55) rectangle (5.0,2.9);% labelwidth
%     \draw[color=lightgray,fill=lightgray,fill opacity=0.15] (5.0,2.55) rectangle (7.0,2.9);% labelsep
%     \draw[color=lightgray,fill=lightgray,fill opacity=0.15] (7.0,2.55) rectangle (17.5,2.9) node[pos=.5,text=black]%
%       {\begin{tabular}{c}
%          \textcolor{blue}{\emph{item text below the keyanspic environment}}
%        \end{tabular}};% item width
%   \end{tikzpicture}
%   \caption{Representation of the \myenv*{keyanspic} environment with \mykey{layout-sty}=\mymarg{3, 2} in \mypkg*[no-index]{enumext}.}
%   \label{fig:keyanspicdoc}
% \end{figure}
%
% This environment cannot be nested and must \emph{always} be at the
% \emph{\enquote{first level}} of the \myenv*[no-index]{enumext} environment, the
% |\item| command is disabled and \mymeta{keys} cannot be set using
% |\setenumext|.
%
% \subsubsection{Keys for \texttt{keyanspic}}
%
% \keyexampcmd{label-pos}{above \textnormal{\textcolor{lightgray}{\textbar}} below}{below}[env]
%
% Set the \emph{position} of \mymeta{label} to be centered
% \enquote{above} or \enquote{below} \emph{drawings} or \emph{tabular}
% when the |\anspic| command is executed.
%
% \keyexampcmd{label-sep}{rubber length \textnormal{\textcolor{lightgray}{\textbar}} rigid length}{internal adjustment}[env]
%
% Set the \emph{vertical spacing} between the \mymeta{label} centered
% \enquote{above} or \enquote{below} and \emph{drawings} or \emph{tabular}
% when running the |\anspic| command.
%
% \keyexampcmd{layout-sty}{n° upper \textnormal{\textcolor{gray}{,}} n° lower}{not set}[env]
%
% Set the \emph{number} of \emph{drawings} or \emph{tabular} that will be
% distributed \enquote{upper} and \enquote{lower} within the environment
% when executing the |\anspic| command. The value must be passed in braces
% and if not set or the \mymeta{nº lower} is omitted the \emph{drawings}
% or \emph{tabular} will be put on a \emph{single line}.
%
% \keyexampcmd{layout-sep}{rubber length \textnormal{\textcolor{lightgray}{\textbar}} rigid length}{adjusted parsep from keyans}[env]
%
% Set the \emph{vertical separation} between the number of
% \emph{drawings} or \emph{tabular} placed at the \enquote{upper} and
% \enquote{lower} within the environment when executing the |\anspic|
% command. Internally adjusts the \mykey{parsep} value taken from the
% \myenv*[no-index]{keyans} environment.
%
% \keyexampcmd{layout-top}{rubber length \textnormal{\textcolor{lightgray}{\textbar}} rigid length}{adjusted topsep from keyans}[env]
%
% Set the \emph{vertical space} added to both the top and bottom of the environment.
% Internally adjust the value of \mykey{topsep} taken from \myenv*[no-index]{keyans} environment.
%
% The keys \mykey{mark-ans*}, \mykey{mark-pos*}, \mykey{mark-sep*}, \mykey{save-sep},
% \mykey{wrap-opt}, \mykey{wrap-ans*}, \mykey{show-ans} and \mykey{show-pos}
% are available for this environment.
%
% \subsubsection{The command \cs[no-index]{anspic}}
%
% \begin{function}{\anspic}
%   \begin{syntax}
%      \cmdexamp{anspic}{drawing or tabular}
%      \cmdexamp[*]{anspic}[content]{drawing or tabular}
%   \end{syntax}
%  The \ics*{anspic} command takes three arguments, the \emph{starred
%  argument} `|*|' store the current \mymeta{label} next to the
%  \emph{optional argument} \mymeta{content} in \emph{sequence} and
%  \emph{prop list} \mymarg{store name} set by \mykey{save-ans} key.
% \end{function}
%
% The \emph{starred argument} `|*|' cannot be separated by spaces `\verb*| |'
% from the command, i.e. |\anspic*| and the \emph{optional argument}
% does \enquote{NOT} support \emph{verbatim content}. By design it
% is assumed that the \emph{starred argument} `|*|' will only
% appear \emph{\enquote{once}} within the environment.
%
% \subsubsection*{Example}
%
% \iffalse
%<*example>
% \fi
\begin{examplecode}[frame=single]
\begin{enumext}[save-ans=test,show-ans=true,nosep]
  \item Question with images and labels below.

    \begin{keyanspic}[layout-sty={3,2}]
      \anspic{\includegraphics[scale=0.15]{example-image-a}}
      \anspic{\includegraphics[scale=0.15]{example-image-b}}
      \anspic{\includegraphics[scale=0.15]{example-image-a}}
      \anspic{\includegraphics[scale=0.15]{example-image-a}}
      \anspic*[note]{\includegraphics[scale=0.15]{example-image-a}}
    \end{keyanspic}

  \item Question with images and labels above.

    \begin{keyanspic}[label-pos=above, layout-sty={3,2},layout-sep=0.25cm]
      \anspic{\includegraphics[scale=0.15]{example-image-a}}
      \anspic{\includegraphics[scale=0.15]{example-image-b}}
      \anspic{\includegraphics[scale=0.15]{example-image-a}}
      \anspic{\includegraphics[scale=0.15]{example-image-a}}
      \anspic*[note]{\includegraphics[scale=0.15]{example-image-a}}
    \end{keyanspic}

  \item Question with images and labels below on a single line.

    \begin{keyanspic}
      \anspic{\includegraphics[scale=0.15]{example-image-a}}
      \anspic{\includegraphics[scale=0.15]{example-image-b}}
      \anspic{\includegraphics[scale=0.15]{example-image-a}}
      \anspic{\includegraphics[scale=0.15]{example-image-a}}
      \anspic*[note]{\includegraphics[scale=0.15]{example-image-a}}
    \end{keyanspic}

\end{enumext}
\end{examplecode}
% \iffalse
%</example>
% \fi
%
% \begin{enumext}[save-ans=test, mark-ans*={\oldtextasteriskcentered},show-ans=true,nosep, before={\setenumext[keyans]{nosep,parsep=5mm}}]
% \item Question with images and labels below.
%    \begin{keyanspic}[layout-sty={3,2}]
%      \anspic{\includegraphics[scale=0.15]{example-image-a}}
%      \anspic{\includegraphics[scale=0.15]{example-image-b}}
%      \anspic{\includegraphics[scale=0.15]{example-image-a}}
%      \anspic{\includegraphics[scale=0.15]{example-image-a}}
%      \anspic*[note]{\includegraphics[scale=0.15]{example-image-a}}
%    \end{keyanspic}
% \item Question with images and labels above.
%    \begin{keyanspic}[label-pos=above, layout-sty={3,2},layout-sep=0.25cm,layout-top=10pt]
%      \anspic{\includegraphics[scale=0.15]{example-image-a}}
%      \anspic{\includegraphics[scale=0.15]{example-image-b}}
%      \anspic{\includegraphics[scale=0.15]{example-image-a}}
%      \anspic{\includegraphics[scale=0.15]{example-image-a}}
%      \anspic*[note]{\includegraphics[scale=0.15]{example-image-a}}
%    \end{keyanspic}
% \item Question with images and labels below on a single line.
%    \begin{keyanspic}[layout-top=10pt]
%      \anspic{\includegraphics[scale=0.15]{example-image-a}}
%      \anspic{\includegraphics[scale=0.15]{example-image-b}}
%      \anspic{\includegraphics[scale=0.15]{example-image-a}}
%      \anspic{\includegraphics[scale=0.15]{example-image-a}}
%      \anspic*[note]{\includegraphics[scale=0.15]{example-image-a}}
%    \end{keyanspic}
% \end{enumext}
%
% \begin{important}
%  Remember to pass the |alt=|\mymarg{description} key to the |\includegraphics|
%  command when creating a \emph{tagged} PDF.
% \end{important}
%
% \subsection{Printing stored content}\label{subsec:printing}
%
% \subsubsection{The command \cs[no-index]{getkeyans}}\label{subsubsec:getkeyans}
%
% \begin{function}{\getkeyans}
%   \begin{syntax}
%      \cmdexamp{getkeyans}{store name \textnormal{\textcolor{red}{:}} position}
%   \end{syntax}
%   The command \ics*{getkeyans} prints the \emph{\enquote{stored content}}
%   in \emph{prop list} \mymarg{store name} defined by
%   \mykey{save-ans} key in the \mymeta{position} returned by the
%   \mykey{show-pos} key.
%
%   The \emph{\enquote{stored content}} can only be accessed \emph{after}
%   it is stored, if \mymarg{store name} does not exist the command
%   will return an error.
% \end{function}
%
% The form taken by the argument \mymarg{store name \textnormal{\textcolor{red}{:}} position}
% is the same as that used to generate the \emph{\enquote{internal
% label and ref}} system when \mykey{save-ref} key are active, so to
% refer to a \emph{\enquote{stored content}}. For example
% |\getkeyans{test:4}| will return the \emph{\enquote{stored content}}
% at position |4| of the environment in which the key |save-ans=test|
% was set.
%
% \subsubsection{The command \cs{foreachkeyans}}\label{subsubsec:foreachkeyans}
%
% \begin{function}{\foreachkeyans}
%   \begin{syntax}
%     \cmdexamp{foreachkeyans}[key \textnormal{\textcolor{gray}{=}} val]{store name}
%   \end{syntax}
% \end{function}
%
% The command \ics*{foreachkeyans} goes through and executes the command |\getkeyans|
% on the contents in \emph{prop list} \mymarg{store name}. If you pass without options run
% |\getkeyans| on all contents in \emph{prop list} \mymarg{store name}.
%
% \subsubsection*{Options for command}\label{subsubsec:optcmdforeach}
%
% \keyexampcmd{sep}{code}{\{; \}}[foreachkeyans]
%
% Establishes the \emph{separation} between \emph{\enquote{each}} \mymarg{content} stored in \emph{prop list} \mymarg{store name}.
% For example, you can use \mykey[foreachkeyans]{sep}|={\\[10pt]}| for vertical separation of stored
% contents.
%
% \keyexampcmd{step}{integer}{\mydim{1}}[foreachkeyans]
%
% Sets the \emph{step} (increment) applied to the value set by key \mykey[foreachkeyans]{start}
% for \emph{\enquote{each}} \mymarg{content} stored in \emph{prop list} \mymarg{store name}. The value must be a
% \mymeta{positive integer}.
%
% \keyexampcmd{start}{integer}{\mydim{1}}[foreachkeyans]
%
% Sets the \emph{position} of the \emph{prop list} \mymarg{store name} from which execution
% will start. The value must be a \mymeta{positive integer}.
%
% \keyexampcmd{stop}{integer}{\mydim{0}}[foreachkeyans]
%
% Sets the \emph{position} of the \emph{prop list} \mymarg{store name} from which execution
% will finish. The value must be a \mymeta{positive integer}.
%
% \keyexampcmd{before}{code}{empty}[foreachkeyans]
%
% Sets the \mymarg{code} that will be executed \mymeta{before} each \mymarg{content} stored
% in \emph{prop list} \mymarg{store name}. The \mymarg[type=tt]{code} must be passed between braces.
%
% \keyexampcmd{after}{code}{empty}[foreachkeyans]
%
% Sets the \mymarg{code} that will be executed \mymeta{after} each \mymarg{content} stored
% in \emph{prop list} \mymarg{store name}. The \mymarg{code} must be passed between braces.
%
% \keyexampcmd{wrapper}{code{} \space\usebox{\hashbox} more code}{empty}[foreachkeyans]
%
% Wraps the \mymarg{content} stored in \emph{prop list} \mymarg{store name} referenced by |{#1}|.
% The \mymarg{code} must be passed between braces. For example
% |\foreachkeyans[wrapper={\makebox[1em][l]{#1}}]|\mymarg{store name}.
%
% \subsubsection{The command \cs[no-index]{printkeyans}}\label{subsubsec:printkeyans}
%
% \begin{function}{\printkeyans}
%   \begin{syntax}
%      \cmdexamp{printkeyans}{store name}
%      \cmdexamp{printkeyans}[keys]{store name}
%      \cmdexamp[*]{printkeyans}[keys]{store name}
%   \end{syntax}
%   The command \ics*{printkeyans} prints \emph{\enquote{all stored content}}
%   in \emph{sequence} \mymarg{store name} defined by \mykey{save-ans}
%   key placing this inside the \myenv*[no-index]{enumext} or \myenv*[no-index]{enumext*}
%   environment if the \emph{starred argument} `|*|' is used.
%
%   The \emph{\enquote{stored content}} can only be accessed \emph{after} it is
%   stored in the \emph{sequence}, if \mymarg{store name} does not exist the
%   command will return an error.
% \end{function}
%
% The \emph{optional argument} allows managing the \mymeta{keys} in the
% \emph{\enquote{first level}} of the environment in which the
% \emph{\enquote{stored content}} of the \emph{sequence} \mymarg{store
% name} will be printed, if the \emph{starred argument} `|*|' is used it will be
% \myenv*[no-index]{enumext*} otherwise \myenv*[no-index]{enumext}.
%
% The default values for the \emph{\enquote{first level}} are the
% same as the default values for the \myenv*[no-index]{enumext} and \myenv*[no-index]{enumext*}
% environments along with the keys |nosep,first=\small,font=\small| and |columns=2|.
% For the inner levels of the environment \myenv*[no-index]{enumext} saved in the
% \emph{sequence} \mymarg{store name} the default values are the same
% as those established for the second, third and fourth levels plus the
% keys |nosep,first=\small,font=\small|. If the environment \myenv*[no-index]{enumext*}
% is saved within the \emph{sequence} \mymarg{store name} it will have
% the same default values plus the keys |nosep,first=\small, font=\small|.
% These are the default values of the package: the values set with |\setenumext|
% for the environments \myenv*[no-index]{enumext} and \myenv*[no-index]{enumext*}
% do not modify them.
%
% Since the command encapsulates by default the \myenv*[no-index]{enumext}
% environment or the \myenv*[no-index]{enumext*} environment, some
% considerations must be taken into account:
%
%\begin{itemize}[topsep=0pt]
% \item If |\printkeyans*|\mymarg{store name} is executed and the \emph{sequence}
%   \mymarg{store name} already contains any \myenv*[no-index]{enumext*} environment
%   an error will be returned, since nesting is not possible.
%
% \item If |\printkeyans*|\mymarg{store name} is executed and the \emph{sequence}
%   \mymarg{store name} contains any \myenv*[no-index]{enumext} environments, they
%   will start with the \mymeta{keys} set for the first level unless they
%   are set in the \emph{optional argument} or \mykey{save-key} is used to modify it.
%
% \item If |\printkeyans|\mymarg{store name} is executed and the \emph{sequence}
%   \mymarg{store name} contains any environment \myenv*[no-index]{enumext*}, they
%   will start with the \mymeta{keys} set by default unless they are set
%   in the \emph{optional argument} or \mykey{save-key} is used to modify it.
%\end{itemize}
%
% The default values for the \emph{\enquote{first level}} of |\printkeyans| commands and
% |\printkeyans*| are established using |\setenumext|\myoarg{print \textnormal{\textcolor{gray}{,}} \textcolor{star}{1}}\mymarg{keys} and
% |\setenumext|\myoarg{print\textnormal{\textcolor{star}{\texttt{*}}}}\mymarg{keys}.
%
% To set the \mymeta{keys} for the environment
% \myenv*[no-index]{enumext} \enquote{saved} in the \emph{sequence} \mymarg{store
% name}, use |\setenumext|\myoarg{print
% \textnormal{\textcolor{gray}{,}} level}\mymarg{keys}, and to
% set the \mymeta{keys} for the environment \myenv*[no-index]{enumext*}
% \enquote{saved} in the \emph{sequence} \mymarg{store name}, use
% |\setenumext|\myoarg{print \textnormal{\textcolor{gray}{,}}
% \textnormal{\textcolor{star}{\texttt{*}}}}\mymarg{keys}.
%
% \subsubsection*{Example}
%
% \iffalse
%<*example>
% \fi
\begin{examplecode}[frame=single]
\begin{enumext}[save-ans=sample,columns=1,show-pos=true,nosep,save-ref=true]
  \item Factor $3x+3y+3z$. \anskey{$3(x+y+z)$}
  \item True False

    \begin{enumext}[nosep]
      \item \LaTeX2e\ is cool? \anskey{Very True!}
    \end{enumext}

  \item Related to Linux

    \begin{enumext}[nosep]
      \item You use Linux? \anskey{Yes}
      \item Rate the following package and class
        \begin{enumext}[nosep]
          \item \texttt{xsim} \anskey{very good}
          \item \texttt{exsheets} \anskey{obsolete}
        \end{enumext}
    \end{enumext}
\end{enumext}

The answer to \ref{sample:4} is \getkeyans{sample:4} and the answers to
all the worksheets are as follows:

\printkeyans{sample}
\end{examplecode}
% \iffalse
%</example>
% \fi
%
% \begin{enumext}[save-ans=sampleXX,columns=1,show-pos=true,save-ref=true]
%  \item Factor $3x+3y+3z$. \anskey{$3(x+y+z)$}
%  \item True False
%    \begin{enumext}[nosep]
%      \item \LaTeX2e\ is cool? \anskey{Very True!}
%    \end{enumext}
%  \item Related to Linux
%    \begin{enumext}[nosep]
%      \item You use Linux? \anskey{Yes}
%      \item Rate the following package and class
%        \begin{enumext}[nosep]
%          \item \texttt{xsim} \anskey{very good}
%          \item \texttt{exsheets} \anskey{obsolete}
%        \end{enumext}
%    \end{enumext}
% \end{enumext}
%
% The answer to \ref{sampleXX:4} is \getkeyans{sampleXX:4} and the answers to
% all the worksheets are as follows:
%
% \printkeyans[columns=1]{sampleXX}
%
% \section{Full examples}
%
% Some adapted questions taken from \href{https://tex.stackexchange.com}{TeX-SX} are included as an example.
% The source files of the examples can be compiled with the excellent \arara\footnote{The cool \TeX\ automation tool:
% \url{https://www.ctan.org/pkg/arara}} tool.
%
% \subsubsection*{Example 1}
%
% Adapted from the response given by Enrico Gregorio in \href{https://tex.stackexchange.com/q/461730/7832}{Squares for answer choice options and perfect alignment to mathematical answers}.
%   \myattachfile*{enumext-exa-1}
%
% \setenumext[keyans]{label=\Alph*,font=\small,nosep,wrap-label={\fbox{\makebox[\height]{##1}}}}
% \begin{enumext}[columns=2,save-ans=sabastiano,itemsep=2pt, topsep=5pt,below*={8pt}]
% \item La velocità di \qty{1,00e2}{m/s} espressa in \unit{km/h} è:
%
%   \begin{keyans}
%     \item \qty{36}{km/h}.
%     \item* \qty{360}{km/h}.
%     \item \qty{27,8}{km/h}.
%     \item \qty{3,60e8}{km/h}.
%   \end{keyans}
%
% \item  In fisica nucleare si usa l'angstrom (simbolo:
% $\qty{1}{\angstrom} = \qty{1e-10}{m}$) e il fermi o femtometro
% ($\qty{1}{fm} = \qty{1e-15}{m}$). Qual è la relazione tra queste due
% unità di misura?
%   \begin{keyans}
%     \item* $\qty{1}{\angstrom}=\qty{1e5}{fm}$.
%     \item $\qty{1}{\angstrom}=\qty{1e-5}{fm}$.
%     \item $\qty{1}{\angstrom}=\qty{1e-15}{fm}$.
%     \item $\qty{1}{\angstrom}=\qty{1e3}{fm}$.
%   \end{keyans}
%
% \item La velocità di \qty{1,00e2}{m/s} espressa in \unit{km/h} è:
%
%   \begin{keyans}
%     \item \qty{36}{km/h}.
%     \item* \qty{360}{km/h}.
%     \item \qty{27,8}{km/h}.
%     \item \qty{3,60e8}{km/h}.
%   \end{keyans}
%
% \item  In fisica nucleare si usa l'angstrom (simbolo:
% $\qty{1}{\angstrom} = \qty{1e-10}{m}$) e il fermi o femtometro
% ($\qty{1}{fm} = \qty{1e-15}{m}$). Qual è la relazione tra queste due
% unità di misura?
%   \begin{keyans}
%     \item* $\qty{1}{\angstrom}=\qty{1e5}{fm}$.
%     \item $\qty{1}{\angstrom}=\qty{1e-5}{fm}$.
%     \item $\qty{1}{\angstrom}=\qty{1e-15}{fm}$.
%     \item $\qty{1}{\angstrom}=\qty{1e3}{fm}$.
%   \end{keyans}
% \end{enumext}
%
% \printkeyans[columns=1]{sabastiano}
%
% \subsubsection*{Example 2}
%
% Adapted from the response given by Florent Rougon in \href{https://tex.stackexchange.com/q/529034/7832}{Multiple choice questions with proposed answers in random order — addition of automatic correction (cross mark)}.
%   \myattachfile*{enumext-exa-2}
% \setenumext[keyans]{label=\Alph*,font=\small,nosep,wrap-label={\mySquared{##1}}}
% \begin{enumext}[columns=2,save-ans=sabastiano2,mark-sep*=2pt,mark-ans*={\textcolor{gray}{✔}},save-ref=true,show-ans=true,itemsep=0pt]
% \item La velocità di \qty{1,00e2}{m/s} espressa in \unit{km/h} è:
%   \begin{keyans}
%     \item \qty{36}{km/h}.
%     \item* \qty{360}{km/h}.
%     \item \qty{27,8}{km/h}.
%     \item \qty{3,60e8}{km/h}.
%   \end{keyans}
%
% \item In fisica nucleare si usa l'angstrom (simbolo:
% $\qty{1}{\angstrom} = \qty{1e-10}{m}$) e il fermi o femtometro
% ($\qty{1}{fm} = \qty{1e-15}{m}$). Qual è la relazione tra queste due
% unità di misura?
%   \begin{keyans}
%     \item* $\qty{1}{\angstrom}=\qty{1e5}{fm}$.
%     \item $\qty{1}{\angstrom}=\qty{1e-5}{fm}$.
%     \item $\qty{1}{\angstrom}=\qty{1e-15}{fm}$.
%     \item $\qty{1}{\angstrom}=\qty{1e3}{fm}$.
%   \end{keyans}
%
% \item La velocità di \qty{1,00e2}{m/s} espressa in \unit{km/h} è:
%   \begin{keyans}
%     \item \qty{36}{km/h}.
%     \item* \qty{360}{km/h}.
%     \item \qty{27,8}{km/h}.
%     \item \qty{3,60e8}{km/h}.
%   \end{keyans}
%
% \item In fisica nucleare si usa l'angstrom (simbolo:
% $\qty{1}{\angstrom} = \qty{1e-10}{m}$) e il fermi o femtometro
% ($\qty{1}{fm} = \qty{1e-15}{m}$). Qual è la relazione tra queste due
% unità di misura?
%   \begin{keyans}
%     \item* $\qty{1}{\angstrom}=\qty{1e5}{fm}$.
%     \item $\qty{1}{\angstrom}=\qty{1e-5}{fm}$.
%     \item $\qty{1}{\angstrom}=\qty{1e-15}{fm}$.
%     \item $\qty{1}{\angstrom}=\qty{1e3}{fm}$.
%   \end{keyans}
% \end{enumext}
%
% \printkeyans*[columns=2]{sabastiano2}
%
% \subsubsection*{Example 3}
%
% A \emph{\enquote{simple multiple choice}} test.
%   \myattachfile*{enumext-exa-3}
% \setenumext[keyans]{labelsep=8pt,label=\Alph*,font=\small,nosep,wrap-label={\myCircled{##1}}}
% \begin{enumext}[mode-box,save-ans=simplemc,columns=1,nosep,save-ref=true]
%   \item First type of questions
%      \begin{keyans*}[labelsep=8pt,label=\Alph*,font=\small,nosep,wrap-label={\myCircled{#1}},columns=2]
%        \item value
%        \item*[$x=5$] correct
%        \item value
%        \item value
%      \end{keyans*}
%   \item Second type of questions
%      \begin{enumext}[no-store,label=\Roman*.,nosep]
%        \item $2\alpha+2\delta=90^{\circ}$
%        \item $\alpha=\delta$
%        \item $\angle EDF=45^{\circ}$
%      \end{enumext}
%      \begin{keyans}[columns=2,noitemsep]
%        \item I only
%        \item II only
%        \item I and II only
%        \item* I and III only
%        \item I, II, and III
%      \end{keyans}
%   \item Third type of questions
%        \begin{enumext}[no-store,list-offset=0.5cm,label=(\arabic*),nosep]
%            \item $2\alpha+2\delta=90^{\circ}$
%            \item $\angle EDF=45^{\circ}$
%        \end{enumext}
%    \begin{keyans}[columns=2,noitemsep]
%        \item value
%        \item value
%        \item*[some note] value
%        \item value
%        \item value
%    \end{keyans}
%
% \item Question with image and label below:
%  \begin{keyanspic}[label-sep=3pt]
%    \anspic{\includegraphics[scale=0.15]{example-image-a}}
%    \anspic{\includegraphics[scale=0.15]{example-image-b}}
%    \anspic{\includegraphics[scale=0.15]{example-image-a}}
%    \anspic{\includegraphics[scale=0.15]{example-image-a}}
%    \anspic*[A duck]{\includegraphics[scale=0.32]{example-image-duck}}
%  \end{keyanspic}
%
% \item Question with image on right side:
%    \begin{keyans}[mini-env=0.5\linewidth, columns=1,noitemsep]
%        \item value
%        \item value
%        \item value
%        \item*[other note] correct
%        \item value
%        \miniright
%        \includegraphics[scale=0.15]{example-image-b}
%    \end{keyans}
% \end{enumext}
%
% Test keys
%
% \printkeyans[columns=2]{simplemc}
%
% \subsubsection*{Example 4}
%
% A \emph{\enquote{simple worksheet}} using ducks :)
%   \myattachfile*{enumext-exa-4}
%
% \begin{enumext}[save-ans=simplews,columns=1,check-ans=true,nosep,save-ref=true,wrap-label={\tikz[scale=0.25]\duck[signpost=\scalebox{0.6}{#1}];}]
%  \item Factor $x^{2}-2x+1$ \anskey{$\left(x-1\right)^{2}$}
%  \item Factor $3x+3y+3z$ \anskey{$3(x+y+z)$}
%  \item[] The following questions need to be cuaqtified :)
%  \item True False
%    \begin{enumext}[nosep]
%      \item $\alpha > \delta$ \anskey{False}
%      \item \LaTeX2e\ is cool? \anskey{Very True!}
%    \end{enumext}
%  \item Related to Linux
%    \begin{enumext}[nosep]
%      \item You use Linux? \anskey{Yes}
%      \item Usually uses the package manager? \anskey{Yes, \texttt{dnf}}
%      \item Rate the following package and class
%        \begin{enumext}[nosep]
%          \item \texttt{xsim-exam} \anskey{doesn't exist for now :(}
%          \item \texttt{xsim} \anskey{very good}
%          \item \texttt{exsheets} \anskey{obsolete}
%        \end{enumext}
%    \end{enumext}
% \end{enumext}
%
% The answer to \ref{simplews:1} is \getkeyans{simplews:1} and the answer
% to \ref{simplews:3} is \getkeyans{simplews:3}.
%
% \printkeyans[columns=2]{simplews}
%
% \subsubsection*{Example 5}
%
% Adapted from the response given by Stephen in \href{https://tex.stackexchange.com/a/691544}{SAT like question format}.
%   \myattachfile*{enumext-exa-5}
% \setenumext[keyans]{label=\Alph*),font=\small,nosep,wrap-label={##1}}
% \begin{enumext}[label=\arabic*,labelwidth=20pt,labelsep=0pt,columns=2,save-ans=stylesat, wrap-label={\stylesat{#1}}]
%  \item Which choice best describes what happens in the passage?
%   \begin{keyans}
%     \item* One character argues with another character who intrudes on her home.
%     \item One character receives a surprising request from another character.
%     \item One character reminisces about choices she has made over the years.
%     \item One character criticizes another character for pursuing an unexpected course of action.
%   \end{keyans}
%
%  \item Which choice best describes what happens in the passage?
%   \begin{keyans}
%    \item One character argues with another character who intrudes on her home.
%    \item One character receives a surprising request from another character.
%    \item* One character reminisces about choices she has made over the years.
%    \item One character criticizes another character for pursuing an unexpected course of action.
%   \end{keyans}
%
%  \item Which choice best describes what happens in the passage?
%   \begin{keyans}
%    \item One character argues with another character who intrudes on her home.
%    \item* One character receives a surprising request from another character.
%    \item One character reminisces about choices she has made over the years.
%    \item One character criticizes another character for pursuing an unexpected course of action.
%   \end{keyans}
%
%  \item Which choice best describes what happens in the passage?
%   \begin{keyans}
%    \item One character argues with another character who intrudes on her home.
%    \item One character receives a surprising request from another character.
%    \item One character reminisces about choices she has made over the years.
%    \item* One character criticizes another character for pursuing an unexpected course of action.
%   \end{keyans}
% \end{enumext}
%
% \printkeyans[columns=4]{stylesat}
%
% \subsubsection*{Example 6}
%
% Adapted from the response to \href{https://tex.stackexchange.com/q/728584}{Environment for enumerate environment}.
%   \myattachfile*{enumext-exa-6}
% \setenumext[keyans]{label={\Alph*},font={\bfseries},wrap-ans*={\textcolor{red}{\bfseries##1}},mark-ans*={\textcolor{red}{✔}},wrap-opt={ },save-sep={ }}
% \showrating{true}
% \begin{enumext}[azetinasty,start=10,wrap-label={\azwrapp{#1}},save-ans=test2024]
%  \item sample
%    \begin{keyans}
%      \item sample
%      \item*[(8.5a, KSC)] answer
%      \item sample
%      \item sample
%    \end{keyans}
%  \RenewDocumentCommand \rating { }{9.5a, KSC}
%  \item sample
%      \begin{keyans}
%      \item sample
%      \item sample
%      \item sample
%      \item*[(9.5a, KSC)] answer
%    \end{keyans}
% \end{enumext}
% \showrating{false}
% \RenewDocumentCommand \rating { }{10.5a, KSC}
% \begin{enumext}[save-ans=test2024,resume*]
%  \item sample
%    \begin{keyans}
%      \item sample
%      \item*[(10.5a, KSC)] answer
%      \item sample
%      \item sample
%    \end{keyans}
%  \RenewDocumentCommand \rating { }{11.5a, KSC}
%  \item sample
%      \begin{keyans}
%      \item sample
%      \item sample
%      \item sample
%      \item*[(11.5a, KSC)] answer
%    \end{keyans}
% \end{enumext}
% \printkeyans[start=10]{test2024}
%
% \section{Tagged PDF examples}
%
% This section is just to show the compatibility of \mypkg*[no-index]{enumext} with
% \emph{tagged} \texttt{PDF} using |lualatex|. The files described below
% are just for testing and are intended as examples and, in a way, to
% simplify the time of Matthew Bertucci (\texttt{@mbertucci}) when he sees this
% excellent package and adds it to \href{https://github.com/latex3/tagging-project}{The LaTeX Tagged PDF repository}.
%
% To compile the tests with |lualatex-dev| the packages \mypkg[no-index]{multicol},
% \mypkg{unicode-math}, \mypkg{geometry}, \mypkg{graphicx}, \mypkg{luamml}
% and \mypkg[no-index]{hyperref} are required along with the line:
%
% \iffalse
%<*example>
% \fi
\begin{examplecode}[frame=single]
\DocumentMetadata{lang = en-US, pdfversion = 2.0, pdfstandard = `\mydim{ua-2}`, tagging=on,}
\end{examplecode}
% \iffalse
%</example>
% \fi
%
% \begin{important}
%  All examples have been checked using \href{https://dev.verapdf-rest.duallab.com}{\texttt{veraPDF}}
%  together with \href{https://ngpdf.com}{\texttt{ngpdf}}.
% \end{important}
%
% \begin{itemize}[topsep=1pt,partopsep=0pt,parsep=1pt,itemsep=1pt]
% \item The file \myfile{enumext-01.tex} contains the basic tests for the
%   \myenv*[no-index]{enumext} and \myenv*[no-index]{enumext*} environments and the nesting
%   between them plus the use of the \mykey{label}, \mykey{labelwidth},
%   \mykey{labelsep}, \mykey{ref}, \mykey{align} and \mykey{wrap-label} keys.
%
% \item The file \myfile{enumext-02.tex} contains the tests for the \myenv*[no-index]{enumext}
%   and \myenv*[no-index]{enumext*} environments and the support for \myenv[no-index]{minipage}
%   and \myenv[no-index]{multicols} environments using the keys \mykey{columns},
%   \mykey{columns-sep}, \mykey{mini-env}, \mykey{mini-right} and
%   |\miniright| command.
%   \myattachfile{enumext-02}
%
% \item The file \myfile{enumext-03.tex} contains the tests for the \myenv*[no-index]{enumext}
%   and \myenv*[no-index]{keyanspic} environments activated by the \mykey{save-ans}
%   key together with the \mykey{save-sep} and \mykey{save-ref} keys and
%   the |\printkeyans| command.
%   \myattachfile{enumext-03}
%
% \item The file \myfile{enumext-04.tex} contains the tests for the |\anskey|
%   command and the \myenv*[no-index]{anskey*} environment activated by the
%   \mykey{save-ans} key along with the |\getkeyans| and |\printkeyans| commands.
%   \myattachfile{enumext-04}
%
% \item The file \myfile{enumext-05.tex} contains the tests for the environments
%   \myenv*[no-index]{keyans}, \myenv*[no-index]{keyans*} and \myenv*[no-index]{keyanspic} activated by the key
%   \mykey{save-ans} together with the keys \mykey{no-store} and \mykey{show-ans}
%   and the commands |\setenumext|, |\setenumextmeta|, |\printkeyans| and |\foreachkeyans|.
%   \myattachfile{enumext-05}
%
% \item The file \myfile{enumext-06.tex} contains the tests for the environments
%   \myenv*[no-index]{enumext} and \myenv*[no-index]{enumext*} for \emph{fake} \myenv[no-index]{itemize} and \myenv[no-index]{description}.
%   \myattachfile{enumext-06}
%
% \item The file \myfile{enumext-07.tex} contains the tests for starting the
%   environments with |\setenumext{resume}|, the |\resetenumext| command
%   and the \mykey{series}, \mykey{resume}, \mykey{resume*} and \mykey{reset} keys.
%   \myattachfile{enumext-07}
% \end{itemize}
%
% \section{The way of non-enumerated lists}
%
% It is possible to use (or abuse) the \myenv*[no-index]{enumext} and \myenv*[no-index]{enumext*} environments to
% mimic \emph{non-enumerated} list environments such as \myenv[no-index]{itemize} and
% \myenv[no-index]{description}, clearly the \mymeta{keys} to \emph{\enquote{store answers}},
% the \myenv*[no-index]{keyans}, \myenv*[no-index]{keyans*} and \myenv*[no-index]{keyanspic} environments lose their
% sense and it is not the focus of \mypkg*[no-index]{enumext} package, but, why not
% to do it?.
%
% Other uses of the \myenv*[no-index]{enumext} environment that can be helpful for
% specific purposes are shown as examples. The \emph{trick}
% to generate these \emph{\enquote{fake environments}} is to set \mykey{label}|={}|
% or \mykey{label}|=|\mymarg{some} and play with the |list-indent|, |list-offset|,
% |font| and |wrap-label| keys.
%
% \subsection*{Fake \texttt{itemize} environment}\label{subsec:fake-itemize}
%
% The \mykey{label} key is set here using the default settings in
% \hologo{LaTeX} for the four levels |\textbullet|, |\textendash|, |\textasteriskcentered|
% and |\textperiodcentered| together with the
% \mykey{nosep} key to reduce the vertical spaces in the left side example and set the
% \mykey{label} key in \emph{mathematical mode} for the right side as
% |\ast|, |\diamond|, |\circ| and |\star| for the four levels
% together with the \mykey{nosep} key
%
% \noindent\begin{minipage}[t]{0.45\linewidth}
% \begin{enumext}[nosep, label=\textbullet]
%  \item First level item
%    \begin{enumext}[nosep,label=\normalfont\bfseries\textendash]
%      \item Second level item
%        \begin{enumext}[nosep,label=\textasteriskcentered]
%          \item Third level item
%            \begin{enumext}[nosep,label=\textperiodcentered]
%                \item Fourth level item
%            \end{enumext}
%        \end{enumext}
%    \end{enumext}
%  \item First level item
% \end{enumext}
% \end{minipage}\hfill
% \begin{minipage}[t]{0.45\linewidth}
% \begin{enumext}[nosep, label=$\ast$]
%  \item First level item
%    \begin{enumext}[nosep,label=$\diamond$]
%      \item Second level item
%        \begin{enumext}[nosep,label=$\circ$]
%          \item Third level item
%            \begin{enumext}[nosep,label=$\star$]
%                \item Fourth level item
%            \end{enumext}
%        \end{enumext}
%    \end{enumext}
%  \item First level item
% \end{enumext}
% \end{minipage}
%
% \subsection*{Fake \texttt{description} environment}\label{subsec:fake-description}
%
% Here \mykey{label}|={}| and |list-indent=2.5em,font=\bfseries| are set.
%
% \begin{enumext}[label={},list-indent=2.5em,font=\bfseries,noitemsep,topsep=3pt]
%  \item[SomeThing] A short one-line description.
%  \item[] This is an entry \emph{without} a label.
%  \item[Something] A short \emph{one-line} description text.
%  \item[Something long] A much \emph{longer} description text may take
%    more than one line or more than one paragraph.
%
%    Lorem ipsum dolor sit amet, consectetur adipiscing elit, sed do eiusmod
%    tempor incididunt ut labore et dolore magna aliqua.
% \end{enumext}
%
% Adding |list-indent=0pt| gives the \emph{widest style}:
%
% \begin{enumext}[label={},list-indent=0pt,font=\bfseries,noitemsep,topsep=3pt]
%  \item[SomeThing] A short one-line description.
%  \item[] This is an entry \emph{without} a label.
%  \item[Something] A short \emph{one-line} description text.
%  \item[Something long] A much \emph{longer} description text may take
%    more than one line or more than one paragraph.
%
%    Lorem ipsum dolor sit amet, consectetur adipiscing elit, sed do eiusmod
%    tempor incididunt ut labore et dolore magna aliqua.
% \end{enumext}
%
% \begin{important}*
% The small space at the beginning of the \emph{\enquote{unlabeled entry}} corresponds to
% |\labelsep| and can be removed using |\hspace{-\labelsep}| at the
% beginning of the line.
% \end{important}
%
% \begin{important}
%  When \emph{tagged} PDF is active the default \myenv[no-index]{description} style
%  is NOT available due to the redefinition of |\makelabel| for the
%  \mykey{align} key which uses |\makebox| in this case, meaning that
%  |\item|\myoarg{content} will not extend beyond |\labelwidth| which
%  causes overlaps,
% \end{important}
%
% \subsection*{Description indented by label}\label{subsec:description-indented}
%
% Here \mykey{label}|={}| is set and a convenient value is given to
% \mykey{labelsep} and \mykey{labelwidth}, for example taking the
% \emph{longest label} as reference and passing it as value using:
%
% \iffalse
%<*example>
% \fi
\begin{examplecode}[frame=none]
\newlength{\descitemwd}
\settowidth{\descitemwd}{\textbf{Something long}}
\end{examplecode}
% \iffalse
%</example>
% \fi
%
% and then use \mykey{labelsep}|=4pt,|\mykey{labelwidth}|=\descitemwd,font=\bfseries|.
%
% \newlength{\descitemwd}
% \settowidth{\descitemwd}{\textbf{Something long}}
% \begin{enumext}[label={},labelsep=4pt,labelwidth=\descitemwd,font=\bfseries,noitemsep,topsep=3pt]
%  \item[SomeThing] A short one-line description.
%  \item[] This is an entry \textit{without} a label.
%  \item[Something] A short one-line description.
%  \item[Something long] A much longer description. Lorem ipsum dolor sit amet, consectetuer adipiscing elit.
%    Ut purus elit, vestibulum ut, placerat ac, adipiscing vitae, felis.
%    Curabitur dictum gravida mauris.
% \end{enumext}
%
% The environment can be translated so that the \mymeta{labels} are on
% the left margin calculating the value passed to the \mykey{list-offset}
% key, in this case it will be equal to the sum of the values set by the
% \mykey{labelwidth} and \mykey{labelsep} keys finally resulting as
% \mykey{list-offset}|={-\descitemwd - 4pt}|.
%
% \begin{enumext}[label={},labelsep=4pt,labelwidth=\descitemwd,list-offset={-\descitemwd-4pt}, font=\bfseries,noitemsep,topsep=3pt]
%  \item[SomeThing] A short one-line description.
%  \item This is an entry \textit{without} a label.
%  \item[Something] A short one-line description.
%  \item[Something long] A much longer description. Lorem ipsum dolor sit amet, consectetuer adipiscing elit.
%    Ut purus elit, vestibulum ut, placerat ac, adipiscing vitae, felis.
%    Curabitur dictum gravida mauris.
% \end{enumext}
%
% Adding \mykey{align}|=right| gives:
%
% \begin{enumext}[label={},align=right,labelsep=4pt,labelwidth=\descitemwd,list-offset={-\descitemwd - 4pt}, font=\bfseries,noitemsep,topsep=3pt]
%  \item[SomeThing] A short one-line description.
%  \item This is an entry \textit{without} a label.
%  \item[Something] A short one-line description.
%  \item[Something long] A much longer description. Lorem ipsum dolor sit amet, consectetuer adipiscing elit.
%    Ut purus elit, vestibulum ut, placerat ac, adipiscing vitae, felis.
%    Curabitur dictum gravida mauris.
% \end{enumext}
%
% \begin{important}*
% At this point \mykey{list-offset}|={-\descitemwd - 4pt}| instead of
% |list-offset={-\labelwidth - \labelsep}|,
% this is because the parameters |\labelwidth| and |\labelsep| take the
% default values, as if \mykey{label} had not been set.
% \end{important}
%
% \subsection*{Description with multi-line labels}\label{subsec:description-multiline}
%
% The \mykey{label} key does not accept \emph{multiline material}, this
% is where the \mykey{wrap-label} and \mykey{wrap-label*} keys comes into
% play. Unlike the \mypkg[no-index]{enumitem} package, the \mykey{align} key only
% supports three options, so a command is created in the
% style |\parleft| of \mypkg[no-index]{enumitem} that allows placing \emph{multiline labels}
% using |\parbox|.
%
% \iffalse
%<*example>
% \fi
\begin{examplecode}[frame=none]
\NewDocumentCommand \labelbx { s +m }
  {%
    \SuspendTagging{\parbox}%
    \IfBooleanTF{#1}
      {\strut\smash{\parbox[t]{\labelwidth}{\raggedright{#2}}}}%
      {\strut\smash{\parbox[t]{\labelwidth}{\raggedleft{#2}}}}%
    \ResumeTagging{\parbox}%
  }
\end{examplecode}
% \iffalse
%</example>
% \fi
%
% Now it is only necessary to set |wrap-label*={\labelbx{#1}}|.
%
% \begin{enumext}[mode-box,label={},labelsep=4pt,labelwidth=\descitemwd,list-offset=-\dimeval{\descitemwd+4pt}, wrap-label*={\labelbx{#1}}, font=\bfseries,noitemsep,topsep=3pt]
%  \item[SomeThing] A short one-line description.
%  \item This is an entry \textit{without} a label.
%  \item[Something] A short one-line description.
%  \item[Something \\ long] A much longer description. Lorem ipsum dolor sit amet, consectetuer adipiscing elit.
%    Ut purus elit, vestibulum ut, placerat ac, adipiscing vitae, felis.
%    Curabitur dictum gravida mauris.
%
%    Lorem ipsum dolor sit amet, consectetuer adipiscing elit.
%    Ut purus elit, vestibulum ut, placerat ac, adipiscing vitae, felis.
%    Curabitur dictum gravida mauris.
%
%  \item[SoMeThInG \\ LoNg] A much longer description. Lorem ipsum dolor sit amet, consectetuer adipiscing elit.
%    Ut purus elit, vestibulum ut, placerat ac, adipiscing vitae, felis.
%    Curabitur dictum gravida mauris.
% \end{enumext}
%
% \section*{Final notes}
%
% The original implementation (if you can call it that) of the ideas that
% led to the creation of \mypkg*[no-index]{enumext} were some macros using the \mypkg{enumerate}\cite{enumerate}
% package for personal use created in early 2003, the code was quite
% questionable, but functional for these simple requirements.
%
% With the great answers given by Christian Hupfer in
% \href{https://tex.stackexchange.com/a/364763/7832}{Create a fake label ref
% using list} and the answer given by David Carlisle in
% \href{https://tex.stackexchange.com/a/363963/7832}{Change the use of label
% ref by data save in an array (list)} I managed to create a more solid
% code than the original version, now using the \mypkg{l3prop}\cite{inter3} and \mypkg{l3seq}\cite{inter3}
% modules together with the \mypkg[no-index]{hyperref}\cite{hyperref} and \mypkg[no-index]{enumitem}\cite{enumitem} packages, which did
% the job, but with some limitations.
%
% As time went by I took these limitations as a personal challenge which
% I called \emph{\enquote{reinventing the wheel}}, since there were packages and classes
% that did more or less what I was looking for, but did not fit my simple
% requirements. This \emph{\enquote{reinventing the wheel}} finally ended up becoming
% \mypkg*[no-index]{enumext}.
%
% \subsubsection*{Why list environments?}
%
% The answer is simple, first I love the beauty of its syntax and many of
% what I had already written used the \myenv[no-index]{enumerate} environment or
% lists created using the \mypkg[no-index]{enumitem} package. In my mind I thought:
% how complicated could it be to write a package that looked like
% \mypkg[no-index]{enumitem}? It seemed simple enough, of course I didn't have in
% mind the mess I was getting into working with \myenv[no-index]{list}
% environments, \myenv[no-index]{minipage} and adding support for the
% \mypkg[no-index]{multicol} and \mypkg[no-index]{hyperref} packages.
%
% Of course, seeing the final result of the experiment
% \emph{\enquote{reinventing the wheel}} I am quite satisfied.
%
% \subsubsection*{Why not random questions and other utilities}
%
% The \emph{\enquote{random}} type questions I love and hate them at the
% same time, although they simplify a lot the work when creating a
% multiple choice test, but you lose the beauty of typesetting a
% document with \hologo{LaTeX}, that is to say the output does not always
% look as nice as it should, even if they are only alternatives these
% must follow a certain order when presented either numerical or
% presentation, that said handling that using \emph{nested lists} is quite
% complicated so I do not classify to be implemented.
%
% \subsubsection*{Why has it taken so long?}
%
% One of the setbacks, beyond my laziness, was including compatibility
% with \emph{tagged} \texttt{PDF}. To be honest, it's something I never
% considered at any point, but I firmly believe that being able to create
% \emph{accessible documents} provides a great opportunity in the world of
% mathematics education. From my perspective as a \emph{high school} teacher,
% beyond theorems and deep mathematics, the use of exercise lists is one
% of the most common things. Being able to open the way to work in
% parallel with those who have different abilities is really important
% and I regret not having looked into this in the past. I hope that
% \mypkg*[no-index]{enumext} serves this purpose and inspires more users and
% authors to follow this path.
%
% \let\stdsection\section
% \def\section*#1{\stdsection{#1}}
% \expandafter\let\expandafter\thebibliography\csname HDorg@thebibliography\endcsname
%
% \begin{thebibliography}{10}
%
% \bibitem{multicol} \textsc{Mittelbach, Frank}. \enquote{An
% environment for \textsf{multicolumn} output}. Available from
% \textsc{ctan}, \url{https://www.ctan.org/pkg/multicol}, 2025.
%
% \bibitem{exam} \textsc{Hirschhorn, Philip}. \enquote{Using the
% \textsf{exam} document class}. Available from
% \textsc{ctan}, \url{https://www.ctan.org/pkg/exam}, 2023.
%
% \bibitem{xsim} \textsc{Niederberger, Clemens}. \enquote{\textsf{xsim} – eXercise
% Sheets IMproved}. Available from
% \textsc{ctan}, \url{https://www.ctan.org/pkg/xsim}, 2023.
%
% \bibitem{latexref} \textsc{Berry, Karl}. \enquote{\hologo{LaTeX2e}: An
% Unofficial Reference Manual}. Available from \textsc{ctan},
% \url{https://ctan.org/pkg/latex2e-help-texinfo}, 2026.
%
% \bibitem{enumitem} \textsc{Bezos, Javier}. \enquote{Customizing lists
% with the \textsf{enumitem} package}. Available from
% \textsc{ctan}, \url{https://www.ctan.org/pkg/enumitem}, 2025.
%
% \bibitem{footnotehyper} \textsc{Burnol, Jean-François}. \enquote{The \textsf{footnotehyper}
% package}. Available from
% \textsc{ctan}, \url{https://www.ctan.org/pkg/footnotehyper}, 2025.
%
% \bibitem{shortlst} \textsc{Lemvig, Mogens}. \enquote{The \textsf{shortlst}
% package}. Available from
% \textsc{ctan}, \url{https://www.ctan.org/pkg/shortlst}, 1998.
%
% \bibitem{tasks} \textsc{Niederberger, Clemens}. \enquote{\textsf{tasks} – Horizontally
% columned lists}. Available from
% \textsc{ctan}, \url{https://www.ctan.org/pkg/tasks}, 2022.
%
% \bibitem{userguide} The \hologo{LaTeX} Project. \enquote{\hologo{LaTeX} for authors
% current version}. Available from
% \textsc{ctan}, \url{https://ctan.org/pkg/latex-base}, 2026.
%
% \bibitem{hyperref} The \hologo{LaTeX} Project. \enquote{Extensive support
% for hypertext in \hologo{LaTeX}}. Available from
% \textsc{ctan}, \url{https://www.ctan.org/pkg/hyperref}, 2026.
%
% \bibitem{ltnews} The \hologo{LaTeX} Project. \enquote{\hologo{LaTeX} News,
% Issue 41, June 2025}. Available from
% \textsc{ctan}, \url{https://ctan.org/tex-archive/macros/latex/base}, 2025.
%
% \bibitem{enumerate} The \hologo{LaTeX} Project. \enquote{\textsf{enumerate} – Enumerate
% with redefinable labels}. Available from
% \textsc{ctan}, \url{https://www.ctan.org/pkg/enumerate}, 2025.
%
% \bibitem{inter3} The \hologo{LaTeX} Project. \enquote{The \hologo{LaTeX}3
% Interfaces}. Available from
% \textsc{ctan}, \url{https://www.ctan.org/pkg/l3kernel}, 2026.
%
% \bibitem{latex2e} The \hologo{LaTeX} Project. \enquote{The \hologo{LaTeX2e}
% sources}. Available from
% \textsc{ctan}, \url{https://ctan.org/tex-archive/macros/latex/base}, 2026.
%
% \bibitem{latexlab} The \hologo{LaTeX} Project. \enquote{\textsf{latex-lab} – \hologo{LaTeX}
% laboratory}. Available from
% \textsc{ctan}, \url{https://www.ctan.org/pkg/latex-lab}, 2026.
%
% \bibitem{luavisual} \textsc{Gundlach, Patrick}. \enquote{The \textsf{lua-visual-debug}
% package}. Available from
% \textsc{ctan}, \url{https://www.ctan.org/pkg/lua-visual-debug}, 2026.
%
% \bibitem{scontents} \textsc{González, Pablo}. \enquote{\textsf{scontents} - Stores
% \hologo{LaTeX} contents in memory or files}. Available from
% \textsc{ctan}, \url{https://www.ctan.org/pkg/scontents}, 2026.
%
% \bibitem{tagpdf} \textsc{Fischer, Ulrike}. \enquote{\textsf{tagpdf} – \hologo{LaTeX}
% kernel code for PDF tagging}. Available from
% \textsc{ctan}, \url{https://www.ctan.org/pkg/tagpdf}, 2026.
%
% \bibitem{ltsockets} \textsc{Mittelbach, Frank}. \enquote{\hologo{LaTeX}’s
% socket management}. Available from
% \textsc{ctan}, \url{https://ctan.org/tex-archive/macros/latex/base}, 2026.
%
% \bibitem{expl3} The \hologo{LaTeX} Project. \enquote{The \textsf{expl3}
% package}. Available from
% \textsc{ctan}, \url{https://www.ctan.org/pkg/l3kernel}, 2026.
%
% \end{thebibliography}
%
% \let\section\stdsection
%
% \section{Change history}
% \label{sec:changes}
%
% \setlist[itemize,1]{label=\textendash,leftmargin=*,labelsep=0.5em,nosep,noitemsep}
% \newlength\descrwidth
% \settowidth{\descrwidth}{\textsf{v1.0, (ctan), 2024-11-01} }
% \begin{description}[font=\small\sffamily,wide=0pt,style=multiline,leftmargin=1.1\descrwidth,nosep,noitemsep]
%  \item [\fileversion{} (ctan), \filedate]
%    \begin{itemize}
%      \item Update requirement to \hologo{LaTeX} release 2026-11-01.
%      \item Update code for \emph{tagged} \texttt{PDF}.
%      \item Fixed bad interaction between |\setenumext| and |\printkeyans|: the
%        values set for the environments no longer modify the printed content.
%      \item The error messages follow a common layout: a short text with the
%        line context and, for errors, an explanation in the extended text.
%      \item Adjustments in the documentation: the examples are attached to the
%        manual, the bibliography follows the order of appearance and the
%        \enquote{References} section is numbered in the bookmarks.
%      \item Internal: cleanup warnings and details returned by |expltools|.
%      \item Internal: replace |\exp_after:wN| by |\exp_after:w|.
%      \item Internal: the implementation is organized by command, with the
%        variables, the messages, the keys and the internal functions first and
%        the public command at the end, and all function bodies follow the same
%        layout.
%    \end{itemize}
%  \item [v2.2 (ctan), 2026-06-01]
%    \begin{itemize}
%      \item Update requirement to \hologo{LaTeX} release 2026-06-01.
%      \item Move |\setcounter| to |\begin{list}|\mymarg{arg one}|{|\mymeta{arg two}|\setcounter}| for |enumext*|.
%      \item Adjustments in the documentation.
%    \end{itemize}
%  \item [v2.1 (ctan), 2026-01-15]
%    \begin{itemize}
%      \item Fixed bad interaction between |\setenumext| and the \texttt{\textcolor{keyname}{ref}} key.
%    \end{itemize}
%  \item [v2.0 (ctan), 2025-12-02]
%    \begin{itemize}
%      \item Update requirement to \mypkg[no-index]{multicol} release 2025-10-21.
%      \item Adjustments in the documentation.
%    \end{itemize}
%  \item [v1.9 (ctan), 2025-11-01]
%    \begin{itemize}
%      \item Update requirement to \hologo{LaTeX} release 2025-11-01.
%      \item Move |\setcounter| to |\begin{list}|\mymarg{arg one}|{|\mymeta{arg two}|\setcounter}|.
%      \item Update code for \emph{tagged} \texttt{PDF}.
%    \end{itemize}
%  \item [v1.8 (ctan), 2025-10-04]
%    \begin{itemize}
%      \item Replacing |\scantokens| (primitive) with |\tl_retokenize:n|.
%      \item Cleanup warnings and details returned by |expltools|.
%    \end{itemize}
%  \item [v1.7 (ctan), 2025-07-10]
%    \begin{itemize}
%      \item Fixed |\setenumext{enumext*}{|\texttt{\textcolor{keyname}{resume}}|}|.
%      \item Fixed bad interaction between |\setenumext| and the \texttt{\textcolor{keyname}{resume}} key.
%      \item The behavior of the key \texttt{\textcolor{keyname}{resume}}|*| has been updated and documented.
%    \end{itemize}
%  \item [v1.6 (ctan), 2025-07-04]
%    \begin{itemize}
%      \item Syntax simplification for |\setenumextmeta|.
%      \item Environments can be started with the key \texttt{\textcolor{keyname}{resume}} \emph{without value}.
%      \item Add |\resetenumext|, \texttt{\textcolor{keyname}{reset}} and \texttt{\textcolor{keyname}{reset}}|*| keys.
%      \item The \texttt{\textcolor{keyname}{resume}}, \texttt{\textcolor{keyname}{resume}}|*| and \texttt{\textcolor{keyname}{series}}
%            keys can now be set per level.
%      \item Fixed bad interaction between |\printkeyans| and the \texttt{\textcolor{keyname}{resume}}, \texttt{\textcolor{keyname}{resume}}|*| keys.
%    \end{itemize}
%  \item [v1.5 (ctan), 2025-06-11]
%    \begin{itemize}
%      \item Replacing |\regex_match:| (deprecated) with |\regex_if_match:|.
%      \item Add keys |beginpenalty|, |midpenalty| and |endpenalty|.
%    \end{itemize}
%  \item [v1.4 (ctan), 2025-06-09]
%    \begin{itemize}
%      \item Improved implementation of the \texttt{\textcolor{keyname}{start}} key for \emph{tagged} \texttt{PDF}.
%      \item Improved implementation of the \texttt{\textcolor{keyname}{ref}} key.
%      \item Fixed the behavior of the |save-sep| key.
%      \item Fixed the behavior of the \texttt{\textcolor{keyname}{resume}}|*| key.
%    \end{itemize}
%  \item [v1.3 (ctan), 2025-06-01]
%    \begin{itemize}
%      \item Removed dependency on the \mypkg{scontents} package.
%      \item The |anskey*| environment has been rewritten using the new
%            \texttt{c-type} argument.
%    \end{itemize}
%  \item [v1.2 (ctan), 2025-03-28]
%    \begin{itemize}
%      \item Replace signature (prevent expansion for optional argument).
%      \item Solve Inconsistent local/global assignment.
%    \end{itemize}
%  \item [v1.1 (ctan), 2024-11-14]
%    \begin{itemize}
%      \item Fixed implementation for |font| and |base-fix| keys.
%      \item Added new keys for symbol marks.
%      \item Update and improvements in the internal code.
%      \item Adjustments in the documentation.
%    \end{itemize}
%  \item [v1.0 (ctan), 2024-11-01]
%    \begin{itemize}
%      \item First public release.
%    \end{itemize}
% \end{description}
%
% \newpage
%
% \indexprologue{
% The italic numbers denote the pages where the corresponding entry is
% described.}
%
% \printindex[userdoc]
%
% \newpage
%
% \StartImplementation
%
% \StopEventually{^^A
% \newgeometry{top=0.5in,bottom=0.3in,left=1.0in,right=0.5in,footskip=0.2in,headheight=1cm,headsep=0.27cm}
% \addtocontents{toc}{\protect\setcounter{tocdepth}{2}}
% \cleardoublepage
% \phantomsection
% \indexprologue{
% The italic numbers denote the pages where the corresponding entry is
% described, the numbers underlined and all others indicate the line on
% which they are implemented in the package code.
% }
% \printindex
% }
%
% \newgeometry{top=0.5in,bottom=0.3in,left=1.95in,right=0.5in,footskip=0.2in,headheight=1cm,headsep=0.27cm}
%
% \section{Implementation}
% \label{sec:implementation}
% \addtocontents{toc}{\protect\setcounter{tocdepth}{0}}
%
% The most recent publicly released version of \mypkg*[no-index]{enumext} is
% available at \textsc{ctan}: \url{https://www.ctan.org/pkg/enumext}.
% While general feedback via email is welcomed, specific bugs or
% feature requests should be reported through the issue
% tracker: \textcolor{gray}{\scriptsize\faIcon[regular]{github}} \url{https://github.com/pablgonz/enumext/issues}.
%
% \begin{important}*
% The documentation presented here is far from professional, it contains
% a lot of obvious information that to the eye of a \hologo{TeX}pert are
% superfluous, but, after so many years developing this project is the
% only way to remember what does what.
% \end{important}
%
% \subsection{General conventions}\label{subsec:general-conventions}
%
% Variables containing \lstinline+i+, \lstinline+ii+, \lstinline+iii+ and
% \lstinline+iv+ are associated by level with the \myenv*[no-index]{enumext}
% environment, variables containing \lstinline+v+ are associated with the
% \myenv*[no-index]{keyans} environment, variables containing \lstinline+vi+ are
% associated with the \myenv*[no-index]{keyanspic} environment, variables
% containing \lstinline+vii+ are associated with the \myenv*{enumext*}
% environment and variables containing \lstinline+viii+ are associated
% with the \myenv*[no-index]{keyans*} environment.
%
% To simplify writing and documentation some variables and functions that
% are common to the different levels of the environments are described
% using a capital \enquote{\texttt{\textcolor{keyname}{X}}}.
%
% The temporary function \myhlc|\@@_tmp:n| is used in
% different parts of the package code for variable creation or execution
% of other functions that are grouped into this one.
%
% All variables and functions defined in this package are private and
% are NOT intended to work or be used by another package or module.
%
%
% \subsection{Initial set up}\label{subsec:initial-setup}
%
% Start the \pkg{DocStrip} guards.
%    \begin{macrocode}
%<*package>
%    \end{macrocode}
% Identify the internal prefix (\LaTeX3 \pkg{DocStrip} convention) for
% \myclass{l3doc} class.
%    \begin{macrocode}
%<@@=enumext>
%    \end{macrocode}
%
% \subsection{Declaration of the package}\label{subsec:declaration}
%
% Checks for a minimum release version of \hologo{LaTeX}; otherwise it
% returns an \enquote{error}.
%    \begin{macrocode}
\NeedsTeXFormat{LaTeX2e}[2026-11-01]
\IfFormatAtLeastF { 2026-11-01 }
  {
    \PackageError { enumext }
      { LaTeX~kernel~too~old~need~2026-11-01~or~later }
      {
        enumext~requires~a~LaTeX2e~release~2026-11-01~or~
        later.~Update~your~TeX~distribution.
      }
  }
%    \end{macrocode}
% Now declare the \mypkg*[no-index]{enumext} package.
%    \begin{macrocode}
\ProvidesExplPackage {enumext} {2026-10-23} {2.3} {Enumerate exercise sheets}
%    \end{macrocode}
% Finally, check if the \mypkg{multicol} package is loaded; if not, it is loaded.
%    \begin{macrocode}
\hook_gput_code:nnn {begindocument} {enumext}
  {
    \IfPackageLoadedTF { multicol }
      {
        \msg_info:nnn { enumext } { package-load } { multicol }
      }
      {
        \msg_info:nnn { enumext } { package-not-load } { multicol }
        \RequirePackage{multicol}[2025-10-21]
      }
  }
%    \end{macrocode}
%
% \subsection{Definition of variables}\label{subsec:variables}
%
% Variables that do not appear in this section are created by means of
% \myhlc|\keys_define:nn| or some function described below.
%
% \begin{variable}[int]{
%   \l_@@_level_int,
%   \l_@@_level_h_int,
%   \l_@@_anskey_level_int,
%   \l_@@_keyans_level_int,
%   \l_@@_keyans_level_h_int,
%   \l_@@_keyans_pic_level_int,
% }
%   Integer variables will control the nesting levels of the environments,
%   \myenv*{anskey*} environment and |\anskey| command.
%    \begin{macrocode}
\int_new:N  \l_@@_level_int
\int_new:N  \l_@@_level_h_int
\int_new:N  \l_@@_anskey_level_int
\int_new:N  \l_@@_keyans_level_int
\int_new:N  \l_@@_keyans_level_h_int
\int_new:N  \l_@@_keyans_pic_level_int
%    \end{macrocode}
% \end{variable}
%
% \begin{variable}[int]{
%   \l_@@_starred_bool,
%   \g_@@_starred_bool,
%   \l_@@_starred_first_bool,
%   \l_@@_standard_bool,
%   \g_@@_standard_bool,
%   \l_@@_standard_first_bool,
%   \l_@@_keyans_env_bool,
%   \g_@@_start_line_tl,
%   \g_@@_envir_name_tl,
%   \l_@@_envir_name_tl,
% }
%   Internal variables used by functions \myhlc|\@@_is_not_nested:|,
%   \myhlc|\@@_is_on_first_level:| and \myhlc|\@@_keyans_name_and_start:| (\S\ref{subsubsec:utilities}).
%    \begin{macrocode}
\bool_new:N \l_@@_starred_bool
\bool_new:N \g_@@_starred_bool
\bool_new:N \l_@@_starred_first_bool
\bool_new:N \l_@@_standard_bool
\bool_new:N \g_@@_standard_bool
\bool_new:N \l_@@_standard_first_bool
\bool_new:N \l_@@_keyans_env_bool
\tl_new:N   \g_@@_start_line_tl
\tl_new:N   \g_@@_envir_name_tl
\tl_new:N   \l_@@_envir_name_tl
%    \end{macrocode}
% \end{variable}
%
% \begin{variable}[int]{
%   \l_@@_counter_i_tl,
%   \l_@@_counter_ii_tl,
%   \l_@@_counter_iii_tl,
%   \l_@@_counter_iv_tl,
%   \l_@@_counter_v_tl,
%   \l_@@_counter_vi_tl,
%   \l_@@_counter_vii_tl,
%   \l_@@_counter_viii_tl,
% }
%
%   Variables to store the \emph{\enquote{name of the counters}}
%   \icounter{enumXi}, \icounter{enumXii}, \icounter{enumXiii} and
%   \icounter{enumXiv} for \myenv*[no-index]{enumext}  environment,
%   \icounter{enumXv} for \myenv*[no-index]{keyans} environment and
%   \icounter{enumXvi} for the \myenv*[no-index]{keyanspic} environment.
%   The counters \icounter{enumXvii} and \icounter{enumXviii} are used
%   by \myenv*[no-index]{enumext*} and \myenv*[no-index]{keyans*} environments.
%
%   The initial values of these variables are set by the function
%   \myhlc|\@@_define_counter:Nn| (\S\ref{subsec:counters}) and then modified by the function
%   \myhlc|\@@_label_style:Nnn| used by \mykey{label} key (\S\ref{key:label-ref}).
%    \begin{macrocode}
\cs_set_protected:Nn \@@_tmp:n
  {
    \tl_new:c { l_@@_counter_ #1 _tl }
  }
\clist_map_inline:nn { i, ii, iii, iv, v, vi, vii, viii } { \@@_tmp:n {#1} }
%    \end{macrocode}
% \end{variable}
%
% \begin{variable}[int]{
%   \l_@@_ref_key_arg_tl,
%   \l_@@_ref_the_count_tl,
%   \l_@@_renew_counter_X_tl,
%   \l_@@_the_counter_X_tl,
% }
%   Internal variables used by \mykey{ref} key (\S\ref{key:label-ref}).
%    \begin{macrocode}
\tl_new:N \l_@@_ref_key_arg_tl
\tl_new:N \l_@@_ref_the_count_tl
\cs_set_protected:Nn \@@_tmp:n
  {
    \tl_new:c  { l_@@_renew_counter_ #1 _tl }
    \tl_new:c  { l_@@_the_counter_ #1 _tl   }
    \tl_set:ce { l_@@_the_counter_ #1 _tl   } { \exp_not:c { theenumX #1 } }
  }
\clist_map_inline:nn { i, ii, iii, iv, v, vi, vii, viii } { \@@_tmp:n {#1} }
%    \end{macrocode}
% \end{variable}
%
% \begin{variable}[int]{
%   \l_@@_series_name_tl,
%   \l_@@_resume_count_bool,
%   \l_@@_resume_count_X_bool,
%   \l_@@_resume_series_X_bool,
%   \l_@@_resume_star_key_X_bool,
%   \g_@@_resume_last_keys_X_tl,
% }
%   Internal variables used by \mykey{resume}, \mykey{resume*} and
%   \mykey{series} keys (\S\ref{key:series}).
%    \begin{macrocode}
\tl_new:N   \l_@@_series_name_tl
\bool_new:N \l_@@_resume_count_bool
\cs_set_protected:Nn \@@_tmp:n
  {
    \bool_new:c { l_@@_resume_count_ #1 _bool    }
    \bool_new:c { l_@@_resume_series_ #1 _bool   }
    \bool_new:c { l_@@_resume_star_key_ #1 _bool }
    \tl_new:c   { g_@@_resume_last_keys_ #1 _tl  }
  }
\clist_map_inline:nn { i, ii, iii, iv, vii } { \@@_tmp:n {#1} }
%    \end{macrocode}
% \end{variable}
%
% \begin{variable}[int]{
%   \l_@@_current_widest_dim,
%   \g_@@_counter_styles_tl,
%   \g_@@_widest_label_tl,
%   \l_@@_label_width_by_box,
% }
%   The variable \myhlc|\l_@@_current_widest_dim| stores the current
%   label width, the variable \myhlc|\g_@@_counter_styles_tl| stores
%   the default \mymeta{label style} and the variable
%   \myhlc|\g_@@_widest_label_tl| the label width. These variables are
%   used by \mykey{widest} (\S\ref{key:start-widest}) and \mykey{label} (\S\ref{key:label}) keys.
%    \begin{macrocode}
\dim_new:N \l_@@_current_widest_dim
\tl_new:N  \g_@@_counter_styles_tl
\tl_new:N  \g_@@_widest_label_tl
\box_new:N \l_@@_label_width_by_box
%    \end{macrocode}
% \end{variable}
%
% \begin{variable}[int]{
%   \l_@@_leftmargin_tmp_X_bool,
%   \l_@@_leftmargin_tmp_X_dim,
%   \l_@@_leftmargin_X_dim,
%   \l_@@_itemindent_X_dim,
% }
%   The boolean variable \myhlc|\l_@@_leftmargin_tmp_X_bool| and the
%   dimensional variable \myhlc|\l_@@_leftmargin_tmp_X_dim|
%   are used by the \mykey{list-indent} key (\S\ref{key:list-indent}).
%   The variables \myhlc|\l_@@_leftmargin_X_dim| and
%   \myhlc|\l_@@_itemindent_X_dim| are used and set by the function
%   \myhlc|\@@_calc_hspace:NNNNNNN| (\S\ref{subsubsec:leftmargin}).
%    \begin{macrocode}
\cs_set_protected:Nn \@@_tmp:n
  {
    \bool_new:c { l_@@_leftmargin_tmp_ #1 _bool }
    \dim_new:c  { l_@@_leftmargin_tmp_ #1 _dim  }
    \dim_new:c  { l_@@_leftmargin_ #1 _dim      }
    \dim_new:c  { l_@@_itemindent_ #1 _dim      }
  }
\clist_map_inline:nn { i, ii, iii, iv, v, vi, vii, viii } { \@@_tmp:n {#1} }
%    \end{macrocode}
% \end{variable}
%
% \begin{variable}[int]{
%   \l_@@_multicols_above_X_skip,
%   \l_@@_multicols_below_X_skip,
%   \l_@@_align_label_pos_X_str,
% }
%   Internal variables used by \mykey{columns} key (\S\ref{subsec:vspace-multicols})
%   and \mykey{align} key (\S\ref{key:label}).
%    \begin{macrocode}
\cs_set_protected:Nn \@@_tmp:n
  {
    \skip_new:c  { l_@@_multicols_above_ #1 _skip }
    \skip_new:c  { l_@@_multicols_below_ #1 _skip }
    \str_new:c   { l_@@_align_label_pos_ #1 _str  }
  }
\clist_map_inline:nn { i, ii, iii, iv, v } { \@@_tmp:n {#1} }
%    \end{macrocode}
% \end{variable}
%
% \begin{variable}[int]{
%   \g_@@_minipage_stat_int,
%   \l_@@_minipage_temp_skip,
%   \l_@@_minipage_left_skip,
%   \l_@@_minipage_right_skip,
%   \l_@@_minipage_after_skip,
%   \g_@@_minipage_right_skip,
%   \g_@@_minipage_after_skip,
%   \l_@@_minipage_left_X_dim,
%   \l_@@_minipage_active_X_bool,
% }
%   Internal variables used by \ics*{miniright} command (\S\ref{cmd:miniright})
%   and the keys \mykey{mini-right}, \mykey{mini-right*}, \mykey{mini-env}
%   and \mykey{mini-sep} (\S\ref{key:mini-multi},~\S\ref{subsec:vspace-minipage}).
%    \begin{macrocode}
\int_new:N  \g_@@_minipage_stat_int
\skip_new:N \l_@@_minipage_temp_skip
\skip_new:N \l_@@_minipage_left_skip
\skip_new:N \l_@@_minipage_right_skip
\skip_new:N \l_@@_minipage_after_skip
\skip_new:N \g_@@_minipage_right_skip
\skip_new:N \g_@@_minipage_after_skip
\cs_set_protected:Nn \@@_tmp:n
  {
    \dim_new:c  { l_@@_minipage_left_ #1 _dim    }
    \bool_new:c { l_@@_minipage_active_ #1 _bool }
  }
\clist_map_inline:nn { i, ii, iii, iv, v, vii, viii } { \@@_tmp:n {#1} }
%    \end{macrocode}
% \end{variable}
%
% \begin{variable}[int]{
%   \l_@@_wrap_label_X_bool,
%   \l_@@_wrap_label_opt_X_bool,
%   \l_@@_start_X_int,
%   \l_@@_fake_item_indent_X_tl,
%   \l_@@_label_fill_left_X_tl,
%   \l_@@_label_fill_right_X_tl,
%   \l_@@_vspace_a_star_X_bool,
%   \l_@@_vspace_b_star_X_bool,
% }
%   The bool vars \myhlc|\l_@@_wrap_label_X_bool| and
%   \myhlc|\l_@@_wrap_label_opt_X_bool| are used by \mykey{wrap-label}
%   and \mykey{wrap-label*} keys (\S\ref{key:label}), the integer
%   \myhlc|\l_@@_start_X_int| are used by the \mykey{start} and \mykey{start*} keys (\S\ref{key:start-widest}),
%   the token list \myhlc|\l_@@_fake_item_indent_X_tl | is used by \mykey{itemindent} key (\S\ref{subsubsec:fake-indent}),
%   the variables \myhlc|\l_@@_label_fill_left_X_tl| and \myhlc|\l_@@_label_fill_left_X_tl| are used by the
%   \mykey{align} key (\S\ref{key:label}). The boolean vars \myhlc|\l_@@_vspace_a_star_X_bool|,
%   \myhlc|\l_@@_vspace_b_star_X_bool| are used by \mykey{above},
%   \mykey{above*}, \mykey{below} and \mykey{below*} keys (\S\ref{key:before}).
%    \begin{macrocode}
\cs_set_protected:Nn \@@_tmp:n
  {
    \bool_new:c { l_@@_wrap_label_ #1 _bool     }
    \bool_new:c { l_@@_wrap_label_opt_ #1 _bool }
    \int_new:c  { l_@@_start_ #1 _int           }
    \tl_new:c   { l_@@_fake_item_indent_ #1 _tl }
    \tl_new:c   { l_@@_label_fill_left_ #1 _tl  }
    \tl_new:c   { l_@@_label_fill_right_ #1 _tl }
    \bool_new:c { l_@@_vspace_a_star_ #1 _bool  }
    \bool_new:c { l_@@_vspace_b_star_ #1 _bool  }
  }
\clist_map_inline:nn { i, ii, iii, iv, v, vii, viii } { \@@_tmp:n {#1} }
%    \end{macrocode}
% \end{variable}
%
% \begin{variable}[int]{
%   \l_@@_store_active_bool,
%   \l_@@_store_name_tl,
%   \g_@@_store_name_tl,
%   \l_@@_store_current_label_tl,
%   \l_@@_store_current_opt_arg_tl,
% }
%
%   The variable \myhlc|\l_@@_store_active_bool| setting by
%   \mykey{save-ans} key (\S\ref{key:save-ans}) activates all the
%   mechanism related to \ics*{anskey}, \myenv*[no-index]{anskey*}, \myenv*[no-index]{keyans},
%   \myenv*[no-index]{keyans*} and \myenv*[no-index]{keyanspic} environments.
%
%   The variable \myhlc|\l_@@_store_name_tl| saves the \mymarg{store name}
%   set by the \mykey{save-ans} key of the \emph{sequence} and \emph{prop
%   list} where the content is stored. The variable \myhlc|\g_@@_store_name_tl|
%   is just a global copy of \mymarg{store name} used by different functions.
%
%   The variables \myhlc|\l_@@_store_current_label_tl| and
%   \myhlc|\l_@@_store_current_opt_arg_tl| save the \emph{current label} and
%   \emph{optional argument} of \ics*{item*} (\S\ref{cmd:item-keyans}) and \ics*{anspic*}
%   (\S\ref{cmd:anspic}) for the \myenv*[no-index]{keyans}, \myenv*[no-index]{keyans*} and
%   \myenv*[no-index]{keyanspic} environments.
%    \begin{macrocode}
\bool_new:N \l_@@_store_active_bool
\tl_new:N   \l_@@_store_name_tl
\tl_new:N   \g_@@_store_name_tl
\tl_new:N   \l_@@_store_current_label_tl
\tl_new:N   \l_@@_store_current_opt_arg_tl
%    \end{macrocode}
% \end{variable}
%
% \begin{variable}[int]{
%   \l_@@_store_anskey_arg_tl,
%   \l_@@_store_anskey_env_tl,
%   \l_@@_write_anskey_env_bool,
%   \l_@@_write_anskey_env_file_name_tl,
%   \l_@@_write_anskey_env_file_iow,
% }
%   The variable \myhlc|\l_@@_store_anskey_arg_tl| save the \emph{argument}
%   of |\anskey| (\S\ref{cmd:anskey}) and the variables
%   \myhlc|\l_@@_store_anskey_env_tl| save the \mymeta{body}
%   of the environment \myenv*[no-index]{anskey*} (\S\ref{env:anskey}).
%
%   The variables \myhlc|\l_@@_write_anskey_env_bool|, \myhlc|\l_@@_write_anskey_env_file_name_tl|
%   and \myhlc|\l_@@_write_anskey_env_file_iow| they are used by the \mykey{write-env}
%   and \mykey{overwrite} keys in the \myenv*[no-index]{anskey*} environment implementation.
%    \begin{macrocode}
\tl_new:N   \l_@@_store_anskey_arg_tl
\tl_new:N   \l_@@_store_anskey_env_tl
\bool_new:N \l_@@_write_anskey_env_bool
\tl_new:N   \l_@@_write_anskey_env_file_name_tl
\iow_new:N  \l_@@_write_anskey_env_file_iow
%    \end{macrocode}
% \end{variable}
%
% \begin{variable}[int]{\c_@@_anskey_env_hidden_space_str}
%   The \myhlc|\c_@@_anskey_env_hidden_space_str| is a constant \emph{string}
%   to used to hide the \mymeta{forced space} added by \hologo{TeX} when
%   recording content in a macro. This \emph{string} contains the \emph{reserved phrase}
%   `|%^^Aenumextheol%|' which is added to the end of the argument stored
%   in \emph{sequence} and \emph{prop list} when the key |force-eol| is false.
%    \begin{macrocode}
\str_const:Ne \c_@@_anskey_env_hidden_space_str
  { \c_percent_str \c_circumflex_str \c_circumflex_str A enumextheol
    \c_percent_str }
%    \end{macrocode}
% \end{variable}
%
% \begin{variable}[int]{
%   \l_@@_setkey_tmpa_tl,
%   \l_@@_setkey_tmpb_tl,
%   \l_@@_setkey_tmpa_int,
%   \l_@@_setkey_tmpa_seq,
%   \l_@@_setkey_tmpb_seq,
% }
%   Internal variables used by the command |\setenumext| (\S\ref{cmd:setenumext}).
%    \begin{macrocode}
\tl_new:N  \l_@@_setkey_tmpa_tl
\tl_new:N  \l_@@_setkey_tmpb_tl
\int_new:N \l_@@_setkey_tmpa_int
\seq_new:N \l_@@_setkey_tmpa_seq
\seq_new:N \l_@@_setkey_tmpb_seq
%    \end{macrocode}
% \end{variable}
%
% \begin{variable}[int]{
%   \l_@@_meta_path_str,
%   \l_@@_foreach_print_seq,
%   \l_@@_foreach_name_prop_tl,
%   \l_@@_foreach_default_keys_tl,
% }
%   Internal variables used by the |\setenumextmeta| command (\S\ref{cmd:setenumextmeta})
%   and |\foreachkeyans| command (\S\ref{cmd:foreachkeyans}).
%    \begin{macrocode}
\str_new:N \l_@@_meta_path_str
\seq_new:N \l_@@_foreach_print_seq
\tl_new:N  \l_@@_foreach_name_prop_tl
\tl_new:N  \l_@@_foreach_default_keys_tl
%    \end{macrocode}
% \end{variable}
%
%
% \begin{variable}[int]{
%   \l_@@_print_keyans_starred_tl,
%   \l_@@_print_keyans_star_bool,
%   \l_@@_print_keyans_cmd_bool,
%   \l_@@_mark_position_str,
%   \l_@@_mark_position_v_str,
%   \l_@@_mark_position_viii_str,
%   \l_@@_mark_sep_tmpa_dim,
%   \l_@@_mark_sep_tmpb_dim,
%   \l_@@_show_pos_tmp_int,
%   \g_@@_item_symbol_aux_tl,
%   \l_@@_print_keyans_X_tl,
%   \l_@@_store_save_key_X_tl,
%   \l_@@_store_save_key_X_bool,
%   \l_@@_store_upper_level_X_bool,
% }
%   Internal variables used by command \ics*{printkeyans}
%   (\S\ref{cmd:printkeyans}), \mykey{show-pos}, \mykey{show-ans}, \mykey{mark-pos}, \mykey{mark-sep} keys (\S\ref{subsec:storage-system}),
%   \mykey{item-sym*} key (\S\ref{key:item-starred}), \mykey{save-key} key
%   (\S\ref{key:save-key}) and \emph{\enquote{storing structure}}.
%    \begin{macrocode}
\tl_new:N   \l_@@_print_keyans_starred_tl
\bool_new:N \l_@@_print_keyans_star_bool
\bool_new:N \l_@@_print_keyans_cmd_bool
\str_new:N  \l_@@_mark_position_str
\str_new:N  \l_@@_mark_position_v_str
\str_new:N  \l_@@_mark_position_viii_str
\dim_new:N  \l_@@_mark_sep_tmpa_dim
\dim_new:N  \l_@@_mark_sep_tmpb_dim
\int_new:N  \l_@@_show_pos_tmp_int
\tl_new:N   \g_@@_item_symbol_aux_tl
\cs_set_protected:Nn \@@_tmp:n
  {
    \tl_new:c   { l_@@_print_keyans_ #1 _tl        }
    \tl_new:c   { l_@@_store_save_key_ #1 _tl      }
    \bool_new:c { l_@@_store_save_key_ #1 _bool    }
    \bool_new:c { l_@@_store_upper_level_ #1 _bool }
  }
\clist_map_inline:nn { i, ii, iii, iv, vii } { \@@_tmp:n {#1} }
%    \end{macrocode}
% \end{variable}
%
% \begin{variable}[int]{
%   \l_@@_anspic_args_seq,
%   \l_@@_anspic_mini_width_dim,
%   \l_@@_anspic_above_int,
%   \l_@@_anspic_below_int,
%   \l_@@_anspic_label_above_bool,
%   \l_@@_anspic_mini_pos_str,
%   \l_@@_anspic_label_box,
%   \l_@@_anspic_body_box,
%   \l_@@_anspic_label_htdp_dim,
%   \l_@@_anspic_body_htdp_dim,
% }
%    Internal variables used by \myenv*[no-index]{keyanspic} environment and
%   \ics*{anspic} command (\S\ref{env:keyanspic}).
%    \begin{macrocode}
\seq_new:N  \l_@@_anspic_args_seq
\dim_new:N  \l_@@_anspic_mini_width_dim
\int_new:N  \l_@@_anspic_above_int
\int_new:N  \l_@@_anspic_below_int
\bool_new:N \l_@@_anspic_label_above_bool
\str_new:N  \l_@@_anspic_mini_pos_str
\box_new:N  \l_@@_anspic_label_box
\box_new:N  \l_@@_anspic_body_box
\dim_new:N  \l_@@_anspic_label_htdp_dim
\dim_new:N  \l_@@_anspic_body_htdp_dim
%    \end{macrocode}
% \end{variable}
%
% \begin{variable}[int]{
%   \l_@@_check_answers_bool,
%   \g_@@_check_ans_key_bool,
%   \l_@@_check_start_line_env_tl,
%   \l_@@_item_wrap_key_bool,
%   \g_@@_check_starred_cmd_int,
%   \g_@@_item_anskey_int,
%   \g_@@_item_number_int,
%   \l_@@_item_number_bool,
%   \g_@@_item_answer_diff_int,
% }
%   Internal variables used by \emph{\enquote{internal check answer}} mechanism
%   (\S\ref{subsubsec:check-answer}) used by the \mykey{check-ans}, \mykey{no-store}, \mykey{wrap-ans*}
%   keys and check for starred commands |\item*| in \myenv*[no-index]{keyans} and
%   \myenv*[no-index]{keyans*} environments and |\anspic*|
%   in \myenv*[no-index]{keyanspic} environment.
%    \begin{macrocode}
\bool_new:N \l_@@_check_answers_bool
\bool_new:N \g_@@_check_ans_key_bool
\tl_new:N   \l_@@_check_start_line_env_tl
\bool_new:N \l_@@_item_wrap_key_bool
\int_new:N  \g_@@_check_starred_cmd_int
\int_new:N  \g_@@_item_anskey_int
\int_new:N  \g_@@_item_number_int
\bool_new:N \l_@@_item_number_bool
\int_new:N  \g_@@_item_answer_diff_int
%    \end{macrocode}
% \end{variable}
%
% \begin{variable}[int]{\l_@@_hyperref_bool, \l_@@_footnotes_key_bool}
%   The boolean variable \myhlc|\l_@@_hyperref_bool| will determine if the
%   \mypkg{hyperref} package is present or load in memory
%   (\S\ref{subsec:hyperref}). The boolean variable \myhlc|\l_@@_footnotes_key_bool|
%   determine if \mypkg[no-index]{hyperref} is load with key |hyperfootnotes=true|.
%    \begin{macrocode}
\bool_new:N \l_@@_hyperref_bool
\bool_new:N \l_@@_footnotes_key_bool
%    \end{macrocode}
% \end{variable}
%
% \begin{variable}[int]{
%   \l_@@_newlabel_arg_one_tl,
%   \l_@@_newlabel_arg_two_tl,
%   \l_@@_write_aux_file_tl,
%   \l_@@_label_copy_X_tl,
% }
%   Internal variables used by \mykey{save-ref} key (\S\ref{subsec:storage-system}).
%   The variables \myhlc|\l_@@_label_copy_X_tl| correspond to temporary
%   copies of the \mymeta{labels} defined by level on which operations will be performed.
%
%   The variables \myhlc|\l_@@_newlabel_arg_one_tl| and \myhlc|\l_@@_newlabel_arg_two_tl|
%   will be used to form the arguments passed to the function
%   \myhlc|\@@_newlabel:nn| (\S\ref{fun:newlabel}) and the variable \myhlc|\l_@@_write_aux_file_tl|
%   will be in charge of executing the writing code in the |.aux| file.
%    \begin{macrocode}
\tl_new:N \l_@@_newlabel_arg_one_tl
\tl_new:N \l_@@_newlabel_arg_two_tl
\tl_new:N \l_@@_write_aux_file_tl
\cs_set_protected:Nn \@@_tmp:n
  {
    \tl_new:c { l_@@_label_copy_ #1 _tl }
  }
\clist_map_inline:nn { i, ii, iii, iv, v, vi, vii, viii } { \@@_tmp:n {#1} }
%    \end{macrocode}
% \end{variable}
%
% \begin{variable}[int]{
%   \g_@@_footnote_standard_int,
%   \g_@@_footnote_starred_int,
%   \g_@@_footnote_standard_arg_seq,
%   \g_@@_footnote_starred_arg_seq,
%   \g_@@_footnote_standard_int_seq,
%   \g_@@_footnote_starred_int_seq,
% }
%   Internal variables used for redefinition of |\footnote| (\S\ref{subsec:footnote}).
%    \begin{macrocode}
\int_new:N \g_@@_footnote_standard_int
\int_new:N \g_@@_footnote_starred_int
\seq_new:N \g_@@_footnote_standard_arg_seq
\seq_new:N \g_@@_footnote_starred_arg_seq
\seq_new:N \g_@@_footnote_standard_int_seq
\seq_new:N \g_@@_footnote_starred_int_seq
%    \end{macrocode}
% \end{variable}
%
% \begin{variable}[int]{
%   \l_@@_item_starred_X_bool,
%   l_@@_item_column_pos_X_int,
%   \g_@@_item_count_all_X_int,
%   \l_@@_joined_item_X_int,
%   \l_@@_joined_item_aux_X_int,
%   \l_@@_tmpa_X_int,
%   \l_@@_item_text_X_box,
%   \l_@@_joined_width_X_dim,
%   \l_@@_item_width_X_dim,
%   \g_@@_item_symbol_aux_X_tl,
%   \l_@@_align_label_X_str,
%   \g_@@_minipage_active_X_bool,
%   \l_@@_miniright_code_X_box,
%   \g_@@_minipage_center_X_bool,
%   \g_@@_minipage_right_X_dim,
% }
%    Internal variables used by \myenv*[no-index]{enumext*} and \myenv*[no-index]{keyans*} environments.
%    \begin{macrocode}
\cs_set_protected:Nn \@@_tmp:n
  {
    \bool_new:c { l_@@_item_starred_ #1 _bool    }
    \int_new:c  { l_@@_item_column_pos_ #1 _int  }
    \int_new:c  { g_@@_item_count_all_ #1 _int   }
    \int_new:c  { l_@@_joined_item_ #1 _int      }
    \int_new:c  { l_@@_joined_item_aux_ #1 _int  }
    \int_new:c  { l_@@_tmpa_ #1 _int             }
    \box_new:c  { l_@@_item_text_ #1 _box        }
    \dim_new:c  { l_@@_joined_width_ #1 _dim     }
    \dim_new:c  { l_@@_item_width_ #1 _dim       }
    \tl_new:c   { g_@@_item_symbol_aux_ #1 _tl   }
    \str_new:c  { l_@@_align_label_ #1 _str      }
    \bool_new:c { g_@@_minipage_active_ #1 _bool }
    \box_new:c  { l_@@_miniright_code_ #1 _box   }
    \bool_new:c { g_@@_minipage_center_ #1 _bool }
    \dim_new:c  { g_@@_minipage_right_ #1 _dim   }
  }
\clist_map_inline:nn { vii, viii } { \@@_tmp:n {#1} }
%    \end{macrocode}
% \end{variable}
%
% \begin{variable}[int]{\c_@@_all_envs_clist}
%    An internal |clist-var| variable to run with \myhlc|\@@_tmp:n|.
%    \begin{macrocode}
\clist_const:Nn \c_@@_all_envs_clist
  {
    {level-1}{i}, {level-2}{ii}, {level-3}{iii}, {level-4}{iv},
    {keyans}{v}, {enumext*}{vii}, {keyans*}{viii}
  }
%    \end{macrocode}
% \end{variable}
%
% \subsection{Some utility functions}\label{subsec:utility-functions}
%
% \subsubsection{Definition of error messages}
%
% \begin{macro}[int]{prop-seq-int-hook, item-answer-hook}
% Messages written to the |.log| file by the internal hooks that show the values of the
% global variables.
%    \begin{macrocode}
\msg_new:nnn { enumext } { prop-seq-int-hook }
  {
    *~Package~enumext:~Elements~in~
    \c_backslash_str g_@@_ #1 _prop~=~#2.\\
    *~Package~enumext:~Elements~in~
    \c_backslash_str g_@@_ #1 _seq~=~#3.\\
    *~Package~enumext:~Value~of~
    \c_backslash_str g_@@_resume_ #1 _int~=~#4.
  }
\msg_new:nnn { enumext } { item-answer-hook }
  {
    *~Package~enumext:~Value~of~
    \c_backslash_str g_@@_item_number_int~=~#1.\\
    *~Package~enumext:~Value~of~
    \c_backslash_str g_@@_item_anskey_int~=~#2.\\
    *~Package~enumext:~Difference~item_number_int~-~item_anskey_int~=~#3.
  }
%    \end{macrocode}
% \end{macro}
%
% \subsubsection{General utility functions}
%
% \begin{macro}[int]{\keys_precompile:neN, \seq_use:NV}
%    Non-standard kernel variants used by the |\printkeyans| command
%    (\S\ref{cmd:printkeyans}) and |\foreachkeyans| command (\S\ref{cmd:foreachkeyans}).
%    \begin{macrocode}
\cs_generate_variant:Nn \keys_precompile:nnN { neN }
\cs_generate_variant:Nn \seq_use:Nn { NV }
%    \end{macrocode}
% \end{macro}
%
% \begin{macro}[int]{\@@_at_begin_document:n}
%    A internal \emph{\enquote{hook}} function used for copying plain
%    \myenv[no-index]{list} and \myenv[no-index]{minipage} environments definition and
%    \mypkg[no-index]{hyperref} detection.
%    \begin{macrocode}
\cs_new_protected:Nn \@@_at_begin_document:n
  {
    \hook_gput_code:nnn {begindocument} {enumext} {#1}
  }
%    \end{macrocode}
% \end{macro}
%
% \begin{macro}[int]{\@@_after_env:nn}
%    An internal \emph{\enquote{hook}} function to execute code \mykey[env]{mini-right}
%    and \mykey[env]{mini-right*} keys outside the \myenv*[no-index]{enumext*} and
%    \myenv*[no-index]{keyans*} environments and print \mykey[env]{check-ans} outside the
%    \myenv*[no-index]{enumext} and \myenv*[no-index]{enumext*} environments.
%    \begin{macrocode}
\cs_new_protected:Nn \@@_after_env:nn
  {
    \hook_gput_code:nnn {env/#1/after} {enumext} {#2}
  }
%    \end{macrocode}
% \end{macro}
%
% \begin{macro}[int]{\@@_level:}
%    Function for check current level in \myenv*[no-index]{enumext}.
%    \begin{macrocode}
\cs_new:Nn \@@_level:
  {
    \int_to_roman:n { \l_@@_level_int }
  }
%    \end{macrocode}
% \end{macro}
%
% \begin{macro}[int]{\@@_if_is_int:nT, \@@_if_is_int:nF, \@@_if_is_int:nTF}
%    A conditional function to determine whether the variable being passed is
%    an integer used by \mykey[env,foreachkeyans]{start} and \mykey[env]{widest} keys.
%    This function is taken directly from the answer given by Henri Menke
%    in \href{https://tex.stackexchange.com/a/427559}{How to test if an expl3 function argument is an integer expression?}.
%    \begin{macrocode}
\prg_new_protected_conditional:Npnn \@@_if_is_int:n #1 { T, F, TF }
  {
    \regex_if_match:nnTF { ^[\+\-]?[\d]+$ } {#1} % $
      { \prg_return_true: }
      { \prg_return_false: }
  }
%    \end{macrocode}
% \end{macro}
%
% \begin{macro}[int]{\@@_show_length:nnn}
%    Internal function used by \mykey[env]{show-length} key to show
%    \emph{\enquote{all lengths}} calculated and used in
%    \myenv*[no-index]{enumext}, \myenv*[no-index]{enumext*}, \myenv*[no-index]{keyans} and
%    \myenv*[no-index]{keyans*} environments.
%    \begin{macrocode}
\cs_new:Nn \@@_show_length:nnn
  {
    *~#2
    \prg_replicate:nn { 14 - \str_count:n {#2} } {~}
      =~\use:c { #1 _use:c } { l_@@_ #2 _ #3 _ #1 } \\
  }
%    \end{macrocode}
% \end{macro}
%
% \begin{macro}[int]{\@@_unskip_unkern:}
%   The function \myhlc|\@@_unskip_unkern:| will remove the last \mymeta{skip}
%   or \mymeta{kern} at execution time using the values \mydim{11} and
%   \mydim{12} of |\lastnodetype| to apply |\unskip| or |\unkern| depending
%   on the case.
%    \begin{macrocode}
\cs_new_protected:Nn \@@_unskip_unkern:
  {
    \int_case:nnT { \lastnodetype }
      {
        { 11 }{ \unskip }
        { 12 }{ \unkern }
      }
  }
%    \end{macrocode}
% \end{macro}
%
% \subsubsection{Utilities for environments and levels}\label{subsubsec:utilities}
%
% \begin{macro}[int]{\@@_is_not_nested:, \@@_is_on_first_level:}
%   The function \myhlc|\@@_is_not_nested:| sets the variables \myhlc|\g_@@_standard_bool|
%   and \myhlc|\g_@@_starred_bool| to |true| only if the
%   environments \myenv*[no-index]{enumext} and \myenv*[no-index]{enumext*} are NOT nested
%   in each other and saves the environment name in \myhlc|\l_@@_envir_name_tl|.
%    \begin{macrocode}
\cs_new_protected:Nn \@@_is_not_nested:
  {
    \str_case:en { \@currenvir }
      {
        {enumext}
          {
            \tl_set:Ne \l_@@_envir_name_tl { enumext }
            \bool_lazy_and:nnT
              { \bool_not_p:n { \g_@@_standard_bool } }
              { \int_compare_p:nNn { \l_@@_level_h_int } = { 0 } }
              {
                \bool_gset_true:N \g_@@_standard_bool
              }
          }
        {enumext*}
          {
            \tl_set:Ne \l_@@_envir_name_tl { enumext* }
            \bool_lazy_and:nnT
              { \bool_not_p:n { \g_@@_starred_bool } }
              { \int_compare_p:nNn { \l_@@_level_int } = { 0 } }
              {
                \bool_gset_true:N \g_@@_starred_bool
              }
          }
      }
  }
%    \end{macrocode}
%   The function \myhlc|\@@_is_on_first_level:| will sets the variables
%   \myhlc|\l_@@_standard_first_bool| (\S\ref{key:save-ans}),
%   \myhlc|\l_@@_starred_first_bool| (\S\ref{key:save-ans}) to |true|
%   only if the environment is not nested and it is at the \emph{\enquote{first level}}.
%   They also save the \emph{start line number} of each environment
%   in the variable \myhlc|\g_@@_start_line_tl| and the \emph{name} of each
%   environment in the variable \myhlc|\g_@@_envir_name_tl| to use in messages
%   related to the \mykey[env]{check-ans} key and |.log| file.
%    \begin{macrocode}
\cs_new_protected:Nn \@@_is_on_first_level:
  {
    \bool_lazy_all:nT
      {
        { \bool_if_p:N \g_@@_standard_bool }
        { \int_compare_p:nNn { \l_@@_level_int } = { 1 } }
        { \int_compare_p:nNn { \l_@@_level_h_int } = { 0 } }
      }
      {
        \bool_set_true:N \l_@@_standard_first_bool
        \tl_gset:Nn \g_@@_envir_name_tl { enumext }
        \tl_gset:Ne \g_@@_start_line_tl
          {
            on~line~\exp_not:V \inputlineno
          }
      }
    \bool_lazy_all:nT
      {
        { \bool_if_p:N \g_@@_starred_bool }
        { \int_compare_p:nNn { \l_@@_level_h_int } = { 1 } }
        { \int_compare_p:nNn { \l_@@_level_int } = { 0 } }
      }
      {
        \bool_set_true:N \l_@@_starred_first_bool
        \tl_gset:Nn \g_@@_envir_name_tl { enumext* }
        \tl_gset:Ne \g_@@_start_line_tl
          {
            on~line~\exp_not:V \inputlineno
          }
      }
  }
%    \end{macrocode}
% \end{macro}
%
% \begin{macro}[int]{\@@_keyans_name_and_start:}
%   The function \myhlc|\@@_keyans_name_and_start:| will save the start line
%   number and name of the environments \myenv*[no-index]{keyans}, \myenv*[no-index]{keyans*} and
%   \myenv*[no-index]{keyanspic} in the variables \myhlc|\l_@@_check_start_line_env_tl|
%   and \myhlc|\l_@@_envir_name_tl| to use in the \myhlc|\@@_check_starred_cmd:n| function.
%    \begin{macrocode}
\cs_new_protected:Nn \@@_keyans_name_and_start:
  {
    \str_case:en { \@currenvir }
      {
        {keyans}
          {
            \tl_set:Ne \l_@@_envir_name_tl { keyans }
            \tl_set:Ne \l_@@_check_start_line_env_tl
              {
                in~'keyans'~start~on~line~\exp_not:V \inputlineno
              }
          }
        {keyans*}
          {
            \tl_set:Ne \l_@@_envir_name_tl { keyans* }
            \tl_set:Ne \l_@@_check_start_line_env_tl
              {
                in~'keyans*'~start~on~line~\exp_not:V \inputlineno
              }
          }
        {keyanspic}
          {
            \tl_set:Ne \l_@@_envir_name_tl { keyanspic }
            \tl_set:Ne \l_@@_check_start_line_env_tl
              {
                in~'keyanspic'~start~on~line~\exp_not:V \inputlineno
              }
          }
      }
  }
%    \end{macrocode}
% \end{macro}
%
% \subsubsection{Utilities for log and terminal}
%
% \begin{macro}[int]{\@@_reset_global_vars:, \@@_reset_global_int:,
%                    \@@_reset_global_bool:, \@@_reset_global_tl:}
%   The function \myhlc|\@@_reset_global_vars:| will be passed to the function
%   \myhlc|\@@_execute_after_env:| and will return the global variables to
%   their default values after being used.
%    \begin{macrocode}
\cs_new_protected:Nn \@@_reset_global_vars:
  {
    \@@_reset_global_int:
    \@@_reset_global_bool:
    \@@_reset_global_tl:
  }
\cs_new_protected:Nn \@@_reset_global_int:
  {
    \int_gzero:N \g_@@_item_number_int
    \int_gzero:N \g_@@_item_anskey_int
    \int_gzero:N \g_@@_item_answer_diff_int
  }
\cs_new_protected:Nn \@@_reset_global_bool:
  {
    \bool_gset_false:N \g_@@_check_ans_key_bool
    \bool_gset_false:N \g_@@_standard_bool
    \bool_gset_false:N \g_@@_starred_bool
  }
\cs_new_protected:Nn \@@_reset_global_tl:
  {
    \tl_gclear:N \g_@@_store_name_tl
    \tl_gclear:N \g_@@_start_line_tl
    \tl_gclear:N \g_@@_envir_name_tl
  }
%    \end{macrocode}
% \end{macro}
%
% \begin{macro}[int]{\@@_log_global_vars:, \@@_log_answer_vars:}
%   The function \myhlc|\@@_log_global_vars:| will be passed to the function
%   \myhlc|\@@_execute_after_env:| and write to the |.log| file the number of
%   elements saved in the \emph{prop list} and \emph{sequence} created by the
%   \mykey[env]{save-ans} key along with the value of the integer variable
%   created for the \mykey[env]{resume} key.
%    \begin{macrocode}
\cs_new_protected:Nn \@@_log_global_vars:
  {
    \msg_log:nneeee { enumext } { prop-seq-int-hook }
      { \g_@@_store_name_tl }
      { \prop_count:c { g_@@_ \g_@@_store_name_tl _prop } }
      { \seq_count:c { g_@@_ \g_@@_store_name_tl _seq } }
      { \int_use:c { g_@@_resume_ \g_@@_store_name_tl _int } }
  }
%    \end{macrocode}
%   The function \myhlc|\@@_log_answer_vars:| will be passed to the function
%   \myhlc|\@@_execute_after_env:| and write to the |.log| file the number
%   of items and answers along with the difference between them.
%    \begin{macrocode}
\cs_new_protected:Nn \@@_log_answer_vars:
  {
    \msg_log:nneee { enumext } { item-answer-hook }
      { \int_use:N \g_@@_item_number_int }
      { \int_use:N \g_@@_item_anskey_int }
      { \int_eval:n { \g_@@_item_number_int - \g_@@_item_anskey_int} }
  }
%    \end{macrocode}
% \end{macro}
%
% \subsection{Copying \env{list} and \env{minipage} environments}\label{subsec:copy-envs}
%
% The \myenv[no-index]{list} environment provided by \hologo{LaTeX} has the
% following plain form\cite{latex2e}:
%
% \iffalse
%<*example>
% \fi
\begin{examplecode}[frame=none]
\list`\mymarg{arg one}\mymarg{arg two}`
  \item`\myoarg{opt}`
\endlist
\end{examplecode}
% \iffalse
%</example>
% \fi
%
% And \myenv[no-index]{minipage} environment provided by \hologo{LaTeX} has the
% following (simplified) plain form:
%
% \iffalse
%<*example>
% \fi
\begin{examplecode}[frame=none]
\minipage`\myoarg{pos}\myoarg{height}\myoarg{inner-pos}\mymarg{width}`
  `\mymeta{internal implement}`
\endminipage
\end{examplecode}
% \iffalse
%</example>
% \fi
%
% As a precaution they are copied using \myhlc|\@@_at_begin_document:n| in
% case any package redefines the \myenv[no-index]{list} environment or a related
% command.
%
% \begin{important}
%  For compatibility with \emph{tagged} \texttt{PDF}, |\NewCommandCopy| should be used and not \myhlc|\cs_new_eq:NN| for |\item|. When
%  \emph{tagged} \texttt{PDF} is active |\item| is redefined using
%  \mypkg{ltcmd} (see \mypkg{latex-lab-block}\cite{latexlab}).
% \end{important}
%
% \begin{macro}[int]{\@@_start_list:nn, \@@_stop_list:, \@@_item_std:w,
%                    \@@_minipage:w, \@@_endminipage:}
%   The functions \myhlc|\@@_start_list:nn| and \myhlc|\@@_stop_list:| correspond
%   to copies of |\list| and |\endlist| from plain definition of \myenv[no-index]{list}
%   environment, the function \myhlc|\@@_item_std:w| is a copy of the |\item| command.
%    \begin{macrocode}
\@@_at_begin_document:n
  {
    \cs_new_eq:NN   \@@_start_list:nn \list
    \cs_new_eq:NN   \@@_stop_list: \endlist
    \NewCommandCopy \@@_item_std:w \item
  }
%    \end{macrocode}
%   The functions \myhlc|\@@_minipage:w| and \myhlc|\@@_endminipage:| correspond
%   to copies of |\minipage| and |\endminipage| from plain definition of
%   \myenv[no-index]{minipage} environment.
%    \begin{macrocode}
\@@_at_begin_document:n
  {
    \cs_new_eq:NN \@@_minipage:w \minipage
    \cs_new_eq:NN \@@_endminipage: \endminipage
  }
%    \end{macrocode}
% \end{macro}
%
% \subsection{Compatibility with hyperref and footnotehyper}\label{subsec:hyperref}
%
% \begin{macro}[int]{\@@_after_hyperref:, \@@_hypertarget:nn, \@@_phantomsection:}
%    First, the necessary rules are defined using \emph{\enquote{hooks}} to
%    determine if the \mypkg[no-index]{hyperref} package is loaded.
%    \begin{macrocode}
\hook_gput_code:nnn { begindocument } { enumext } { \@@_after_hyperref: }
\hook_gset_rule:nnnn { begindocument } { enumext } { after } { hyperref }
%    \end{macrocode}
%   The function \myhlc|\@@_after_hyperref:| sets the state of the boolean
%   variable \myhlc|\l_@@_hyperref_bool| to |true| if the package is
%   loaded. The public macro |\IfHyperBoolean| determines if the
%   |hyperfootnotes=true| key is present, and if so, the boolean variable
%   \myhlc|\l_@@_footnotes_key_bool| is set to |true|.
%    \begin{macrocode}
\cs_new_protected:Nn \@@_after_hyperref:
  {
    \IfPackageLoadedT { hyperref }
      {
        \msg_info:nnn { enumext } { package-load } { hyperref }
        \bool_set_true:N \l_@@_hyperref_bool
        \IfHyperBoolean{hyperfootnotes}
          {
            \bool_set_true:N \l_@@_footnotes_key_bool
          }
          {  }
      }
%    \end{macrocode}
%   If the variable \myhlc|\l_@@_footnotes_key_bool| is true, it checks
%   whether the package \mypkg{footnotehyper} is loaded; if it is not,
%   \myhlc|\l_@@_footnotes_key_bool| is set to false and |\footnote| is
%   redefined.
%    \begin{macrocode}
\bool_if:NT \l_@@_footnotes_key_bool
  {
    \IfPackageLoadedTF { footnotehyper }
      {
        \msg_info:nnn { enumext } { package-load } { footnotehyper }
      }
      {
        \bool_set_false:N \l_@@_footnotes_key_bool
      }
  }
%    \end{macrocode}
%   The functions \myhlc|\@@_hypertarget:nn| and \myhlc|\@@_phantomsection:|
%   correspond to the internal copies of |\hypertarget| and
%   |\phantomsection|. If the boolean variable \myhlc|\l_@@_hyperref_bool|
%   is false the functions \myhlc|\@@_hypertarget:nn| and \myhlc|\@@_phantomsection:|
%   will be disabled.
%    \begin{macrocode}
    \bool_if:NTF \l_@@_hyperref_bool
      {
        \cs_new_eq:NN \@@_hypertarget:nn \hypertarget
        \cs_new_eq:NN \@@_phantomsection: \phantomsection
      }
      {
        \cs_new_eq:NN \@@_hypertarget:nn \use_none:nn
        \cs_new_eq:NN \@@_phantomsection: \prg_do_nothing:
      }
  }
%    \end{macrocode}
% \end{macro}
%
% \begin{macro}[int]{\@@_newlabel:nn}\label{fun:newlabel}
%   The function \myhlc|\@@_newlabel:nn| write the information to the |.aux|
%   file when using the \mykey[env]{save-ref} key. The arguments taken by the
%   function are:
%
% \begin{arguments}
%    \item \myhlc|\l_@@_newlabel_arg_one_tl|
%    \item \myhlc|\l_@@_newlabel_arg_two_tl|
% \end{arguments}
%
% \begin{important}*
%   The trick here is to manage the number of arguments passed to
%   |\newlabel{#1}{#2}| according to the presence of the \mypkg[no-index]{hyperref}
%   package.
% \end{important}
%
%    \begin{macrocode}
\cs_new_protected:Nn \@@_newlabel:nn
  {
    \protected@write \@auxout { }
      {
        \token_to_str:N \newlabel {#1}
          {
            {#2}
            \bool_if:NT \l_@@_hyperref_bool
              { { \thepage } {#2} {#1} }
            { }
          }
      }
    \@@_hypertarget:nn {#1} { }
    \@@_phantomsection:
  }
%    \end{macrocode}
% \end{macro}
%
% \subsection{Internal redefining \cs{footnote} command}\label{subsec:footnote}
%
% To keep the correct numbering of |\footnote| and to make it work
% correctly in the \myenv*[no-index]{enumext*} and \myenv*[no-index]{keyans*} environments
% and \mykey{mini-env} key it is necessary to redefine the |\footnote| command.
% This implementation is adapted from the answer given by Clea F. Rees (\texttt{@cfr})
% in \href{https://tex.stackexchange.com/a/700092}{footnotes in boxes compatible with hyperref}.
%
% \begin{macro}[int]{\@@_footnotetext:nn, \@@_renew_footnote:,
%                    \@@_print_footnote:, \@@_renew_footnote_mini:,
%                    \@@_print_footnote_mini:}
%   Redefinition of the |\footnote| command using |\footnotetext| and
%   |\footnotemark| for the \mykey[env]{mini-env} key in the \myenv*[no-index]{enumext} and
%   \myenv*[no-index]{keyans} environments.
%    \begin{macrocode}
\cs_new_protected:Nn \@@_footnotetext:nn
  {
    \footnotetext[#1]{#2}
  }
\cs_new_protected:Nn \@@_renew_footnote:
  {
    \RenewDocumentCommand \footnote { o +m }
      {
        \tl_if_novalue:nTF {##1}
          {
            \stepcounter{footnote}
            \int_gset_eq:Nc \g_@@_footnote_standard_int { c@footnote }
          }
          {
            \int_gset:Nn \g_@@_footnote_standard_int {##1}
          }
        \footnotemark [ \g_@@_footnote_standard_int ]
        \seq_gput_right:Nn \g_@@_footnote_standard_arg_seq {##2}
        \seq_gput_right:NV
          \g_@@_footnote_standard_int_seq \g_@@_footnote_standard_int
      }
  }
\cs_new_protected:Nn \@@_print_footnote:
  {
    \seq_if_empty:NF \g_@@_footnote_standard_int_seq
      {
        \seq_map_pairwise_function:NNN
          \g_@@_footnote_standard_int_seq
          \g_@@_footnote_standard_arg_seq
          \@@_footnotetext:nn
      }
    \seq_gclear:N \g_@@_footnote_standard_arg_seq
    \seq_gclear:N \g_@@_footnote_standard_int_seq
  }
%    \end{macrocode}
%   The \myenv*[no-index]{enumext*} and \myenv*[no-index]{keyans*} environments are implemented
%   using \myenv[no-index]{minipage}, so |\footnote| must also be redefined to keep
%   these numbering as if it were part of the document.
%    \begin{macrocode}
\cs_new_protected:Nn \@@_renew_footnote_mini:
  {
    \RenewDocumentCommand \footnote { o +m }
      {
        \tl_if_novalue:nTF {##1}
          {
            \stepcounter{footnote}
            \int_gset_eq:Nc \g_@@_footnote_starred_int { c@footnote }
          }
          {
            \int_gset:Nn \g_@@_footnote_starred_int {##1}
          }
        \footnotemark [ \g_@@_footnote_starred_int ]
        \seq_gput_right:Nn \g_@@_footnote_starred_arg_seq {##2}
        \seq_gput_right:NV
          \g_@@_footnote_starred_int_seq \g_@@_footnote_starred_int
      }
  }
\cs_new_protected:Nn \@@_print_footnote_mini:
  {
    \seq_if_empty:NF \g_@@_footnote_starred_int_seq
      {
        \seq_map_pairwise_function:NNN
          \g_@@_footnote_starred_int_seq
          \g_@@_footnote_starred_arg_seq
          \@@_footnotetext:nn
      }
    \seq_gclear:N \g_@@_footnote_starred_arg_seq
    \seq_gclear:N \g_@@_footnote_starred_int_seq
  }
%    \end{macrocode}
% \end{macro}
%
% \begin{macro}[int]{\@@_renew_footnote_standard:, \@@_print_footnote_standard:,
%                    \@@_renew_footnote_starred:, \@@_print_footnote_starred:}
%   The redefinition of |\footnote| is encapsulated to pass it to the internal
%   \myhlc|\@@_mini_page| environment used by the \mykey[env]{mini-env}
%   key in the \myenv*[no-index]{enumext} and \myenv*[no-index]{keyans} environments. It is
%   executed when \emph{tagged} \texttt{PDF} is active or when
%   the \mypkg[no-index]{footnotehyper} package is not loaded.
%    \begin{macrocode}
\cs_new_protected:Nn \@@_renew_footnote_standard:
  {
    \bool_if:NT \g_@@_standard_bool
      {
        \IfDocumentMetadataTF
          {
            \@@_renew_footnote:
          }
          {
            \bool_if:NF \l_@@_footnotes_key_bool
              {
                \@@_renew_footnote:
              }
          }
      }
  }
\cs_new_protected:Nn \@@_print_footnote_standard:
  {
    \bool_if:NT \g_@@_standard_bool
      {
        \IfDocumentMetadataTF
          {
            \@@_print_footnote:
          }
          {
            \bool_if:NF \l_@@_footnotes_key_bool
              {
                \@@_print_footnote:
              }
          }
      }
  }
%    \end{macrocode}
%   The redefinition of |\footnote| is encapsulated to pass it to the
%   \myenv*[no-index]{enumext*} and \myenv*[no-index]{keyans*} environments. It is executed
%   when \emph{tagged} \texttt{PDF} is active or when the
%   \mypkg[no-index]{footnotehyper} package is not loaded.
%    \begin{macrocode}
\cs_new_protected:Nn \@@_renew_footnote_starred:
  {
    \IfDocumentMetadataTF
      {
        \@@_renew_footnote_mini:
      }
      {
        \bool_if:NF \l_@@_footnotes_key_bool
          {
            \@@_renew_footnote_mini:
          }
      }
  }
\cs_new_protected:Nn \@@_print_footnote_starred:
  {
    \IfDocumentMetadataTF
      {
        \@@_print_footnote_mini:
      }
      {
        \bool_if:NF \l_@@_footnotes_key_bool
          {
            \@@_print_footnote_mini:
          }
      }
  }
%    \end{macrocode}
%   In \myenv*[no-index]{enumext*} and \myenv*[no-index]{keyans*} environments it is necessary to use
%   \emph{\enquote{hooks}} to print |\footnote| with support for
%   \emph{tagged} \texttt{PDF}.
%    \begin{macrocode}
\@@_after_env:nn { enumext* }
  {
    \@@_print_footnote_starred:
  }
\@@_after_env:nn { keyans* }
  {
    \@@_print_footnote_starred:
  }
%    \end{macrocode}
% \end{macro}
%
% \subsection{The internal \env{minipage} environment}\label{subsec:internal-minipage}
%
% \begin{macro}[int]{\@@_internal_mini_page:, \@@_mini_page}
%   The function \myhlc|\@@_internal_mini_page:| creates an internal
%   \myhlc|\@@_mini_page| environment (\emph{custom version} of
%   \myenv[no-index]{minipage}) setting the |\if@minipage| switch to |false|
%   to allow spaces at the \emph{\enquote{above}} of the environment,
%   plus \myhlc|\skip_vertical:N \c_zero_skip| is added to maintain alignment
%   on \emph{\enquote{top}} in the first part and  \myhlc|\skip_vertical:N \c_zero_skip|
%   in the second part to allow spaces \emph{\enquote{below}}. This environment
%   will be used internally by the \mykey[env]{mini-env} key, it is NOT documented
%   in the user interface and is for internal use only. Within this environment
%   |\footnote| is redefined to make the notes look the same as if they were
%   elsewhere in the document. This implementation is adapted from the answer
%   given by Max Chernoff (\texttt{@MaxChernoff}) in \href{https://tex.stackexchange.com/a/720966/7832}{Customize \texttt{minipage} to support spaces below it}.
%
% \begin{important}*
%   This function is passed to the function \myhlc|\@@_safe_exec:| in the
%   \myenv*[no-index]{enumext} environment definition (\S\ref{env:enumext})
%   and \myhlc|\@@_safe_exec_vii:| in the \myenv*[no-index]{enumext*} environment
%   definition (\S\ref{env:starred}).
% \end{important}
%    \begin{macrocode}
\cs_new_protected:Nn \@@_internal_mini_page:
  {
    \int_compare:nNnT { \l_@@_level_int } = { 0 }
      {
        \DeclareDocumentEnvironment{@@_mini_page}{ m }
          {
            \@@_renew_footnote_standard:
            \@@_minipage:w [ t ] {##1}
              \legacy_if_gset_false:n { @minipage }
              \skip_vertical:N \c_zero_skip
          }
          {
              \skip_vertical:N \c_zero_skip
            \@@_endminipage:
            \@@_print_footnote_standard:
          }
      }
  }
%    \end{macrocode}
% \end{macro}
%
% \subsection{Definition of public dimension}\label{subsec:public-dimensions}
%
% The package \mypkg*[no-index]{enumext} only provides a single public dimension
% |\itemwidth| and is intended for user convenience only and is not
% for internal use as such. This dimension is set in all environments
% and is only used by the \mykey{wrap-ans} key at its default value.
%    \begin{macrocode}
\dim_zero_new:N \itemwidth
%    \end{macrocode}
%
% \subsection{Definition of counters}\label{subsec:counters}
%
% \subsubsection{Definition of error messages}
%
% \begin{macro}[int]{counters}
% Message used in the creation of counters by \mypkg*[no-index]{enumext} package.
%    \begin{macrocode}
\msg_new:nnn { enumext } { counters }
  {
    The~counter~'#1'~is~already~defined~by~some~package~or~macro,~
    it~cannot~be~continued.
  }
%    \end{macrocode}
% \end{macro}
%
% \subsubsection{Internal functions for the counters}
%
% \begin{macro}[int]{\@@_define_counter:Nn, enumXi, enumXii, enumXiii,
%                    enumXiv, enumXv, enumXvi, enumXvii, enumXviii}
%   To create the necessary \emph{\enquote{counters}}, the function first
%   checks that they are not already defined by the user or a package
%   such as \mypkg{enumitem}; otherwise it returns an error and aborts
%   the package loading. The arguments taken by the function are:
%
% \begin{arguments}
%    \item A token list \myhlc|\l_@@_counter_X_tl| for \emph{\enquote{store}} the counter's name.
%    \item The counter's name.
% \end{arguments}
%
%    \begin{macrocode}
\cs_new_protected:Nn \@@_define_counter:Nn
  {
    \cs_if_exist:cTF { c@ #2 }
      { \msg_fatal:nnn { enumext } { counters }{#2} }
      {
        \tl_set:Nn #1 {#2}
        \newcounter {#2}
      }
  }
%    \end{macrocode}
%   The counters created here are \icounter[no-index]{enumXi}, \icounter[no-index]{enumXii},
%   \icounter[no-index]{enumXiii} and \icounter[no-index]{enumXiv} for \myenv*[no-index]{enumext}
%   environment, \icounter[no-index]{enumXv} for \myenv*[no-index]{keyans} environment,
%   \icounter[no-index]{enumXvi} for \myenv*[no-index]{keyanspic} environment, \icounter[no-index]{enumXvii} for
%   \myenv*[no-index]{enumext*} and \icounter[no-index]{enumXviii} for the \myenv*[no-index]{keyans*}
%   environments.
%    \begin{macrocode}
\@@_define_counter:Nn \l_@@_counter_i_tl    { enumXi    }
\@@_define_counter:Nn \l_@@_counter_ii_tl   { enumXii   }
\@@_define_counter:Nn \l_@@_counter_iii_tl  { enumXiii  }
\@@_define_counter:Nn \l_@@_counter_iv_tl   { enumXiv   }
\@@_define_counter:Nn \l_@@_counter_v_tl    { enumXv    }
\@@_define_counter:Nn \l_@@_counter_vi_tl   { enumXvi   }
\@@_define_counter:Nn \l_@@_counter_vii_tl  { enumXvii  }
\@@_define_counter:Nn \l_@@_counter_viii_tl { enumXviii }
%    \end{macrocode}
% \end{macro}
%
% \begin{macro}[int]{\c@__@@_resume_i_int, \c@__@@_resume_ii_int,
%                    \c@__@@_resume_iii_int, \c@__@@_resume_iv_int,
%                    \c@__@@_resume_vii_int}
%  In version 1.6 the command |\resetenumext| (\S\ref{cmd:resetenumext})
%  was added which internally uses |\counterwithin*| so for its correct
%  operation, \emph{\enquote{real counters}} are created instead of
%  the \emph{\enquote{integer variables}} for the keys \mykey[env]{resume}
%  and \mykey[env]{resume*}.
%    \begin{macrocode}
\cs_set_protected:Nn \@@_tmp:n
  {
    \cs_if_exist:cTF { c@ @@_resume_ #1 _int }
      { \msg_fatal:nne { enumext } { counters }{ @@_resume_ #1 _int } }
      {
        \newcounter { @@_resume_ #1 _int }
      }
  }
\clist_map_inline:nn {i,ii,iii,iv,vii} { \@@_tmp:n {#1} }
%    \end{macrocode}
% \end{macro}
%
% \subsection{Definition of labels}\label{subsec:definition-labels}
%
% This part of the code is inspired by the \mypkg[no-index]{enumitem} package. The
% idea is to be able to access the counters using \ilabel{arabic*},
% \ilabel{Alph*}, \ilabel{alph*}, \ilabel{Roman*} and \ilabel{roman*}
% to use them in the \mykey{label} key.
%
% \begin{important}*
%  Direct support for this is provided since \hologo{LaTeX} release
%  2025-06-01\cite{ltnews}, but the original implementation is kept
%  so as not to hinder the internal \emph{\enquote{label and ref}} system.
% \end{important}
%
%
% \begin{macro}[int]{\@@_register_default_label_wd:Nn}
%   These \mymeta{counters} will be used as default \mymeta{labels} if
%   the \mykey[env]{label} key is not used for the different levels of the
%   \myenv*[no-index]{enumext}, \myenv*[no-index]{enumext*}, \myenv*[no-index]{keyans} and \myenv*[no-index]{keyans*}
%   environments, so it is necessary to get a default value for \mykey[env]{labelwidth}
%   from these \mymeta{labels} at the same time.
%    \begin{macrocode}
\cs_new_protected:Nn \@@_register_default_label_wd:Nn
  {
    \tl_const:cn { c_@@_widest_ \cs_to_str:N #1 _tl } {#2}
    \tl_gput_right:Nn \g_@@_counter_styles_tl {#1}
  }
\@@_register_default_label_wd:Nn \arabic { 0 }
\@@_register_default_label_wd:Nn \Alph   { M }
\@@_register_default_label_wd:Nn \alph   { m }
\@@_register_default_label_wd:Nn \Roman  { VIII }
\@@_register_default_label_wd:Nn \roman  { viii }
%    \end{macrocode}
% \end{macro}
%
% \begin{macro}[int]{\@@_label_width_by_box:Nn, \@@_label_width_by_box:cv}
%    The function \myhlc|\@@_label_width_by_box:Nn| set the default |\labelwidth|
%    using a box width if no \mykey[env]{labelwidth} key is passed.
%    \begin{macrocode}
\cs_new_protected:Nn \@@_label_width_by_box:Nn
  {
    \hbox_set:Nn \l_@@_label_width_by_box {#2}
    \dim_set:Nn #1 { \box_wd:N \l_@@_label_width_by_box }
  }
\cs_generate_variant:Nn \@@_label_width_by_box:Nn { cv }
%    \end{macrocode}
% \end{macro}
%
% \begin{macro}[int]{\@@_label_style:Nnn, \@@_label_style:cvn}
%   The function \myhlc|\@@_label_style:Nnn| is used by the \mykey[env]{label}
%   key to create the variables containing the \mymeta{label style}
%   and allows using \ilabel[no-index]{arabic*}, \ilabel[no-index]{Alph*}, \ilabel[no-index]{alph*},
%   \ilabel[no-index]{Roman*} and \ilabel[no-index]{roman*} as arguments.
%
%   It loops through the defined counter styles in \myhlc|\g_@@_counter_styles_tl|
%   (|\arabic|, |\alph|, |\Alph|, |\roman| and |\Roman|) for example,
%   looking for \ilabel[no-index]{roman*} and replacing that by |\roman|\mymarg{counter},
%   and doing the same for the \myhlc|\g_@@_widest_label_tl| to keep both in sync.
%    \begin{macrocode}
\cs_new_protected:Nn \@@_label_style:Nnn
  {
    \tl_clear_new:N #1
    \tl_put_right:Ne #1 { \tl_trim_spaces:n {#3} }
    \tl_gset_eq:NN \g_@@_widest_label_tl #1
    \tl_map_inline:Nn \g_@@_counter_styles_tl
      {
        \tl_replace_all:Nne #1 { ##1* } { \exp_not:N ##1 {#2} }
        \tl_greplace_all:Nne \g_@@_widest_label_tl { ##1* }
          { \tl_use:c { c_@@_widest_ \cs_to_str:N ##1 _tl } }
      }
    \@@_label_width_by_box:Nn \l_@@_current_widest_dim
      { \tl_use:N \g_@@_widest_label_tl }
    \tl_set_eq:cN { the #2 } #1
  }
\cs_generate_variant:Nn \@@_label_style:Nnn { cvn }
%    \end{macrocode}
% \end{macro}
%
% \subsection{Setting keys associated with label}\label{key:label}
%
% When \emph{tagged} \texttt{PDF} is active |\makelabel| is redefined using
% |\makebox| to work correctly (\S\ref{cmd:item-enumext}). From
% the user side it is convenient to have a key that allows using this
% redefinition with |\makebox| without having \myhlc|\IfDocumentMetadataTF| active.
%
% \begin{macro}[int]{mode-box}
%   Defines the key \mykey[env]{mode-box} only for the \emph{\enquote{first level}}
%   of \myenv*[no-index]{enumext} and \myenv*[no-index]{enumext*} environments.
%    \begin{macrocode}
\cs_set_protected:Nn \@@_tmp:n
  {
    \keys_define:nn { enumext /  #1 }
      {
        mode-box .bool_set:N = \l_@@_mode_box_bool,
        mode-box .initial:n = false,
        mode-box .value_forbidden:n = true,
      }
  }
\clist_map_inline:nn { level-1, enumext* } { \@@_tmp:n {#1} }
%    \end{macrocode}
% \end{macro}
%
% \begin{macro}[int]{font, labelsep, labelwidth, wrap-label, wrap-label*}
%   Definition of keys \mykey[env]{font}, \mykey[env]{labelsep}, \mykey[env]{labelwidth},
%   \mykey[env]{wrap-label} and \mykey[env]{wrap-label*} keys for \myenv*[no-index]{enumext}
%   and \myenv*[no-index]{keyans} environments.
%    \begin{macrocode}
\cs_set_protected:Nn \@@_tmp:nn
  {
    \keys_define:nn { enumext / #1 }
      {
        font        .tl_set:c = { l_@@_label_font_style_ #2 _tl },
        font        .value_required:n = true,
        labelsep    .dim_set:c = { l_@@_labelsep_ #2 _dim },
        labelsep    .initial:n = {0.3333em},
        labelsep    .value_required:n = true,
        labelwidth  .dim_set:c = { l_@@_labelwidth_ #2 _dim },
        labelwidth  .value_required:n = true,
        wrap-label  .cs_set_protected:cp = { @@_wrapper_label_ #2 :n } ##1,
        wrap-label  .initial:n = {##1},
        wrap-label  .value_required:n = true,
        wrap-label* .code:n =
          {
            \bool_set_true:c { l_@@_wrap_label_opt_ #2 _bool }
            \keys_set:nn { enumext / #1 } { wrap-label = {##1} }
          },
        wrap-label* .value_required:n = true,
      }
  }
\clist_map_inline:Nn \c_@@_all_envs_clist { \@@_tmp:nn #1 }
%    \end{macrocode}
% \end{macro}
%
% \begin{macro}[int]{align}
%   The \mykey[env]{align} key is implemented differently for \emph{\enquote{starred}}
%   and \emph{\enquote{non starred}} environments. For compatibility with
%   \emph{tagged} \texttt{PDF} it is necessary to set \myhlc|\l_@@_align_label_pos_X_str|.
%    \begin{macrocode}
\cs_set_protected:Nn \@@_tmp:nn
  {
    \keys_define:nn { enumext / #1 }
      {
        align .choices:nn =
          { left , right , center }
          {
            \int_case:nn { \l_keys_choice_int }
              {
                {1}
                  {
                    \tl_clear:c { l_@@_label_fill_left_ #2 _tl  }
                    \tl_set:cn  { l_@@_label_fill_right_ #2 _tl } { \hfill }
                    \str_set:cn { l_@@_align_label_pos_ #2 _str } { l }
                  }
                {2}
                  {
                    \tl_set:cn  { l_@@_label_fill_left_ #2 _tl  } { \hfill }
                    \tl_clear:c { l_@@_label_fill_right_ #2 _tl }
                    \str_set:cn { l_@@_align_label_pos_ #2 _str } { r }
                  }
                {3}
                  {
                    \tl_set:cn { l_@@_label_fill_left_ #2 _tl  } { \hfill }
                    \tl_set:cn { l_@@_label_fill_right_ #2 _tl } { \hfill }
                    \str_set:cn { l_@@_align_label_pos_ #2 _str } { c }
                  }
              }
          },
        align .initial:n = left,
        align .value_required:n = true,
        align / unknown .code:n =
          \msg_error:nneee { enumext } { unknown-choice }
            { align } { left,~right,~ center } { \exp_not:n {##1} },
      }
  }
\clist_map_inline:nn
  {
    {level-1}{i}, {level-2}{ii}, {level-3}{iii}, {level-4}{iv}, {keyans}{v}
  }
  { \@@_tmp:nn #1 }
%    \end{macrocode}
%    \begin{macrocode}
\cs_set_protected:Nn \@@_tmp:nn
  {
    \keys_define:nn { enumext / #1 }
      {
        align .choices:nn =
          { left , right , center }
          {
            \int_case:nn { \l_keys_choice_int }
              {
                {1} { \str_set:cn { l_@@_align_label_ #2 _str } { l } }
                {2} { \str_set:cn { l_@@_align_label_ #2 _str } { r } }
                {3} { \str_set:cn { l_@@_align_label_ #2 _str } { c } }
              }
          },
        align .initial:n = left,
        align .value_required:n = true,
        align / unknown .code:n =
          \msg_error:nneee { enumext } { unknown-choice }
            { align } { left,~right,~ center } { \exp_not:n {##1} },
      }
  }
\clist_map_inline:nn { {enumext*}{vii}, {keyans*}{viii} } { \@@_tmp:nn #1 }
%    \end{macrocode}
% \end{macro}
%
% \subsection{Setting \texttt{label} and \texttt{ref} keys}\label{key:label-ref}
%
% The implementation of the keys \mykey{label} and \mykey{ref} are part
% of the core of the package \mypkg*[no-index]{enumext}, here the default values
% for \mymeta{label}, the value of the variables \myhlc|\l_@@_label_X_tl|,
% the default values for |\labelwidth| and the \emph{\enquote{label and ref}} system.
%
% \subsubsection{Definition of error messages}
%
% \begin{macro}[int]{key-ref-empty}
% Messages used by \mykey[env]{ref} key.
%    \begin{macrocode}
\msg_new:nnn { enumext } { key-ref-empty }
  {
    Key~'ref'~needs~a~value~in~'#1'~\msg_line_context:.
  }
%    \end{macrocode}
% \end{macro}
%
% \subsubsection{Define and set \texttt{label} and \texttt{ref} keys for \env{enumext} environment}
%
% \begin{macro}[int]{label, ref, \l_@@_label_i_tl, \l_@@_label_ii_tl, \l_@@_label_iii_tl, \l_@@_label_iv_tl}
%   Sets the default \mymeta{labels} of the \emph{four levels} of
%   \myenv*[no-index]{enumext} environment, along with the default values for
%   \mykey[env]{labelwidth}, \mykey[env]{widest} and \mykey[env]{ref} keys.
%    \begin{macrocode}
\cs_set_protected:Nn \@@_tmp:nnn
  {
    \tl_const:cn { c_@@_default_label_ #2 _tl } {#3}
    \keys_define:nn { enumext / #1 }
      {
        label .code:n =
          {
            \@@_label_style:cvn { l_@@_label_ #2 _tl }
              { l_@@_counter_ #2 _tl } {##1}
            \dim_set_eq:cN  { l_@@_labelwidth_ #2 _dim }
              \l_@@_current_widest_dim
          },
        label .initial:n = #3,
        label .value_required:n = true,
        ref   .code:n = { \@@_standard_ref:nn {#2} {##1} },
        ref   .value_required:n = true,
      }
  }
\@@_tmp:nnn { level-1 } {   i } { \arabic*.}
\@@_tmp:nnn { level-2 } {  ii } { (\alph*) }
\@@_tmp:nnn { level-3 } { iii } { \roman*. }
\@@_tmp:nnn { level-4 } {  iv } { \Alph*.  }
%    \end{macrocode}
% \end{macro}
%
% \begin{macro}[int]{\@@_standard_ref:nn, \@@_standard_ref:}
%   The \myhlc|\@@_standard_ref:nn| function will first pass the key
%   \emph{argument} \texttt{\textcolor{keyname}{ref}} to the variable
%   \myhlc|\l_@@_ref_key_arg_tl| and analyze its state, if it is not \emph{empty}
%   it will set a copy of of the \emph{current counter style} save in
%   \myhlc|\l_@@_the_counter_X_tl| to \myhlc|\l_@@_ref_the_count_tl| and
%   then set the variable \myhlc|\l_@@_renew_counter_X_tl| which will
%   modify |\theenumX|.
%    \begin{macrocode}
\cs_new_protected:Nn \@@_standard_ref:nn
  {
    \tl_set:Nn \l_@@_ref_key_arg_tl {#2}
    \tl_if_empty:NTF \l_@@_ref_key_arg_tl
      {
        \msg_error:nnn { enumext } { key-ref-empty } { enumext }
      }
      {
        \tl_set_eq:Nc \l_@@_ref_the_count_tl { l_@@_the_counter_ #1 _tl }
        \tl_set:ce { l_@@_renew_counter_ #1 _tl }
          {
            \exp_not:N \renewcommand { \exp_not:V \l_@@_ref_the_count_tl }
              { \exp_not:V \l_@@_ref_key_arg_tl }
          }
      }
  }
%    \end{macrocode}
%    Finally the function \myhlc|\@@_standard_ref:| will execute the
%    modification for the reference system in the second argument of the
%    environment definition \myenv*[no-index]{enumext}.
%    \begin{macrocode}
\cs_new_protected:Nn \@@_standard_ref:
  {
    \tl_if_empty:cF { l_@@_renew_counter_ \@@_level: _tl }
      {
        \tl_use:c { l_@@_renew_counter_ \@@_level: _tl }
      }
  }
%    \end{macrocode}
% \end{macro}
%
% \subsubsection{Define and set \texttt{label} and \texttt{ref} keys for \env{enumext*} and \env{keyans*} environments}
%
% \begin{macro}[int]{label, ref, \l_@@_label_vii_tl, \l_@@_label_viii_tl}
%    Sets the default \mymeta{labels} for \myenv*[no-index]{enumext*} and
%    \myenv*[no-index]{keyans*} environments, along with the default value for
%    \mykey[env]{labelwidth} key and \mykey[env]{ref} key.
%    \begin{macrocode}
\cs_set_protected:Nn \@@_tmp:nnn
  {
    \tl_const:cn { c_@@_default_label_ #2 _tl } {#3}
    \keys_define:nn { enumext / #1 }
      {
        label .code:n =
          {
            \@@_label_style:cvn { l_@@_label_ #2 _tl }
              { l_@@_counter_ #2 _tl } {##1}
            \dim_set_eq:cN  { l_@@_labelwidth_ #2 _dim }
                \l_@@_current_widest_dim
          },
        label .initial:n = #3,
        label .value_required:n = true,
        ref   .code:n = { \@@_starred_ref:n {##1} },
        ref   .value_required:n = true,
      }
  }
\@@_tmp:nnn { enumext* } {  vii } { \arabic*.}
\@@_tmp:nnn { keyans*  } { viii } { \Alph*) }
%    \end{macrocode}
% \end{macro}
%
% \begin{macro}[int]{\@@_starred_ref:n, \@@_starred_ref:}
%   The implementation of \myhlc|\@@_starred_ref:n| is the same as that
%   used for the environment \myenv*[no-index]{enumext}.
%    \begin{macrocode}
\cs_new_protected:Nn \@@_starred_ref:n
  {
    \tl_set:Nn \l_@@_ref_key_arg_tl {#1}
    \int_compare:nNnT { \l_@@_level_h_int } = { 1 }
      {
        \tl_if_empty:NTF \l_@@_ref_key_arg_tl
          {
            \msg_error:nnn { enumext } { key-ref-empty } { enumext* }
          }
          {
            \tl_set_eq:NN \l_@@_ref_the_count_tl \l_@@_the_counter_vii_tl
            \tl_set:Ne \l_@@_renew_counter_vii_tl
              {
                \exp_not:N \renewcommand { \exp_not:V \l_@@_ref_the_count_tl }
                  { \exp_not:V \l_@@_ref_key_arg_tl }
              }
          }
      }
    \int_compare:nNnT { \l_@@_keyans_level_h_int } = { 1 }
      {
        \tl_if_empty:NTF \l_@@_ref_key_arg_tl
          {
            \msg_error:nnn { enumext } { key-ref-empty } { keyans* }
          }
          {
            \tl_set_eq:NN \l_@@_ref_the_count_tl \l_@@_the_counter_viii_tl
            \tl_set:Ne \l_@@_renew_counter_viii_tl
              {
                \exp_not:N \renewcommand  { \exp_not:V \l_@@_ref_the_count_tl }
                  { \exp_not:V \l_@@_ref_key_arg_tl }
              }
          }
      }
  }
%    \end{macrocode}
%    Finally the function \myhlc|\@@_starred_ref:| will execute the
%    modification for the reference system in the second argument of the
%    \myenv*[no-index]{enumext*} and \myenv*[no-index]{keyans*} environment definition.
%    \begin{macrocode}
\cs_new_protected:Nn \@@_starred_ref:
  {
    \int_compare:nNnT { \l_@@_level_h_int } = { 1 }
      {
        \tl_if_empty:NF \l_@@_renew_counter_vii_tl
          {
            \tl_use:N \l_@@_renew_counter_vii_tl
          }
      }
    \int_compare:nNnT { \l_@@_keyans_level_h_int } = { 1 }
      {
        \tl_if_empty:NF \l_@@_renew_counter_viii_tl
          {
            \tl_use:N \l_@@_renew_counter_viii_tl
          }
      }
  }
%    \end{macrocode}
% \end{macro}
%
% \subsubsection{Define and set \texttt{label} and \texttt{ref} keys for \env{keyans} and \env{keyanspic} environments}
%
% \begin{macro}[int]{label, ref, \l_@@_label_v_tl, \l_@@_label_vi_tl}
%   Sets the default \mymeta{label} for \myenv*[no-index]{keyans}
%   and \myenv*[no-index]{keyanspic} environment, along with the default value for
%   \mykey[env]{labelwidth} if it has not been established and \mykey[env]{ref} key.
%   The \myenv*[no-index]{keyanspic} environment use the same \mymeta{label} as
%   the \myenv*[no-index]{keyans} environment.
%    \begin{macrocode}
\keys_define:nn { enumext / keyans }
  {
    label .code:n =
      {
        \@@_label_style:cvn { l_@@_label_v_tl }
          { l_@@_counter_v_tl } {#1}
        \@@_label_style:cvn { l_@@_label_vi_tl }
          { l_@@_counter_vi_tl } {#1}
        \dim_set_eq:NN
          \l_@@_labelwidth_v_dim \l_@@_current_widest_dim
      },
    label .initial:n = \Alph*),
    label .value_required:n = true,
    ref   .code:n = { \@@_keyans_ref:n {#1} },
    ref   .value_required:n = true,
  }
%    \end{macrocode}
% \end{macro}
%
% \begin{macro}[int]{\@@_keyans_ref:n, \@@_keyans_ref:}
%   The implementation of \myhlc|\@@_keyans_ref:n| is the same as that
%   used for the environment \myenv*[no-index]{enumext}.
%    \begin{macrocode}
\cs_new_protected:Nn \@@_keyans_ref:n
  {
    \tl_set:Nn \l_@@_ref_key_arg_tl {#1}
    \tl_if_empty:NTF \l_@@_ref_key_arg_tl
      {
        \msg_error:nnn { enumext } { key-ref-empty } { keyans }
      }
      {
        \tl_set_eq:NN \l_@@_ref_the_count_tl \l_@@_the_counter_v_tl
        \tl_put_right:Ne \l_@@_renew_counter_v_tl
          {
            \exp_not:N \renewcommand { \exp_not:V \l_@@_ref_the_count_tl }
              { \exp_not:V \l_@@_ref_key_arg_tl }
          }
      }
  }
%    \end{macrocode}
%    Finally the function \myhlc|\@@_keyans_ref:| will execute the
%    modification for the reference system in the second argument of the
%    \myenv*[no-index]{keyans*} environment definition.
%    \begin{macrocode}
\cs_new_protected:Nn \@@_keyans_ref:
  {
    \tl_if_empty:NF \l_@@_renew_counter_v_tl
      {
        \tl_use:N \l_@@_renew_counter_v_tl
      }
  }
%    \end{macrocode}
% \end{macro}
%
% \subsection{Setting \texttt{start}, \texttt{start*} and \texttt{widest} keys}\label{key:start-widest}
%
% \begin{macro}[int]{\@@_start_from:NNn, \@@_start_from:ccn, \@@_start_from:cce}
%    The function \myhlc|\@@_start_from:NNn| used by \mykey[env,foreachkeyans]{start} and \mykey[env]{start*}
%    keys take three arguments:
%
% \begin{arguments}
%    \item \myhlc|\l_@@_label_X_tl|
%    \item \myhlc|\l_@@_start_X_int|
%    \item \meta{integer or string}
% \end{arguments}
%    The argument `|#1|' is the \emph{\enquote{counter style}} set by
%    \mykey[env]{label} key, `|#2|' is returned by the function,
%    `|#3|' can be an \mymeta{integer} or \mymeta{string} of the
%    form |\Alph|, |\alph|, |\Roman| or |\roman|. This effectively
%    allows \mykey{start*}|=A| or \mykey{start}|=1| to be used.
%
% \begin{important}*
%  In version 1.6 it is allowed to pass the \mykey[env]{resume} key \emph{without
%  value} by means of the command |\setenumext|, for the correct operation
%  of this, the boolean variable \myhlc|\l_@@_resume_count_bool|
%  set by the \mykey{resume} key \emph{without value} must be set to |false| (\S\ref{key:series}).
%  This is necessary to be able to \enquote{reset} the \emph{start value}
%  by means of the \mykey{start} or \mykey{start*} keys.
% \end{important}
%
%    \begin{macrocode}
\cs_new_protected:Nn \@@_start_from:NNn
  {
    \bool_set_false:N \l_@@_resume_count_bool
    \@@_if_is_int:nTF {#3}
      {
        \int_set:Nn #2 {#3}
      }
      {
        \regex_if_match:nVT { \c{Alph} | \c{alph} } #1
          { \int_set:Nn #2 { \int_from_alph:n {#3} } }
        \regex_if_match:nVT { \c{Roman} | \c{roman} } #1
          { \int_set:Nn #2  { \int_from_roman:n {#3} } }
      }
  }
\cs_generate_variant:Nn \@@_start_from:NNn { ccn, cce }
%    \end{macrocode}
% \end{macro}
%
% \begin{macro}[int]{\@@_widest_from:nNNn, \@@_widest_from:nccn}
%    The function \myhlc|\@@_widest_from:nNNn| used by the \mykey[env]{widest}
%    key take four arguments:
%
% \begin{arguments}
%    \item The counter associated with the environment level
%    \item \myhlc|\l_@@_label_X_tl|
%    \item \myhlc|\l_@@_labelwidth_X_dim|
%    \item \meta{integer or string}
% \end{arguments}
%
%    The arguments `|#2|' and `|#3|' are the values set by
%    \mykey[env]{label} and \mykey[env]{labelwidth} keys, `|#4|' can
%    be an \mymeta{integer} or \mymeta{string} of the
%    form |\Alph|, |\alph|, |\Roman| or |\roman|.
%    The value of `|#4|' is set temporarily for the
%    identified counter in this point (level), then the value is expanded
%    into a \emph{\enquote{box}} and the \emph{\enquote{width}} of the
%    \emph{\enquote{box}} is returned.
%    \begin{macrocode}
\cs_new_protected:Nn \@@_widest_from:nNNn
  {
    \@@_if_is_int:nTF {#4}
      {
        \setcounter{enumX #1} {#4}
      }
      {
        \regex_if_match:nVT { \c{Alph} | \c{alph} } #2
          { \setcounter{enumX #1} { \int_from_alph:n {#4} } }
        \regex_if_match:nVT { \c{Alph} | \c{alph} } #2
          { \setcounter{enumX #1} { \int_from_roman:n {#4} } }
      }
    \@@_label_width_by_box:cv
      { l_@@_labelwidth_ #1 _dim } { l_@@_label_ #1 _tl }
  }
\cs_generate_variant:Nn \@@_widest_from:nNNn { nccn }
%    \end{macrocode}
% \end{macro}
%
% \begin{macro}[int]{start, start*, widest}
%    Now define and set \mykey[env]{start*}, \mykey[env,foreachkeyans]{start} and \mykey[env]{widest}
%    keys for \myenv*[no-index]{enumext}, \myenv*[no-index]{enumext*}, \myenv*[no-index]{keyans} and
%    \myenv*[no-index]{keyans*} environments.
%
%    \begin{macrocode}
\cs_set_protected:Nn \@@_tmp:nn
  {
    \keys_define:nn { enumext / #1 }
      {
        start* .code:n =
          {
            \@@_start_from:ccn
              { l_@@_label_ #2 _tl }
              { l_@@_start_ #2 _int } {##1}
          },
        start* .value_required:n = true,
        start  .code:n =
          {
            \@@_start_from:cce
              { l_@@_label_ #2 _tl }
              { l_@@_start_ #2 _int } { \int_eval:n {##1} }
          },
        start  .initial:n = 1,
        start  .value_required:n = true,
        widest .code:n =
          {
            \@@_widest_from:nccn {#2}
              { l_@@_label_ #2 _tl }
              { l_@@_labelwidth_ #2 _dim } {##1}
          },
        widest .value_required:n = true,
      }
  }
\clist_map_inline:Nn \c_@@_all_envs_clist { \@@_tmp:nn #1 }
%    \end{macrocode}
% \end{macro}
%
% \subsection{Setting keys for penalties}\label{subsec:set-penalties}
%
% \begin{macro}[int]{beginpenalty, midpenalty, endpenalty}
%   The three parameters \textcolor{red}{\texttt{\textbackslash}}|@beginparpenalty|,
%   \textcolor{red}{\texttt{\textbackslash}}|@itempenalty| and
%   \textcolor{red}{\texttt{\textbackslash}}|@endparpenalty| work together
%   to ensure that list environments look good, avoiding unsightly page
%   breaks that can break the flow of the \myenv[no-index]{list}, so it's a good
%   idea to have a \mymeta{keys} to access these.
%    \begin{macrocode}
\cs_set_protected:Nn \@@_tmp:nn
  {
    \keys_define:nn { enumext / #1 }
      {
        beginpenalty .int_set:c = { l_@@_beginparpenalty_ #2 _int },
        beginpenalty .initial:n = -51,
        beginpenalty .value_required:n = true,
        midpenalty   .int_set:c = { l_@@_itempenalty_ #2 _int },
        midpenalty   .initial:n = -51,
        midpenalty   .value_required:n = true,
        endpenalty   .int_set:c = { l_@@_endparpenalty_ #2 _int },
        endpenalty   .initial:n = -51,
        endpenalty   .value_required:n = true,
      }
  }
\clist_map_inline:Nn \c_@@_all_envs_clist { \@@_tmp:nn #1 }
%    \end{macrocode}
% \end{macro}
%
% \subsection{Setting keys for vertical spaces}\label{subsec:set-vspaces}
%
% \begin{macro}[int]{topsep, partopsep, parsep, itemsep, noitemsep, nosep}
%   Define and set \mykey[env]{topsep}, \mykey[env]{partopsep}, \mykey[env]{parsep},
%   \mykey[env]{itemsep}, \mykey[env]{noitemsep} and \mykey[env]{nosep} keys for
%   \myenv*[no-index]{enumext}, \myenv*[no-index]{enumext*}, \myenv*[no-index]{keyans} and
%   \myenv*[no-index]{keyans*} environments.
%
%    \begin{macrocode}
\cs_set_protected:Nn \@@_tmp:nnnnnn
  {
    \keys_define:nn { enumext / #1 }
      {
        topsep    .skip_set:c = { l_@@_topsep_ #2 _skip },
        topsep    .initial:n = {#3},
        topsep    .value_required:n = true,
        partopsep .skip_set:c = { l_@@_partopsep_ #2 _skip },
        partopsep .initial:n = {#4},
        partopsep .value_required:n = true,
        parsep    .skip_set:c = { l_@@_parsep_ #2 _skip },
        parsep    .initial:n = {#5},
        parsep    .value_required:n = true,
        itemsep   .skip_set:c = { l_@@_itemsep_ #2 _skip },
        itemsep   .initial:n = {#6},
        itemsep   .value_required:n = true,
        noitemsep .meta:n = { itemsep = 0pt, parsep = 0pt },
        noitemsep .value_forbidden:n = true,
        nosep     .meta:n =
          { itemsep = 0pt, parsep = 0pt, topsep = 0pt, partopsep = 0pt },
        nosep     .value_forbidden:n = true,
      }
  }
%    \end{macrocode}
%    Now the values are set based on the standard \myclass{article} class in
%    \mydim{10pt}.
%    \begin{macrocode}
\@@_tmp:nnnnnn { level-1 } { i } { 8.0pt plus 2.0pt minus 4.0pt }
  { 2.0pt plus 1.0pt minus 1.0pt } { 4.0pt plus 2.0pt minus 1.0pt }
  { 4.0pt plus 2.0pt minus 1.0pt }
\@@_tmp:nnnnnn { level-2 } { ii } { 4.0pt plus 2.0pt minus 1.0pt }
  { 2.0pt plus 1.0pt minus 1.0pt } { 2.0pt plus 1.0pt minus 1.0pt }
  { 2.0pt plus 1.0pt minus 1.0pt }
\@@_tmp:nnnnnn { level-3 } { iii } { 2.0pt plus 1.0pt minus 1.0pt }
  { 1.0pt minus 1.0pt }{ 0pt }{ 2.0pt plus 1.0pt minus 1.0pt }
\@@_tmp:nnnnnn { level-4 } { iv } { 2.0pt plus 1.0pt minus 1.0pt }
  { 1.0pt minus 1.0pt }{ 0pt }{ 2.0pt plus 1.0pt minus 1.0pt }
\@@_tmp:nnnnnn { keyans  } { v }{ 4.0pt plus 2.0pt minus 1.0pt }
  { 2.0pt plus 1.0pt minus 1.0pt }{ 2.0pt plus 1.0pt minus 1.0pt }
  { 2.0pt plus 1.0pt minus 1.0pt }
\@@_tmp:nnnnnn { enumext* } { vii } { 8.0pt plus 2.0pt minus 4.0pt }
  { 2.0pt plus 1.0pt minus 1.0pt } { 4.0pt plus 2.0pt minus 1.0pt }
  { 4.0pt plus 2.0pt minus 1.0pt }
\@@_tmp:nnnnnn { keyans* } { viii } { 4.0pt plus 2.0pt minus 1.0pt }
  { 2.0pt plus 1.0pt minus 1.0pt } { 2.0pt plus 1.0pt minus 1.0pt }
  { 2.0pt plus 1.0pt minus 1.0pt }
%    \end{macrocode}
% \end{macro}
%
% \subsection{Setting \texttt{base-fix} key}\label{key:base-fix}
%
% When nesting starting right after |\item| (without material between them) there
% is a problem with the alignment of the \emph{baseline} between the two environments.
% One way to get around this problem is to place \myhlc|\mode_leave_vertical:|
% apply |\vspace{-\baselineskip}| and set |\topsep=|\mydim{0pt} for
% the \emph{\enquote{first level}} of the nested \myenv*[no-index]{enumext} environment.
%
% \begin{macro}[int]{base-fix}
%   Defines the key \mykey[env]{base-fix} only for the \emph{\enquote{first level}}
%   of \myenv*[no-index]{enumext} environment.
%    \begin{macrocode}
\keys_define:nn { enumext / level-1 }
  {
    base-fix .bool_set:N = \l_@@_base_line_fix_bool,
    base-fix .initial:n = false,
    base-fix .value_forbidden:n = true,
  }
%    \end{macrocode}
% \end{macro}
%
% \begin{macro}[int]{\@@_nested_base_line_fix:}
%   The function \myhlc|\@@_nested_base_line_fix:| passed to the \myhlc|\@@_parse_keys:n|
%   function in the definition of the \myenv*[no-index]{enumext} environment (\S\ref{env:enumext})
%   will be responsible for applying the \emph{baseline correction} and
%   adjusting the \mymeta{keys} for the \myenv*[no-index]{enumext} environment and
%   the |\printkeyans| with \emph{starred argument} `|*|' (\S\ref{cmd:printkeyans}).
%
%   First, the function code is implemented from the user side of the
%   \mykey{base-fix} key, that is, only the user knows when it is necessary
%   to apply it within the document in which case the variable
%   \myhlc|\l_@@_print_keyans_star_bool| set by the |\printkeyans| command is
%   false and the variable \myhlc|\l_@@_base_line_fix_bool| is true.
%
%   Sets the values of the keys \mykey[env]{topsep}, \mykey[env]{above}
%   and \mykey[env]{above*} for the \emph{\enquote{first level}} of \myenv*[no-index]{enumext}
%   environment equal to \mydim{0pt} and finally sets the variable \myhlc|\l_@@_base_line_fix_bool|
%   to false.
%    \begin{macrocode}
\cs_new_protected:Nn \@@_nested_base_line_fix:
  {
    \bool_lazy_all:nT
      {
        { \bool_if_p:N \l_@@_starred_first_bool }
        { \bool_if_p:N \l_@@_base_line_fix_bool }
        { \bool_not_p:n { \l_@@_print_keyans_star_bool } }
      }
      {
        \mode_leave_vertical:
        \vspace { -\dim_eval:n { \baselineskip + \parsep } }
        \keys_set:nn { enumext / level-1 }
          {
            topsep = 0pt, above = 0pt, above* = 0pt,
          }
      }
%    \end{macrocode}
%   When the |\printkeyans| command is run with the \emph{starred
%   argument} `|*|' the variable \myhlc|\l_@@_print_keyans_star_bool| is true
%   and a simplified version of |\vspace| can be run using \myhlc|\skip_vertical:n|.
%    \begin{macrocode}
    \bool_lazy_and:nnT
      { \bool_if_p:N \l_@@_starred_first_bool }
      { \bool_if_p:N \l_@@_print_keyans_star_bool }
      {
        \mode_leave_vertical:
        \skip_vertical:n { -\baselineskip }
        \skip_vertical:N \c_zero_skip
        \keys_set:nn { enumext / level-1 }
          {
            topsep = 0pt, above = 0pt, above* = 0pt,
          }
      }
    \bool_set_false:N \l_@@_base_line_fix_bool
  }
%    \end{macrocode}
% \end{macro}
%
% \subsection{Setting keys for horizontal spaces}\label{key:list-indent}
%
% \begin{macro}[int]{itemindent, rightmargin, listparindent, list-offset, list-indent}
%    Define and set \mykey[env]{itemindent}, \mykey[env]{rightmargin}, \mykey[env]{listparindent}, \mykey[env]{list-offset}
%    and \mykey[env]{list-indent} keys for \myenv*[no-index]{enumext}, \myenv*[no-index]{enumext*}, \myenv*[no-index]{keyans} and
%    \myenv*[no-index]{keyans*} environments.
%    \begin{macrocode}
\cs_set_protected:Nn \@@_tmp:nn
  {
    \keys_define:nn { enumext / #1 }
      {
        itemindent    .dim_set:c = { l_@@_fake_item_indent_ #2 _dim },
        itemindent    .value_required:n = true,
        rightmargin   .dim_set:c = { l_@@_rightmargin_ #2 _dim },
        rightmargin   .value_required:n = true,
        listparindent .dim_set:c = { l_@@_listparindent_ #2 _dim },
        listparindent .value_required:n = true,
        list-offset   .dim_set:c = { l_@@_listoffset_ #2 _dim },
        list-offset   .value_required:n = true,
        list-indent   .code:n =
          {
            \bool_set_true:c { l_@@_leftmargin_tmp_ #2 _bool }
            \dim_set:cn { l_@@_leftmargin_tmp_ #2 _dim } {##1}
          },
        list-indent   .value_required:n = true,
      }
  }
\clist_map_inline:nn
  {
    {level-1}{i}, {level-2}{ii}, {level-3}{iii}, {level-4}{iv}, {keyans}{v}
  }
  { \@@_tmp:nn #1 }
%    \end{macrocode}
%
% For \myenv*[no-index]{enumext*} and \myenv*[no-index]{keyans*} environments the situation
% is a bit different, the \mykey{list-indent} key behaves like the \mykey{list-offset} key.
%    \begin{macrocode}
\cs_set_protected:Nn \@@_tmp:nn
  {
    \keys_define:nn { enumext / #1 }
      {
        itemindent    .dim_set:c = { l_@@_fake_item_indent_ #2 _dim },
        itemindent    .value_required:n = true,
        rightmargin   .dim_set:c = { l_@@_rightmargin_ #2 _dim },
        rightmargin   .value_required:n = true,
        listparindent .dim_set:c = { l_@@_listparindent_ #2 _dim },
        listparindent .value_required:n = true,
        list-offset   .dim_set:c = { l_@@_listoffset_ #2 _dim },
        list-offset   .value_required:n = true,
        list-indent   .meta:n = { list-offset = ##1 },
        list-indent   .value_required:n = true,
      }
  }
\clist_map_inline:nn
  {
    {enumext*}{vii}, {keyans*}{viii}
  }
  { \@@_tmp:nn #1 }
%    \end{macrocode}
% \end{macro}
%
% \subsubsection{Functions for setting the fake \texttt{itemindent}}\label{subsubsec:fake-indent}
%
% \begin{macro}[int]{\@@_fake_item_indent:, \@@_keyans_fake_item_indent:,
%                    \@@_fake_item_indent_vii:, \@@_fake_item_indent_viii:}
%
%    The \mykey[env]{itemindent} key does not set the value of |\itemindent|, it only
%    sets the value of the \emph{horizontal space} applied using \myhlc|\skip_horizontal:N|. This
%    value is stored in the variable and only applied when it is greater
%    than \mydim{0pt}. Here it is necessary to place
%    \myhlc|\mode_leave_vertical:| and the plain \hologo{TeX} macro
%    |\ignorespaces| to avoid unwanted extra space when using the
%    \mykey{itemindent} key.
%    \begin{macrocode}
\cs_set_protected:Nn \@@_fake_item_indent:
  {
    \dim_compare:nNnT
      { \dim_use:c { l_@@_fake_item_indent_ \@@_level: _dim } }
      >
      { \c_zero_dim }
      {
        \tl_set:ce { l_@@_fake_item_indent_ \@@_level: _tl }
          {
            \exp_not:N \mode_leave_vertical:
            \exp_not:n { \skip_horizontal:n }
              { \dim_use:c { l_@@_fake_item_indent_ \@@_level: _dim } }
            \exp_not:N \ignorespaces
          }
      }
  }
\cs_set_protected:Nn \@@_keyans_fake_item_indent:
  {
    \dim_compare:nNnT
      { \l_@@_fake_item_indent_v_dim } > { \c_zero_dim }
      {
        \tl_set:Ne \l_@@_fake_item_indent_v_tl
          {
            \exp_not:N \mode_leave_vertical:
            \exp_not:N \skip_horizontal:N \l_@@_fake_item_indent_v_dim
            \exp_not:N \ignorespaces
          }
      }
  }
\cs_set_protected:Nn \@@_fake_item_indent_vii:
  {
    \dim_compare:nNnT
      { \l_@@_fake_item_indent_vii_dim } > { \c_zero_dim }
      {
        \tl_set:Ne \l_@@_fake_item_indent_vii_tl
          {
            \exp_not:N \skip_horizontal:N \l_@@_fake_item_indent_vii_dim
            \exp_not:N \ignorespaces
          }
      }
  }
\cs_set_protected:Nn \@@_fake_item_indent_viii:
  {
    \dim_compare:nNnT
      { \l_@@_fake_item_indent_viii_dim } > { \c_zero_dim }
      {
        \tl_set:Ne \l_@@_fake_item_indent_viii_tl
          {
            \exp_not:N \skip_horizontal:N \l_@@_fake_item_indent_viii_dim
            \exp_not:N \ignorespaces
          }
      }
  }
%    \end{macrocode}
% \end{macro}
%
% \subsection{Setting \texttt{show-length} key}\label{subsec:set-show-length}
%
% \begin{macro}[int]{show-length}
%    Define and set \mykey[env]{show-length} key for \myenv*[no-index]{enumext},
%    \myenv*[no-index]{enumext*}, \myenv*[no-index]{keyans} and \myenv*[no-index]{keyans*} environments.
%    The function sets the boolean variable \myhlc|\l_@@_show_length_X_bool|
%    used in the definition of all environments to |true|
%    and calls the function \myhlc|\@@_show_length:nnn| which prints all the
%    values of the \emph{\enquote{vertical}} and \emph{\enquote{horizontal}}
%    parameters calculated and used.
%    \begin{macrocode}
\cs_set_protected:Nn \@@_tmp:nn
  {
    \keys_define:nn { enumext / #1 }
      {
        show-length .bool_set:c = { l_@@_show_length_ #2 _bool },
        show-length .initial:n = false,
      }
  }
\clist_map_inline:Nn \c_@@_all_envs_clist { \@@_tmp:nn #1 }
%    \end{macrocode}
% \end{macro}
%
% \subsection{Setting \texttt{before}, \texttt{after} and \texttt{first} keys}\label{key:before}
%
% \begin{macro}[int]{before, before*, after, first}
%    Define and set \mykey[env,foreachkeyans]{before}, \mykey[env]{before*}, \mykey[env,foreachkeyans]{after} and \mykey[env]{first} keys
%    for \myenv*[no-index]{enumext}, \myenv*[no-index]{enumext*}, \myenv*[no-index]{keyans} and  \myenv*[no-index]{keyans*} environments.
%    \begin{macrocode}
\cs_set_protected:Nn \@@_tmp:nn
  {
    \keys_define:nn { enumext / #1 }
      {
        before  .tl_set:c = { l_@@_before_no_starred_key_ #2 _tl },
        before  .value_required:n = true,
        before* .tl_set:c = { l_@@_before_starred_key_ #2 _tl },
        before* .value_required:n = true,
        after   .tl_set:c = { l_@@_after_stop_list_ #2 _tl },
        after   .value_required:n = true,
        first   .tl_set:c = { l_@@_after_list_args_ #2 _tl },
        first   .value_required:n = true,
      }
  }
\clist_map_inline:Nn \c_@@_all_envs_clist { \@@_tmp:nn #1 }
%    \end{macrocode}
% \end{macro}
%
% \subsubsection{Functions for \texttt{before}, \texttt{after} and \texttt{first} keys in \env{enumext}}
%
% \begin{macro}[int]{\@@_before_args_exec:, \@@_before_keys_exec:,
%                    \@@_after_stop_list:, \@@_after_args_exec:}
%    The function \myhlc|\@@_before_args_exec:| executes the \mymarg{code} set by the
%    \mykey[env]{before*} key \emph{\enquote{before}} the \myenv*[no-index]{enumext}
%    environment is started. The \mymarg{code} is executed \emph{\enquote{without}}
%    knowing any definition of the \mymarg{arg two} of the list:
%    \mymarg[cbc=module_name,ac=araracolor]{code}|\list|\mymarg{arg one}\mymarg{arg two}.
%    \begin{macrocode}
\cs_new_protected:Nn \@@_before_args_exec:
  {
    \tl_use:c { l_@@_before_starred_key_ \@@_level: _tl }
  }
%    \end{macrocode}
%    The function \myhlc|\@@_before_keys_exec:| executes the \mymarg{code} set by the
%    \mykey[env,foreachkeyans]{before} key \emph{\enquote{before}} the \myenv*[no-index]{enumext}
%    environment is started in \emph{second argument} of the list. The
%    \mymarg{code} is executed \emph{\enquote{knowing}}
%    all definition and values provides by \mymeta{keys}:
%    |\list|\mymarg{arg one}|{|\mymeta{arg two}\mymarg[cbc=module_name,ac=araracolor]{code}|}|
%    \begin{macrocode}
\cs_new_protected:Nn \@@_before_keys_exec:
  {
    \tl_use:c { l_@@_before_no_starred_key_ \@@_level: _tl }
  }
%    \end{macrocode}
%    The function \myhlc|\@@_after_stop_list:| executes the \mymarg{code} set by the
%    \mykey[env,foreachkeyans]{after} key \emph{\enquote{after}} the \myenv*[no-index]{enumext}
%    environment has finished: |\endlist|\mymarg[cbc=module_name,ac=araracolor]{code}.
%    \begin{macrocode}
\cs_new_protected:Nn \@@_after_stop_list:
  {
    \tl_use:c { l_@@_after_stop_list_ \@@_level: _tl }
  }
%    \end{macrocode}
%    The function \myhlc|\@@_after_args_exec:| executes the \mymarg{code} set by the
%    \mykey[env]{first} key after the end of the second argument of the list
%    defining the \myenv*[no-index]{enumext} environment, just before the first
%    occurrence of |\item|: |\list|\mymarg{arg one}\mymarg{arg two}\mymarg[cbc=module_name,ac=araracolor]{code}|\item|.
%    \begin{macrocode}
\cs_new_protected:Nn \@@_after_args_exec:
  {
    \tl_use:c { l_@@_after_list_args_ \@@_level: _tl }
  }
%    \end{macrocode}
% \end{macro}
%
% \subsubsection{Functions for \texttt{before}, \texttt{after} and \texttt{first} keys in \env{keyans}}
%
% \begin{macro}[int]{\@@_before_args_exec_v:, \@@_before_keys_exec_v:,
%                    \@@_after_stop_list_v:, \@@_after_args_exec_v:}
% Same implementation as the one used in the \myenv*[no-index]{enumext} environment.
%    \begin{macrocode}
\cs_new_protected:Nn \@@_before_args_exec_v:
  {
    \tl_use:N \l_@@_before_starred_key_v_tl
  }
\cs_new_protected:Nn \@@_before_keys_exec_v:
  {
    \tl_use:N \l_@@_before_no_starred_key_v_tl
  }
\cs_new_protected:Nn \@@_after_stop_list_v:
  {
    \tl_use:N \l_@@_after_stop_list_v_tl
  }
\cs_new_protected:Nn \@@_after_args_exec_v:
  {
    \tl_use:N \l_@@_after_list_args_v_tl
  }
%    \end{macrocode}
% \end{macro}
%
% \subsubsection{Functions for \texttt{before}, \texttt{after} and \texttt{first} keys in \env{enumext*} and \env{keyans*}}
%
% \begin{macro}[int]{\@@_before_args_exec_vii:, \@@_before_args_exec_viii:,
%                    \@@_before_keys_exec_vii:, \@@_before_keys_exec_viii:,
%                    \@@_after_stop_list_vii:, \@@_after_stop_list_viii:}
% Same implementation as the one used in the \myenv*[no-index]{enumext} environment,
% except that there are no \myhlc|\@@_after_args_exec_vii:| and
% \myhlc|\@@_after_args_exec_viii:| functions: the \mymarg{code} set by the
% \mykey[env]{first} key is used directly by \myhlc|\@@_start_item_vii:w| and
% \myhlc|\@@_start_item_viii:w|, inside the \env{minipage} of each \cs{item}.
%    \begin{macrocode}
\cs_new_protected:Nn \@@_before_args_exec_vii:
  {
    \tl_use:N \l_@@_before_starred_key_vii_tl
  }
\cs_new_protected:Nn \@@_before_args_exec_viii:
  {
    \tl_use:N \l_@@_before_starred_key_viii_tl
  }
\cs_new_protected:Nn \@@_before_keys_exec_vii:
  {
    \tl_use:N \l_@@_before_no_starred_key_vii_tl
  }
\cs_new_protected:Nn \@@_before_keys_exec_viii:
  {
    \tl_use:N \l_@@_before_no_starred_key_viii_tl
  }
\cs_new_protected:Nn \@@_after_stop_list_vii:
  {
    \tl_use:N \l_@@_after_stop_list_vii_tl
  }
\cs_new_protected:Nn \@@_after_stop_list_viii:
  {
    \tl_use:N \l_@@_after_stop_list_viii_tl
  }
%    \end{macrocode}
% \end{macro}
%
% \subsection{Setting keys for \env{multicols} and \env{minipage}}\label{key:mini-multi}
%
% \begin{macro}[int]{mini-env, mini-sep, columns-sep, columns, mini-right, mini-right*}
%    The default value of the \mykey[env]{columns-sep} key is handled by
%    the state of the boolean variable \myhlc|\l_@@_columns_sep_X_bool|
%    which is handled in the internal definition of the \myenv*[no-index]{enumext}
%    and \myenv*[no-index]{keyans} environments.
%    Define and set \mykey[env]{mini-env}, \mykey[env]{mini-sep}, \mykey{columns-sep}
%    and \mykey[env]{columns} keys for \myenv*[no-index]{enumext}, \myenv*[no-index]{enumext*}, \myenv*[no-index]{keyans} and
%    \myenv*[no-index]{keyans*} environments.
%    \begin{macrocode}
\cs_set_protected:Nn \@@_tmp:nn
  {
    \keys_define:nn { enumext / #1 }
      {
        mini-env    .dim_set:c = { l_@@_minipage_right_ #2 _dim },
        mini-env    .value_required:n = true,
        mini-sep    .dim_set:c = { l_@@_minipage_hsep_ #2 _dim },
        mini-sep    .initial:n = 0.3333em,
        mini-sep    .value_required:n = true,
        columns-sep .dim_set:c = { l_@@_columns_sep_ #2 _dim },
        columns-sep .value_required:n = true,
        columns     .int_set:c = { l_@@_columns_ #2 _int },
        columns     .initial:n = 1,
        columns     .value_required:n = true,
      }
  }
\clist_map_inline:Nn \c_@@_all_envs_clist { \@@_tmp:nn #1 }
%    \end{macrocode}
%
% For \myenv*[no-index]{enumext*} and \myenv*[no-index]{keyans*} environments the situation
% is a bit different, the command |\miniright| is not available, so the
% keys \mykey[env]{mini-right} and \mykey[env]{mini-right*} are added to implement support for
% \myenv[no-index]{minipage} environment.
%    \begin{macrocode}
\cs_set_protected:Nn \@@_tmp:nn
  {
    \keys_define:nn { enumext / #1 }
      {
        mini-right  .tl_gset:c = { g_@@_miniright_code_ #2 _tl },
        mini-right  .value_required:n = true,
        mini-right* .code:n =
          {
            \bool_gset_true:c { g_@@_minipage_center_ #2 _bool }
            \keys_set:nn { enumext / #1 } { mini-right = {##1} }
          },
        mini-right* .value_required:n = true,
      }
  }
\clist_map_inline:nn { {enumext*}{vii}, {keyans*}{viii} } { \@@_tmp:nn #1 }
%    \end{macrocode}
% \end{macro}
%
% \subsection{Adjustment of vertical spaces for \env{multicols}} \label{subsec:vspace-multicols}
%
% When nesting a \emph{\enquote{list environment}} inside the \myenv[no-index]{multicols} environment,
% the values of the \emph{\enquote{vertical spaces}} are lost, basically the
% \myenv[no-index]{multicols} environment takes control over them. Graphically, this can be seen in figure~\ref{fig:six}.
%
% \begin{figure}[!htbp]
%   \centering
%   \begin{tikzpicture}[line cap=round,line join=round,x=0.8cm,y=0.8cm,every node/.style={font=\footnotesize}]
%     \draw[color=white] (7.0,6.10) rectangle (17.5,6.50) node[pos=.5,text=black] { \textcolor{blue}{\texttt{\textbackslash{}multicolsep [±\textbackslash{}topskip]}} };% space above
%     \draw[color=white] (7.0,4.85) rectangle (17.5,5.25) node[pos=.5,text=black] { \textcolor{blue}{\texttt{\textbackslash{}multicolsep [±\textbackslash{}prevdepth?]}} };% space below
%     \draw[color=white] (2.0,6.10) rectangle (5,6.50) node[pos=.5,text=keyname,font=\footnotesize\ttfamily] { topsep };% topsep above
%     \draw[color=white] (2.0,4.85) rectangle (5,5.25) node[pos=.5,text=keyname,font=\footnotesize\ttfamily] { topsep };% topsep below
%     \draw[<->,color=gray] (2.0,7.0) --node[above,text=keyname]{\vphantom{p}\texttt{labelwidth}} (5.0,7.0);
%     \draw[<->,color=gray] (5.0,7.0) --node[above,text=keyname]{\texttt{labelsep}} (7.0,7.0);
%     \draw[color=lightgray,fill=lightgray,fill opacity=0.15] (2.0,6.5) rectangle (5.0,6.85);% labelwidth
%     \draw[->,color=lightgray,shorten >=3pt] (4.5,6.3) -- (7.0,6.3); % glue space = topsep (above)
%     \draw[->,color=lightgray,shorten >=3pt] (4.5,5.05) -- (7.0,5.05); % glue space = topsep (above)
%     \draw[color=lightgray,fill=lightgray,fill opacity=0.15] (5.0,6.5) rectangle (7.0,6.85);% labelsep
%     \draw[color=lightgray,fill=lightgray,fill opacity=0.15] (7.0,6.5) rectangle (17.5,6.85);% item width
%     \draw[color=lightgray,fill=lightgray,fill opacity=0.15] (7.0,5.25) rectangle (11.5,6.10) node[pos=.5,text=black]
%        {\begin{tabular}{c}
%          \texttt{column one} \\
%          \textcolor{blue}{\emph{nested list or items}}
%        \end{tabular}};
%     \draw[color=lightgray,fill=lightgray,fill opacity=0.15] (13.5,5.25) rectangle (17.5,6.10) node[pos=.5,text=black]
%        {\begin{tabular}{c}
%          \texttt{column two} \\
%          \textcolor{blue}{\emph{nested list or items}}
%        \end{tabular}};
%     \draw[color=lightgray,fill=lightgray,fill opacity=0.15] (2.0,4.5) rectangle (5.0,4.85);% labelwidth
%     \draw[color=lightgray,fill=lightgray,fill opacity=0.15] (5.0,4.5) rectangle (7.0,4.85);% labelsep
%     \draw[color=lightgray,fill=lightgray,fill opacity=0.15] (7.0,4.5) rectangle (17.5,4.85);% item width
%   \end{tikzpicture}
%   \caption{Representation of the vertical space in \myenv[no-index]{multicols} for a nested level.}
%   \label{fig:six}
% \end{figure}
%
% To keep the desired spaces \emph{above} and \emph{below} in the \emph{\enquote{list environment}}
% (|\topsep + [\partopsep]|) it is necessary to \emph{\enquote{adjust}} the spaces added
% by the \myenv[no-index]{multicols} environment.
% The most appropriate option in this case is to use a \emph{\enquote{context sensitive}}
% vertical space with |\addvspace|.
%
% \begin{important}*
% It should be made clear that the implementation here is a \emph{\enquote{bit
% questionable}}. At first glance doing |\multicolsep=\topsep| seemed right, but the results
% were not always as expected. An almost \emph{imperceptible} detail is that in some cases the
% |\itemsep| values of are \emph{\enquote{stretched}}, possibly due to
% the use of |\raggedcolumns| and this affects the lower space when
% closing the environment, which is \emph{\enquote{smaller}} than
% expected. My attempts to find the correct values using |\showoutput| and
% |\showboxdepth| absolutely failed.
% \end{important}
%
% \subsubsection{Adjustment of vertical spaces for \env{multicols} in \env{enumext}}
%
%
% \begin{macro}[int]{\@@_multi_set_vskip:}
%    The function \myhlc|\@@_multi_set_vskip:| will take care of determining the
%    \emph{\enquote{adjusted spaces}} applied \emph{\enquote{above}} and
%    \emph{\enquote{below}} the \myenv[no-index]{multicols} environment in
%    \myenv*[no-index]{enumext}.
%
%    The default values are set taking into account that \hologo{TeX} is
%    in \mymeta{horizontal mode}, then the settings for the
%    \mymeta{vertical mode}, in which |\partopsep| comes into play, are made.
%
%    Set the values of \myhlc|\l_@@_multicols_above_X_skip| and \myhlc|\l_@@_multicols_below_X_skip|
%    equal to the value of |\topsep| in the \emph{current level}.
%    \begin{macrocode}
\cs_new_protected:Nn \@@_multi_set_vskip:
  {
    \skip_set:cn { l_@@_multicols_above_ \@@_level: _skip }
      {
        \skip_use:c { l_@@_topsep_ \@@_level: _skip }
      }
    \skip_set:cn { l_@@_multicols_below_ \@@_level: _skip }
      {
        \skip_use:c { l_@@_topsep_ \@@_level: _skip }
      }
    \@@_add_pre_parsep:
  }
%    \end{macrocode}
% \end{macro}
%
% \begin{macro}[int]{\@@_add_pre_parsep:}
%    The function \myhlc|\@@_add_pre_parsep:| \emph{\enquote{adjusted}} the value
%    of \myhlc|\l_@@_multicols_above_X_skip| detecting the value of
%    |\parsep| from the previous level. This is necessary since
%    |\parsep| from the previous level affects the \emph{vertical spaces}.
%    \begin{macrocode}
\cs_new_protected:Nn \@@_add_pre_parsep:
  {
    \int_case:nn { \l_@@_level_int }
      {
        { 2 }
          {
            \skip_if_eq:nnF { \l_@@_parsep_i_skip } { \c_zero_skip }
              {
                \skip_add:Nn \l_@@_multicols_above_ii_skip
                  {
                    \l_@@_parsep_i_skip
                  }
              }
          }
        { 3 }
          {
            \skip_if_eq:nnF { \l_@@_parsep_ii_skip } { \c_zero_skip }
              {
                \skip_add:Nn \l_@@_multicols_above_iii_skip
                  {
                    \l_@@_parsep_ii_skip
                  }
              }
          }
        { 4 }
          {
            \skip_if_eq:nnF { \l_@@_parsep_iii_skip } { \c_zero_skip }
              {
                \skip_add:Nn \l_@@_multicols_above_iv_skip
                  {
                    \l_@@_parsep_iii_skip
                  }
              }
          }
      }
  }
%    \end{macrocode}
% \end{macro}
%
% \begin{macro}[int]{\@@_multi_addvspace:}
%   The function \myhlc|\@@_multi_addvspace:| will apply the spaces set using
%   |\addvspace| \emph{\enquote{above}} the \myenv[no-index]{multicols} environment
%   in \myenv*[no-index]{enumext}, taking into account whether  \hologo{TeX} is in
%   \mymeta{horizontal mode} or \mymeta{vertical mode}.
%    \begin{macrocode}
\cs_new_protected:Nn \@@_multi_addvspace:
  {
    \@@_multi_set_vskip:
    \mode_if_vertical:T
      {
        \skip_add:cn { l_@@_multicols_above_ \@@_level: _skip }
          {
            \skip_use:c { l_@@_partopsep_ \@@_level: _skip }
          }
        \skip_add:cn { l_@@_multicols_below_ \@@_level: _skip }
          {
            \skip_use:c { l_@@_partopsep_ \@@_level: _skip }
          }
      }
    \par\nopagebreak
    \addvspace{ \skip_use:c { l_@@_multicols_above_ \@@_level: _skip } }
  }
%    \end{macrocode}
% \end{macro}
%
% \subsubsection{Adjustment of vertical spaces for \env{multicols} in \env{keyans}}
%
% \begin{macro}[int]{\@@_keyans_multi_set_vskip:, \@@_keyans_multi_addvspace:}
%    The function \myhlc|\@@_keyans_multi_set_vskip:| will take care of determining the
%    \emph{\enquote{adjusted spaces}} applied \emph{\enquote{above}} and
%    \emph{\enquote{below}} the \myenv[no-index]{multicols} environment in
%    \myenv*[no-index]{keyans}. The implementation of this function is the same
%    as the one used in \myenv*[no-index]{enumext}.
%    \begin{macrocode}
\cs_new_protected:Nn \@@_keyans_multi_set_vskip:
  {
    \skip_set:Nn \l_@@_multicols_above_v_skip
      {
        \l_@@_topsep_v_skip
      }
    \skip_set:Nn \l_@@_multicols_below_v_skip
      {
        \l_@@_topsep_v_skip
      }
  }
\cs_new_protected:Nn \@@_keyans_multi_addvspace:
  {
    \@@_keyans_multi_set_vskip:
    \mode_if_vertical:T
      {
        \skip_add:Nn \l_@@_multicols_above_v_skip
          {
            \skip_use:N \l_@@_partopsep_v_skip
          }
        \skip_add:Nn \l_@@_multicols_below_v_skip
          {
            \skip_use:N \l_@@_partopsep_v_skip
          }
      }
    \par\nopagebreak
    \addvspace{ \l_@@_multicols_above_v_skip }
  }
%    \end{macrocode}
% \end{macro}
%
% \subsection{Adjustment of vertical spaces for \env{minipage}}\label{subsec:vspace-minipage}
%
% When nesting a \emph{\enquote{list environment}} within the \myenv[no-index]{minipage} environment,
% the values of the \emph{\enquote{vertical spaces}} are lost. Graphically it can
% be seen like in the figure~\ref{fig:seven}.
%
% \begin{figure}[!htbp]
% \centering
% \begin{tikzpicture}[line cap=round,line join=round,x=0.8cm,y=0.8cm,every node/.style={font=\footnotesize}]
%  \draw[color=white] (7.0,6.10) rectangle (12.5,6.5) node[pos=.5,text=black]%
%   { \textcolor{blue}{\texttt{space above left}} };% space above left
% \draw[color=white] (13.5,6.10) rectangle (17.5,6.5) node[pos=.5,text=black]%
%   { \textcolor{blue}{\texttt{space above right}} };% space above right
% \draw[color=white] (7.0,4.85) rectangle (17.5,5.25) node[pos=.5,text=black]%
%   { \textcolor{blue}{\texttt{space below minipage's}} };% space below after second
% \draw[color=white] (2.0,6.10) rectangle (5,6.50) node[pos=.5,text=keyname,font=\footnotesize\ttfamily] { topsep };% topsep above
% \draw[color=white] (2.0,4.85) rectangle (5,5.25) node[pos=.5,text=keyname,font=\footnotesize\ttfamily] { topsep };% topsep below
% \draw[->,color=lightgray,shorten >=3pt] (4.5,6.3) -- (7.0,6.3); % glue space = topsep (above)
% \draw[->,color=lightgray,shorten >=3pt] (4.5,5.05) -- (7.0,5.05); % glue space = topsep (above)
% \draw[<->,color=gray] (2.0,7.0) --node[above,text=keyname]{\vphantom{p}\texttt{labelwidth}} (5.0,7.0);
% \draw[<->,color=gray] (5.0,7.0) --node[above,text=keyname]{\texttt{labelsep}} (7.0,7.0);
% \draw[color=lightgray,fill=lightgray,fill opacity=0.15] (2.0,6.5) rectangle (5.0,6.85);% labelwidth
% \draw[color=lightgray,fill=lightgray,fill opacity=0.15] (5.0,6.5) rectangle (7.0,6.85);% labelsep
% \draw[color=lightgray,fill=lightgray,fill opacity=0.15] (7.0,6.5) rectangle (17.5,6.85);% item width
% \draw[color=lightgray,fill=lightgray,fill opacity=0.15] (7.0,5.25) rectangle (12.5,6.10) node[pos=.5,text=black]%
%   {\begin{tabular}{c}
%      \texttt{minipage left [t]} \\
%      \textcolor{blue}{\emph{nested list}}
%    \end{tabular}};
% \draw[color=lightgray,fill=lightgray,fill opacity=0.15] (13.5,5.25) rectangle (17.5,6.10) node[pos=.5,text=black]%
%  {\begin{tabular}{c}
%      \texttt{minipage right [t]} \\
%      \textcolor{blue}{\emph{drawing or tabular}}
%    \end{tabular}};
% \draw[color=lightgray,fill=lightgray,fill opacity=0.15] (2.0,4.5) rectangle (5.0,4.85);% labelwidth
% \draw[color=lightgray,fill=lightgray,fill opacity=0.15] (5.0,4.5) rectangle (7.0,4.85);% labelsep
% \draw[color=lightgray,fill=lightgray,fill opacity=0.15] (7.0,4.5) rectangle (17.5,4.85);% item width
% \end{tikzpicture}
% \caption{Representation of the \myenv[no-index]{minipage} spacing adjustment for a nested level.}
% \label{fig:seven}
% \end{figure}
%
% To keep the \emph{\enquote{left}} and \emph{\enquote{right}}
% environments \emph{\enquote{aligned on top}}, preserving the |\baselineskip|
% and keeping the desired \emph{\enquote{spaces}} (|\topsep + [\partopsep]|) it is necessary
% to \emph{\enquote{adjust}} the \emph{\enquote{vertical spaces}}
% for \myenv[no-index]{minipage} environments.
%
% Here there are several complications that must be circumvented: the
% \myenv[no-index]{minipage} environment eliminates the \enquote{top} spaces, the \myenv[no-index]{multicols}
% environment can be nested in the \myenv[no-index]{minipage} environment, the \enquote{top} and
% \enquote{bottom} spaces are affected when |topsep=|\mydim{0pt} and to this is added the
% |\partopsep| parameter that comes into action according to whether \hologo{TeX}
% is in \mymeta{horizontal mode} or \mymeta{vertical mode}. Depending on these cases, small
% adjustments must be made using |\vspace| and |\addvspace| to
% obtain the \emph{\enquote{desired vertical spacing}}.
%
% \begin{important}*
% Again, it must be made clear that the implementation here is a \emph{\enquote{bit
% questionable}}, but hunting the spaces (|glue|) produced by the \myenv[no-index]{minipage}
% environment is quite complicated, even more if \myenv[no-index]{multicols} it is nested.
% The setting of the values was more \emph{\enquote{trial and error}}
% (approx to |\strutbox|), using the help of the \mypkg{lua-visual-debug}\cite{luavisual}
% package, again the attempts to find the correct values using |\showoutput| and
% |\showboxdepth| absolutely failed.
% \end{important}
%
% \subsubsection{Definition of error messages}
%
% \begin{macro}[int]{wrong-miniright-place, wrong-miniright-use,
%                    wrong-miniright-starred, many-miniright-used}
% Messages used by |\miniright| command.
%    \begin{macrocode}
\msg_new:nnnn { enumext } { wrong-miniright-place }
  {
    Wrong~place~for~'\token_to_str:N \miniright'~\msg_line_context:.
  }
  {
    Works~in~'enumext'~and~'keyans'~with~key~'mini-env'.
  }
\msg_new:nnnn { enumext } { wrong-miniright-use }
  {
    Wrong~use~for~'\token_to_str:N \miniright'~\msg_line_context:.
  }
  {
    '\token_to_str:N \miniright'~needs~a~key~'mini-env'.
  }
\msg_new:nnn { enumext } { wrong-miniright-starred }
  {
    Can't~use~'\token_to_str:N \miniright'~in~starred~environments~
    \msg_line_context:.
  }
\msg_new:nnn { enumext } { many-miniright-used }
  {
    Can't~use~'\token_to_str:N \miniright'~more~than~once~
    \msg_line_context:.
  }
%    \end{macrocode}
% \end{macro}
%
% \subsubsection{Adjustment of vertical spaces for \env{minipage} in \env{enumext}}
%
% \begin{macro}[int]{\@@_minipage_set_skip:, \@@_minipage_add_space:}
%    The function \myhlc|\@@_minipage_set_skip:| will take care of determining the
%    \emph{\enquote{adjust}} spaces applied \emph{\enquote{above}} and
%    \emph{\enquote{below}} the \myhlc|\@@_mini_page| environment in
%    \myenv*[no-index]{enumext}.
%
%    First it sets the value of \myhlc|\l_@@_minipage_right_skip| equal to
%    |\topsep|, then checks if \hologo{TeX} is in \mymeta{vertical mode} and, if so,
%    adds |\partopsep|, followed by setting the value of \myhlc|\l_@@_minipage_after_skip|.
%    \begin{macrocode}
\cs_new_protected:Nn \@@_minipage_set_skip:
  {
    \skip_set:Nn \l_@@_minipage_right_skip
      {
        \skip_use:c { l_@@_topsep_ \@@_level: _skip }
      }
    \mode_if_vertical:T
      {
        \skip_add:Nn \l_@@_minipage_right_skip
          {
            \skip_use:c { l_@@_partopsep_ \@@_level: _skip }
          }
      }
    \skip_set_eq:NN \l_@@_minipage_after_skip \l_@@_minipage_right_skip
%    \end{macrocode}
%   Adjusts the values \myhlc|\l_@@_multicols_above_X_skip|
%   and \myhlc|\l_@@_multicols_below_X_skip| and calls the function
%   \myhlc|\@@_pre_itemsep_skip:|.
%    \begin{macrocode}
    \skip_set_eq:cN
      { l_@@_multicols_above_ \@@_level: _skip } \l_@@_minipage_right_skip
    \skip_set_eq:cN
      { l_@@_multicols_below_ \@@_level: _skip } \l_@@_minipage_right_skip
    \@@_pre_itemsep_skip:
%    \end{macrocode}
%   If the environment \myenv[no-index]{multicols} is active, it sets |\topskip=|\mydim{0pt}
%   and then makes |\multicolsep| have the same value as \myhlc|\l_@@_multicols_above_X_skip|.
%    \begin{macrocode}
    \int_compare:nNnT
      { \int_use:c { l_@@_columns_ \@@_level: _int } } > { 1 }
      {
        \skip_zero:N \topskip
        \skip_set_eq:Nc \multicolsep { l_@@_multicols_above_ \@@_level: _skip }
      }
  }
%    \end{macrocode}
%   The function \myhlc|\@@_minipage_add_space:| will apply the spaces on the \emph{\enquote{left side}} using
%   |\addvspace| \emph{\enquote{above}} the \myhlc|\@@_mini_page|
%   environment, taking into account whether \hologo{TeX} is in
%   \mymeta{horizontal mode} or \mymeta{vertical mode}. Here the plain \hologo{TeX} macro
%   |\nointerlineskip| is used to prevent baseline \emph{\enquote{glue}} being
%   added between the next pair of boxes in a \emph{vertical list}. For the latter
%   some adjustments are made since the |\partopsep| parameter comes into
%   play and this affects the \emph{vertical spacing}.
%    \begin{macrocode}
\cs_new_protected:Nn \@@_minipage_add_space:
  {
    \@@_minipage_set_skip:
    \@@_unskip_unkern:
    \mode_if_vertical:TF
      {
        \nopagebreak\nointerlineskip
      }
      {
        \par\nopagebreak\nointerlineskip
        \skip_zero:c { l_@@_partopsep_ \@@_level: _skip }
      }
    \int_compare:nNnTF
      { \int_use:c { l_@@_columns_ \@@_level: _int } } > { 1 }
      {
        \addvspace{ 0.445\box_ht:N \strutbox }
      }
      {
        \addvspace{ 0.250\box_ht:N \strutbox }
      }
  }
%    \end{macrocode}
% \end{macro}
%
% \begin{macro}[int]{\@@_pre_itemsep_skip:}
%   The function \myhlc|\@@_pre_itemsep_skip:| will adjust the spaces below the
%   environment \myenv[no-index]{minipage} and the environment \myenv[no-index]{multicols} if
%   it is nested in it, taking into account the value of |\itemsep| from
%   the previous level.
%    \begin{macrocode}
\cs_new_protected:Nn \@@_pre_itemsep_skip:
  {
    \int_case:nn { \l_@@_level_int }
      {
        { 2 }
          {
            \skip_if_eq:nnTF
              { \l_@@_itemsep_i_skip } { \l_@@_minipage_after_skip }
              {
                \skip_set:Nn \l_@@_minipage_after_skip
                  { 0.150\box_ht:N \strutbox }
                \skip_set:Nn \l_@@_multicols_below_ii_skip
                  { 0.350\box_ht:N \strutbox }
              }
              {
                \dim_compare:nNnT
                  { \l_@@_itemsep_i_skip } < { \l_@@_minipage_after_skip }
                  {
                    \skip_sub:Nn
                      \l_@@_minipage_after_skip { \l_@@_itemsep_i_skip }
                    \skip_sub:Nn
                      \l_@@_multicols_below_ii_skip { \l_@@_itemsep_i_skip }
                    \skip_add:Nn
                      \l_@@_minipage_after_skip { 0.150\box_ht:N \strutbox }
                    \skip_add:Nn
                      \l_@@_multicols_below_ii_skip
                        { 0.350\box_ht:N \strutbox }
                  }
                \dim_compare:nNnT
                  { \l_@@_itemsep_i_skip } > { \l_@@_minipage_after_skip }
                  {
                    \skip_set:Nn \l_@@_minipage_temp_skip
                      {
                        \l_@@_itemsep_i_skip - \l_@@_minipage_after_skip
                      }
                    \skip_sub:Nn
                      \l_@@_minipage_after_skip { \l_@@_itemsep_i_skip }
                    \skip_sub:Nn
                      \l_@@_multicols_below_ii_skip { \l_@@_itemsep_i_skip }
                    \skip_add:Nn
                      \l_@@_minipage_after_skip
                      { 0.150\box_ht:N \strutbox + \l_@@_minipage_temp_skip }
                    \skip_add:Nn
                      \l_@@_multicols_below_ii_skip
                      { 0.350\box_ht:N \strutbox + \l_@@_minipage_temp_skip }
                  }
              }
          }
        { 3 }
          {
            \skip_if_eq:nnTF
              { \l_@@_itemsep_ii_skip } { \c_zero_skip }
              {
                \skip_set:Nn \l_@@_minipage_after_skip
                  { 0.150\box_ht:N \strutbox }
                \skip_set:Nn \l_@@_multicols_below_iii_skip
                  { 0.350\box_ht:N \strutbox }
              }
              {
                \dim_compare:nNnT
                  { \l_@@_itemsep_ii_skip } < { \l_@@_minipage_after_skip }
                  {
                    \skip_sub:Nn
                      \l_@@_minipage_after_skip { \l_@@_itemsep_ii_skip }
                    \skip_sub:Nn
                      \l_@@_multicols_below_iii_skip
                        { \l_@@_itemsep_ii_skip }
                    \skip_add:Nn
                      \l_@@_minipage_after_skip { 0.150\box_ht:N \strutbox }
                    \skip_add:Nn
                      \l_@@_multicols_below_iii_skip
                        { 0.350\box_ht:N \strutbox }
                  }
                \dim_compare:nNnT
                  { \l_@@_itemsep_ii_skip } > { \l_@@_minipage_after_skip }
                  {
                    \skip_set:Nn \l_@@_minipage_temp_skip
                      {
                        \l_@@_itemsep_ii_skip - \l_@@_minipage_after_skip
                      }
                    \skip_sub:Nn
                      \l_@@_minipage_after_skip { \l_@@_itemsep_ii_skip }
                    \skip_sub:Nn
                      \l_@@_multicols_below_iii_skip
                        { \l_@@_itemsep_ii_skip }
                    \skip_add:Nn
                      \l_@@_minipage_after_skip
                      { 0.150\box_ht:N \strutbox + \l_@@_minipage_temp_skip }
                    \skip_add:Nn
                      \l_@@_multicols_below_iii_skip
                      { 0.350\box_ht:N \strutbox + \l_@@_minipage_temp_skip }
                  }
              }
          }
        { 4 }
          {
            \skip_if_eq:nnTF { \l_@@_itemsep_iii_skip } { \c_zero_skip }
              {
                \skip_set:Nn \l_@@_minipage_after_skip
                  { 0.150\box_ht:N \strutbox }
                \skip_set:Nn \l_@@_multicols_below_iv_skip
                  { 0.350\box_ht:N \strutbox }
              }
              {
                \dim_compare:nNnT
                  { \l_@@_itemsep_iii_skip } < { \l_@@_minipage_after_skip }
                  {
                    \skip_sub:Nn
                      \l_@@_minipage_after_skip { \l_@@_itemsep_iii_skip }
                    \skip_sub:Nn
                      \l_@@_multicols_below_iv_skip
                        { \l_@@_itemsep_iii_skip }
                    \skip_add:Nn
                      \l_@@_minipage_after_skip { 0.150\box_ht:N \strutbox }
                    \skip_add:Nn
                      \l_@@_multicols_below_iv_skip
                        { 0.350\box_ht:N \strutbox }
                  }
                \dim_compare:nNnT
                  { \l_@@_itemsep_iii_skip } > { \l_@@_minipage_after_skip }
                  {
                    \skip_set:Nn \l_@@_minipage_temp_skip
                      {
                        \l_@@_itemsep_iii_skip - \l_@@_minipage_after_skip
                      }
                    \skip_sub:Nn
                      \l_@@_minipage_after_skip { \l_@@_itemsep_iii_skip }
                    \skip_sub:Nn
                      \l_@@_multicols_below_iv_skip
                        { \l_@@_itemsep_iii_skip }
                    \skip_add:Nn
                      \l_@@_minipage_after_skip
                      { 0.150\box_ht:N \strutbox + \l_@@_minipage_temp_skip }
                    \skip_add:Nn
                      \l_@@_multicols_below_iv_skip
                      { 0.350\box_ht:N \strutbox + \l_@@_minipage_temp_skip }
                  }
              }
          }
      }
  }
%    \end{macrocode}
% \end{macro}
%
% \subsubsection{Adjustment of vertical spaces for \env{minipage} in \env{keyans}}
%
% \begin{macro}[int]{\@@_keyans_minipage_set_skip:, \@@_keyans_minipage_add_space:, \@@_keyans_pre_itemsep_skip:}
%    The function \myhlc|\@@_keyans_mini_set_vskip:| will take care of determining the
%    \enquote{adjusted} spaces applied \emph{\enquote{above}} and
%    \emph{\enquote{below}} the \myhlc|\@@_mini_page| environment in
%    \myenv*[no-index]{keyans}. The implementation of this function is the same
%    as the one used in \myenv*[no-index]{enumext}.
%    \begin{macrocode}
\cs_new_protected:Nn \@@_keyans_minipage_set_skip:
  {
    \skip_zero:N \l_@@_minipage_after_skip
    \skip_zero:N \l_@@_minipage_left_skip
    \skip_zero:N \l_@@_minipage_right_skip
    \skip_set:Nn \l_@@_minipage_right_skip
      {
        \l_@@_topsep_v_skip
      }
    \mode_if_vertical:T
      {
        \skip_add:Nn \l_@@_minipage_right_skip
          {
            \l_@@_partopsep_v_skip
          }
      }
    \skip_set_eq:NN \l_@@_minipage_after_skip \l_@@_minipage_right_skip
    \skip_set_eq:NN \l_@@_multicols_above_v_skip \l_@@_minipage_right_skip
    \skip_set_eq:NN \l_@@_multicols_below_v_skip \l_@@_minipage_right_skip
    \@@_keyans_pre_itemsep_skip:
    \int_compare:nNnT { \l_@@_columns_v_int } > { 1 }
      {
        \skip_zero:N \topskip
        \skip_set_eq:NN \multicolsep \l_@@_minipage_right_skip
      }
  }
\cs_new_protected:Nn \@@_keyans_minipage_add_space:
  {
    \@@_keyans_minipage_set_skip:
    \@@_unskip_unkern:
    \mode_if_vertical:TF
      {
        \nopagebreak\nointerlineskip
      }
      {
        \par\nopagebreak\nointerlineskip
        \skip_zero:N \l_@@_partopsep_v_skip
      }
    \int_compare:nNnTF { \l_@@_columns_v_int } > { 1 }
      {
        \addvspace{ 0.445\box_ht:N \strutbox }
      }
      {
        \addvspace{ 0.250\box_ht:N \strutbox }
      }
  }
\cs_new_protected:Nn \@@_keyans_pre_itemsep_skip:
  {
    \skip_if_eq:nnTF
      { \l_@@_itemsep_i_skip } { \l_@@_minipage_after_skip }
      {
        \skip_set:Nn \l_@@_minipage_after_skip { 0.150\box_ht:N \strutbox }
        \skip_set:Nn \l_@@_multicols_below_v_skip { 0.350\box_ht:N \strutbox }
      }
      {
        \dim_compare:nNnT
          { \l_@@_itemsep_i_skip } < { \l_@@_minipage_after_skip }
          {
            \skip_sub:Nn \l_@@_minipage_after_skip { \l_@@_itemsep_i_skip }
            \skip_sub:Nn \l_@@_multicols_below_v_skip { \l_@@_itemsep_i_skip }
            \skip_add:Nn \l_@@_minipage_after_skip { 0.150\box_ht:N \strutbox }
            \skip_add:Nn \l_@@_multicols_below_v_skip
              { 0.350\box_ht:N \strutbox }
          }
        \dim_compare:nNnT
          { \l_@@_itemsep_i_skip } > { \l_@@_minipage_after_skip }
          {
            \skip_set:Nn \l_@@_minipage_temp_skip
              {
                \l_@@_itemsep_i_skip - \l_@@_minipage_after_skip
              }
            \skip_sub:Nn \l_@@_minipage_after_skip { \l_@@_itemsep_i_skip }
            \skip_sub:Nn \l_@@_multicols_below_v_skip { \l_@@_itemsep_i_skip }
            \skip_add:Nn \l_@@_minipage_after_skip
              { 0.150\box_ht:N \strutbox + \l_@@_minipage_temp_skip }
            \skip_add:Nn \l_@@_multicols_below_v_skip
              { 0.350\box_ht:N \strutbox + \l_@@_minipage_temp_skip }
          }
      }
  }
%    \end{macrocode}
% \end{macro}
%
% \subsubsection{Adjustment of vertical spaces for \env{minipage} in \env{enumext*} and \env{keyans*}}
%
% \begin{macro}[int]{\@@_mini_set_vskip_vii:, \@@_mini_set_vskip_viii:}
%    The functions \myhlc|\@@_mini_set_vskip_vii:| and \myhlc|\@@_mini_set_vskip_viii:| will take care of determining the
%    \enquote{adjusted} spaces applied \emph{\enquote{above}} and
%    \emph{\enquote{below}} the \myhlc|\@@_mini_page| environment in
%    \myenv*[no-index]{enumext*} and \myenv*[no-index]{keyans*}.
%    \begin{macrocode}
\cs_new_protected:Nn \@@_mini_set_vskip_vii:
  {
    \skip_zero:N \l_@@_minipage_left_skip
    \skip_gzero:N \g_@@_minipage_right_skip
    \skip_gzero:N \g_@@_minipage_after_skip
    \skip_if_eq:nnTF { \l_@@_topsep_vii_skip } { \c_zero_skip }
      {
        \skip_set:Nn \l_@@_minipage_left_skip { 0.5\box_dp:N \strutbox }
        \skip_gset:Nn \g_@@_minipage_right_skip { 0.325\box_dp:N \strutbox }
      }
      {
        \skip_set:Nn \l_@@_minipage_left_skip { 0.5875\box_dp:N \strutbox }
        \skip_gset:Nn \g_@@_minipage_right_skip
          {
            \l_@@_topsep_vii_skip
          }
        \skip_gset:Nn \g_@@_minipage_after_skip
          {
            0.325\box_dp:N \strutbox + \l_@@_topsep_vii_skip
          }
      }
  }
\cs_new_protected:Nn \@@_mini_set_vskip_viii:
  {
    \skip_zero:N \l_@@_minipage_after_skip
    \skip_zero:N \l_@@_minipage_left_skip
    \skip_zero:N \l_@@_minipage_right_skip
    \skip_if_eq:nnTF { \l_@@_topsep_viii_skip } { \c_zero_skip }
      {
        \skip_set:Nn \l_@@_minipage_left_skip
          {
            0.5\box_dp:N \strutbox
          }
        \skip_set:Nn \l_@@_minipage_right_skip
          {
            \l_@@_partopsep_viii_skip
          }
        \skip_set:Nn \l_@@_minipage_after_skip
          {
            1.6\box_dp:N \strutbox
          }
      }
      {
        \skip_set:Nn \l_@@_minipage_left_skip
          {
            0.5875\box_dp:N \strutbox
          }
        \skip_set:Nn \l_@@_minipage_right_skip
          {
            \l_@@_topsep_viii_skip
          }
        \skip_set:Nn \l_@@_minipage_after_skip
          {
            0.325\box_dp:N \strutbox + \l_@@_topsep_viii_skip
          }
      }
  }
%    \end{macrocode}
% \end{macro}
%
% \begin{macro}[int]{\@@_mini_addvspace_vii:, \@@_mini_addvspace_viii:}
%   The functions \myhlc|\@@_mini_addvspace_vii:| and
%   \myhlc|\@@_mini_addvspace_viii:| will apply the vertical space
%   \emph{\enquote{only above}} the \myhlc|\@@_mini_page| environment on
%   the \emph{left side} when the \mykey[env]{mini-right} key is active in the
%   \myenv*[no-index]{enumext*} and \myenv*[no-index]{keyans*} environments.
%
%   Here it does NOT take into account whether \hologo{TeX} is in
%   \mymeta{horizontal mode} or \mymeta{vertical mode}, since |\partopsep|
%   is equal to \mydim{0pt} in both environments.
%    \begin{macrocode}
\cs_new_protected:Nn \@@_mini_addvspace_vii:
  {
    \@@_mini_set_vskip_vii:
    \par\nopagebreak
    \addvspace { \l_@@_minipage_left_skip }
  }
\cs_new_protected:Nn \@@_mini_addvspace_viii:
  {
    \@@_mini_set_vskip_viii:
    \par\nopagebreak
    \addvspace { \l_@@_minipage_left_skip }
  }
%    \end{macrocode}
% \end{macro}
%
% \subsubsection{The command \cs{miniright}}\label{cmd:miniright}
%
% The command |\miniright| will close the \myhlc|\@@_mini_page|
% environment on the \emph{\enquote{left side}}, open the \myhlc|\@@_mini_page|
% environment on the \emph{\enquote{right side}} adding the \emph{adjusted vertical space}.
% By default |\centering| is added when starting the \emph{\enquote{right side}}
% environment. The \emph{starred argument} `|*|' inhibits the use of |\centering|
% command i.e. the usual \hologo{LaTeX} justification is maintained in
% the \myhlc|\@@_mini_page| on the \emph{\enquote{right side}}.
%
% \begin{macro}{\miniright}
%   First, some checks prevent the command from being executed outside the
%   \myenv*[no-index]{enumext} environment or somewhere inappropriate; then the
%   internal functions are called to execute it in the \myenv*[no-index]{enumext} and
%   \myenv*[no-index]{keyans} environments.
%    \begin{macrocode}
\NewDocumentCommand \miniright { s }
  {
    \int_compare:nNnT { \l_@@_keyans_pic_level_int } = { 1 }
      {
        \msg_error:nnn { enumext } { wrong-miniright-place }
      }
    % outside
    \bool_lazy_and:nnT
      { \int_compare_p:nNn { \l_@@_level_int } = { 0 } }
      { \int_compare_p:nNn { \l_@@_level_h_int } = { 0 } }
      {
        \msg_error:nnn { enumext } { wrong-miniright-place }
      }
    % starred env
    \bool_lazy_and:nnT
      { \bool_if_p:N \g_@@_starred_bool }
      { \bool_not_p:n { \l_@@_standard_bool } }
      {
        \msg_error:nnn { enumext } { wrong-miniright-starred }
      }
    % exec
    \int_compare:nNnTF { \l_@@_keyans_level_int } = { 1 }
      {
        \@@_keyans_mini_right_cmd:n {#1}
      }
      { \@@_mini_right_cmd:n {#1} }
  }
%    \end{macrocode}
% \end{macro}
%
% \begin{macro}[int]{\@@_mini_right_cmd:n}
%    The function \myhlc|\@@_mini_right_cmd:n| takes as argument the \emph{starred} `|*|'
%    of the |\miniright| command in the \myenv*[no-index]{enumext} environment.
%    It checks if the \mykey[env]{mini-env} key is active via the variable
%    \myhlc|\l_@@_minipage_right_X_dim|; if so, it closes the
%    \myenv[no-index]{multicols} environment with the \myhlc|\@@_mini_page|
%    environment on the \emph{\enquote{left side}}, then opens the
%    \myhlc|\@@_mini_page| environment on the \emph{\enquote{right
%    side}}, applies the adjusted \emph{\enquote{vertical spaces}}, followed by
%    adding the |\centering| command when the \emph{starred argument} `|*|' is not
%    present, and sets \myhlc|\g_@@_minipage_stat_int| to zero; otherwise it returns an error.
%    \begin{macrocode}
\cs_new_protected:Nn \@@_mini_right_cmd:n
  {
    \dim_compare:nNnTF
      { \dim_use:c { l_@@_minipage_right_ \@@_level: _dim } } > { \c_zero_dim }
      {
        \@@_multicols_stop:
        \int_compare:nNnT
          { \int_use:c { l_@@_columns_ \@@_level: _int } } = { 1 }
          {
            \par\addvspace{ \l_@@_minipage_after_skip }
          }
        \end@@_mini_page
        \hfill
        \@@_mini_page{ \dim_use:c { l_@@_minipage_right_ \@@_level: _dim } }
          \par\nointerlineskip
          \addvspace { \l_@@_minipage_right_skip }
          \bool_if:nF {#1}
            {
              \centering
            }
          \int_gzero:N \g_@@_minipage_stat_int
      }
      { \msg_error:nnn { enumext } { wrong-miniright-use } }
    % paranoia
    \RenewDocumentCommand \miniright { s }
      {
        \msg_error:nn { enumext } { many-miniright-used }
      }
  }
%    \end{macrocode}
% \end{macro}
%
% \begin{macro}[int]{\@@_keyans_mini_right_cmd:n}
%    The function \myhlc|\@@_keyans_mini_right_cmd:n| takes as argument the \emph{starred} `|*|'
%    of the |\miniright| command in the \myenv*[no-index]{keyans}
%    environment. The implementation of this function is the same as that of
%    the \myhlc|\@@_mini_right_cmd:n| function of the \myenv*[no-index]{enumext}
%    environment.
%    \begin{macrocode}
\cs_new_protected:Nn \@@_keyans_mini_right_cmd:n
  {
    \dim_compare:nNnTF { \l_@@_minipage_right_v_dim } > { \c_zero_dim }
      {
        \@@_keyans_multicols_stop:
        \int_compare:nNnT { \l_@@_columns_v_int } = { 1 }
          {
            \par\addvspace{ \l_@@_minipage_after_skip }
          }
        \end@@_mini_page
        \hfill
        \@@_mini_page{ \l_@@_minipage_right_v_dim }
          \par\nointerlineskip
          \addvspace { \l_@@_minipage_right_skip }
          \bool_if:nF {#1}
            {
              \centering
            }
          \int_gzero:N \g_@@_minipage_stat_int
      }
      { \msg_error:nnn { enumext } { wrong-miniright-use } }
    % paranoia
    \RenewDocumentCommand \miniright { s }
      {
        \msg_error:nn { enumext } { many-miniright-used }
      }
  }
%    \end{macrocode}
% \end{macro}
%
% \subsection{Setting \texttt{above} and \texttt{below} keys}\label{subsec:set-above-below}
%
% While having controlled the \emph{vertical spaces} within
% the \myenv*[no-index]{enumext} and \myenv*[no-index]{keyans} environments when using the
% \mykey{columns} or \mykey{mini-env} keys, sometimes the
% \emph{\enquote{vertical spaces above}} or \emph{\enquote{vertical spaces below}} the environments are
% not as expected and it is necessary to be able to apply a
% \emph{\enquote{fine correction}} to these. Since these \emph{glitches} cannot be
% corrected completely, the best option is to provide a couple of
% \mymeta{keys} dedicated to this purpose, in this case it is best to use
% |\vspace| or |\vspace*| when convenient.
%
% \begin{macro}[int]{above, above*, below, below*}
%     Define \mykey[env]{above}, \mykey[env]{above*}, \mykey[env]{below} and \mykey[env]{below*} keys for
%     \myenv*[no-index]{enumext} and \myenv*[no-index]{keyans} environments.
%    \begin{macrocode}
\cs_set_protected:Nn \@@_tmp:nn
  {
    \keys_define:nn { enumext / #1 }
      {
        above  .skip_set:c = { l_@@_vspace_above_ #2 _skip },
        above  .value_required:n = true,
        above* .code:n =
          {
            \bool_set_true:c { l_@@_vspace_a_star_ #2 _bool }
              \keys_set:nn { enumext / #1 } { above = {##1} }
          },
        above* .value_required:n = true,
        below  .skip_set:c = { l_@@_vspace_below_ #2 _skip },
        below  .value_required:n = true,
        below* .code:n =
          {
            \bool_set_true:c { l_@@_vspace_b_star_ #2 _bool }
              \keys_set:nn { enumext / #1 } { below = {##1} }
          },
        below* .value_required:n = true,
      }
  }
\clist_map_inline:Nn \c_@@_all_envs_clist { \@@_tmp:nn #1 }
%    \end{macrocode}
% \end{macro}
%
% \subsubsection{Functions for \texttt{above} and \texttt{below} keys in \env{enumext}}
%
% \begin{macro}[int]{\@@_vspace_above:}
%    The function \myhlc|\@@_vspace_above:| apply the \emph{vertical space above}
%    the \myenv*[no-index]{enumext} environment set by the \mykey[env]{above*} and \mykey[env]{above} keys.
%    \begin{macrocode}
\cs_new_protected:Nn \@@_vspace_above:
  {
    \skip_if_eq:nnF
      { \skip_use:c { l_@@_vspace_above_ \@@_level: _skip } } { \c_zero_skip }
      {
        \bool_if:cTF { l_@@_vspace_a_star_ \@@_level: _bool }
          {
            \vspace*{ \skip_use:c { l_@@_vspace_above_ \@@_level: _skip } }
          }
          {
            \vspace { \skip_use:c { l_@@_vspace_above_ \@@_level: _skip } }
          }
      }
  }
%    \end{macrocode}
% \end{macro}
%
% \begin{macro}[int]{\@@_vspace_below:}
%    The function \myhlc|\@@_vspace_below:| apply the \emph{vertical space below}
%    the \myenv*[no-index]{enumext} environment set by the \mykey[env]{below*} and \mykey[env]{below} keys.
%    \begin{macrocode}
\cs_new_protected:Nn \@@_vspace_below:
  {
    \skip_if_eq:nnF
      { \skip_use:c { l_@@_vspace_below_ \@@_level: _skip } } { \c_zero_skip }
      {
        \bool_if:cTF { l_@@_vspace_b_star_ \@@_level: _bool }
          {
            \vspace*{ \skip_use:c { l_@@_vspace_below_ \@@_level: _skip } }
          }
          {
            \vspace { \skip_use:c { l_@@_vspace_below_ \@@_level: _skip } }
          }
      }
  }
%    \end{macrocode}
% \end{macro}
%
% \subsubsection{Functions for \texttt{above} and \texttt{below} keys in \env{keyans}}
%
% \begin{macro}[int]{\@@_vspace_above_v:}
%    The function \myhlc|\@@_vspace_above_v:| apply the \emph{vertical space above}
%    the \myenv*[no-index]{keyans} environment set by the \mykey[env]{above} and \mykey[env]{above*} keys.
%    \begin{macrocode}
\cs_new_protected:Nn \@@_vspace_above_v:
  {
    \skip_if_eq:nnF { \l_@@_vspace_above_v_skip } { \c_zero_skip }
      {
        \bool_if:NTF \l_@@_vspace_a_star_v_bool
          {
            \vspace*{ \l_@@_vspace_above_v_skip }
          }
          { \vspace { \l_@@_vspace_above_v_skip } }
      }
  }
%    \end{macrocode}
% \end{macro}
%
% \begin{macro}[int]{\@@_vspace_below_v:}
%    The function \myhlc|\@@_vspace_below_v:| apply the \emph{vertical space below}
%    the \myenv*[no-index]{keyans} environment set by the \mykey[env]{below*} and \mykey[env]{below} keys.
%    \begin{macrocode}
\cs_new_protected:Nn \@@_vspace_below_v:
  {
    \skip_if_eq:nnF { \l_@@_vspace_below_v_skip } { \c_zero_skip }
      {
        \bool_if:NTF \l_@@_vspace_b_star_v_bool
          {
            \vspace*{ \l_@@_vspace_below_v_skip }
          }
          { \vspace { \l_@@_vspace_below_v_skip } }
      }
  }
%    \end{macrocode}
% \end{macro}
%
% \subsubsection{Functions for \texttt{above} and \texttt{below} keys in \env{enumext*} \env{keyans*}}
%
% \begin{macro}[int]{\@@_vspace_above_vii:, \@@_vspace_above_viii:}
%    The functions \myhlc|\@@_vspace_above_vii:| and \myhlc|\@@_vspace_above_viii:| apply the \emph{vertical space above}
%    the \myenv*[no-index]{enumext*} and \myenv*[no-index]{keyans*} environments set by the \mykey[env]{above} and \mykey[env]{above*} keys.
%    \begin{macrocode}
\cs_new_protected:Nn \@@_vspace_above_vii:
  {
    \skip_if_eq:nnF { \l_@@_vspace_above_vii_skip } { \c_zero_skip }
      {
        \bool_if:NTF \l_@@_vspace_a_star_vii_bool
          {
            \vspace*{ \l_@@_vspace_above_vii_skip }
          }
          { \vspace { \l_@@_vspace_above_vii_skip } }
      }
  }
\cs_new_protected:Nn \@@_vspace_above_viii:
  {
    \skip_if_eq:nnF { \l_@@_vspace_above_viii_skip } { \c_zero_skip }
      {
        \bool_if:NTF \l_@@_vspace_a_star_viii_bool
          {
            \vspace*{ \l_@@_vspace_above_viii_skip }
          }
          { \vspace { \l_@@_vspace_above_viii_skip } }
      }
  }
%    \end{macrocode}
% \end{macro}
%
% \begin{macro}[int]{\@@_vspace_below_vii:, \@@_vspace_below_viii:}
%    The functions \myhlc|\@@_vspace_below_vii:| and  \myhlc|\@@_vspace_below_viii:|
%    apply the \emph{vertical space below} the \myenv*[no-index]{enumext*} and
%    \myenv*[no-index]{keyans*} environments set by the \mykey[env]{below*} and \mykey[env]{below} keys.
%    \begin{macrocode}
\cs_new_protected:Nn \@@_vspace_below_vii:
  {
    \skip_if_eq:nnF { \l_@@_vspace_below_vii_skip } { \c_zero_skip }
      {
        \bool_if:NTF \l_@@_vspace_b_star_vii_bool
          {
            \vspace*{ \l_@@_vspace_below_vii_skip }
          }
          { \vspace { \l_@@_vspace_below_vii_skip } }
      }
  }
\cs_new_protected:Nn \@@_vspace_below_viii:
  {
    \skip_if_eq:nnF { \l_@@_vspace_below_viii_skip } { \c_zero_skip }
      {
        \bool_if:NTF \l_@@_vspace_b_star_viii_bool
          {
            \vspace*{ \l_@@_vspace_below_viii_skip }
          }
          { \vspace { \l_@@_vspace_below_viii_skip } }
      }
  }
%    \end{macrocode}
% \end{macro}
%
% \subsection{Setting \texttt{series}, \texttt{resume} and \texttt{resume*} keys}\label{key:series}
%
% The \mykey{series} key is responsible for the whole process of the
% \mykey{resume} and \mykey{resume*} keys. The idea behind this is to be
% able to absorb the \mymeta{keys} passed to the \emph{optional argument}
% of the environments \myenv*[no-index]{enumext} and \myenv*[no-index]{enumext*}, but, discarding
% some specific \mymeta{keys}.
%
% \subsubsection{Definition of error messages}
%
% \begin{macro}[int]{unknown-series-starred, unknown-series-standard}
% Messages used by \mykey[env]{resume} key.
%    \begin{macrocode}
\msg_new:nnn { enumext } { unknown-series-starred }
  {
    The~series~'#1'~for~the~resume~key~does~not~exist~in~the~
    enumext*~environment~\msg_line_context:.
  }
\msg_new:nnn { enumext } { unknown-series-standard }
  {
    The~series~'#1'~for~the~resume~key~does~not~exist~at~level~
    \int_use:N \l_@@_level_int
    \c_space_tl of~enumext~environment~\msg_line_context:.
  }
%    \end{macrocode}
% \end{macro}
%
% \subsubsection{Definition of keys for \texttt{series}, \texttt{resume} and \texttt{resume*}}
%
% \begin{macro}[int]{series, resume, resume*}
%   Defines the keys \mykey[env]{series}, \mykey[env]{resume}
%   and \mykey[env]{resume*} for the \emph{\enquote{all levels}}
%   of \myenv*[no-index]{enumext} and \myenv*[no-index]{enumext*}.
%   It is not necessary to make sure that |\printkeyans| is not running
%   otherwise the \emph{start value} of the environments would be increased when
%   using \mykey{resume} or \mykey{resume}|*| keys.
%
%   In version 1.6 it is allowed to pass the key \mykey{resume} \emph{without
%   value} by means of the command |\setenumext|, for the correct operation
%   of this, the boolean variable \myhlc|\l_@@_resume_count_bool|
%   set by the key \mykey{resume} \emph{without value} must be set to |true| to
%   be later processed by the function \myhlc|\@@_parse_series:n| in the definition
%   of the environments \myenv*[no-index]{enumext} and \myenv*[no-index]{enumext*}.
%    \begin{macrocode}
\cs_set_protected:Nn \@@_tmp:nn
  {
    \keys_define:nn { enumext / #1 }
      {
        series  .str_set:N = \l_@@_series_name_str,
        series  .value_required:n = true,
        resume  .code:n =
          {
            \bool_if:NF \l_@@_print_keyans_cmd_bool
              {
                \tl_set:Nn \l_@@_series_name_tl {##1}
                \tl_if_empty:NTF \l_@@_series_name_tl
                  {
                    \bool_set_true:c { l_@@_resume_count_ #2 _bool }
                    \bool_set_eq:Nc
                      \l_@@_resume_count_bool
                      { l_@@_resume_count_ #2 _bool }
                  }
                  {
                    \@@_resume:n {##1}
                  }
              }
          },
        resume* .code:n =
          {
            \bool_if:NF \l_@@_print_keyans_cmd_bool
              {
                \bool_set_true:c { l_@@_resume_count_ #2 _bool }
                \bool_set_eq:Nc
                  \l_@@_resume_count_bool { l_@@_resume_count_ #2 _bool }
                \bool_set_true:c { l_@@_resume_star_key_ #2 _bool }
                \@@_resume_star:
              }
          },
        resume* .value_forbidden:n = true,
      }
  }
\clist_map_inline:nn
  {
    {level-1}{i}, {level-2}{ii}, {level-3}{iii}, {level-4}{iv},{enumext*}{vii},
  }
  { \@@_tmp:nn #1 }
%    \end{macrocode}
% \end{macro}
%
% \subsubsection{Internal function to save counter and integer values}
%
% \begin{macro}[int]{\@@_standard_save_counter:, \@@_standard_save_counter_aux:,
%                    \@@_starred_save_counter:, \@@_starred_save_counter_aux:}
%   The \myhlc|\@@_standard_save_counter:| and \myhlc|\@@_starred_save_counter:|
%   functions will save the last counter value to \myhlc|\g_@@_series_#1_int|
%   if the \mykey[env]{series}|=|\mymarg{series name} key has been passed, to
%   \myhlc|\c@ @@_resume_X_int| if it has passed the key \mykey[env]{resume}
%   \emph{without value} and the key \mykey{series} is not active, in
%   \myhlc|\g_@@_series_#1_X_int| if the key
%   \mykey{resume}|=|\mymarg{series name} has been passed and in
%   \myhlc|\g_@@_series_#1_X_int| if the key has
%   been passed \mykey[env]{save-ans}|=|\mymarg{store name}.
%
% \begin{important}*
%   The variables \myhlc|\l_@@_series_name_str| and \myhlc|\l_@@_series_name_tl|
%   contain the same \mymarg{series name} but are executed at different
%   moments, the integer variable with \myhlc|\l_@@_series_name_str| sets the value
%   when execute \mykey{series}|=|\mymarg{series name} and the
%   integer variable with \myhlc|\l_@@_series_name_tl| sets the subsequent
%   values when use \mykey{resume}|=|\mymarg{series name}.
%   This function is passed to the \myenv*[no-index]{enumext} environment definition (\S\ref{env:enumext})
%   and the \myenv*[no-index]{enumext*} environment definition (\S\ref{env:starred}).
% \end{important}
%    \begin{macrocode}
\cs_new_protected:Nn \@@_standard_save_counter:
  {
    \bool_if:NTF \g_@@_standard_bool
      {
        \@@_standard_save_counter_aux:
        \int_compare:nNnT { \l_@@_level_int } = { 1 }
          {
            \int_if_exist:cT { g_@@_resume_ \l_@@_store_name_tl _int }
              {
                \int_gset_eq:cN
                  { g_@@_resume_ \l_@@_store_name_tl _int } \value{enumXi}
              }
          }
      }
      {
        \@@_standard_save_counter_aux:
      }
  }
\cs_new_protected:Nn \@@_standard_save_counter_aux:
  {
    \str_if_empty:NF \l_@@_series_name_str
      {
        \int_gset_eq:cc
          { g_@@_series_ \l_@@_series_name_str _ \@@_level: _int }
          { c@enumX \@@_level: }
      }
    \tl_if_empty:NTF \l_@@_series_name_tl
      {
        \str_if_empty:NT \l_@@_series_name_str
          {
            \tl_if_empty:NT \l_@@_store_name_tl
              {
                \int_gset_eq:cc
                  { c@ @@_resume_ \@@_level: _int } { c@enumX \@@_level: }
              }
          }
      }
      {
        \int_if_exist:cT
          { g_@@_series_ \l_@@_series_name_tl _ \@@_level: _int }
          {
            \int_gset_eq:cc
              { g_@@_series_ \l_@@_series_name_tl _ \@@_level: _int }
              { c@enumX \@@_level: }
          }
      }
  }
\cs_new_protected:Nn \@@_starred_save_counter:
  {
    \bool_if:NTF \g_@@_starred_bool
      {
        \@@_starred_save_counter_aux:
        \int_if_exist:cT { g_@@_resume_ \l_@@_store_name_tl _int }
          {
            \int_gset_eq:cN
              { g_@@_resume_ \l_@@_store_name_tl _int } \value{enumXvii}
          }
      }
      {
        \@@_starred_save_counter_aux:
      }
  }
\cs_new_protected:Nn \@@_starred_save_counter_aux:
  {
    \str_if_empty:NF \l_@@_series_name_str
      {
        \int_gset_eq:cN
          { g_@@_series_ \l_@@_series_name_str _vii_int } \value{enumXvii}
      }
    \tl_if_empty:NTF \l_@@_series_name_tl
      {
        \str_if_empty:NT \l_@@_series_name_str
          {
            \tl_if_empty:NT \l_@@_store_name_tl
              {
                \int_gset_eq:cc { c@ @@_resume_vii_int } { c@enumXvii }
              }
          }
      }
      {
        \int_if_exist:cT { g_@@_series_ \l_@@_series_name_tl _vii_int }
          {
            \int_gset_eq:cN
              { g_@@_series_ \l_@@_series_name_tl _vii_int } \value{enumXvii}
          }
      }
  }
%    \end{macrocode}
% \end{macro}
%
% \subsubsection{Internal function for resume counters}
%
% \begin{macro}[int]{\@@_resume_counter:}
%   The \myhlc|\@@_resume_counter:| function is executed by the \mykey[env]{resume*}
%   key and \mykey[env]{resume} key \emph{without value}, only the
%   \emph{\enquote{counters}} for the \emph{\enquote{levels}} of the
%   environments in which it is executed will be set. If the
%   \mykey[env]{save-ans} key is active it will set the \emph{\enquote{counter}}
%   according to the value of the integer variable created by that key.
%    \begin{macrocode}
\cs_new_protected:Nn \@@_resume_counter:
  {
    \cs_set:Nn \@@_tmp:n
      {
        \bool_if:cT { l_@@_resume_count_ \int_to_roman:n {##1} _bool }
          {
            \exp_args:Ne \int_set:cn { l_@@_start_ \int_to_roman:n {##1} _int }
              {
                \int_use:c { c@ @@_resume_ \int_to_roman:n {##1} _int } + 1
              }
          }
      }
    \int_compare:nNnT { \l_@@_level_int } > { 0 }
      {
        \bool_lazy_and:nnTF
          { \bool_if_p:N \l_@@_standard_first_bool }
          { \bool_if_p:N \l_@@_store_active_bool }
          {
            \int_set:Nn \l_@@_start_i_int
              {
                \int_use:c { g_@@_resume_ \l_@@_store_name_tl _int } + 1
              }
          }
          {
            \int_step_function:nN { \l_@@_level_int } \@@_tmp:n
          }
      }
    \int_compare:nNnT { \l_@@_level_h_int } = { 1 }
      {
        \bool_lazy_and:nnTF
          { \bool_if_p:N \l_@@_starred_first_bool }
          { \bool_if_p:N \l_@@_store_active_bool }
          {
            \int_set:Nn \l_@@_start_vii_int
              {
                \int_use:c { g_@@_resume_ \l_@@_store_name_tl _int } + 1
              }
          }
          {
            \bool_if:NT \l_@@_resume_count_vii_bool
              {
                \int_set:Nn \l_@@_start_vii_int
                  {
                    \int_use:c { c@ @@_resume_vii_int } + 1
                  }
              }
          }
      }
  }
%    \end{macrocode}
% \end{macro}
%
% \subsubsection{Internal functions for \texttt{series} key}
%
% \begin{macro}[int]{\@@_filter_series:n, \@@_filter_series_key:n, \@@_filter_series_pair:nn}
%   The function \myhlc|\@@_filter_series:n| will be in charge of filtering the
%   \mymeta{keys} to be stored, where `|#1|' represents the \emph{optional argument}
%   passed to the environment. This implementation is adapted directly from
%   the code provided by Jonathan P. Spratte (\texttt{@Skillmon}) in \href{https://chat.stackexchange.com/transcript/41?m=65647962#65647962}{chat-TeX-SX}
%    \begin{macrocode}
\cs_new:Nn \@@_filter_series:n
  {
    \use:e
      {
        \keyval_parse:NNn
          \@@_filter_series_key:n
          \@@_filter_series_pair:nn {#1}
      }
  }
%    \end{macrocode}
%   The function \myhlc|\@@_filter_series_key:n| will be responsible for
%   filtering the \mymeta{keys} that are passed \emph{\enquote{without value}} by
%   excluding the \mykey[env]{resume}, \mykey[env]{resume*}, \mykey[env]{reset}, \mykey[env]{reset*}
%   and \mykey[env]{base-fix} keys.
%    \begin{macrocode}
\cs_new:Nn \@@_filter_series_key:n
  {
    \str_case:nnF {#1}
      {
        { resume } { } { resume* } { } { reset } { } { reset* } { } { base-fix } { }
      }
      { , { \exp_not:n {#1} } }
  }
%    \end{macrocode}
%   The function \myhlc|\@@_filter_series_pair:nn| will be responsible for
%   filtering the \mymeta{keys} that are passed \emph{\enquote{with value}} by
%   excluding the \mykey[env]{series}, \mykey{resume}, \mykey[env,foreachkeyans]{start}, \mykey[env]{start*}, \mykey[env]{save-ans} and
%   \mykey[env]{save-key} keys.
%    \begin{macrocode}
\cs_new:Nn \@@_filter_series_pair:nn
  {
    \str_case:nnF {#1}
      {
        { series } { } { resume } { } { start } { }
        { start* } { }  { save-ans } { } { save-key } { }
      }
      { , { \exp_not:n {#1} } = { \exp_not:n {#2} } }
  }
%    \end{macrocode}
% \end{macro}
%
% \begin{macro}[int]{\@@_save_last_keys:n, \@@_resume_last_counter:}
%   The function \myhlc|\@@_save_last_keys:n| will be in charge of saving the
%   filtering \mymeta{keys} when the keys \mykey[env]{series}|=|\mymarg{series name}
%   or \mykey[env]{resume}|=|\mymarg{series name} or \mykey{resume}|*| are NOT
%   active and will save them in the variable \myhlc|\g_@@_resume_last_keys_X_tl| for the
%   \myenv*[no-index]{enumext} environment and in the variable \myhlc|\g_@@_resume_last_keys_vii_tl|
%   for the \myenv*[no-index]{enumext*} environment.
%
% \begin{important}*
%   The boolean variable \myhlc|\l_@@_resume_series_X_bool|
%   is set to |true| by the key \mykey{resume}|=|\mymarg{series name},
%   the boolean variable \myhlc|\l_@@_resume_star_key_X_bool|
%   is set to |true| by the key \mykey{resume}|*|, in this case
%   both variables must be set to |false| so that
%   they don't override the default filtered \mymeta{keys}.
% \end{important}
%    \begin{macrocode}
\cs_new_protected:Nn \@@_save_last_keys:n
  {
    \int_compare:nNnT { \l_@@_level_int } > { 0 }
      {
        \bool_if:cF { l_@@_resume_series_ \@@_level: _bool }
          {
            \bool_if:cF { l_@@_resume_star_key_ \@@_level: _bool }
              {
                \tl_gclear:c { g_@@_resume_last_keys_ \@@_level: _tl }
                \tl_gset:ce
                  { g_@@_resume_last_keys_ \@@_level: _tl }
                  { \@@_filter_series:n {#1} }
              }
          }
      }
    \int_compare:nNnT { \l_@@_level_h_int } = { 1 }
      {
        \bool_if:NF \l_@@_resume_series_vii_bool
          {
            \tl_gclear:N \g_@@_resume_last_keys_vii_tl
            \tl_gset:Ne \g_@@_resume_last_keys_vii_tl
              { \@@_filter_series:n {#1} }
          }
      }
  }
%    \end{macrocode}
%   The \myhlc|\@@_resume_last_counter:| function will be in charge of setting
%   the \emph{\enquote{counters}} when the keys \mykey{series}|=|\mymarg{series name}
%   or \mykey{resume}|=|\mymarg{series name} are NOT active and the
%   \mykey{resume} key is being used \emph{without value} either in the
%   \emph{optional argument} of the environments or through the
%   |\setenumext| command.
%
% \begin{important}*
%   The boolean variable \myhlc|\l_@@_resume_count_bool|
%   is set to |true| by the keys \mykey{resume} \emph{without value}
%   and \mykey{resume}|*|; and set to |false| by the keys
%   \mykey[env,foreachkeyans]{start} and \mykey[env]{start*} (\S\ref{key:start-widest}).
% \end{important}
%    \begin{macrocode}
\cs_new_protected:Nn \@@_resume_last_counter:
  {
    \bool_lazy_and:nnT
      { \bool_if_p:N \l_@@_resume_count_bool }
      { \tl_if_empty_p:N \l_@@_series_name_tl }
      {
        \@@_resume_counter:
      }
  }
%    \end{macrocode}
% \end{macro}
%
% \begin{macro}[int]{\@@_parse_series:n}
%   The \myhlc|\@@_parse_series:n| function handled by the \mykey[env]{series} key
%   will be responsible for \emph{storing} the \emph{filtered}
%   \mymeta{keys} from the \emph{optional arguments} of the
%   \myenv*[no-index]{enumext} and \myenv*[no-index]{enumext*} environments for the
%   \mykey[env]{resume} and \mykey{resume}|*| keys.
%   If the \mykey{series} key is NOT active it will call the
%   \myhlc|\@@_save_last_keys:n| function to \emph{store} the \emph{filtered}
%   \mymeta{keys} that will be used by the \mykey[env]{resume*} key and then the
%   \myhlc|\@@_resume_last_counter:| function used by the \mykey{resume} key
%   \emph{without value} if it is active, otherwise \emph{store} the
%   \emph{filtered} \mymeta{keys} in the global variable \myhlc|\g_@@_series_#1_X_tl|
%   along with the creation of the integer variable \myhlc|\g_@@_series_#1_X_int|
%   used by the \mykey{resume} key \emph{with value}.
%
% \begin{important}*
%   This function is passed to the function \myhlc|\@@_parse_keys:n|
%   in the \myenv*[no-index]{enumext} environment definition (\S\ref{env:enumext})
%   and to the function \myhlc|\@@_parse_keys_vii:n| in the \myenv*[no-index]{enumext*}
%   environment definition (\S\ref{env:starred}).
% \end{important}
%    \begin{macrocode}
\cs_new_protected:Nn \@@_parse_series:n
  {
    \str_if_empty:NTF \l_@@_series_name_str
      {
        \int_compare:nNnT { \l_@@_level_int } > { 0 }
          {
            \@@_save_last_keys:n {#1}
            \@@_resume_last_counter:
          }
        \int_compare:nNnT { \l_@@_level_h_int } = { 1 }
          {
            \@@_save_last_keys:n {#1}
            \@@_resume_last_counter:
          }
      }
      {
        \int_compare:nNnT { \l_@@_level_int } > { 0 }
          {
            \tl_gclear_new:c
              { g_@@_series_ \l_@@_series_name_str _ \@@_level: _tl }
            \tl_gset:ce
              { g_@@_series_ \l_@@_series_name_str _ \@@_level: _tl }
              { \@@_filter_series:n {#1} }
            \int_if_exist:cF
              { g_@@_series_ \l_@@_series_name_str _ \@@_level: _int }
              {
                \int_new:c
                  { g_@@_series_ \l_@@_series_name_str _ \@@_level: _int }
              }
          }
        \int_compare:nNnT { \l_@@_level_h_int } = { 1 }
          {
            \tl_gclear_new:c { g_@@_series_ \l_@@_series_name_str _vii_tl }
            \tl_gset:ce
              { g_@@_series_ \l_@@_series_name_str _vii_tl }
              { \@@_filter_series:n {#1} }
            \int_if_exist:cF { g_@@_series_ \l_@@_series_name_str _vii_int }
              {
                \int_new:c { g_@@_series_ \l_@@_series_name_str _vii_int }
              }
          }
      }
  }
%    \end{macrocode}
% \end{macro}
%
% \subsubsection{Internal functions for \texttt{resume} key with value}
%
% \begin{macro}[int]{\@@_resume:n}
%   The function \myhlc|\@@_resume:n| will handle the \emph{argument} \mymarg{series name}
%   passed to the \mykey[env]{resume} key in \myenv*[no-index]{enumext} and \myenv*[no-index]{enumext*} environments.
%   First it checks if the global variable \myhlc|\g_@@_series_#1_X_tl|
%   exists; if so, it calls the function \myhlc|\@@_resume_series:n| and passes the \mymeta{keys}
%   stored in \myhlc|\g_@@_series_#1_X_tl| to the environments,
%   otherwise it returns an error.
%    \begin{macrocode}
\cs_new_protected:Nn \@@_resume:n
  {
    \int_compare:nNnT { \l_@@_level_int } > { 0 }
      {
        \tl_if_exist:cTF { g_@@_series_ \tl_to_str:n {#1} _ \@@_level: _tl }
          {
            \@@_resume_series:n {#1}
            \exp_args:Ne \keys_set:nv { enumext /
              level-\int_use:N \l_@@_level_int }
              { g_@@_series_ \tl_to_str:n {#1} _ \@@_level: _tl }
          }
          {
            \msg_error:nnn { enumext } { unknown-series-standard } {#1}
          }
      }
    \int_compare:nNnT { \l_@@_level_h_int } = { 1 }
      {
        \tl_if_exist:cTF { g_@@_series_ \tl_to_str:n {#1} _vii_tl }
          {
            \@@_resume_series:n {#1}
            \keys_set:nv { enumext / enumext* }
              { g_@@_series_ \tl_to_str:n {#1} _vii_tl }
          }
          {
            \msg_error:nnn { enumext } { unknown-series-starred } {#1}
          }
      }
  }
%    \end{macrocode}
% \end{macro}
%
% \begin{macro}[int]{\@@_resume_series:n, \@@_resume_integer_series:}
%   The function \myhlc|\@@_resume_series:n| will set the variable
%   \myhlc|\l_@@_resume_series_X_bool| to |true| and pass the \mymarg{argument}
%   to the variable \myhlc|\l_@@_series_name_tl| then call the function
%   \myhlc|\@@_resume_integer_series:|.
%    \begin{macrocode}
\cs_new_protected:Nn \@@_resume_series:n
  {
    \int_compare:nNnT { \l_@@_level_int } > { 0 }
      {
        \bool_set_true:c { l_@@_resume_series_ \@@_level: _bool }
        \tl_clear:N \l_@@_series_name_tl
        \tl_set:Nn \l_@@_series_name_tl {#1}
        \@@_resume_integer_series:
      }
    \int_compare:nNnT { \l_@@_level_h_int } = { 1 }
      {
        \bool_set_true:N \l_@@_resume_series_vii_bool
        \tl_clear:N \l_@@_series_name_tl
        \tl_set:Nn \l_@@_series_name_tl {#1}
        \@@_resume_integer_series:
      }
  }
%    \end{macrocode}
%   The function \myhlc|\@@_resume_integer_series:| will be executed when the
%   \mykey[env]{resume}|=|\mymarg{series name} key is active, setting the \emph{start value}
%   for the \emph{\enquote{counter}} of the \emph{\enquote{current level}} of the environments
%   in which it is run according to the value of the \emph{\enquote{integer variables}} created by the
%   \mykey[env]{series} key. If the \mykey[env]{save-ans} key is active it will
%   set the \emph{start value} for the \emph{\enquote{counter}} according
%   to the value of the integer variable created by that key.
%    \begin{macrocode}
\cs_new_protected:Nn \@@_resume_integer_series:
  {
    \cs_set:Nn \@@_tmp:n
      {
        \int_if_exist:cT { g_@@_series_ \l_@@_series_name_tl _
          \int_to_roman:n {##1} _int }
          {
            \exp_args:Ne \int_set:cn { l_@@_start_ \int_to_roman:n {##1} _int }
              {
                \int_use:c { g_@@_series_ \l_@@_series_name_tl _
                  \int_to_roman:n {##1} _int } + 1
              }
          }
      }
    \int_compare:nNnT { \l_@@_level_int } > { 0 }
      {
        \bool_lazy_and:nnTF
          { \bool_if_p:N \l_@@_standard_first_bool }
          { \bool_if_p:N \l_@@_store_active_bool }
          {
            \int_set:Nn \l_@@_start_i_int
              {
                \int_use:c { g_@@_resume_ \l_@@_store_name_tl _int } + 1
              }
          }
          {
            \int_step_function:nN { \l_@@_level_int } \@@_tmp:n
          }
      }
    \int_compare:nNnT { \l_@@_level_h_int } = { 1 }
      {
        \bool_lazy_and:nnTF
          { \bool_if_p:N \l_@@_starred_first_bool }
          { \bool_if_p:N \l_@@_store_active_bool }
          {
            \int_set:Nn \l_@@_start_vii_int
              {
                \int_use:c { g_@@_resume_ \l_@@_store_name_tl _int } + 1
              }
          }
          {
            \int_set:Nn \l_@@_start_vii_int
              {
                \int_use:c { g_@@_series_ \l_@@_series_name_tl _vii_int } + 1
              }
          }
      }
  }
%    \end{macrocode}
% \end{macro}
%
% \subsubsection{Internal function for \texttt{resume*} key}
%
% \begin{macro}[int]{\@@_resume_star:}
%   The function \myhlc|\@@_resume_star:| will handle the \mykey[env]{resume*}
%   key in the \myenv*[no-index]{enumext} and \myenv*[no-index]{enumext*} environments. This
%   function will execute the filtered \mymeta{keys} in the last one and
%   will continue with the numbering and \mymeta{keys} according to the
%   last execution of the environment \myenv*[no-index]{enumext} or \myenv*[no-index]{enumext*}
%   in which the keys \mykey[env]{resume}|=|\mymarg{series name} or \mykey[env]{series}|=|\mymarg{series name}
%   were NOT active.
%    \begin{macrocode}
\cs_new_protected:Nn \@@_resume_star:
  {
    \cs_set:Nn \@@_tmp:n
      {
        \tl_if_empty:cTF { g_@@_resume_last_keys_ \int_to_roman:n {##1} _tl }
          {
            \@@_resume_counter:
          }
          {
            \@@_resume_counter:
            \exp_args:Ne \keys_set:nv
              { enumext / level-\int_use:N \l_@@_level_int }
              { g_@@_resume_last_keys_ \int_to_roman:n {##1} _tl }
          }
      }
    \int_compare:nNnT { \l_@@_level_int } > { 0 }
      {
        \int_step_function:nN { \l_@@_level_int } \@@_tmp:n
      }
    \int_compare:nNnT { \l_@@_level_h_int } = { 1 }
      {
        \tl_if_empty:NTF \g_@@_resume_last_keys_vii_tl
          {
            \@@_resume_counter:
          }
          {
            \@@_resume_counter:
            \keys_set:nV { enumext / enumext* } \g_@@_resume_last_keys_vii_tl
          }
      }
  }
%    \end{macrocode}
% \end{macro}
%
% \subsection{The \cs{resetenumext} command}\label{cmd:resetenumext}
%
% Sometimes it is necessary to be able to reset the \emph{\enquote{counters}}
% of the environments according to some value, for example |\chapter|. Since
% \emph{\enquote{internal counters}} are used for the \mykey{resume} and \mykey{resume*}
% keys which set the \emph{start value}, but are not accessible by the user,
% it is necessary to provide a public command for this. This implementation is an
% adaptation of the answers given by Clea F. Rees (\texttt{@cfr}) and
% Jonathan P. Spratte (\texttt{@Skillmon}) in \href{https://tex.stackexchange.com/q/746644}{Correct
% implementation of optional argument (comma-separated) in expl3}.
%
% \begin{macro}[int]{out-of-range}
% Message used by |\resetenumext| command.
%    \begin{macrocode}
\msg_new:nnnn { enumext } { out-of-range }
  {
    The~number~must~be~exactly~1,~2,~3~or~4.
  }
  {
    Received:~'#1'.
  }
%    \end{macrocode}
% \end{macro}
%
% \begin{macro}[int]{\@@_standard_reset:nn, \@@_starred_reset:n,
%                    \@@_reset_count_resume:nn, \@@_reset_count_resume:en,
%                    \@@_reset_count_resume_all:n, \@@_reset_count_resume_levels:n}
%   Internal functions used by the \ics*{resetenumext} command.
%    \begin{macrocode}
\cs_new_protected:Npn \@@_standard_reset:nn #1 %#2
  {
    \@@_reset_count_resume:en { \int_to_roman:n {#1} } %{#2}
  }
\cs_new_protected:Nn \@@_starred_reset:n
  {
    \@@_reset_count_resume:nn { vii } {#1}
  }
\cs_new_protected:Nn \@@_reset_count_resume:nn
  {
    \counterwithin*{enumX #1}{#2}
    \counterwithin*{@@_resume_ #1 _int}{#2}
  }
\cs_generate_variant:Nn \@@_reset_count_resume:nn { e }
\cs_new_protected:Nn \@@_reset_count_resume_all:n
  {
    \clist_map_inline:nn { i,ii,iii,iv,vii }
      {
        \@@_reset_count_resume:nn {##1} {#1}
      }
  }
\cs_new_protected:Nn \@@_reset_count_resume_levels:n
  {
    \clist_map_inline:nn { i,ii,iii,iv }
      {
        \@@_reset_count_resume:nn {##1} {#1}
      }
  }
%    \end{macrocode}
% \end{macro}
%
% \begin{macro}{\resetenumext}
%   The \ics*{resetenumext} command \emph{\enquote{resets}} the \emph{start value}
%   of the \emph{\enquote{counters}} for the \myenv*[no-index]{enumext} and \myenv*[no-index]{enumext*}
%   environments along with the \emph{\enquote{internal counters}} used by
%   the keys \mykey[env]{resume} \emph{without value} and \mykey[env]{resume*}
%   according to the value of \mymarg{some counter}.
%    \begin{macrocode}
\NewDocumentCommand \resetenumext { s o m }
  {
    \bool_if:nTF {#1}
      {
        \@@_reset_count_resume_all:n {#3}
      }
      {
        \tl_if_novalue:nTF {#2}
          {
            \@@_reset_count_resume_levels:n {#3}
          }
          {
            \str_if_eq:nnTF {#2} { * }
              { \@@_starred_reset:n {#3} }
              {
                \bool_lazy_and:nnTF
                  { \int_compare_p:nNn {#2} > { 0 } }
                  { \int_compare_p:nNn {#2} < { 5 } }
                  { \@@_standard_reset:nn {#2} {#3} }
                  {
                    \msg_error:nne { enumext } { out-of-range }
                      { \int_eval:n {#2} }
                  }
              }
          }
      }
  }
%    \end{macrocode}
% \end{macro}
%
% \subsection{The \texttt{reset} and \texttt{reset*} keys}\label{key:reset}
%
% The |\resetenumext| command does not work, for example, after an
% unnumbered chapter, so it is preferable to provide a pair of \mymeta{keys} that
% adjust the internal variables if necessary.
%
% \begin{macro}[int]{reset, reset*}
%     Defines the keys \mykey[env]{reset} and \mykey[env]{reset*} for the \emph{\enquote{all levels}}
%     of \myenv*[no-index]{enumext} and \myenv*[no-index]{enumext*}.
%    \begin{macrocode}
\cs_set_protected:Nn \@@_tmp:n
  {
    \keys_define:nn { enumext / #1 }
      {
        reset  .code:n = { \@@_standard_reset_key: },
        reset  .value_forbidden:n = true,
        reset* .code:n = { \@@_standard_reset_key_star: },
        reset* .value_forbidden:n = true,
      }
  }
\clist_map_inline:nn {level-1, level-2, level-3, level-4} { \@@_tmp:n {#1} }
\keys_define:nn { enumext / enumext* }
  {
    reset  .code:n = { \@@_starred_reset_key: },
    reset  .value_forbidden:n = true,
    reset* .code:n = { \@@_starred_reset_key: },
    reset* .value_forbidden:n = true,
  }
%    \end{macrocode}
% \end{macro}
%
% \subsubsection{Internal functions for \texttt{reset} and \texttt{reset*} keys}
%
% \begin{macro}[int]{\@@_standard_reset_key:, \@@_standard_reset_key_star:, \@@_starred_reset_key:}
%   The function \myhlc|\@@_standard_reset_key:| will be handled by the
%   \mykey[env]{reset} key and will \emph{\enquote{reset}} the counter \myhlc|\c@ @@_resume_X_int|
%   to \emph{\enquote{zero}} according to the \emph{level} at
%   which it is executed within the \myenv*[no-index]{enumext} environment.
%    \begin{macrocode}
\cs_new_protected:Nn \@@_standard_reset_key:
  {
    \int_compare:nNnT { \l_@@_level_int } > { 0 }
      {
        \int_if_exist:cT { c@ @@_resume_ \int_to_roman:n { \l_@@_level_int }
          _int }
          {
            \int_gzero:c { c@ @@_resume_ \int_to_roman:n { \l_@@_level_int }
              _int }
          }
      }
  }
%    \end{macrocode}
%   The function \myhlc|\@@_standard_reset_key_star:| will be handled by the
%   \mykey[env]{reset*} key and will \emph{\enquote{reset}} the counters \myhlc|\c@ @@_resume_X_int|
%   to \emph{\enquote{zero}} from the \emph{level} at which it
%   is executed within the \myenv*[no-index]{enumext} environment to the \emph{lower levels}.
%    \begin{macrocode}
\cs_new_protected:Nn \@@_standard_reset_key_star:
  {
    \cs_set:Nn \@@_tmp:n
      {
        \int_if_exist:cT { c@ @@_resume_ \int_to_roman:n {##1} _int }
          {
            \int_gzero:c { c@ @@_resume_ \int_to_roman:n {##1} _int }
          }
      }
    \int_compare:nNnT { \l_@@_level_int } > { 0 }
      {
        \int_step_function:nnN { \l_@@_level_int } { 4 } \@@_tmp:n
      }
  }
%    \end{macrocode}
%   The function \myhlc|\@@_starred_reset_key:| will be handled by
%   \mykey{reset} keys and \mykey{reset*} will \emph{\enquote{reset}} the
%   counter \myhlc|\c@ @@_resume_vii_int| to \emph{\enquote{zero}} when
%   executed in the \myenv*[no-index]{enumext*} environment.
%    \begin{macrocode}
\cs_new_protected:Nn \@@_starred_reset_key:
  {
    \int_gzero:c { c@ @@_resume_vii_int }
  }
%    \end{macrocode}
% \end{macro}
%
% \subsection{Setting \texttt{save-ans}, \texttt{check-ans} and  \texttt{no-store} keys}\label{subsec:set-save-check-nostore}
%
% The key \mykey{save-ans} is directly associated with the keys
% \mykey{check-ans}, \mykey{no-store}, \mykey{resume} and
% \mykey{resume*}, this will activate the entire \emph{\enquote{storage
% system}} in the \mypkg*[no-index]{enumext} package.
%
% \subsubsection{Definition of error messages}
%
% \begin{macro}[int]{save-ans-empty, save-ans-log, save-ans-log-hook}
% Messages used by \mykey[env]{save-ans} key.
%    \begin{macrocode}
\msg_new:nnn { enumext } { save-ans-empty }
  {
    Key~'save-ans'~needs~a~value~in~'#1'~\msg_line_context:.
  }
\msg_new:nnn { enumext } { save-ans-log }
  {
    *~Package~enumext:~Start~#1~with~save-ans=#2~\msg_line_context:.
  }
\msg_new:nnn { enumext } { save-ans-log-hook }
  {
    *~Package~enumext:~Stop~#1~with~save-ans=#2~\msg_line_context:.
  }
%    \end{macrocode}
% \end{macro}
%
% \begin{macro}[int]{items-same-answer, item-greater-answer, item-less-answer}
%  Messages used by the internal system to check answer used by
%  \mykey[env]{check-ans} key.
%    \begin{macrocode}
\msg_new:nnn { enumext } { items-same-answer }
  {
    ****************************************\\
    *~Package~enumext:~Checking~answers~in~'#1'~
    for~\c_left_brace_str #2 \c_right_brace_str\\
    *~started~#3~and~close~\msg_line_context: :~
    'OK',~all~items~with~answer.\\
    ****************************************
  }
\msg_new:nnn { enumext } { item-greater-answer }
  {
    Checking~answers~in~'#1'~for~\c_left_brace_str #2 \c_right_brace_str\\
    started~#3~and~close~\msg_line_context: :~'NOT~OK'\\
    Items~>~Answers.
  }
\msg_new:nnn { enumext } { item-less-answer }
  {
    Checking~answers~in~'#1'~for~\c_left_brace_str #2 \c_right_brace_str\\
    started~#3~and~close~\msg_line_context: :~'NOT~OK'\\
    Items~<~Answers.
  }
%    \end{macrocode}
% \end{macro}
%
% \begin{macro}[int]{missing-starred, many-starred}
%  Messages used by the internal system to check for
%  \emph{\enquote{starred}} |\item*| and |\anspic*| commands.
%    \begin{macrocode}
\msg_new:nnn { enumext } { missing-starred }
  {
    Missing~'\c_backslash_str #1*'~#2.
  }
\msg_new:nnn { enumext } { many-starred }
  {
    Many~'\c_backslash_str #1*'~#2.
  }
%    \end{macrocode}
% \end{macro}
%
% \begin{macro}[int]{store-prop, store-seq, store-int}
% Messages written to the |.log| file when the \emph{prop list}, \emph{sequence} and
% integer variables used by the \mykey[env]{save-ans} key are created.
%    \begin{macrocode}
\msg_new:nnn { enumext } { store-prop }
  {
    *~Package~enumext:~Creating~
    \c_backslash_str g_@@_ #1 _prop~\msg_line_context:.
  }
\msg_new:nnn { enumext } { store-seq }
  {
    *~Package~enumext:~Creating~
    \c_backslash_str g_@@_ #1 _seq~\msg_line_context:.
  }
\msg_new:nnn { enumext } { store-int }
  {
    *~Package~enumext:~Creating~
    \c_backslash_str g_@@_resume_ #1 _int~\msg_line_context:.
  }
%    \end{macrocode}
% \end{macro}
%
% \subsubsection{Setting \texttt{save-ans} key}\label{key:save-ans}
%
% \begin{macro}[int]{save-ans}
%     Defines the key \mykey[env]{save-ans} only for the \emph{\enquote{first level}}
%     of \myenv*[no-index]{enumext} and \myenv*[no-index]{enumext*}.
%    \begin{macrocode}
\cs_set_protected:Nn \@@_tmp:n
  {
    \keys_define:nn { enumext / #1 }
      {
        save-ans .code:n = { \@@_storing_set:n {##1} },
        save-ans .value_required:n = true,
      }
  }
\clist_map_inline:nn { level-1, enumext* } { \@@_tmp:n {#1} }
%    \end{macrocode}
% \end{macro}
%
% \subsubsection{Internal functions for \texttt{save-ans} key}
%
% \begin{macro}[int]{\@@_start_save_ans_msg:, \@@_stop_save_ans_msg:}
%   The functions \myhlc|\@@_start_save_ans_msg:| and \myhlc|\@@_stop_save_ans_msg:|
%   will display in the terminal and |.log| file the environment in which
%   the \mykey[env]{save-ans} key was executed along with the line at the
%   beginning and end of it. The function \myhlc|\@@_start_save_ans_msg:| will
%   be passed to \myhlc|\@@_storing_set:n| and the function
%   \myhlc|\@@_stop_save_ans_msg:| will be passed to the function
%   \myhlc|\@@_execute_after_env:|.
%    \begin{macrocode}
\cs_new_protected:Nn \@@_start_save_ans_msg:
  {
    \msg_term:nnVV { enumext } { save-ans-log }
      \g_@@_envir_name_tl \l_@@_store_name_tl
  }
\cs_new_protected:Nn \@@_stop_save_ans_msg:
  {
    \msg_term:nnVV { enumext } { save-ans-log-hook }
      \g_@@_envir_name_tl \g_@@_store_name_tl
  }
%    \end{macrocode}
% \end{macro}
%
% \begin{macro}[int]{\@@_storing_set:n, \@@_storing_exec:}
%    The function \myhlc|\@@_storing_set:n| first passes the value of the
%    \mykey[env]{save-ans} key to the variable \myhlc|\l_@@_store_name_tl|
%    which contains the \mymarg{store name} of the \emph{sequence}
%    and \emph{prop list} to be used. If \myhlc|\l_@@_store_name_tl| is
%    \emph{empty} it returns an error message, otherwise it returns the
%    appropriate message \myhlc|\@@_start_save_ans_msg:| and proceed to execute
%    the function \myhlc|\@@_storing_exec:| for \myenv*[no-index]{enumext} and
%    \myenv*[no-index]{enumext*} environments.
%    \begin{macrocode}
\cs_new_protected:Nn \@@_storing_set:n
  {
    \tl_set:Ne \l_@@_store_name_tl {#1}
    \tl_if_empty:NTF \l_@@_store_name_tl
      {
        \bool_lazy_or:nnT
          { \l_@@_standard_first_bool } { \l_@@_starred_first_bool }
          {
            \msg_error:nnV { enumext } { save-ans-empty } \g_@@_envir_name_tl
          }
      }
      {
        \bool_lazy_or:nnT
          { \l_@@_standard_first_bool } { \l_@@_starred_first_bool }
          {
            \@@_start_save_ans_msg:
            \@@_storing_exec:
          }
      }
  }
%    \end{macrocode}
%   The function \myhlc|\@@_storing_exec:| will set to true the variable
%   \myhlc|\l_@@_store_active_bool| which activates the use of the
%   |\anskey| command and the \myenv*[no-index]{anskey*}, \myenv*[no-index]{keyans},
%   \myenv*[no-index]{keyans*} and \myenv*[no-index]{keyanspic} environments and will set
%   to |true| the variable \myhlc|\l_@@_check_answers_bool| used for
%   internal checking answers mechanism set by the \mykey[env]{check-ans} and \mykey[env]{no-store}
%   keys, copy \mymarg{store name} into the variable \myhlc|\g_@@_store_name_tl|.
%    \begin{macrocode}
\cs_new_protected:Nn \@@_storing_exec:
  {
    \bool_set_true:N \l_@@_store_active_bool
    \bool_set_true:N \l_@@_check_answers_bool
    \tl_gset:NV \g_@@_store_name_tl \l_@@_store_name_tl
%    \end{macrocode}
%   The \emph{prop list} \myhlc|\g_@@_series_#1_prop| and the \emph{sequence} \myhlc|\g_@@_series_#1_seq|
%   will be created globally to \emph{\enquote{store content}} in case
%   they do not exist together with the integer variable
%   \myhlc|\g_@@_series_#1_int| used by the keys
%   \mykey[env]{resume} and \mykey[env]{resume*}.
%    \begin{macrocode}
    \prop_if_exist:cF { g_@@_ \l_@@_store_name_tl _prop }
      {
        \msg_log:nnV { enumext } { store-prop } \l_@@_store_name_tl
        \prop_new:c { g_@@_ \l_@@_store_name_tl _prop }
      }
    \seq_if_exist:cF { g_@@_ \l_@@_store_name_tl _seq }
      {
        \msg_log:nnV { enumext } { store-seq } \l_@@_store_name_tl
        \seq_new:c { g_@@_ \l_@@_store_name_tl _seq }
      }
    \int_if_exist:cF { g_@@_resume_ \l_@@_store_name_tl _int }
      {
        \msg_log:nnV { enumext } { store-int } \l_@@_store_name_tl
        \int_new:c { g_@@_resume_ \l_@@_store_name_tl _int }
      }
  }
%    \end{macrocode}
% \end{macro}
%
% \subsubsection{The check answer mechanism}\label{subsubsec:check-answer}
%
% The internal mechanism for \emph{\enquote{checking answers}} follows this
% logic:
%
% \begin{quote}
%  If the line begins with |\item| or |\item*| and does NOT
%  \emph{open a nested environment}, each |\item| or |\item*| must contain a
%  \emph{single} execution of the |\anskey| command, i.e. the counter of the
%  executions of the |\anskey| command must be equal to the counter
%  associated with the sum of executions of |\item| and |\item*|.
%
%  If the line begins with |\item| or |\item*| and \emph{opens a nested
%  environment} each |\item| or |\item*| in the nested environment must
%  have a \emph{single} execution of the |\anskey| command and the counter
%  associated to the sum of |\item| and |\item*| executions must decrementing
%  by \emph{\enquote{one}} to maintain equality.
%
% \end{quote}
%
% In order for the mechanism for the check-answer to work (not counting
% \myenv*[no-index]{keyans}, \myenv*[no-index]{keyans*} and \myenv*[no-index]{keyanspic}) it is necessary to:
%
% \begin{enumerate}[left=0pt,noitemsep, partopsep=0pt, topsep=5pt]
%   \item Keep track of the total number of |\item| and |\item*| (enumerated) that appear within
%         the environment including the nested levels.
%   \item Keep track of the total number of |\item| and |\item*| (enumerated) that appear per
%         level of nesting.
%   \item Keep track of the number of times the environment nests.
% \end{enumerate}
%
%
% The integer variable associated to the sum of each |\item| and |\item*|
% in the environment \myhlc|\g_@@_item_number_int| must match the
% integer variable \myhlc|\g_@@_item_anskey_int| associated
% to the execution of the command |\anskey|. The cases are analyzed below:
%
% \begin{enumerate}[label=\alph*),left=0pt,noitemsep, partopsep=0pt, topsep=5pt]
%   \item If the list only has one level the number of |\item + \item*| = |\anskey|
%
%   \item If the list has \emph{nested levels}, for each level of nesting the count
%    is decremented by one (for the |\item| or |\item*| that opens the nest) so that the account
%    remains the same.
% \end{enumerate}
%
% With \myenv*[no-index]{keyans}, \myenv*[no-index]{keyans*} and \myenv*[no-index]{keyanspic} it is
% enough to increase in one the integer of |\anskey|. The integers
% created must be global if they are not lost in the interior levels of
% nesting and to execute the test a \emph{\enquote{hook}}
% function is used after closing the \emph{first level} of the environment.
%
% \subsubsection{Setting \texttt{check-ans} and \texttt{no-store} keys}\label{key:check-ans}
%
% \begin{macro}[int]{check-ans, no-store}
%    Defines the keys \mykey[env]{check-ans} and \mykey[env]{no-store} for all
%    levels of \myenv*[no-index]{enumext} and \myenv*[no-index]{enumext*} environments.
%    \begin{macrocode}
\cs_set_protected:Nn \@@_tmp:n
  {
    \keys_define:nn { enumext / #1 }
      {
        check-ans .bool_set:N = \l_@@_check_ans_key_bool,
        check-ans .initial:n = false,
        check-ans .value_required:n = true,
        no-store  .code:n =
          {
            \bool_set_false:N \l_@@_check_answers_bool
            \bool_set_false:N \l_@@_check_ans_key_bool
          },
        no-store  .value_forbidden:n = true,
      }
  }
\clist_map_inline:nn
  {
    level-1, level-2, level-3, level-4, enumext*
  }
  { \@@_tmp:n {#1} }
%    \end{macrocode}
% \end{macro}
%
% \subsubsection{Set-up check answer mechanism}\label{subsubsec:set-check-answer}
%
% \begin{macro}[int]{\@@_check_ans_active:, \@@_check_ans_level:}
%   The function \myhlc|\@@_check_ans_active:| will first check the state of
%   the variable \myhlc|\l_@@_store_name_tl|, that is, the \mykey[env]{save-ans} key
%   is active, if so it will check the state of the variable
%   \myhlc|\l_@@_check_answers_bool| handled by the key \mykey[env]{no-store} and
%   will execute the function \myhlc|\@@_check_ans_level:| only if
%   |true|, i.e. the key \mykey{no-store} is not active.
%    \begin{macrocode}
\cs_new_protected:Nn \@@_check_ans_active:
  {
    \tl_if_empty:NF \l_@@_store_name_tl
      {
        \bool_if:NT \l_@@_check_answers_bool
          {
            \@@_check_ans_level:
          }
      }
  }
%    \end{macrocode}
%   The function \myhlc|\@@_check_ans_level:| will decrement by \emph{\enquote{one}}
%   the value of the variable \myhlc|\g_@@_item_number_int| which keeps track
%   of the executions of |\item| and |\item*| for each level of nesting
%   of the environment \myenv*[no-index]{enumext}, taking into account whether it
%   is nested within \myenv*[no-index]{enumext*} or the opposite and set \myhlc|\l_@@_item_number_bool|
%   to |false|.
%    \begin{macrocode}
\cs_new_protected:Nn \@@_check_ans_level:
  {
    \int_case:nn { \l_@@_level_int }
      {
        { 1 }
          {
            \bool_lazy_all:nT
              {
                { \bool_if_p:N \g_@@_starred_bool }
                { \int_compare_p:nNn { \l_@@_level_h_int } = { 1 } }
              }
              {
                \int_gdecr:N \g_@@_item_number_int
                \bool_set_false:N \l_@@_item_number_bool
              }
          }
        { 2 }
          {
            \int_gdecr:N \g_@@_item_number_int
            \bool_set_false:N \l_@@_item_number_bool
          }
        { 3 }
          {
            \int_gdecr:N \g_@@_item_number_int
            \bool_set_false:N \l_@@_item_number_bool
          }
        { 4 }
          {
            \int_gdecr:N \g_@@_item_number_int
            \bool_set_false:N \l_@@_item_number_bool
          }
      }
%    \end{macrocode}
%   This is executed only if \myenv*[no-index]{enumext*} is nested in the
%   \emph{\enquote{first level}} of \myenv*[no-index]{enumext}, for the rest of the
%   cases the value of \myhlc|\g_@@_item_number_int| is already decreased.
%    \begin{macrocode}
    \int_case:nn { \l_@@_level_h_int }
      {
        { 1 }
          {
            \bool_lazy_all:nT
              {
                { \bool_if_p:N \g_@@_standard_bool }
                { \int_compare_p:nNn { \l_@@_level_int } = { 1 } }
              }
              {
                \int_gdecr:N \g_@@_item_number_int
                \bool_set_false:N \l_@@_item_number_bool
              }
          }
      }
  }
%    \end{macrocode}
% \end{macro}
%
% \begin{macro}[int]{\@@_check_ans_key_hook:}
%   The function \myhlc|\@@_check_ans_key_hook:| will \emph{export} the status
%   of the local variable \myhlc|\l_@@_check_ans_key_bool| to the global variable
%   \myhlc|\g_@@_check_ans_key_bool| only if the key \mykey[env]{check-ans} is active.
%    \begin{macrocode}
\cs_new_protected:Nn \@@_check_ans_key_hook:
  {
    \bool_lazy_and:nnT
      { \bool_if_p:N \l_@@_check_ans_key_bool }
      { \bool_if_p:N \g_@@_standard_bool }
      {
        \bool_gset_true:N \g_@@_check_ans_key_bool
      }
    \bool_lazy_and:nnT
      { \bool_if_p:N \l_@@_check_ans_key_bool }
      { \bool_if_p:N \g_@@_starred_bool }
      {
        \bool_gset_true:N \g_@@_check_ans_key_bool
      }
  }
%    \end{macrocode}
% \end{macro}
%
% \begin{macro}[int]{\@@_item_answer_diff:}
%   The function \myhlc|\@@_item_answer_diff:| will set the value of the
%   variable \myhlc|\g_@@_item_answer_diff_int| which is used by the functions
%   \myhlc|\@@_check_ans_show:| for the key \mykey[env]{save-ans} and by the
%   function \myhlc|\@@_check_ans_log:| by the internal \emph{\enquote{check answer}}
%   mechanism. This function will be passed to the function
%   \myhlc|\@@_execute_after_env:|.
%    \begin{macrocode}
\cs_new_protected:Nn \@@_item_answer_diff:
  {
    \int_gset:Nn \g_@@_item_answer_diff_int
      {
        \int_sign:n { \g_@@_item_number_int - \g_@@_item_anskey_int }
      }
  }
%    \end{macrocode}
% \end{macro}
%
% \begin{macro}[int]{\@@_check_ans_show:, \@@_check_ans_msg_less:,
%                    \@@_check_ans_msg_same_ok:, \@@_check_ans_msg_greater:}
%   The function \myhlc|\@@_check_ans_show:| will be executed within the
%   function \myhlc|\@@_execute_after_env:| when the key \mykey[env]{check-ans} is
%   active, that is, when \myhlc|\g_@@_check_ans_key_bool| is |true|
%   and will return the appropriate message according to the value of
%   \myhlc|\g_@@_item_answer_diff_int| set by the function \myhlc|\@@_item_answer_diff:|.
%    \begin{macrocode}
\cs_new_protected:Nn \@@_check_ans_show:
  {
    \int_case:nn { \g_@@_item_answer_diff_int }
      {
        { -1 }{ \@@_check_ans_msg_less:    }
        {  0 }{ \@@_check_ans_msg_same_ok: }
        {  1 }{ \@@_check_ans_msg_greater: }
      }
  }
\cs_new_protected:Nn \@@_check_ans_msg_less:
  {
    \msg_warning:nneee { enumext } { item-less-answer } { \g_@@_store_name_tl }
      { \g_@@_envir_name_tl } { \g_@@_start_line_tl }
  }
\cs_new_protected:Nn \@@_check_ans_msg_same_ok:
  {
    \msg_term:nneee { enumext } { items-same-answer } { \g_@@_store_name_tl }
      { \g_@@_envir_name_tl } { \g_@@_start_line_tl }
  }
\cs_new_protected:Nn \@@_check_ans_msg_greater:
  {
    \msg_warning:nneee { enumext } { item-greater-answer }
      { \g_@@_store_name_tl }
      { \g_@@_envir_name_tl } { \g_@@_start_line_tl }
  }
%    \end{macrocode}
% \end{macro}
% \begin{macro}[int]{\@@_check_ans_log:, \@@_check_ans_log_msg_less:,
%                    \@@_check_ans_log_msg_same_ok:, \@@_check_ans_log_msg_greater:}
%   The function \myhlc|\@@_check_ans_log:| will be executed within the
%   function \myhlc|\@@_execute_after_env:| when the key \mykey[env]{check-ans} is
%   not active, that is, when \myhlc|\g_@@_check_ans_key_bool| is |false|
%   and write in the log the appropriate message according to the value of
%   \myhlc|\g_@@_item_answer_diff_int| set by the function
%   \myhlc|\@@_item_answer_diff:|.
%    \begin{macrocode}
\cs_new_protected:Nn \@@_check_ans_log:
  {
    \int_case:nn { \g_@@_item_answer_diff_int }
      {
        { -1 }{ \@@_check_ans_log_msg_less:    }
        {  0 }{ \@@_check_ans_log_msg_same_ok: }
        {  1 }{ \@@_check_ans_log_msg_greater: }
      }
  }
\cs_new_protected:Nn \@@_check_ans_log_msg_less:
  {
    \msg_log:nneee { enumext } { item-less-answer } { \g_@@_store_name_tl }
      { \g_@@_envir_name_tl } { \g_@@_start_line_tl }
  }
\cs_new_protected:Nn \@@_check_ans_log_msg_same_ok:
  {
    \msg_log:nneee { enumext } { items-same-answer }  { \g_@@_store_name_tl }
      { \g_@@_envir_name_tl } { \g_@@_start_line_tl }
  }
\cs_new_protected:Nn \@@_check_ans_log_msg_greater:
  {
    \msg_log:nneee { enumext } { item-greater-answer } { \g_@@_store_name_tl }
      { \g_@@_envir_name_tl } { \g_@@_start_line_tl }
  }
%    \end{macrocode}
% \end{macro}
%
% \subsubsection{Check for \cs{item*} and \cs{anspic*} commands}
%
% \begin{macro}[int]{\@@_check_starred_cmd:n}
%    The function \myhlc|\@@_check_starred_cmd:n| performs an \emph{extra check}
%    for the \myenv*[no-index]{keyans}, \myenv*[no-index]{keyans*} and \myenv*[no-index]{keyanspic} environments.
%    Unlike the \emph{check} executed by \mykey[env]{check-ans} key this one
%    is not controlled by any key, it is intended to prevent the forgetting
%    of |\item*| or |\anspic*| in these environments.
%    \begin{macrocode}
\cs_new_protected:Nn \@@_check_starred_cmd:n
  {
    \int_compare:nNnT
      { \g_@@_check_starred_cmd_int } = { 0 }
      {
        \msg_warning:nnnV
          { enumext } { missing-starred }{#1} \l_@@_check_start_line_env_tl
      }
    \int_compare:nNnT
      { \g_@@_check_starred_cmd_int } > { 1 }
      {
        \msg_warning:nnnV
          { enumext } { many-starred }{#1} \l_@@_check_start_line_env_tl
      }
    \int_gzero:N \g_@@_check_starred_cmd_int
    \tl_clear:N \l_@@_check_start_line_env_tl
  }
%    \end{macrocode}
% \end{macro}
%
% \subsection{Keys and functions associated with storage}\label{subsec:storage-system}
%
% \subsubsection{Keys for \texttt{marks}, \texttt{wrap} and \texttt{show}}
%
% The \mypkg*[no-index]{enumext} package provides a set of \mymeta{keys} for
% manipulating \emph{\enquote{symbol marks}} associated with \emph{\enquote{answers}}
% and how they are displayed and stored in the \emph{sequence} and \emph{prop list}
% as well as an internal \emph{\enquote{label and ref}} system.
%
% \begin{macro}[int]{mark-ans*, mark-pos*, mark-sep*, wrap-ans*, wrap-opt,
%                    save-sep, show-ans, show-pos}
%   For the \myenv*[no-index]{keyans} and \myenv*[no-index]{keyans*} environments only the keys
%   \mykey[env]{mark-ans*}, \mykey[env]{mark-pos*}, \mykey[env]{mark-sep*},
%   \mykey[env]{wrap-ans*}, \mykey[env]{wrap-opt}, \mykey[env]{save-sep}, \mykey[env]{show-ans}
%   and \mykey[env]{show-pos} are added.
%    \begin{macrocode}
\cs_set_protected:Nn \@@_tmp:nn
  {
    \keys_define:nn { enumext / #1 }
      {
        mark-ans* .tl_set:c = { l_@@_mark_answer_sym_ #2 _tl },
        mark-ans* .initial:n = \textasteriskcentered,
        mark-ans* .value_required:n = true,
        mark-pos* .choices:nn =
          { left , right , center }
          {
            \int_case:nn { \l_keys_choice_int }
              {
                {1} { \str_set:cn { l_@@_mark_position_ #2 _str } { l } }
                {2} { \str_set:cn { l_@@_mark_position_ #2 _str } { r } }
                {3} { \str_set:cn { l_@@_mark_position_ #2 _str } { c } }
              }
          },
        mark-pos* .initial:n = right,
        mark-pos* .value_required:n = true,
        mark-pos* / unknown .code:n =
          \msg_error:nneee { enumext } { unknown-choice }
            { mark-pos } { left,~right,~center } { \exp_not:n {##1} },
        mark-sep* .dim_set:c = { l_@@_mark_sym_sep_ #2 _dim },
        mark-sep* .value_required:n = true,
        wrap-ans* .cs_set_protected:cp = { @@_keyans_wrapper_item_ #2 :n } ##1,
        wrap-ans* .value_required:n = true,
        wrap-opt  .cs_set_protected:cp = { @@_keyans_wrapper_opt_ #2 :n } ##1,
        wrap-opt  .initial:n = [{##1}],
        wrap-opt  .value_required:n = true,
        save-sep  .tl_set:c = { l_@@_store_keyans_item_opt_sep_ #2 _tl },
        save-sep  .initial:n = {,~},
        save-sep  .value_required:n = true,
        show-ans  .bool_set:N = \l_@@_show_answer_bool,
        show-ans  .initial:n = false,
        show-ans  .value_required:n = true,
        show-pos  .bool_set:N = \l_@@_show_position_bool,
        show-pos  .initial:n = false,
        show-pos  .value_required:n = true,
      }
  }
\clist_map_inline:nn { {keyans}{v}, {keyans*}{viii} } { \@@_tmp:nn #1 }
%    \end{macrocode}
% \end{macro}
%
% \begin{macro}[int]{mark-ref, save-ref, show-ans, show-pos, mark-ans,
%                    mark-pos, mark-sep, wrap-ans, mark-ans*, mark-pos*,
%                    mark-sep*, wrap-ans*, wrap-opt, save-sep}
%   Adds the \mymeta{keys} \mykey[env]{mark-ref} and \mykey[env]{save-ref} related to the
%   \emph{\enquote{storage system}} and internal mechanism of
%   \emph{\enquote{label and ref}} along with the \mymeta{keys} \mykey[env]{show-ans},
%   \mykey[env]{show-pos} and the\mymeta{keys} \mykey[env]{mark-ans}, \mykey[env]{mark-pos},
%   \mykey[env]{mark-sep} and \mykey[env]{wrap-ans} for the command |\anskey|,
%   the environment \myenv*[no-index]{anskey*} and the \mymeta{keys} for environments
%   \myenv*[no-index]{keyans} and \myenv*[no-index]{keyans*} only at the \emph{first level}
%   of \myenv*[no-index]{enumext} and \myenv*[no-index]{enumext*}.
%    \begin{macrocode}
\cs_set_protected:Nn \@@_tmp:n
  {
    \keys_define:nn { enumext / #1 }
      {
        mark-ref .tl_set:N = \l_@@_mark_ref_sym_tl,
        mark-ref .initial:n = \textreferencemark,
        mark-ref .value_required:n = true,
        save-ref .bool_set:N = \l_@@_store_ref_key_bool,
        save-ref .initial:n = false,
        save-ref .value_required:n = true,
        show-ans .bool_set:N = \l_@@_show_answer_bool,
        show-ans .initial:n = false,
        show-ans .value_required:n = true,
        show-pos .bool_set:N = \l_@@_show_position_bool,
        show-pos .initial:n = false,
        show-pos .value_required:n = true,
        mark-ans .tl_set:N = \l_@@_mark_answer_sym_tl,
        mark-ans .initial:n = \textasteriskcentered,
        mark-ans .value_required:n = true,
        mark-sep .dim_set:N = \l_@@_mark_sym_sep_dim,
        mark-sep .value_required:n = true,
        mark-pos .choices:nn =
          { left , right , center }
          {
            \int_case:nn { \l_keys_choice_int }
              {
                {1} { \str_set:Nn \l_@@_mark_position_str { l } }
                {2} { \str_set:Nn \l_@@_mark_position_str { r } }
                {3} { \str_set:Nn \l_@@_mark_position_str { c } }
              }
          },
        mark-pos .initial:n = right,
        mark-pos .value_required:n = true,
        mark-pos / unknown .code:n =
          \msg_error:nneee { enumext } { unknown-choice }
            { mark-pos } { left,~right,~center } { \exp_not:n {##1} },

        wrap-ans  .cs_set_protected:Np = \@@_anskey_wrapper:n ##1,
        wrap-ans  .initial:n =
                    {
                      \fbox{\parbox[t]{\dimeval{\itemwidth -2\fboxsep
                        -2\fboxrule}}{##1}}
                    },
        wrap-ans  .value_required:n = true,
        mark-ans* .code:n =
          {
            \keys_set:nn { enumext / keyans  } { mark-ans* = {##1} }
            \keys_set:nn { enumext / keyans* } { mark-ans* = {##1} }
          },
        mark-ans* .value_required:n = true,
        mark-pos* .code:n =
          {
            \keys_set:nn { enumext / keyans  } { mark-pos* = {##1} }
            \keys_set:nn { enumext / keyans* } { mark-pos* = {##1} }
          },
        mark-pos* .value_required:n = true,
        mark-sep* .code:n =
          {
            \keys_set:nn { enumext / keyans  } { mark-sep* = {##1} }
            \keys_set:nn { enumext / keyans* } { mark-sep* = {##1} }
          },
        mark-sep* .value_required:n = true,
        wrap-ans* .code:n =
          {
            \keys_set:nn { enumext / keyans  } { wrap-ans* = {##1} }
            \keys_set:nn { enumext / keyans* } { wrap-ans* = {##1} }
          },
        wrap-ans* .value_required:n = true,
        wrap-opt  .code:n =
          {
            \keys_set:nn { enumext / keyans  } { wrap-opt = {##1} }
            \keys_set:nn { enumext / keyans* } { wrap-opt = {##1} }
          },
        wrap-opt  .value_required:n = true,
        save-sep  .code:n =
          {
            \keys_set:nn { enumext / keyans  } { save-sep = {##1} }
            \keys_set:nn { enumext / keyans* } { save-sep = {##1} }
          },
        save-sep  .value_required:n = true,
      }
  }
\clist_map_inline:nn { level-1, enumext* } { \@@_tmp:n {#1} }
%    \end{macrocode}
% \end{macro}
%
% \subsubsection{Storing structure of the environments}
%
% The idea behind \emph{\enquote{storing structure}} in the \emph{sequence} is
% to have a copy of the \emph{structure of the environment} in which the key
% \mykey{save-ans} is being executed, so the \emph{optional
% argument} passed to the levels of the environment in which it is
% executed must be captured and \emph{\enquote{stored}} in the \emph{sequence}.
%
% \begin{macro}[int]{\@@_store_active_keys:n, \@@_store_active_keys_vii:n}
%   The functions \myhlc|\@@_store_active_keys:n| and \myhlc|\@@_store_active_keys_vii:n|
%   will be responsible for the \emph{\enquote{storing keys}} filtered
%   from the \emph{optional argument} of the environment in which the
%   key \mykey[env]{save-ans} is executed and the levels within this for the \myenv*[no-index]{enumext} and
%   \myenv*[no-index]{enumext*} environments. This function is executed only if
%   the variable \myhlc|\l_@@_store_save_key_X_bool| is false, that is, the key
%   \mykey{store-key} is not active, establishing the variable
%   \myhlc|\l_@@_store_save_key_X_tl| with the filtered \mymeta{keys}.
%    \begin{macrocode}
\cs_new_protected:Nn \@@_store_active_keys:n
  {
    \bool_if:cF { l_@@_store_save_key_ \@@_level: _bool }
      {
        \tl_clear:c { l_@@_store_save_key_ \@@_level: _tl }
        \tl_set:ce
          { l_@@_store_save_key_ \@@_level: _tl }
          { \@@_filter_save_key:n {#1} }
      }
  }
\cs_new_protected:Nn \@@_store_active_keys_vii:n
  {
    \bool_if:NF \l_@@_store_save_key_vii_bool
      {
        \tl_clear:N \l_@@_store_save_key_vii_tl
        \tl_set:Ne \l_@@_store_save_key_vii_tl { \@@_filter_save_key:n {#1} }
      }
  }
%    \end{macrocode}
% \end{macro}
%
% \subsubsection{Setting \texttt{save-key} key}\label{key:save-key}
%
% Since this \emph{\enquote{storing structure}} in the \emph{sequence}
% established by the \mykey{save-ans} key when executing |\anskey| or \myenv*[no-index]{anskey*},
% it cannot be modified. The best option here is to have a key
% that allows modifying the \emph{optional argument} of the \emph{\enquote{storing structure}} in
% the \emph{sequence}.
%
% \begin{macro}[int]{save-key}
%   The values set by this key passed in the \emph{optional argument} of the
%   \myenv*[no-index]{enumext} and \myenv*[no-index]{enumext*} environments will override the
%   values of the \myhlc|\l_@@_store_save_key_X_tl| variable set by the
%   functions \myhlc|\@@_store_active_keys:n| and \myhlc|\@@_store_active_keys_vii:n|.
%   Now define the key \mykey[env]{save-key} for all levels of
%   \myenv*[no-index]{enumext} and \myenv*[no-index]{enumext*} environments.
%    \begin{macrocode}
\cs_set_protected:Nn \@@_tmp:n
  {
    \keys_define:nn { enumext / enumext* }
      {
        save-key .code:n = { \@@_parse_save_key_vii:n {##1} },
        save-key .value_required:n = true,
      }
    \keys_define:nn { enumext / #1 }
      {
        save-key .code:n = { \@@_parse_save_key:n {##1} },
        save-key .value_required:n = true,
      }
  }
\clist_map_inline:nn { level-1, level-2, level-3, level-4 } { \@@_tmp:n {#1} }
%    \end{macrocode}
% \end{macro}
%
% \begin{macro}[int]{\@@_parse_save_key:n, \@@_parse_save_key_vii:n}
%   The functions \myhlc|\@@_parse_save_key:n| and \myhlc|\@@_parse_save_key_vii:n|
%   will be responsible for \emph{\enquote{storing keys}} in the variable
%   \myhlc|\l_@@_store_save_key_X_tl| for \myenv*[no-index]{enumext} and \myenv*[no-index]{enumext*}.
%    \begin{macrocode}
\cs_new_protected:Nn \@@_parse_save_key:n
  {
    \bool_set_true:c { l_@@_store_save_key_ \@@_level: _bool }
    \tl_clear:c { l_@@_save_key_ \@@_level: _tl }
    \tl_set:ce
      { l_@@_store_save_key_ \@@_level: _tl }
      { \@@_filter_save_key:n {#1} }
  }
\cs_new_protected:Nn \@@_parse_save_key_vii:n
  {
    \bool_set_true:N \l_@@_store_save_key_vii_bool
    \tl_clear:N \l_@@_store_save_key_vii_tl
    \tl_set:Ne \l_@@_store_save_key_vii_tl { \@@_filter_save_key:n {#1} }
  }
%    \end{macrocode}
% \end{macro}
%
% \subsubsection{Internal functions to store optional arguments}
%
% \begin{macro}[int]{\@@_filter_save_key:n, \@@_filter_save_key_key:n,
%                    \@@_filter_save_key_pair:nn}
%   The function \myhlc|\@@_filter_save_key:n| will be in charge of
%   \emph{\enquote{filtering keys}} to be \emph{stored} in
%   \emph{sequence} where |{#1}| represents the \emph{optional argument}
%   passed to the environment.
%    \begin{macrocode}
\cs_new:Nn \@@_filter_save_key:n
  {
    \use:e
      {
        \keyval_parse:NNn
          \@@_filter_save_key_key:n
          \@@_filter_save_key_pair:nn {#1}
      }
  }
%    \end{macrocode}
%   The function \myhlc|\@@_filter_save_key_key:n| will be responsible for
%   \emph{\enquote{filtering keys}} that are passed \emph{\enquote{without value}}
%   by excluding the \mykey[env]{resume}, \mykey[env]{resume*}, \mykey[env]{reset}, \mykey[env]{reset*},
%   \mykey[env]{no-store} and \mykey[env]{base-fix} keys.
%    \begin{macrocode}
\cs_new:Nn \@@_filter_save_key_key:n
  {
    \str_case:nnF {#1}
      {
        { resume } { } { resume* } { } { reset } { } { reset* } { } { no-store } { }
          { base-fix } { }
      }
      { , { \exp_not:n {#1} } }
  }
%    \end{macrocode}
%   The function \myhlc|\@@_filter_save_key_pair:nn| will be responsible for
%   \emph{\enquote{filtering keys}} that are passed \emph{\enquote{with value}} by
%   excluding the \mykey[env]{series}, \mykey{resume}, \mykey[env]{save-ans}, \mykey[env]{save-ref},
%   \mykey[env]{save-key}, \mykey[env]{check-ans}, \mykey[env]{show-ans}, \mykey{save-pos},
%   \mykey[env]{mark-ans}, \mykey[env]{mark-pos}, \mykey[env]{mark-sep}, \mykey[env]{wrap-ans},
%   \mykey[env]{mark-ans*}, \mykey[env]{mark-pos*}, \mykey[env]{mark-sep*}, \mykey[env]{wrap-ans*},
%   \mykey[env]{wrap-opt}, \mykey[env]{save-sep}, \mykey[env]{mark-ref}, \mykey[env]{mini-env},
%   \mykey[env]{mini-sep}, \mykey[env]{mini-right} and \mykey[env]{mini-right*} keys.
%    \begin{macrocode}
\cs_new:Nn \@@_filter_save_key_pair:nn
  {
    \str_case:nnF {#1}
      {
        { series    } { } { resume     } { } { save-ans    } { } { save-ref  } { }
        { save-key  } { } { check-ans  } { } { show-ans    } { } { show-pos  } { }
        { mark-ans  } { } { mark-pos   } { } { mark-sep    } { } { wrap-ans  } { }
        { mark-ans* } { } { mark-pos*  } { } { mark-sep*   } { } { wrap-ans* } { }
        { wrap-opt  } { } { save-sep   } { } { mark-ref    } { } { mini-env  } { }
        { mini-sep  } { } { mini-right } { } { mini-right* } { }
      }
      { , { \exp_not:n {#1} } = { \exp_not:n {#2} } }
  }
%    \end{macrocode}
% \end{macro}
%
% \subsubsection{Function for storing content in prop list}
%
% \begin{macro}[int]{\@@_store_addto_prop:n, \@@_store_addto_prop:V}
%   The function \myhlc|\@@_store_addto_prop:n| stores the \mymarg{content} in
%   \emph{prop list} defined by \mykey[env]{save-ans} key. The \emph{\enquote{stored
%   content}} is retrieved by means of the |\getkeyans| command.
%
%   The form in which the \mymarg{content} is \emph{\enquote{stored}} in the
%   \emph{prop list} is \mymarg{position}\mymarg{content}. This
%   function is used by |\anskey| in \myenv*[no-index]{enumext} and \myenv*[no-index]{enumext*}
%   environments, |\item*| in \myenv*[no-index]{keyans} and \myenv*[no-index]{keyans*}
%   environments and |\anspic*| in \myenv*[no-index]{keyanspic} environment.
%    \begin{macrocode}
\cs_new_protected:Nn \@@_store_addto_prop:n
  {
    \prop_gput_if_not_in:cen { g_@@_ \l_@@_store_name_tl _prop }
      {
        \int_eval:n { \prop_count:c { g_@@_ \l_@@_store_name_tl _prop } + 1 }
      }
      {#1}
  }
\cs_generate_variant:Nn \@@_store_addto_prop:n { V }
%    \end{macrocode}
% \end{macro}
%
% \subsubsection{Function for storing content in sequence}
%
% \begin{macro}[int]{\@@_store_addto_seq:n, \@@_store_addto_seq:v, \@@_store_addto_seq:V}
%   The function \myhlc|\@@_store_addto_seq:n| stores the \mymarg{content} in
%   \emph{sequence} defined by \mykey[env]{save-ans} key. This function
%   is used by |\anskey| in \myenv*[no-index]{enumext}, |\item*| in
%   \myenv*[no-index]{keyans} and |\anspic| in \myenv*[no-index]{keyanspic}.
%
%   The form in which the \mymarg{content} is stored in \emph{sequence} is in
%   an internal \myenv*[no-index]{enumext} or \myenv*[no-index]{enumext*} environments with
%   the \emph{\enquote{same structure}} in which the command was executed.
%
%   The \emph{\enquote{stored content}} is retrieved by means of the
%   |\printkeyans| command.
%    \begin{macrocode}
\cs_new_protected:Nn \@@_store_addto_seq:n
  {
    \seq_gput_right:cn { g_@@_ \l_@@_store_name_tl _seq } {#1}
  }
\cs_generate_variant:Nn \@@_store_addto_seq:n { v, V }
%    \end{macrocode}
% \end{macro}
%
% \subsubsection{Functions for storing structure in the sequence}
%
% \begin{macro}[int]{\@@_store_level_open:, \@@_store_level_close:}
%   The \emph{\enquote{storing structure}} is handled by the functions
%   \myhlc|\@@_store_level_open:| and \myhlc|\@@_store_level_close:| which are
%   executed per level within the \myenv*[no-index]{enumext} environment.
%    \begin{macrocode}
\cs_new_protected:Nn \@@_store_level_open:
  {
    \bool_if:NT \l_@@_check_answers_bool
      {
        \tl_if_empty:cTF { l_@@_store_save_key_ \@@_level: _tl }
          {
            \@@_store_addto_seq:n
              {
                \item \begin{enumext}
              }
          }
          {
            \tl_put_left:cn { l_@@_store_save_key_ \@@_level: _tl }
              {
                \item \begin{enumext} [
              }
            \tl_put_right:cn { l_@@_store_save_key_ \@@_level: _tl }
              {
                ]
              }
            \@@_store_addto_seq:v { l_@@_store_save_key_ \@@_level: _tl }
          }
      }
  }
\cs_new_protected:Nn \@@_store_level_close:
  {
    \bool_if:NT \l_@@_check_answers_bool
      {
        \@@_store_addto_seq:n { \end{enumext} }
      }
  }
%    \end{macrocode}
% \end{macro}
%
% \begin{macro}[int]{\@@_store_level_open_vii:, \@@_store_level_close_vii:}
%   The \emph{\enquote{storing structure}} is handled by the functions
%   \myhlc|\@@_store_level_open_vii:| and \myhlc|\@@_store_level_close_vii:|
%   which are executed in the \myenv*[no-index]{enumext*} environment.
%
%    \begin{macrocode}
\cs_new_protected:Nn \@@_store_level_open_vii:
  {
    \bool_if:NT \l_@@_check_answers_bool
      {
        \tl_if_empty:NTF \l_@@_store_save_key_vii_tl
          {
            \@@_store_addto_seq:n
              {
                \item \begin{enumext*}
              }
          }
          {
            \tl_put_left:Nn \l_@@_store_save_key_vii_tl
              {
                \item \begin{enumext*}[
              }
            \tl_put_right:Nn \l_@@_store_save_key_vii_tl
              {
                ]
              }
            \@@_store_addto_seq:V \l_@@_store_save_key_vii_tl
          }
      }
  }
\cs_new_protected:Nn \@@_store_level_close_vii:
  {
    \bool_if:NT \l_@@_check_answers_bool
      {
        \@@_store_addto_seq:n { \end{enumext*} }
      }
  }
%    \end{macrocode}
% \end{macro}
%
% \subsubsection{Function for show marks and position}
%
% \begin{macro}[int]{\@@_print_keyans_box:NN, \@@_print_keyans_box:cc}
%   The function \myhlc|\@@_print_keyans_box:NN| print a box in the left
%   margin with \myhlc|\l_@@_mark_answer_sym_tl| used by the \mykey[env]{wrap-ans},
%   \mykey[env]{show-ans} and \mykey[env]{show-pos} keys. The function takes two
%   arguments:
%
% \begin{arguments}
%   \item \myhlc|\l_@@_labelwidth_X_dim|
%   \item \myhlc|\l_@@_labelsep_X_dim|
% \end{arguments}
%
%    \begin{macrocode}
\cs_new_protected:Nn \@@_print_keyans_box:NN
  {
    \mode_leave_vertical:
    \skip_horizontal:n { -\dim_use:N #2 }
    \hbox_overlap_left:n
      {
        \makebox[ \dim_use:N #1 ][ \l_@@_mark_position_str ]
          {
            \tl_use:N \l_@@_mark_answer_sym_tl
          }
      }
    \skip_horizontal:n { \dim_use:N #2 }
  }
\cs_generate_variant:Nn \@@_print_keyans_box:NN { cc }
%    \end{macrocode}
% \end{macro}
%
% \subsection{The internal label and ref}\label{subsec:internal-label-ref}
%
% The function \myhlc|\@@_store_internal_ref:| handles the
% \emph{\enquote{internal label and ref}} system used by the \mykey{save-ref}
% and \mykey{mark-ref} keys for |\anskey| allow executing
% |\ref|\mymarg{store name \textnormal{\textcolor{red}{:}} position} and
% will return |1||.(a)||.i||.A|.
%
% \begin{macro}[int]{\@@_store_internal_ref:}
%   First the dots \enquote{\textcolor{red}{.}} are removed from
%   the current \mymeta{labels}, since double dots in references are not
%   wanted, then the result is placed in the variable
%   \myhlc|\l_@@_newlabel_arg_two_tl|.
%    \begin{macrocode}
\cs_new_protected:Nn \@@_store_internal_ref:
  {
    \cs_set_protected:Nn \@@_tmp:n
      {
        \tl_set_eq:cc { l_@@_label_copy_ ##1 _tl } { l_@@_label_ ##1 _tl }
        \tl_reverse:c { l_@@_label_copy_ ##1 _tl }
        \tl_remove_once:cn { l_@@_label_copy_ ##1 _tl } { . }
        \tl_reverse:c { l_@@_label_copy_ ##1 _tl }
      }
    \clist_map_inline:nn { i, ii, iii, iv, vii } { \@@_tmp:n {##1} }
    \cs_set:Nn \@@_tmp:n
      { . \tl_use:c { l_@@_label_copy_ \int_to_roman:n {##1} _tl } }
%    \end{macrocode}
%   Here it is necessary to analyze the cases where the environment is started with
%   \myenv*[no-index]{enumext*} and if |\anskey| or \myenv*[no-index]{anskey*} is running
%   alone in it or if it is running in a nested \myenv*[no-index]{enumext} environment
%   within the starting environment.
%    \begin{macrocode}
    \bool_lazy_all:nT
      {
        { \bool_if_p:N \g_@@_starred_bool }
        { \int_compare_p:nNn { \l_@@_level_int } = { 0 } }
      }
      {
        \tl_put_right:Ne \l_@@_newlabel_arg_two_tl
          { \tl_use:N \l_@@_label_copy_vii_tl }
      }
    \bool_lazy_all:nT
      {
        { \bool_not_p:n { \g_@@_standard_bool } }
        { \bool_if_p:N \l_@@_standard_bool }
        { \int_compare_p:nNn { \l_@@_level_int } > { 0 } }
      }
      {
        \tl_put_right:Ne \l_@@_newlabel_arg_two_tl
          {
            \tl_use:N \l_@@_label_copy_vii_tl
            \int_step_function:nnN { 1 } { \l_@@_level_int } \@@_tmp:n
          }
      }
%    \end{macrocode}
%   If started with \myenv*[no-index]{enumext} and if |\anskey| or \myenv*[no-index]{anskey*}
%   is running alone in it or if it is running in a nested \myenv*[no-index]{enumext*}
%   environment within the starting environment.
%    \begin{macrocode}
    \bool_lazy_all:nT
      {
        { \bool_if_p:N \g_@@_standard_bool }
        { \int_compare_p:nNn { \l_@@_level_int } > { 0 } }
        { \int_compare_p:nNn { \l_@@_level_h_int } = { 0 } }
      }
      {
        \tl_put_right:Ne \l_@@_newlabel_arg_two_tl
          {
            \tl_use:N \l_@@_label_copy_i_tl
            \int_step_function:nnN { 2 } { \l_@@_level_int } \@@_tmp:n
          }
      }
    \cs_set:Nn \@@_tmp:n
      { \tl_use:c { l_@@_label_copy_ \int_to_roman:n {##1} _tl } . }
    \bool_lazy_all:nT
      {
        { \bool_if_p:N \g_@@_standard_bool }
        { \bool_if_p:N \l_@@_starred_bool }
        { \int_compare_p:nNn { \l_@@_level_int } > { 0 } }
      }
      {
        \tl_put_right:Ne \l_@@_newlabel_arg_two_tl
          {
            \int_step_function:nnN { 1 } { \l_@@_level_int } \@@_tmp:n
            \tl_use:N \l_@@_label_copy_vii_tl
          }
      }
%    \end{macrocode}
%   Now the variable \myhlc|\l_@@_newlabel_arg_one_tl| is set, which
%   contains \mymarg{store name \textnormal{\textcolor{red}{:}} position}.
%    \begin{macrocode}
    \tl_put_right:Ne \l_@@_newlabel_arg_one_tl
      {
        \l_@@_store_name_tl \c_colon_str
        \int_eval:n { \prop_count:c { g_@@_ \l_@@_store_name_tl _prop } }
      }
%    \end{macrocode}
%   Now execute the function \myhlc|\@@_newlabel:nn| and save the result in
%   the variable \myhlc|\l_@@_write_aux_file_tl| and finally write it in the
%   |.aux| file.
%    \begin{macrocode}
    \tl_put_right:Ne \l_@@_write_aux_file_tl
      {
        \@@_newlabel:nn
          { \exp_not:V \l_@@_newlabel_arg_one_tl }
          { \l_@@_newlabel_arg_two_tl }
      }
    \l_@@_write_aux_file_tl
  }
%    \end{macrocode}
% \end{macro}
%
% \subsection{Common functions for \cs{anskey} and \texttt{anskey*} environment}\label{subsec:common-anskey}
%
% \subsubsection{Definition of error messages}
%
% \begin{macro}[int]{anskey-unnumber-item, anskey-empty-arg, anskey-math-mode}
% Messages used by |\anskey| and \myenv*[no-index]{anskey*}.
%    \begin{macrocode}
\msg_new:nnn { enumext } { anskey-unnumber-item }
  {
    Can't~store~with~an~unnumbered~'\token_to_str:N \item'~\msg_line_context:.
  }
\msg_new:nnn { enumext } { anskey-empty-arg }
  {
    Can't~store~empty~content~\msg_line_context:.
  }
\msg_new:nnn { enumext } { anskey-math-mode }
  {
    '#1'~can't~work~in~math~mode~\msg_line_context:.
  }
%    \end{macrocode}
% \end{macro}
%
% \subsubsection{Internal functions common to \cs[no-index]{anskey} and \env{anskey*}}
%
% \begin{macro}[int]{\@@_store_anskey_arg:n}
%   The internal function \myhlc|\@@_store_anskey_arg:n| first passes the
%   \mymarg{argument} to the \emph{prop list}, then checks the
%   state of the variable \myhlc|\l_@@_store_ref_key_bool| handled by the
%   \mykey[env]{save-ref} key and calls the function \myhlc|\@@_store_internal_ref:|
%   for the \emph{\enquote{internal label and ref}} system. After
%   this, if the \mykey[env]{show-ans} or \mykey[env]{show-pos} keys are active, it
%   shows the \emph{\enquote{wrapped}} \mymarg{argument}.
%    \begin{macrocode}
\cs_new_protected:Nn \@@_store_anskey_arg:n
  {
    \int_gincr:N \g_@@_item_anskey_int
    \@@_store_addto_prop:n {#1}
    \bool_if:NT \l_@@_store_ref_key_bool
      {
        \@@_store_internal_ref:
      }
    \@@_anskey_show_wrap_left:n {#1}
%    \end{macrocode}
%   Now the processing of the \myoarg{key \textnormal{\textcolor{gray}{=}} val}
%   passed to the command starts, building the |\item| in the variable \myhlc|\l_@@_store_anskey_arg_tl|
%   which is \emph{\enquote{stored}} in the \emph{sequence}.
%   First the variable \myhlc|\l_@@_store_anskey_arg_tl| is cleared and the
%   \mymeta{keys} are processed; if the \mykey[anskey,anskey*]{break-col} key is present and
%   the command is running under \myenv*[no-index]{enumext} (not in \myenv*[no-index]{enumext*}),
%   |\columnbreak| and then |\item| are added.
%    \begin{macrocode}
    \tl_clear:N \l_@@_store_anskey_arg_tl
    \bool_lazy_and:nnT
      { \bool_if_p:N \l_@@_store_columns_break_bool }
      { \bool_not_p:n { \l_@@_starred_bool } }
      {
        \tl_put_left:Nn \l_@@_store_anskey_arg_tl { \columnbreak }
      }
    \tl_put_right:Nn \l_@@_store_anskey_arg_tl { \item }
%    \end{macrocode}
%   If the \mykey[anskey,anskey*]{item-join} key is present and the command is running under
%   \myenv*[no-index]{enumext*}, \myparg{number} is added to \myhlc|\l_@@_store_anskey_arg_tl|.
%    \begin{macrocode}
    \bool_lazy_and:nnT
      { \bool_not_p:n { \l_@@_starred_bool } }
      { \int_compare_p:nNn { \l_@@_store_item_join_int } > { 1 } }
      {
        \tl_put_right:Ne \l_@@_store_anskey_arg_tl
          {
            ( \exp_not:V \l_@@_store_item_join_int )
          }
      }
%    \end{macrocode}
%   Then the keys \mykey[anskey,anskey*]{item-star}, \mykey[anskey,anskey*]{item-sym*} and
%   \mykey[anskey,anskey*]{item-pos*} are reviewed and passed to \myhlc|\l_@@_store_anskey_arg_tl|
%   along with the \mymarg{argument} for |\anskey| or \mymeta{body} for |anskey*|.
%    \begin{macrocode}
    \bool_if:NTF \l_@@_store_item_star_bool
      {
        \tl_put_right:Nn \l_@@_store_anskey_arg_tl { * }
        \tl_if_empty:NF \l_@@_store_item_symbol_tl
          {
            \tl_put_right:Ne \l_@@_store_anskey_arg_tl
              {
                [ \exp_not:V \l_@@_store_item_symbol_tl ]
              }
          }
        \dim_compare:nT
          {
            \l_@@_store_item_symbol_sep_dim != \c_zero_dim
          }
          {
            \tl_put_right:Ne \l_@@_store_anskey_arg_tl
              {
                [ \exp_not:V \l_@@_store_item_symbol_sep_dim ]
              }
          }
        \tl_put_right:Nn \l_@@_store_anskey_arg_tl {#1}
      }
      {
        \tl_put_right:Nn \l_@@_store_anskey_arg_tl {#1}
      }
%    \end{macrocode}
%   Finally it checks if the \mykey{save-ref} key is active along with the
%   \mypkg[no-index]{hyperref} package load, if both conditions are met, it will
%   create the |\hyperlink| with \emph{\enquote{symbol}} set by \mykey[env]{mark-ref}
%   key and then store in \emph{sequence}.
%    \begin{macrocode}
    \bool_lazy_and:nnT
      { \bool_if_p:N \l_@@_store_ref_key_bool }
      { \bool_if_p:N \l_@@_hyperref_bool }
      {
        \tl_put_right:Ne \l_@@_store_anskey_arg_tl
          {
            \hfill \exp_not:N \hyperlink
              { \exp_not:V \l_@@_newlabel_arg_one_tl }
              { \exp_not:V \l_@@_mark_ref_sym_tl }
          }
      }
    \@@_store_addto_seq:V \l_@@_store_anskey_arg_tl
  }
%    \end{macrocode}
% \end{macro}
%
% \begin{macro}[int]{\@@_anskey_show_wrap_arg:n}
%   The function \myhlc|\@@_anskey_show_wrap_arg:n| \emph{\enquote{wraps}}
%   the \mymarg{argument} passed to |\anskey| and the \mymeta{body}
%   for \myenv*[no-index]{anskey*} when using the \mykey[env]{wrap-ans} and \mykey{wrap-sep} keys.
%    \begin{macrocode}
\cs_new_protected:Nn \@@_anskey_show_wrap_arg:n
  {
    \par
    \bool_if:NTF \l_@@_starred_bool
      {
        \dim_compare:nNnT { \l_@@_mark_sym_sep_dim } = { \c_zero_dim }
          {
            \dim_set:Nn \l_@@_mark_sym_sep_dim { \l_@@_labelsep_vii_dim }
          }
        \@@_print_keyans_box:NN
          \l_@@_labelwidth_vii_dim \l_@@_mark_sym_sep_dim
      }
      {
        \dim_compare:nNnT { \l_@@_mark_sym_sep_dim } = { \c_zero_dim }
          {
            \dim_set:Nn \l_@@_mark_sym_sep_dim
              {
                \dim_use:c {l_@@_labelsep_ \@@_level: _dim }
              }
          }
        \@@_print_keyans_box:cc
          { l_@@_labelwidth_ \@@_level: _dim } { l_@@_mark_sym_sep_dim }
      }
    \@@_anskey_wrapper:n {#1}
  }
%    \end{macrocode}
% \end{macro}
%
% \begin{macro}[int]{\@@_anskey_show_wrap_left:n}
%   The function \myhlc|\@@_anskey_show_wrap_left:n| will show the
%   \emph{\enquote{mark}} defined by the \mykey[env]{mark-ans} key or the
%   \emph{\enquote{position}} of the \mymarg{content} stored in the \emph{prop list}
%   when using the \mykey[env]{show-pos} key on the left margin next to the
%   \emph{\enquote{wraps}} \mymarg{argument} passed to |\anskey| and
%   the \mymeta{body} in \myenv*[no-index]{anskey*} on the right side when using
%   the \mykey[env]{show-ans} key.
%    \begin{macrocode}
\cs_new_protected:Nn \@@_anskey_show_wrap_left:n
  {
    \bool_if:NT \l_@@_show_answer_bool
      {
        \@@_anskey_show_wrap_arg:n {#1}
      }
    \bool_if:NT \l_@@_show_position_bool
      {
        \tl_set:Ne \l_@@_mark_answer_sym_tl
          {
            \group_begin:
            \exp_not:N \normalfont
            \exp_not:N \footnotesize [ \int_eval:n
              {
                \prop_count:c { g_@@_ \l_@@_store_name_tl _prop }
              }
              ]
            \group_end:
          }
        \@@_anskey_show_wrap_arg:n {#1}
      }
  }
%    \end{macrocode}
% \end{macro}
%
% \subsection{The command \cs{anskey}}\label{cmd:anskey}
%
% Since \emph{\enquote{content}} is stored in a \myenv[no-index]{list} environment
% within \emph{sequences} and the options passed to each level can be managed
% (more or less), it is necessary to have a little more
% control over |\item| when storing.
%
% The |\anskey| command will cover this point and give it similar
% behaviour to that of |\item| in the \myenv*[no-index]{enumext} and \myenv*[no-index]{enumext*}
% environments executed as follows |\anskey|\myoarg{key \textnormal{\textcolor{gray}{=}} val}\mymarg{content}.
%
% \subsubsection{Definition of error messages for \cs[no-index]{anskey}}
%
% \begin{macro}[int]{anskey-wrong-place, anskey-nested,
%                    anskey-cmd-key-unknown, anskey-cmd-key-value-unknown}
% Messages used by |\anskey| command.
%    \begin{macrocode}
\msg_new:nnnn { enumext } { anskey-wrong-place }
  {
    Wrong~place~for~command~'\c_backslash_str #1'~\msg_line_context:.
  }
  {
    '\c_backslash_str #1'~works~in~the~environment~'#2'.
  }
\msg_new:nnn { enumext } { anskey-nested }
  {
    The~command~'\token_to_str:N \anskey'~can't~be~nested~
    \msg_line_context:.
  }
\msg_new:nnnn { enumext } { anskey-cmd-key-unknown }
  {
    The~key~'#1'~is~unknown~by~'\token_to_str:N \anskey'~and~is~being~
    ignored~\msg_line_context:.
  }
  {
    The~command~'\token_to_str:N \anskey'~does~not~have~a~key~called~'#1'.~
    Check~that~you~have~spelled~the~key~name~correctly.
  }
\msg_new:nnnn { enumext } { anskey-cmd-key-value-unknown }
  {
    The~key~'#1=#2'~is~unknown~by~'\token_to_str:N \anskey'~and~is~being~
    ignored~\msg_line_context:.
  }
  {
    The~command~'\token_to_str:N \anskey'~does~not~have~a~key~called~'#1'.~
    Check~that~you~have~spelled~the~key~name~correctly.
  }
%    \end{macrocode}
% \end{macro}
%
% \subsubsection{Definition of keys for \cs[no-index]{anskey}}
%
% \begin{macro}[int]{break-col, item-join, item-star, item-sym*, item-pos*, unknown}
%   First, the keys \mykey[anskey,anskey*]{break-col}, \mykey[anskey,anskey*]{item-join},
%   \mykey[anskey,anskey*]{item-star}, \mykey[anskey,anskey*]{item-sym*} and \mykey[anskey,anskey*]{item-pos*} are added.
%    \begin{macrocode}
\keys_define:nn { enumext / anskey }
  {
    break-col .bool_set:N = \l_@@_store_columns_break_bool,
    break-col .default:n = true,
    break-col .value_forbidden:n = true,
    item-join .int_set:N = \l_@@_store_item_join_int,
    item-join .value_required:n = true,
    item-star .bool_set:N = \l_@@_store_item_star_bool,
    item-star .default:n = true,
    item-star .value_forbidden:n = true,
    item-sym* .tl_set:N = \l_@@_store_item_symbol_tl,
    item-sym* .value_required:n = true,
    item-pos* .dim_set:N = \l_@@_store_item_symbol_sep_dim,
    item-pos* .value_required:n = true,
    unknown   .code:n = { \@@_anskey_unknown:n {#1} },
  }
%    \end{macrocode}
% \end{macro}
%
% \subsubsection{Internal functions for the command}
%
% \begin{macro}[int]{\@@_anskey_unknown:n, \@@_anskey_unknown:nn}
%   The \mymeta{keys} are stored in \cs{l_keys_key_str} and the value
%   (if any) is passed as an argument to the function \myhlc|\@@_anskey_unknown:n|.
%    \begin{macrocode}
\cs_new_protected:Nn \@@_anskey_unknown:n
  {
    \exp_args:NV \@@_anskey_unknown:nn \l_keys_key_str {#1}
  }
\cs_new_protected:Nn \@@_anskey_unknown:nn
  {
    \tl_if_blank:nTF {#2}
      {
        \msg_error:nnn { enumext } { anskey-cmd-key-unknown } {#1}
      }
      {
        \msg_error:nnnn { enumext } { anskey-cmd-key-value-unknown } {#1} {#2}
      }
  }
%    \end{macrocode}
% \end{macro}
%
% \begin{macro}[int]{\@@_anskey_safe_outer:, \@@_anskey_safe_inner:}
%   The \myhlc|\@@_store_anskey_safe_outer:| function will return the
%   appropriate messages when the command is executed outside the
%   environment in which the \mykey[env]{save-ans} key was activated.
%    \begin{macrocode}
\cs_new_protected:Nn \@@_anskey_safe_outer:
  {
    \bool_if:NF \l_@@_store_active_bool
      {
        \msg_error:nnnn { enumext } { anskey-wrong-place }{ anskey }{ enumext }
      }
    \int_compare:nNnT { \l_@@_keyans_level_int } = { 1 }
      {
        \msg_error:nnnn { enumext } { command-wrong-place }{ anskey }{ keyans }
      }
    \int_compare:nNnT { \l_@@_keyans_level_h_int } = { 1 }
      {
        \msg_error:nnnn { enumext } { command-wrong-place }{ anskey }{ keyans* }
      }
    \int_compare:nNnT { \l_@@_keyans_pic_level_int } = { 1 }
      {
        \msg_error:nnnn { enumext }
          { command-wrong-place }{ anskey }{ keyanspic }
      }
  }
%    \end{macrocode}
%   The \myhlc|\@@_anskey_safe_inner:| function will first check if the
%   command is nested, if preceded by a not numbered |\item| or if it is
%   in \emph{math mode} returning the appropriate messages.
%    \begin{macrocode}
\cs_new_protected:Nn \@@_anskey_safe_inner:
  {
    \int_incr:N \l_@@_anskey_level_int
    \int_compare:nNnT { \l_@@_anskey_level_int } > { 1 }
      {
        \msg_error:nn { enumext } { anskey-nested }
      }
    \bool_if:NF \l_@@_item_number_bool
      {
        \msg_error:nn { enumext } { anskey-unnumber-item }
      }
    \mode_if_math:T
      {
        \msg_error:nne { enumext } { anskey-math-mode }
          { \token_to_str:N \anskey }
      }
  }
%    \end{macrocode}
% \end{macro}
%
% \begin{important}*
%   The |\anskey| command will only be present when using the
%   \mykey{save-ans} key in \myenv*[no-index]{enumext} and \myenv*[no-index]{enumext*}
%   environments, otherwise it will return an error.
% \end{important}
%
% \begin{macro}{\anskey}
%   The function \myhlc|\@@_anskey_safe_outer:| is called first to be
%   sure where the command is executed, then the state of the
%   variable \myhlc|\l_@@_check_answers_bool| set by the key \mykey[env]{no-store} is checked;
%   if it is true, \myhlc|\g_@@_item_anskey_int| is incremented for the internal
%   \emph{\enquote{check answer}} system and execute the function
%   \myhlc|\@@_anskey_safe_inner:n| to ensure that the command is not
%   nested and that the argument is not empty, finally search the
%   \myoarg{key \textnormal{\textcolor{gray}{=}} val} and call the
%   function \myhlc|\@@_store_anskey_arg:n|.
%    \begin{macrocode}
\NewDocumentCommand \anskey { o +m }
  {
    \@@_anskey_safe_outer:
    \group_begin:
      \bool_if:NT \l_@@_check_answers_bool
        {
          \tl_if_novalue:nF {#1}
            {
              \keys_set:nn { enumext / anskey } {#1}
            }
          \tl_if_blank:nTF {#2}
            {
              \msg_error:nn { enumext } { anskey-empty-arg }
            }
            {
              \@@_anskey_safe_inner:
              \@@_store_anskey_arg:n {#2}
            }
        }
    \group_end:
  }
%    \end{macrocode}
% \end{macro}
%
% \subsection{The environment \env{anskey*}}\label{env:anskey}
%
% The original implementation of the |anskey*| environment used
% non-public functions from the \mypkg{scontents}\cite{scontents} package,
% which was not the best approach. Fortunately \hologo{LaTeX} release 2025-06-01
% implemented the new \texttt{c-type} argument in the \mypkg[no-index]{ltcmd}\cite{ltnews},
% with which the \mymeta{body} of the environment can be recorded in
% \emph{verbatim mode}, and \myhlc|\tl_retokenize:n| (wrapper around the |\scantokens|)
% provided by \hologo{LaTeX}3 release 2025-07-08 does the work as the
% original implementation.
%
% \subsubsection{Definition of error messages for \env{anskey*}}
%
% \begin{macro}[int]{anskey-env-error}
% Message used by reserved \myenv*[no-index]{anskey*} environment by \mypkg*[no-index]{enumext}
% package.
%    \begin{macrocode}
\msg_new:nnnn { enumext } { anskey-env-error }
  {
    The~environment~'#1'~is~reserved~by~'enumext'~package,~
    it~is~already~defined~\msg_line_context:.
  }
  {
    The~environment~'#1'~is~defined~internally~
    for~the~'save-ans'~key~with~save-ans~key~active.~See~documentation.
  }
%    \end{macrocode}
% \end{macro}
%
% \begin{macro}[int]{anskey-env-wrong, anskey-env-key-unknown,
%                    anskey-env-key-value-unknown, overwrite-file,
%                    writing-file, not-writing}
% Messages used by \myenv*[no-index]{anskey*} environment.
%    \begin{macrocode}
\msg_new:nnn { enumext } { anskey-env-wrong }
  {
    The~environment~anskey*~cannot~be~used~in~'#1'~\msg_line_context:.
  }
\msg_new:nnnn { enumext } { anskey-env-key-unknown }
  {
    The~key~'#1'~is~unknown~by~environment~
    'anskey*'~and~is~being~ignored~\msg_line_context:.
  }
  {
    The~environment~'anskey*'~does~not~have~a~key~called~'#1'.~
    Check~that~you~have~spelled~the~key~name~correctly.
  }
\msg_new:nnnn { enumext } { anskey-env-key-value-unknown }
  {
    The~key~'#1=#2'~is~unknown~by~environment~
    'anskey*'~and~is~being~ignored~\msg_line_context:.
  }
  {
    The~environment~'anskey*'~does~not~have~a~key~called~'#1'.~
    Check~that~you~have~spelled~the~key~name~correctly.
  }
\msg_new:nnn { enumext } { overwrite-file }
  {
    Overwriting~file~'#1'.
  }
\msg_new:nnn { enumext } { writing-file }
  {
    Writing~file~'#1'.
  }
\msg_new:nnn { enumext } { not-writing }
  {
    File~'#1'~already~exists.~Not~writing.
  }
%    \end{macrocode}
% \end{macro}
%
% \subsubsection{Definition of keys for \env{anskey*}}
%
% \begin{macro}[int]{break-col, item-join, item-star, item-sym*, item-pos*,
%                    force-eol, write-env, overwrite, unknown}
%   First, the same keys from the |\anskey| command are added along with the
%   \mykey[anskey,anskey*]{force-eol}, \mykey[anskey,anskey*]{write-env} and \mykey[anskey,anskey*]{overwrite}
%   keys that were in the original implementation that used the
%   \mypkg[no-index]{scontents} support package for these.
%    \begin{macrocode}
\keys_define:nn { enumext / anskey* }
  {
    break-col .bool_set:N = \l_@@_store_columns_break_bool,
    break-col .default:n = true,
    break-col .value_forbidden:n = true,
    item-join .int_set:N = \l_@@_store_item_join_int,
    item-join .value_required:n = true,
    item-star .bool_set:N = \l_@@_store_item_star_bool,
    item-star .default:n = true,
    item-star .value_forbidden:n = true,
    item-sym* .tl_set:N = \l_@@_store_item_symbol_tl,
    item-sym* .value_required:n = true,
    item-pos* .dim_set:N = \l_@@_store_item_symbol_sep_dim,
    item-pos* .value_required:n = true,
    force-eol .bool_set:N = \l_@@_anskey_env_force_eol_bool,
    force-eol .initial:n = false,
    force-eol .default:n = true,
    write-env .code:n =
      {
        \bool_set_true:N \l_@@_write_anskey_env_bool
        \tl_set:Nn \l_@@_write_anskey_env_file_name_tl {#1}
      },
    write-env .value_required:n = true,
    overwrite .bool_set:N = \l_@@_anskey_env_overwrite_bool,
    overwrite .initial:n = false,
    overwrite .default:n = true,
    unknown   .code:n = { \@@_anskey_env_unknown:n {#1} },
  }
%    \end{macrocode}
% \end{macro}
%
% \subsubsection{Internal functions for the environment}
%
% \begin{macro}[int]{\@@_anskey_env_unknown:n, \@@_anskey_env_unknown:nn}
%   The \mymeta{keys} are stored in \cs{l_keys_key_str} and the value
%   (if any) is passed as an argument to the function \myhlc|\@@_anskey_env_unknown:n|.
%    \begin{macrocode}
\cs_new_protected:Nn \@@_anskey_env_unknown:n
  {
    \exp_args:NV \@@_anskey_env_unknown:nn \l_keys_key_str {#1}
  }
\cs_new_protected:Nn \@@_anskey_env_unknown:nn
  {
    \tl_if_blank:nTF {#2}
      {
        \msg_error:nnn { enumext } { anskey-env-key-unknown } {#1}
      }
      {
        \msg_error:nnnn { enumext } { anskey-env-key-value-unknown } {#1} {#2}
      }
  }
%    \end{macrocode}
% \end{macro}
%
% \begin{macro}[int]{\@@_anskey_env_file_if_writable:n, \@@_anskey_env_file_if_writable:nT,
%                    \@@_anskey_env_file_if_writable:nF, \@@_anskey_env_file_if_writable:nTF,
%                    \@@_anskey_env_file_write:nn, \@@_anskey_env_file_write:VV}
%   The conditional function \myhlc|\@@_anskey_env_file_if_writable:n| used by
%   the \mykey[anskey*]{write-env} and \mykey[anskey*]{overwrite} keys in the \myenv*[no-index]{anskey*}
%   environment to determine whether the output file is written or overwritten.
%    \begin{macrocode}
\prg_new_protected_conditional:Npnn \@@_anskey_env_file_if_writable:n #1
  { T, F, TF }
  {
    \bool_if:NTF \l_@@_write_anskey_env_bool
      {
        \file_if_exist:nTF {#1}
          {
            \bool_if:NTF \l_@@_anskey_env_overwrite_bool
              {
                \msg_warning:nne { enumext } { overwrite-file } {#1}
                \prg_return_true:
              }
              {
                \msg_warning:nne { enumext } { not-writing } {#1}
                \prg_return_false:
              }
          }
          {
            \msg_warning:nne { enumext } { writing-file } {#1}
            \prg_return_true:
          }
      }
      { \prg_return_false: }
  }
%    \end{macrocode}
%   The \myhlc|\@@_anskey_env_file_write:nn| function is used by the
%   \mykey{write-env} key in the \myenv*[no-index]{anskey*} environment to write the
%   output file with the \mymeta{body} of the environment.
%    \begin{macrocode}
\cs_new_protected:Nn \@@_anskey_env_file_write:nn
  {
    \@@_anskey_env_file_if_writable:nT {#1}
      {
        \iow_open:Nn \l_@@_write_anskey_env_file_iow {#1}
        \iow_now:Nn  \l_@@_write_anskey_env_file_iow {#2}
        \iow_close:N \l_@@_write_anskey_env_file_iow
      }
  }
\cs_generate_variant:Nn \@@_anskey_env_file_write:nn { VV }
%    \end{macrocode}
% \end{macro}
%
% \begin{macro}[int]{\@@_anskey_env_safe_outer:, \@@_anskey_env_safe_inner:,
%                    \@@_store_anskey_env:n}
%   The function \myhlc|\@@_store_anskey_safe_outer:| will return the
%   appropriate messages when |anskey*| is executed outside the
%   environment in which the \mykey[env]{save-ans} key was activated
%   or within the \myenv*[no-index]{keyans}, \myenv*[no-index]{keyans*} or \myenv*[no-index]{keyanspic}
%   environments.
%    \begin{macrocode}
\cs_new_protected:Nn \@@_anskey_env_safe_outer:
  {
    \bool_if:NF \l_@@_store_active_bool
      {
        \msg_error:nnn { enumext } { anskey-env-error } { anskey* }
      }
    \int_compare:nNnT { \l_@@_keyans_level_int } = { 1 }
      {
        \msg_error:nnn { enumext } { anskey-env-wrong }{ keyans }
      }
    \int_compare:nNnT { \l_@@_keyans_level_h_int } = { 1 }
      {
        \msg_error:nnn { enumext } { anskey-env-wrong } { keyans* }
      }
    \int_compare:nNnT { \l_@@_keyans_pic_level_int } = { 1 }
      {
        \msg_error:nnn { enumext } { anskey-env-wrong } { keyanspic }
      }
  }
%    \end{macrocode}
%   The function \myhlc|\@@_anskey_env_safe_inner:| will first check if
%   preceded by a not numbered |\item| or if it is in \emph{math mode}
%   returning the appropriate messages.
%    \begin{macrocode}
\cs_new_protected:Nn \@@_anskey_env_safe_inner:
  {
    \bool_if:NF \l_@@_item_number_bool
      {
        \msg_error:nn { enumext } { anskey-unnumber-item }
      }
    \mode_if_math:T
      {
        \msg_error:nnn { enumext } { anskey-math-mode } { anskey* }
      }
  }
%    \end{macrocode}
%   The \myhlc|\@@_store_anskey_env:n| function will first pass the
%   \texttt{c-type} argument \mymeta{body} to the variable \myhlc|\l_@@_store_anskey_env_tl| and
%   replace the macro |\obeyedline| with |^^J| and then execute the
%   |write-env| and |overwrite| keys, check the state of the
%   variable \myhlc|\l_@@_anskey_env_force_eol_bool| managed by the |force-eol|
%   key, and adds \myhlc|\c_@@_anskey_env_hidden_space_str| if
%   necessary. Finally it uses \myhlc|\exp_args:Ne| on the
%   \myhlc|\@@_store_anskey_arg:n| to expand the \myhlc|\tl_retokenize:n|
%   function which rescans the \myhlc|\l_@@_store_anskey_env_tl| variable
%   before processing it.
%    \begin{macrocode}
\cs_new_protected:Nn \@@_store_anskey_env:n
  {
    \tl_set:Nn \l_@@_store_anskey_env_tl {#1}
    \RenewDocumentCommand \obeyedline { } { \iow_char:N \^^J }
    \tl_replace_all:Nee \l_@@_store_anskey_env_tl { \obeyedline }
      { \iow_char:N \^^J }
    \@@_anskey_env_file_write:VV
      \l_@@_write_anskey_env_file_name_tl \l_@@_store_anskey_env_tl
    \bool_if:NF \l_@@_anskey_env_force_eol_bool
      {
        \tl_put_right:Ne \l_@@_store_anskey_env_tl
          {
            \c_@@_anskey_env_hidden_space_str
          }
      }
    \exp_args:Ne
      \@@_store_anskey_arg:n
        {
          \tl_retokenize:n { \l_@@_store_anskey_env_tl }
        }
  }
%    \end{macrocode}
%
% \begin{important}*
%  Since |\obeyedline| can be redefined by the user, for example to
%  |\mbox{}\par|, it is necessary to redefine it to |^^J| in order to use
%  \myhlc|\tl_replace_all:Nee| otherwise it returns an error.
% \end{important}
% \end{macro}
%
% \begin{macro}{anskey*}
%   First, the function \myhlc|\@@_anskey_env_safe_outer:| is called to make
%   sure where the environment is running, then the state
%   of the variable \myhlc|\l_@@_check_answers_bool| set by the key
%   \mykey[env]{no-store} is checked. If it is true, it looks for \myoarg{key \textnormal{\textcolor{gray}{=}} val}
%   and verifies that the \emph{argument} \texttt{c-type} \mymeta{body} is not empty.
%   Finally, it runs the internal check function \myhlc|\@@_anskey_env_safe_inner:n|
%   and call the function \myhlc|\@@_store_anskey_arg:n|.
%    \begin{macrocode}
\NewDocumentEnvironment{anskey*}{ o c }
  {
    \@@_anskey_env_safe_outer:
    \bool_if:NT \l_@@_check_answers_bool
      {
        \tl_if_novalue:nF {#1}
          {
            \keys_set:nn { enumext / anskey* } {#1}
          }
        \tl_if_blank:nTF {#2}
          {
            \msg_error:nn { enumext } { anskey-empty-arg }
          }
          {
            \@@_anskey_env_safe_inner:
            \@@_store_anskey_env:n {#2}
          }
      }
  } { }
%    \end{macrocode}
% \end{macro}
%
% \subsection{Executing \texttt{check-ans} system and write \texttt{.log}}\label{subsec:exec-after-env}
%
% \begin{macro}[int]{\@@_execute_after_env:}
%   The \myhlc|\@@_execute_after_env:| function will first return the
%   appropriate message for the end of the environment in which the
%   \mykey[env]{save-ans} key is being executed, then call the
%   \myhlc|\@@_item_answer_diff:| function and then writes the values of
%   the global variables used to the |.log| file. If the key
%   \mykey[env]{check-ans} is active, it executes the function
%   \myhlc|\@@_check_ans_show:| and shows the result in the terminal, otherwise
%   it executes the function \myhlc|\@@_check_ans_log:| and writes the
%   results in the |.log| file, and finally executes the function
%   \myhlc|\@@_reset_global_vars:| returning the used variables to their
%   original state.
%    \begin{macrocode}
\cs_new_protected:Nn \@@_execute_after_env:
  {
    \int_compare:nNnT { \l_@@_level_int } = { 0 }
      {
        \tl_if_empty:NF \g_@@_store_name_tl
          {
            \@@_stop_save_ans_msg:
            \@@_item_answer_diff:
            \@@_log_global_vars:
            \@@_log_answer_vars:
            \bool_if:NTF \g_@@_check_ans_key_bool
              {
                \@@_check_ans_show:
              }
              { \@@_check_ans_log: }
          }
        \@@_reset_global_vars:
      }
  }
%    \end{macrocode}
%
% \begin{important}*
%  This function is passed to the  function \myhlc|\@@_after_env:nn| for the
%  environments \myenv*[no-index]{enumext} (\S\ref{env:enumext}) and \myenv*[no-index]{enumext*} (\S\ref{env:starred})
%  and it is executed only when the environments are not nested or at some level of these.
% \end{important}
% \end{macro}
%
% \subsection{Common functions for \env{keyans}, \env{keyans*} and \env{keyanspic}}\label{subsec:common-keyans}
%
% \subsubsection{Definition of error messages}
%
% \begin{macro}[int]{keyans-wrong-level, wrong-place}
% Messages used by \myenv*[no-index]{keyans}, \myenv*[no-index]{keyans*} and
% \myenv*[no-index]{keyanspic} environments.
%    \begin{macrocode}
\msg_new:nnnn { enumext } { keyans-wrong-level }
  {
    Wrong~level~position~for~'keyans'~\msg_line_context:.
  }
  {
    The~environment~'keyans'~can~only~be~in~the~first~level.
  }
\msg_new:nnnn { enumext } { wrong-place }
  {
    Wrong~place~for~'#1'~environment~\msg_line_context:.
  }
  {
    '#1'~is~only~found~with~'#2'~in~'enumext'.
  }
%    \end{macrocode}
% \end{macro}
%
% \subsubsection{Storing content in prop list}
%
% \begin{macro}[int]{\@@_keyans_addto_prop:n}
%   The function \myhlc|\@@_keyans_addto_prop:n| will pass the current \mymeta{label} for |\item*|
%   in \myenv*[no-index]{keyans} environment and the current \mymeta{label} for
%   |\anspic*| in \myenv*[no-index]{keyanspic} environment followed by the
%   \mymeta{contents} of the \emph{optional argument} of both commands
%   to the \myhlc|\l_@@_store_current_label_tl| variable, which will be stored
%   to the \emph{prop list} defined by the \mykey[env]{save-ans} key using the
%   function \myhlc|\@@_store_addto_prop:V|.
%    \begin{macrocode}
\cs_new_protected:Nn \@@_keyans_addto_prop:n
  {
    \tl_clear:N \l_@@_store_current_label_tl
    \int_compare:nNnTF { \l_@@_keyans_pic_level_int } = { 1 }
      {
        \tl_put_right:Ne \l_@@_store_current_label_tl { \l_@@_label_vi_tl }
      }
      {
        \tl_put_right:Ne \l_@@_store_current_label_tl { \l_@@_label_v_tl }
      }
%    \end{macrocode}
%   If the \emph{optional argument} is present and the \mykey[env]{save-sep}
%   key is not empty, it is saved.
%    \begin{macrocode}
    \tl_if_novalue:nF {#1}
      {
        \tl_if_empty:NF \l_@@_store_keyans_item_opt_sep_v_tl
          {
            \tl_put_right:NV \l_@@_store_current_label_tl
              \l_@@_store_keyans_item_opt_sep_v_tl
          }
        \tl_put_right:Nn \l_@@_store_current_label_tl {#1}
      }
    \@@_store_addto_prop:V \l_@@_store_current_label_tl
  }
%    \end{macrocode}
% \end{macro}
%
% \subsubsection{The \texttt{save-ref} key for \env{keyans}, \env{keyans*} and \env{keyanspic}}\label{key:save-ref-keyans}
%
% The \emph{\enquote{internal label and ref}} system for the \myenv*[no-index]{keyans},
% \myenv*[no-index]{keyans*} and \myenv*[no-index]{keyanspic} environments has \emph{slight differences}
% with the one implemented for |\anskey| basically because in this
% environments the interest is in the current \mymeta{label} for |\item*|
% and |\anspic*| with the \mymeta{contents} of the \emph{optional argument}. The mechanism
% defined here allows executing |\ref|\mymarg{store name \textcolor{red}{:} position} and
% will return |1.(A)|.
%
% \begin{macro}[int]{\@@_keyans_store_ref:, \@@_keyans_store_ref_aux_i:,
%                    \@@_keyans_store_ref_aux_ii:}
%   The function \myhlc|\@@_keyans_store_ref:| handles the
%   \emph{\enquote{internal label and ref}} system used by the
%   \mykey[env]{save-ref} key for |\item*| and |\anspic*| commands.
%   First, copies of the current \mymeta{labels} are created and the
%   dots \enquote{\textcolor{red}{\texttt{.}}} are removed from them, since
%   double dots in references are not wanted.
%    \begin{macrocode}
\cs_new_protected:Nn \@@_keyans_store_ref:
  {
    \bool_if:NT \l_@@_store_ref_key_bool
      {
        \cs_set_protected:Nn \@@_tmp:n
          {
            \tl_set_eq:cc { l_@@_label_copy_ ##1 _tl } { l_@@_label_ ##1 _tl }
            \tl_reverse:c { l_@@_label_copy_ ##1 _tl }
            \tl_remove_once:cn { l_@@_label_copy_ ##1 _tl } { . }
            \tl_reverse:c { l_@@_label_copy_ ##1 _tl }
          }
        \clist_map_inline:nn { i, v, vi, vii, viii } { \@@_tmp:n {##1} }
        \@@_keyans_store_ref_aux_i:
      }
  }
%    \end{macrocode}
%   The auxiliary function \myhlc|\@@_keyans_store_ref_aux_i:| set the
%   variable \myhlc|\l_@@_newlabel_arg_one_tl| which will  contain
%   \mymarg{store name \textcolor{red}{:} position} analyzing whether
%   the environment in which they are executed is \myenv*[no-index]{enumext*} or
%   \myenv*[no-index]{enumext}.
%    \begin{macrocode}
\cs_new_protected:Nn \@@_keyans_store_ref_aux_i:
  {
    \bool_if:NT \g_@@_starred_bool
      {
        \tl_set_eq:NN \l_@@_label_copy_i_tl \l_@@_label_copy_vii_tl
      }
    \int_compare:nNnT { \l_@@_keyans_pic_level_int } = { 1 }
      {
        \tl_put_right:Ne \l_@@_newlabel_arg_two_tl
          { \l_@@_label_copy_i_tl . \l_@@_label_copy_vi_tl }
      }
    \int_compare:nNnT { \l_@@_keyans_level_int } = { 1 }
      {
        \tl_put_right:Ne \l_@@_newlabel_arg_two_tl
          { \l_@@_label_copy_i_tl . \l_@@_label_copy_v_tl }
      }
    \int_compare:nNnT { \l_@@_keyans_level_h_int } = { 1 }
      {
        \tl_put_right:Ne \l_@@_newlabel_arg_two_tl
          { \l_@@_label_copy_i_tl . \l_@@_label_copy_viii_tl }
      }
    \tl_put_right:Ne \l_@@_newlabel_arg_one_tl
      {
        \l_@@_store_name_tl \c_colon_str
        \int_eval:n { \prop_count:c { g_@@_ \l_@@_store_name_tl _prop } }
      }
    \@@_keyans_store_ref_aux_ii:
  }
%    \end{macrocode}
%   Now the auxiliary function \myhlc|\@@_keyans_store_ref_aux_ii:| saves the
%   result in the variable \myhlc|\l_@@_write_aux_file_tl| and finally
%   writes it in the |.aux| file.
%    \begin{macrocode}
\cs_new_protected:Nn \@@_keyans_store_ref_aux_ii:
  {
    \tl_put_right:Ne \l_@@_write_aux_file_tl
      {
        \@@_newlabel:nn
          { \exp_not:V \l_@@_newlabel_arg_one_tl }
          { \l_@@_newlabel_arg_two_tl }
      }
    \l_@@_write_aux_file_tl
  }
%    \end{macrocode}
% \end{macro}
%
% \subsubsection{Storing content in sequence}
%
% \begin{macro}[int]{\@@_keyans_addto_seq:n, \@@_keyans_addto_seq_link:}
%   The function \myhlc|\@@_keyans_addto_seq:n| will pass the
%   contents of the current \mymeta{label} \myhlc|\l_@@_label_v_tl| for the
%   \myenv*[no-index]{keyans} environment and the \myhlc|\l_@@_label_vi_tl| for the \myenv*[no-index]{keyanspic} environment when using
%   |\item*| and |\anspic*|, followed by the \mymeta{contents} of the
%   \emph{optional argument} of both commands to the \myhlc|\l_@@_store_current_label_tl|
%   variable to the sequence defined by the \mykey[env]{save-ans} key.
%    \begin{macrocode}
\cs_new_protected:Nn \@@_keyans_addto_seq:n
  {
    \tl_clear:N \l_@@_store_current_label_tl
    \int_compare:nNnTF { \l_@@_keyans_pic_level_int } = { 1 }
      {
        \tl_put_right:Ne \l_@@_store_current_label_tl
          { \item \l_@@_label_vi_tl }
      }
      {
        \tl_put_right:Ne \l_@@_store_current_label_tl { \item \l_@@_label_v_tl }
      }
    \tl_if_novalue:nF {#1}
      {
        \tl_if_empty:NF \l_@@_store_keyans_item_opt_sep_v_tl
          {
            \tl_put_right:NV \l_@@_store_current_label_tl
              \l_@@_store_keyans_item_opt_sep_v_tl
          }
        \tl_put_right:Nn \l_@@_store_current_label_tl {#1}
      }
    \@@_keyans_addto_seq_link:
  }
%    \end{macrocode}
%   Checks if the \mykey[env]{save-ref} key is active along with the
%   \mypkg[no-index]{hyperref} package load, if both conditions are met, it will
%   create the |\hyperlink| and then store using the
%   \myhlc|\@@_store_addto_seq:V| function.
%   Finally, copy the contents of the variable \myhlc|\l_@@_store_current_label_tl|
%   into the global variable \myhlc|\g_@@_check_ans_item_tl| to be used
%   by the function \myhlc|\@@_check_starred_cmd:n| and increment
%   the value of the integer variable \myhlc|\g_@@_item_anskey_int| handled
%   by the \mykey[env]{check-ans} key.
%    \begin{macrocode}
\cs_new_protected:Nn \@@_keyans_addto_seq_link:
  {
    \bool_lazy_and:nnT
      { \bool_if_p:N \l_@@_store_ref_key_bool }
      { \bool_if_p:N \l_@@_hyperref_bool }
      {
        \tl_put_right:Ne \l_@@_store_current_label_tl
          {
            \hfill \exp_not:N \hyperlink
              {
                \exp_not:V \l_@@_newlabel_arg_one_tl
              }
              { \exp_not:V \l_@@_mark_ref_sym_tl }
          }
      }
    \@@_store_addto_seq:V \l_@@_store_current_label_tl
    \bool_if:NT \l_@@_check_answers_bool
      {
        \int_gincr:N \g_@@_item_anskey_int
      }
  }
%    \end{macrocode}
% \end{macro}
%
% \subsubsection{The \texttt{show-ans} and \texttt{show-pos} keys for \env{keyans} and \env{keyanspic}}
%
% \begin{macro}[int]{\@@_keyans_save_item_opt:n, \@@_keyans_show_item_opt:,
%                    \@@_keyans_show_item_opt_viii:}
%   The function \myhlc|\@@_keyans_save_item_opt:n| will save the optional
%   argument of |\item*| and |\anspic*| in the variable
%   \myhlc|\l_@@_store_current_opt_arg_tl|.
%    \begin{macrocode}
\cs_new_protected:Nn \@@_keyans_save_item_opt:n
  {
    \tl_if_novalue:nF {#1}
      {
        \tl_set:Nn \l_@@_store_current_opt_arg_tl {#1}
      }
  }
%    \end{macrocode}
%   The function \myhlc|\@@_keyans_show_item_opt:|  will print the optional
%   arguments of |\item*| and |\anspic*| when the \mykey[env]{show-ans} or
%   \mykey[env]{show-pos} keys are set next to the key \mykey[env]{wrap-opt} in
%   \myenv*[no-index]{keyans} and \myenv*[no-index]{keyanspic} environments.
%    \begin{macrocode}
\cs_new_protected:Nn \@@_keyans_show_item_opt:
  {
    \tl_if_empty:NF \l_@@_store_current_opt_arg_tl
      {
        \bool_lazy_or:nnT
          { \bool_if_p:N \l_@@_show_answer_bool }
          { \bool_if_p:N \l_@@_show_position_bool }
          {
            \@@_keyans_wrapper_opt_v:n
              { \l_@@_store_current_opt_arg_tl } \c_space_tl
          }
      }
  }
%    \end{macrocode}
%   The function \myhlc|\@@_keyans_show_item_opt_viii:|  will print the optional
%   argument of |\item*| when the \mykey{show-ans} or \mykey{show-pos}
%   keys are set next to the key \mykey{wrap-opt} in \myenv*[no-index]{keyans*}
%   environment.
%    \begin{macrocode}
\cs_new_protected:Nn \@@_keyans_show_item_opt_viii:
  {
    \tl_if_empty:NF \l_@@_store_current_opt_arg_tl
      {
        \bool_lazy_or:nnT
          { \bool_if_p:N \l_@@_show_answer_bool }
          { \bool_if_p:N \l_@@_show_position_bool }
          {
            \@@_keyans_wrapper_opt_viii:n
              { \l_@@_store_current_opt_arg_tl } \c_space_tl
          }
      }
  }
%    \end{macrocode}
% \end{macro}
%
% \begin{macro}[int]{\@@_keyans_pos_mark_set:, \@@_keyans_show_ans:, \@@_keyans_show_pos:}
%   The function \myhlc|\@@_keyans_pos_mark_set:| adjusts the horizontal
%   spaces for the \mykey[env]{mark-sep*} key taking into account the value of
%   the \mykey[env]{align} key and the width of \mymeta{label}.
%    \begin{macrocode}
\cs_new_protected:Nn \@@_keyans_pos_mark_set:
  {
    \@@_label_width_by_box:Nn
      \l_@@_mark_sep_tmpa_dim { \l_@@_label_v_tl }
    \str_case:Vn \l_@@_align_label_pos_v_str
      {
        { l }
            {
              \dim_set:Nn \l_@@_mark_sep_tmpb_dim { \c_zero_dim }
            }
        { r }
            {
              \dim_set:Nn \l_@@_mark_sep_tmpb_dim
                { \l_@@_labelwidth_v_dim - \l_@@_mark_sep_tmpa_dim }
            }
        { c }
            {
              \dim_set:Nn \l_@@_mark_sep_tmpb_dim
                { 0.5\l_@@_labelwidth_v_dim - 0.5\l_@@_mark_sep_tmpa_dim }
            }
      }
%    \end{macrocode}
%   Sets the default values for the keys \mykey[env]{mark-ans*}, \mykey{mark-sep*} and \mykey[env]{mark-pos*}.
%    \begin{macrocode}
    \dim_compare:nNnT { \l_@@_mark_sym_sep_v_dim } = { \c_zero_dim }
      {
        \dim_set:Nn \l_@@_mark_sym_sep_v_dim { \l_@@_labelsep_v_dim }
      }
    \tl_set_eq:NN \l_@@_mark_answer_sym_tl \l_@@_mark_answer_sym_v_tl
    \dim_add:Nn \l_@@_mark_sym_sep_v_dim { \l_@@_mark_sep_tmpb_dim }
    \str_set_eq:NN \l_@@_mark_position_str \l_@@_mark_position_v_str
  }
%    \end{macrocode}
%   The function \myhlc|\@@_keyans_show_ans:| will print the \mymeta{symbol}
%   set by the \mykey{mark-ans*} key when the \mykey[env]{show-ans} key is active.
%    \begin{macrocode}
\cs_new_protected:Nn \@@_keyans_show_ans:
  {
    \bool_lazy_all:nT
      {
        { \bool_if_p:N \l_@@_show_answer_bool   }
        { \bool_if_p:N \l_@@_item_wrap_key_bool }
      }
      {
        \@@_keyans_pos_mark_set:
        \@@_print_keyans_box:NN
          \l_@@_labelwidth_v_dim \l_@@_mark_sym_sep_v_dim
      }
  }
%    \end{macrocode}
%   The function \myhlc|\@@_keyans_show_pos:| will print the \mymeta{position}
%   of the stored content in \emph{prop list}. Need add |1| to
%   \myhlc|\g_@@_#1_prop| for \myenv*[no-index]{keyans}
%   environment.
%    \begin{macrocode}
\cs_new_protected:Nn \@@_keyans_show_pos:
  {
    \int_compare:nNnTF { \l_@@_keyans_level_int } = { 1 }
      {
        \int_incr:N \l_@@_show_pos_tmp_int
      }
      {
        \int_zero:N \l_@@_show_pos_tmp_int
      }
    \bool_lazy_all:nT
      {
        { \bool_if_p:N \l_@@_show_position_bool }
        { \bool_if_p:N \l_@@_item_wrap_key_bool }
      }
      {
        \tl_set:Ne \l_@@_mark_answer_sym_v_tl
          {
            \group_begin:
              \exp_not:N \normalfont
              \exp_not:N \footnotesize [ \int_eval:n
                {
                  \prop_count:c { g_@@_ \l_@@_store_name_tl _prop }
                  + \l_@@_show_pos_tmp_int
                }
              ]
            \group_end:
          }
        \@@_keyans_pos_mark_set:
        \@@_print_keyans_box:NN
          \l_@@_labelwidth_v_dim \l_@@_mark_sym_sep_v_dim
      }
  }
%    \end{macrocode}
% \end{macro}
%
% \subsection{Redefining \cs{item} and \cs{makelabel} in \env{enumext}}\label{cmd:item-enumext}
%
% Redefining the |\item| command is not as simple as it seems. This
% command works in conjunction with the |\makelabel| command, so both
% of them have to be redefined; in addition, a couple of \emph{global}
% variables are needed to pass the values from one command to the other.
%
% When \emph{labeling} \texttt{PDF} is active |\makelabel| is redefined
% as |\hss #1| and the only way to get the \mykey{align} key to work
% correctly is to redefine |\makelabel| using |\makebox|. The best way to
% implement this is to use the conditional command \myhlc|\IfDocumentMetadataTF|
% to force this redefinition and the dedicated \mykey{mode-box} key to
% manually activate it by the user.
%
% The \ics*{item} and |\item|\myoarg{custom} commands work in the
% usual way on \myenv*[no-index]{enumext}, and |\item*|, |\item*|\myoarg{symbol} and
% |\item*|\myoarg{symbol}\myoarg{offset} are added.
%
% \begin{macro}[int]{\@@_default_item:n}
%   First it checks if the \emph{optional argument} is present; if it is NOT
%   present, it checks the state of the variable \myhlc|\l_@@_check_answers_bool|
%   set by the key \mykey[env]{no-store}, sets the boolean variable
%   \myhlc|\l_@@_wrap_label_X_bool| to |true| for the key \mykey[env]{wrap-label}
%   and executes \myhlc|\@@_item_std:w| and the key \mykey[env]{itemindent}; otherwise it checks the state
%   of the boolean variable \myhlc|\l_@@_wrap_label_opt_X_bool| set
%   by the key \mykey[env]{wrap-label*} and executes \myhlc|\@@_item_std:w| with
%   the \emph{optional argument} and the key \mykey{itemindent}.
%    \begin{macrocode}
\cs_new_protected:Nn \@@_default_item:n
  {
    \tl_if_novalue:nTF {#1}
      {
        \bool_if:NT \l_@@_check_answers_bool
          {
            \int_gincr:N \g_@@_item_number_int
            \bool_set_true:N \l_@@_item_number_bool
          }
        \bool_set_true:c { l_@@_wrap_label_ \@@_level: _bool }
        \@@_item_std:w \tl_use:c { l_@@_fake_item_indent_ \@@_level: _tl }
      }
      {
        \bool_set_eq:cc
          { l_@@_wrap_label_ \@@_level: _bool }
          { l_@@_wrap_label_opt_ \@@_level: _bool }
        \@@_item_std:w [#1] \tl_use:c { l_@@_fake_item_indent_ \@@_level: _tl }
      }
  }
%    \end{macrocode}
% \end{macro}
%
% \begin{macro}[int]{\@@_item_starred_exec:nn, \@@_item_starred_exec:}
%    The |\item*|, |\item*|\myoarg{symbol} and |\item*|\myoarg{symbol}\myoarg{offset}
%    work like the \emph{numbered} |\item|, but placing a \mymeta{symbol}
%    to the \emph{\enquote{left}} of the \mymeta{label} separated from
%    it by the value the second \emph{optional argument} \mymeta{offset}.
%
% \begin{arguments}
%   \item \myhlc|\l_@@_item_symbol_X_tl|
%   \item \myhlc|\l_@@_item_symbol_sep_X_dim|
% \end{arguments}
%
%   First a copy of \myhlc|\l_@@_item_symbol_X_tl| is made, which is
%   set by the key \mykey[env,anskey,anskey*]{item-sym*} or passed as \emph{\enquote{first}} \emph{optional argument} in
%   the global variable \myhlc|\g_@@_item_symbol_aux_tl|, followed by setting the
%   variable \myhlc|\l_@@_item_symbol_sep_X_dim| set by the key
%   \mykey[env,anskey,anskey*]{item-pos*} or by the \emph{\enquote{second}} \emph{optional argument},
%   then the state of the variable \myhlc|\l_@@_check_answers_bool|
%   set by the key \mykey[env]{no-store} is checked, the boolean variable
%   \myhlc|\l_@@_wrap_label_X_bool| is set to |true| for the key \mykey[env]{wrap-label}
%   and \myhlc|\@@_item_std:w| and the key \mykey[env]{itemindent} are executed.
%    \begin{macrocode}
\cs_new_protected:Nn \@@_item_starred_exec:nn
  {
    \tl_if_novalue:nTF {#1}
      {
        \tl_gset_eq:Nc
          \g_@@_item_symbol_aux_tl { l_@@_item_symbol_ \@@_level: _tl }
      }
      {
        \tl_gset:Nn \g_@@_item_symbol_aux_tl {#1}
      }
    \tl_if_novalue:nTF {#2}
      {
        \dim_set_eq:cc
          { l_@@_item_symbol_sep_ \@@_level: _dim }
          { l_@@_labelsep_ \@@_level: _dim }
      }
      {
        \dim_set:cn { l_@@_item_symbol_sep_ \@@_level: _dim } {#2}
      }
    \bool_if:NT \l_@@_check_answers_bool
      {
        \int_gincr:N \g_@@_item_number_int
        \bool_set_true:N \l_@@_item_number_bool
      }
    \bool_set_true:c { l_@@_wrap_label_ \@@_level: _bool }
    \@@_item_std:w \tl_use:c { l_@@_fake_item_indent_ \@@_level: _tl }
  }
%    \end{macrocode}
%    The function \myhlc|\@@_item_starred_exec:| will be responsible for
%    executing |\item*| for the \myenv*[no-index]{enumext} environment.
%    \begin{macrocode}
\cs_new_protected:Nn \@@_item_starred_exec:
  {
    \tl_if_empty:cF { l_@@_item_symbol_ \@@_level: _tl }
      {
        \mode_leave_vertical:
        \skip_horizontal:n { -\dim_use:c { l_@@_item_symbol_sep_
          \@@_level: _dim } }
        \hbox_overlap_left:n { \g_@@_item_symbol_aux_tl }
        \skip_horizontal:n { \dim_use:c { l_@@_item_symbol_sep_
          \@@_level: _dim } }
      }
  }
%    \end{macrocode}
% \end{macro}
%
% \begin{macro}[int]{\@@_redefine_item:}
%   The function \myhlc|\@@_redefine_item:| will redefine the |\item|
%   command in the \myenv*[no-index]{enumext} environment adding |\item*|.
%   This function is passed to \myhlc|\@@_list_arg_two_X:| used
%   in the definition of the \myenv*[no-index]{enumext} environment (\S\ref{env:enumext}).
%    \begin{macrocode}
\cs_new_protected:Nn \@@_redefine_item:
  {
    \RenewDocumentCommand \item { s o o }
      {
        \bool_if:nTF {##1}
          {
            \@@_item_starred_exec:nn {##2} {##3}
          }
          { \@@_default_item:n {##2} }
      }
  }
%    \end{macrocode}
% \end{macro}
%
% \begin{macro}[int]{\@@_make_label:, \@@_make_label_std:, \@@_make_label_box:}
%   The function \myhlc|\@@_make_label:| redefines |\makelabel| for the keys
%   \mykey[env]{mode-box}, \mykey[env]{align}, \mykey[env]{font}, \mykey[env]{wrap-label},
%   \mykey[env]{wrap-label*} and |\item*| for \myenv*[no-index]{enumext} environment.
%   This function is passed to \myhlc|\@@_list_arg_two_X:| used
%   in the definition of the \myenv*[no-index]{enumext} environment (\S\ref{env:enumext}).
%    \begin{macrocode}
\cs_new_protected:Nn \@@_make_label:
  {
    \IfDocumentMetadataTF
      {
        \@@_make_label_box:
      }
      {
        \bool_if:NTF \l_@@_mode_box_bool
          {
            \@@_make_label_box:
          }
          {
            \@@_make_label_std:
          }
      }
  }
%    \end{macrocode}
%   Standard definition when \myhlc|\DocumentMetadata| is not active.
%    \begin{macrocode}
\cs_new_protected:Nn \@@_make_label_std:
  {
    \RenewDocumentCommand \makelabel { m }
      {
        \tl_use:c { l_@@_label_fill_left_ \@@_level: _tl }
        \@@_item_starred_exec:
        \tl_use:c { l_@@_label_font_style_ \@@_level: _tl }
        \bool_if:cTF { l_@@_wrap_label_ \@@_level: _bool }
          {
            \use:c { @@_wrapper_label_ \@@_level: :n } {##1}
          }
          {##1}
        \tl_use:c { l_@@_label_fill_right_ \@@_level: _tl }
        \tl_gclear:N \g_@@_item_symbol_aux_tl
      }
  }
%    \end{macrocode}
%   Definition using |\makebox| when \myhlc|\DocumentMetadata| is active or
%   \mykey{mode-box} is active.
%
% \begin{important}
%  Here it is necessary to use |\strut\smash| to maintain text \emph{alignment}
%  in case the user wants to use |\labelbx| for example. In experiments
%  with \emph{mimicking} the \myenv[no-index]{description} environment it was the only way
%  out and it seems to have no adverse effects and may serve in the future
%  as a basis for a more generic \myenv[no-index]{list} environment package than
%  \mypkg*[no-index]{enumext}.
% \end{important}
%
%    \begin{macrocode}
\cs_new_protected:Nn \@@_make_label_box:
  {
    \RenewDocumentCommand \makelabel { m }
      {
        \strut\smash
          {
            \makebox
              [ \dim_use:c { l_@@_labelwidth_ \@@_level: _dim } ]
              [ \str_use:c { l_@@_align_label_pos_ \@@_level: _str } ]
              {
                \@@_item_starred_exec:
                \tl_use:c { l_@@_label_font_style_ \@@_level: _tl }
                \bool_if:cTF { l_@@_wrap_label_ \@@_level: _bool }
                  {
                    \use:c { @@_wrapper_label_ \@@_level: :n } {##1}
                  }
                  {##1}
                \tl_gclear:N \g_@@_item_symbol_aux_tl
              }
          } % close smash
      }
  }
%    \end{macrocode}
% \end{macro}
%
% \subsection{Setting \texttt{item-sym*} and \texttt{item-pos*} keys}\label{key:item-starred}
%
% In order to have a cleaner implementation of |\item*| for the
% \myenv*[no-index]{enumext} and \myenv*[no-index]{enumext*} environments it is best to
% define a couple of keys that allow controlling and setting by default the
% \mymeta{symbol} and its \mymeta{offset}.
%
% \begin{macro}[int]{item-sym*, item-pos*}
%   Define and set \mykey[env,anskey,anskey*]{item-sym*} and \mykey[env,anskey,anskey*]{item-pos*} keys for
%   \myenv*[no-index]{enumext} and \myenv*[no-index]{enumext*}.
%    \begin{macrocode}
\cs_set_protected:Nn \@@_tmp:nn
  {
    \keys_define:nn { enumext / #1 }
      {
        item-sym* .tl_set:c = { l_@@_item_symbol_ #2 _tl },
        item-sym* .value_required:n = true,
        item-sym* .initial:n = {\textborn},
        item-pos* .dim_set:c = { l_@@_item_symbol_sep_ #2 _dim },
        item-pos* .value_required:n = true,
      }
  }
\clist_map_inline:nn
  {
    {level-1}{i}, {level-2}{ii}, {level-3}{iii}, {level-4}{iv}, {enumext*}{vii}
  }
  { \@@_tmp:nn #1 }
%    \end{macrocode}
% \end{macro}
%
% \subsection{Handling unknown keys}\label{key:unknown-keys}
%
% At this point in the code no more \mymeta{keys} will be added. Since the
% definitions of environments and commands have already been quite
% \emph{restrictive}, the only thing left to do is to be equally strict
% with the \mymeta{keys} (consistency matters).
%
% \begin{important}*
%  Well, the paragraph above did not hold entirely: more
%  \mymeta{keys} had to be added than planned, and not everything turns out
%  as expected.
% \end{important}
%
% \subsubsection{Definition of error messages}
%
% \begin{macro}[int]{starred-unknown-key, starred-unknown-key-value}
% Message used by unknown \mymeta{keys} in \myenv*[no-index]{enumext*}.
% environment.
%    \begin{macrocode}
\msg_new:nnnn { enumext } { starred-unknown-key }
  {
    The~key~'#1'~is~unknown~by~environment~
    '\l_@@_envir_name_tl'~and~is~being~ignored~\msg_line_context:.
  }
  {
    The~environment~'\l_@@_envir_name_tl'~does~not~
    have~a~key~called~'#1'.~
    Check~that~you~have~spelled~the~key~name~correctly.
  }
\msg_new:nnnn { enumext } { starred-unknown-key-value }
  {
    The~key~'#1=#2'~is~unknown~by~environment~
    '\l_@@_envir_name_tl'~and~is~being~ignored~\msg_line_context:.
  }
  {
    The~environment~'\l_@@_envir_name_tl'~does~not~
    have~a~key~called~'#1'.~
    Check~that~you~have~spelled~the~key~name~correctly.
  }
%    \end{macrocode}
% \end{macro}
%
% \begin{macro}[int]{standard-unknown-key, standard-unknown-key-value}
% Message used by unknown \mymeta{keys} in \myenv*[no-index]{enumext}
% environment.
%    \begin{macrocode}
\msg_new:nnnn { enumext } { standard-unknown-key }
  {
    The~key~'#1'~is~unknown~by~environment~'\l_@@_envir_name_tl'~on~level~
    \int_use:N \l_@@_level_int \c_space_tl and~is~being~ignored~
    \msg_line_context:.
  }
  {
    The~environment~'\l_@@_envir_name_tl'~does~not~
    have~a~key~called~'#1'~on~level~\int_use:N \l_@@_level_int.~
    Check~that~you~have~spelled~the~key~name~correctly.
  }
\msg_new:nnnn { enumext } { standard-unknown-key-value }
  {
    The~key~'#1=#2'~is~unknown~by~environment~'\l_@@_envir_name_tl'~on~level~
    \int_use:N \l_@@_level_int \c_space_tl and~is~being~ignored~
    \msg_line_context:.
  }
  {
    The~environment~'\l_@@_envir_name_tl'~does~not~
    have~a~key~called~'#1'~on~level~\int_use:N \l_@@_level_int.~
    Check~that~you~have~spelled~the~key~name~correctly.
  }
%    \end{macrocode}
% \end{macro}
%
% \begin{macro}[int]{keyans-unknown-key, keyans-unknown-key-value}
% Message used by unknown \mymeta{keys} in \myenv*[no-index]{keyans},
% \myenv*[no-index]{keyans*} and \myenv*[no-index]{keyanspic} environments.
%    \begin{macrocode}
\msg_new:nnnn { enumext } { keyans-unknown-key }
  {
    The~key~'#1'~is~unknown~by~environment~
    '\l_@@_envir_name_tl'~and~is~being~ignored~\msg_line_context:.
  }
  {
    The~environment~'\l_@@_envir_name_tl'~does~not~
    have~a~key~called~'#1'.~
    Check~that~you~have~spelled~the~key~name~correctly.
  }
\msg_new:nnnn { enumext } { keyans-unknown-key-value }
  {
    The~key~'#1=#2'~is~unknown~by~environment~
    '\l_@@_envir_name_tl'~and~is~being~ignored~\msg_line_context:.
  }
  {
    The~environment~'\l_@@_envir_name_tl'~does~not~
    have~a~key~called~'#1'.~
    Check~that~you~have~spelled~the~key~name~correctly.
  }
%    \end{macrocode}
% \end{macro}
%
% \subsubsection{Handling unknown keys for \env{keyans}, \env{keyans*} and \env{keyanspic}}\label{key:unknown-keys-keyans}
%
% \begin{macro}[int]{unknown, \@@_keyans_unknown_keys:n, \@@_keyans_unknown_keys:nn}
%   Define and set |unknown| key for  \myenv*[no-index]{keyans}, \myenv*[no-index]{keyans*}
%   and \myenv*[no-index]{keyanspic} environments. Here it is necessary to set
%   \myhlc|\l_@@_envir_name_tl| in case an |unknown| key is passed using |\setenumext|.
%    \begin{macrocode}
\cs_set_protected:Nn \@@_tmp:n
  {
    \keys_define:nn { enumext / #1 }
      {
        unknown .code:n =
          {
            \tl_set:Nn \l_@@_envir_name_tl {#1}
            \@@_keyans_unknown_keys:n {##1}
          },
      }
  }
\clist_map_inline:nn { keyans, keyans*, keyanspic } { \@@_tmp:n {#1} }
%    \end{macrocode}
%  Internal functions for handling |unknown| key.
%    \begin{macrocode}
\cs_new_protected:Nn \@@_keyans_unknown_keys:n
  {
    \exp_args:NV \@@_keyans_unknown_keys:nn \l_keys_key_str {#1}
  }
\cs_new_protected:Nn \@@_keyans_unknown_keys:nn
  {
    \tl_if_blank:nTF {#2}
      {
        \msg_error:nne { enumext } { keyans-unknown-key } {#1}
      }
      {
        \msg_error:nnee { enumext } { keyans-unknown-key-value } {#1} {#2}
      }
  }
%    \end{macrocode}
% \end{macro}
%
% \subsubsection{Handling unknown keys for \env{enumext*}}\label{key:unknown-keys-starred}
%
% \begin{macro}[int]{unknown}
%   Define and set |unknown| key for  \myenv*[no-index]{enumext*} environment.
%    \begin{macrocode}
\keys_define:nn { enumext / enumext* }
  {
    unknown .code:n =
      {
        \tl_set:Ne \l_@@_envir_name_tl { enumext* }
        \@@_starred_unknown_keys:n {#1}
      },
  }
%    \end{macrocode}
% \end{macro}
%
% \begin{macro}[int]{\@@_starred_unknown_keys:n, \@@_starred_unknown_keys:nn}
%   Internal functions for handling |unknown| key.
%    \begin{macrocode}
\cs_new_protected:Nn \@@_starred_unknown_keys:n
  {
    \exp_args:NV \@@_starred_unknown_keys:nn \l_keys_key_str {#1}
  }
\cs_new_protected:Nn \@@_starred_unknown_keys:nn
  {
    \tl_if_blank:nTF {#2}
      {
        \msg_error:nne { enumext } { starred-unknown-key } {#1}
      }
      {
        \msg_error:nnee { enumext } { starred-unknown-key-value } {#1} {#2}
      }
  }
%    \end{macrocode}
% \end{macro}
%
% \subsubsection{Handling unknown keys for \env{enumext}}\label{key:unknown-keys-standard}
%
% \begin{macro}[int]{unknown, \@@_standard_unknown_keys:n, \@@_standard_unknown_keys:nn}
%   Defines and set the key |unknown| for \myenv*[no-index]{enumext} environment.
%    \begin{macrocode}
\cs_set_protected:Nn \@@_tmp:n
  {
    \keys_define:nn { enumext / level-#1 }
      {
        unknown .code:n =
          {
            \int_set:Nn \l_@@_level_int {#1}
            \tl_set:Ne \l_@@_envir_name_tl { enumext }
            \@@_standard_unknown_keys:n {##1}
          },
      }
  }
\clist_map_inline:nn {1, 2, 3, 4} { \@@_tmp:n {#1} }
%    \end{macrocode}
%   Internal functions for handling |unknown| key.
%    \begin{macrocode}
\cs_new_protected:Nn \@@_standard_unknown_keys:n
  {
    \exp_args:NV \@@_standard_unknown_keys:nn \l_keys_key_str {#1}
  }
\cs_new_protected:Nn \@@_standard_unknown_keys:nn
  {
    \tl_if_blank:nTF {#2}
      {
        \msg_error:nne { enumext } { standard-unknown-key } {#1}
      }
      {
        \msg_error:nnee { enumext } { standard-unknown-key-value } {#1} {#2}
      }
  }
%    \end{macrocode}
% \end{macro}
%
% \subsection{Redefining \cs{item} and \cs{makelabel} in \env{keyans}}\label{cmd:item-keyans}
%
% The |\item| and |\item|\myoarg{custom} commands work in the
% usual way in \myenv*[no-index]{keyans}, but the |\item*| and |\item*|\myoarg{content}
% commands \emph{store} the current \mymeta{label} next to the \mymeta{content}
% if it is present in the \emph{sequence} and \emph{prop list} defined by
% \mykey{save-ans} key.
%
% \begin{macro}[int]{\@@_keyans_default_item:n}
%   The function \myhlc|\@@_keyans_default_item:n| executes the original
%   behavior of the |\item| along with the keys \mykey[env]{wrap-label},
%   \mykey[env]{wrap-label*} and \mykey[env]{itemindent}.
%    \begin{macrocode}
\cs_new_protected:Nn \@@_keyans_default_item:n
  {
    \tl_if_novalue:nTF {#1}
      {
        \bool_set_true:N \l_@@_wrap_label_v_bool
        \@@_item_std:w \tl_use:N \l_@@_fake_item_indent_v_tl
      }
      {
        \bool_set_eq:NN \l_@@_wrap_label_v_bool \l_@@_wrap_label_opt_v_bool
        \@@_item_std:w [#1] \tl_use:N \l_@@_fake_item_indent_v_tl
      }
  }
%    \end{macrocode}
% \end{macro}
%
% \begin{macro}[int]{\@@_keyans_starred_item:n}
%   The function \myhlc|\@@_keyans_starred_item:n| will take as argument `|#1|'
%   the \emph{optional argument} \myoarg{content} passed to |\item*| and
%   save it via the \myhlc|\@@_keyans_save_item_opt:n| function, then activate
%   the \mykey[env]{wrap-label} key, execute |\item| using \myhlc|\@@_item_std:w|,
%   the \mykey[env]{itemindent} key and print the \emph{optional argument} using
%   the \myhlc|\@@_keyans_show_item_opt:| function handled by the
%   \mykey[env]{wrap-opt} key.
%    \begin{macrocode}
\cs_new_protected:Nn \@@_keyans_starred_item:n
  {
    \@@_keyans_save_item_opt:n {#1}
    \bool_set_true:N \l_@@_wrap_label_v_bool
    \@@_item_std:w \tl_use:N \l_@@_fake_item_indent_v_tl
    \@@_keyans_show_item_opt:
%    \end{macrocode}
%   Now \emph{store} the current \mymeta{label} first in the \emph{prop list}
%   (including the \emph{optional argument}), run the internal \emph{\enquote{label and ref}}
%   system if the \mykey[env]{save-ref} key is active, then \emph{store} in the \emph{sequence}
%   and finally increments \myhlc|\g_@@_check_starred_cmd_int| for internal
%   check system.
%    \begin{macrocode}
    \@@_keyans_addto_prop:n {#1}
    \@@_keyans_store_ref:
    \@@_keyans_addto_seq:n {#1}
    \int_gincr:N \g_@@_check_starred_cmd_int
  }
%    \end{macrocode}
% \end{macro}
%
% \begin{macro}{\item*, \@@_keyans_redefine_item:}
%   The function \myhlc|\@@_keyans_redefine_item:| is responsible for adding
%   the \emph{starred argument} and \emph{optional argument} by the
%   \myhlc|\@@_list_arg_two_v:| function in the definition of the
%   \myenv*[no-index]{keyans} environment. Here the variable
%   \myhlc|\l_@@_item_wrap_key_bool| used by the \mykey[env]{wrap-ans*} key is set
%   to true only when |\item*| is executed and, additionally, it is necessary to use
%   \myhlc|\peek_remove_spaces:n| to avoid an unwanted space when using
%   |\item*| together with the \mykey[env]{itemindent} key.
%   This function is passed to \myhlc|\@@_list_arg_two_v:| used
%   in the definition of the \myenv*[no-index]{keyans} environment (\S\ref{subsec:second-arg}).
%    \begin{macrocode}
\cs_new_protected:Nn \@@_keyans_redefine_item:
  {
    \RenewDocumentCommand \item { s o }
      {
        \bool_if:nTF {##1}
          {
            \bool_set_true:N \l_@@_item_wrap_key_bool % wrap-ans*
            \peek_remove_spaces:n
              {
                \@@_keyans_starred_item:n {##2}
              }
          }
          {
            \bool_set_false:N \l_@@_item_wrap_key_bool
            \@@_keyans_default_item:n {##2}
          }
      }
  }
%    \end{macrocode}
% \end{macro}
%
% \begin{macro}[int]{\@@_keyans_make_label:, \@@_keyans_wrapper_label:n,
%                    \@@_keyans_make_label_std:, \@@_keyans_make_label_box:,
%                    \@@_wrapper_label_v:n}
%   The function \myhlc|\@@_keyans_make_label:| redefines |\makelabel|
%   for the keys \mykey[env]{mode-box}, \mykey[env]{align}, \mykey[env]{font},
%   \mykey[env]{wrap-label}, \mykey[env]{wrap-label*}, \mykey[env]{wrap-ans*} and |\item*|
%   for \myenv*[no-index]{keyans} environment.
%   This function is passed to \myhlc|\@@_list_arg_two_v:| used
%   in the definition of the \myenv*[no-index]{keyans} environment (\S\ref{subsec:second-arg}).
%    \begin{macrocode}
\cs_new_protected:Nn \@@_keyans_make_label:
  {
    \IfDocumentMetadataTF
      {
        \@@_keyans_make_label_box:
      }
      {
        \bool_if:NTF \l_@@_mode_box_bool
          {
            \@@_keyans_make_label_box:
          }
          {
            \@@_keyans_make_label_std:
          }
      }
  }
%    \end{macrocode}
%   Conditionals are added to the \myhlc|\@@_keyans_wrapper_label:n| function
%   to handle the keys \mykey{wrap-ans*}, \mykey{wrap-label} and \mykey{wrap-label*}.
%    \begin{macrocode}
\cs_new_protected:Nn \@@_keyans_wrapper_label:n
  {
    \bool_lazy_all:nT
      {
        { \bool_if_p:N \l_@@_wrap_label_v_bool         }
        { \bool_if_p:N \l_@@_show_answer_bool          }
        { \bool_if_p:N \l_@@_item_wrap_key_bool        }
        { \cs_if_exist_p:N \@@_keyans_wrapper_item_v:n }
      }
      {
        \cs_set_eq:NN \@@_wrapper_label_v:n \@@_keyans_wrapper_item_v:n
      }
    \bool_if:NTF \l_@@_wrap_label_v_bool
      {
        \@@_wrapper_label_v:n {#1}
      }
      {#1}
  }
%    \end{macrocode}
%   Standard definition when \myhlc|\DocumentMetadata| is not active.
%    \begin{macrocode}
\cs_new_protected:Nn \@@_keyans_make_label_std:
  {
    \RenewDocumentCommand \makelabel { m }
      {
        \tl_use:N \l_@@_label_fill_left_v_tl
        \@@_keyans_show_ans:
        \@@_keyans_show_pos:
        \tl_use:N \l_@@_label_font_style_v_tl
        \@@_keyans_wrapper_label:n {##1}
        \tl_use:N \l_@@_label_fill_right_v_tl
      }
  }
%    \end{macrocode}
%   Definition using |\makebox| when \myhlc|\DocumentMetadata| is active or
%   \mykey{mode-box} is active.
%    \begin{macrocode}
\cs_new_protected:Nn \@@_keyans_make_label_box:
  {
    \RenewDocumentCommand \makelabel { m }
      {
        \strut\smash
          {
            \makebox[ \l_@@_labelwidth_v_dim ][ \l_@@_align_label_pos_v_str ]
              {
                \@@_keyans_show_ans:
                \@@_keyans_show_pos:
                \tl_use:N \l_@@_label_font_style_v_tl
                \@@_keyans_wrapper_label:n {##1}
              }
          }
      }
  }
%    \end{macrocode}
% \end{macro}
%
% \subsection{Second argument of the lists}\label{subsec:second-arg}
%
% At this point in the code most of the tools needed to create a
% \emph{custom} \myenv[no-index]{list} environment have already been programmed. Remember that the
% \myhlc|\@@_start_list:nn| function takes two arguments: the \enquote{first}
% one is ready, and the \enquote{second} one is defined for all levels of the
% \myenv*[no-index]{enumext} environment, the \myenv*[no-index]{keyans} environment and the
% \myenv*[no-index]{enumext*} and \myenv*[no-index]{keyans*} environments.
%
% Here the \myhlc|\@@_list_arg_two_X:| function is implemented,
% which is responsible for setting all the list parameters, the
% counter, the redefinition of |\item|, |\makelabel| along with the keys
% \mykey{ref}, \mykey{itemindent} and \mykey{show-length}.
%
% \subsubsection{Definition of error messages}
%
% \begin{macro}[int]{width-negative}
% Messages used in length calculation.
%    \begin{macrocode}
\msg_new:nnn { enumext } { width-negative }
  {
    Ignoring~negative~value~'#1=#2'~\msg_line_context:.\\
    The~key~'#1'~accepts~values~>=~0pt.
  }
%    \end{macrocode}
% \end{macro}
%
% \begin{macro}[int]{list-lengths}
% Messages used by \mykey[env]{show-length} key in \myenv*[no-index]{enumext}.
%    \begin{macrocode}
\msg_new:nnn { enumext } { list-lengths }
  {
    ****~Lengths~used~by~'enumext'~level~'#2'~\msg_line_context:~
    \c_space_tl ****\\
    \@@_show_length:nnn { dim  } { labelsep      } {#1}
    \@@_show_length:nnn { dim  } { labelwidth    } {#1}
    \@@_show_length:nnn { dim  } { itemindent    } {#1}
    \@@_show_length:nnn { dim  } { leftmargin    } {#1}
    \@@_show_length:nnn { dim  } { rightmargin   } {#1}
    \@@_show_length:nnn { dim  } { listparindent } {#1}
    \@@_show_length:nnn { skip } { topsep    } {#1}
    \@@_show_length:nnn { skip } { parsep    } {#1}
    \@@_show_length:nnn { skip } { partopsep } {#1}
    \@@_show_length:nnn { skip } { itemsep   } {#1}
    ****************************************************
  }
%    \end{macrocode}
% \end{macro}
%
% \begin{macro}[int]{list-lengths-not-nested}
%  Messages used by \mykey[env]{show-length} key in  \myenv*[no-index]{enumext*},
%  \myenv*[no-index]{keyans*} and \myenv*[no-index]{keyans}.
%    \begin{macrocode}
\msg_new:nnn { enumext } { list-lengths-not-nested }
  {
    ****~Lengths~used~by~'#2'~environment~\msg_line_context:~\c_space_tl ****\\
    \@@_show_length:nnn { dim  } { labelsep      } {#1}
    \@@_show_length:nnn { dim  } { labelwidth    } {#1}
    \@@_show_length:nnn { dim  } { itemindent    } {#1}
    \@@_show_length:nnn { dim  } { leftmargin    } {#1}
    \@@_show_length:nnn { dim  } { rightmargin   } {#1}
    \@@_show_length:nnn { dim  } { listparindent } {#1}
    \@@_show_length:nnn { skip } { topsep    } {#1}
    \@@_show_length:nnn { skip } { parsep    } {#1}
    \@@_show_length:nnn { skip } { partopsep } {#1}
    \@@_show_length:nnn { skip } { itemsep   } {#1}
    ****************************************************
  }
%    \end{macrocode}
% \end{macro}
%
% \subsubsection{Calculation of \cs{leftmargin} and \cs{itemindent}}\label{subsubsec:leftmargin}
%
% Consider the figure~\ref{fig:eight} where the default margins (on the
% left) of a list are represented.
%
% \begin{figure}[!htbp]
%   \centering
%   \begin{tikzpicture}[line cap=round,line join=round,x=0.8cm,y=0.8cm,every node/.style={font=\footnotesize}]
%     \draw[<->,color=gray] (-0.5,6.175) --node[below,text=lightgray,font=\footnotesize\ttfamily]{margin page} (2,6.175);
%     \draw[line width=0.5pt,dash pattern=on 2pt off 2pt,color=gray] (2,4.85) -- (2,7.5);
%     \draw[<->,color=gray] (2.0,7.0) --node[above,text=keyname]{\vphantom{p}\texttt{list-offset}} (4.5,7.0);
%     \draw[<->,color=gray] (4.5,7.0) --node[above,text=blue!75]{\vphantom{p}\texttt{\textbackslash{}labelwidth}} (7.5,7.0);
%     \draw[<->,color=gray] (7.5,7.0) --node[above,text=blue!75]{\texttt{\textbackslash{}labelsep}} (9.5,7.0);
%     \draw[color=lightgray,fill=lightgray,fill opacity=0.15] (4.5,6.5) rectangle (7.5,6.85);% labelwidth
%     \draw[color=lightgray,fill=lightgray,fill opacity=0.15] (7.5,6.5) rectangle (9.5,6.85);% labelsep
%     \draw[color=lightgray,fill=lightgray,fill opacity=0.15] (9.5,6.5) rectangle (15,6.85);% item width
%     \draw[color=lightgray,fill=lightgray,fill opacity=0.15] (9.5,6.0) rectangle (15,6.35);
%     \draw[color=lightgray,fill=lightgray,fill opacity=0.15] (9.5,5.5) rectangle (15,5.85);
%     \draw[<->,color=gray] (1.25,5.35) --node[below,text=blue!75]%
%       {\begin{tabular}{c}
%          \texttt{\textbackslash{}leftmargin + \textbackslash{}itemindent} \\
%          \textcolor{keyname}{\texttt{list-indent}}
%        \end{tabular}} (9.5,5.35);
%   \end{tikzpicture}
%   \caption{Representation of standard horizontal lengths in \myenv[no-index]{list} environment.}
%   \label{fig:eight}
% \end{figure}
%
% The idea is to have control over these margins so that the list does
% not overlap the left margin of the page. The \emph{key} relationship is that
% the \emph{\enquote{right edge}} of the |\labelsep| equals the \emph{\enquote{right edge}} of the |\itemindent|,
% so that the left edge of the \emph{\enquote{label box}} is at |\leftmargin+\itemindent| minus |\labelwidth+\labelsep|.
% Thus, the handling of the margins by the package will be as shown in
% the figure~\ref{fig:nine}.
%
% \begin{figure}[!htbp]
%   \centering
%   \begin{tikzpicture}[line cap=round,line join=round,x=0.8cm,y=0.8cm,every node/.style={font=\footnotesize}]
%     \draw[<->,color=gray] (-0.5,6.175) --node[below,text=lightgray,font=\footnotesize\ttfamily]{margin page} (2,6.175);
%     \draw[line width=0.5pt,dash pattern=on 2pt off 2pt,color=gray] (2,4.85) -- (2,7.5);
%     \draw[<->,color=gray] (2.0,7.0) --node[above,text=keyname]{\vphantom{p}\texttt{list-offset}} (4.5,7.0);
%     \draw[<->,color=gray] (4.5,7.0) --node[above,text=blue!75]{\vphantom{p}\texttt{\textbackslash{}labelwidth}} (7.5,7.0);
%     \draw[<->,color=gray] (7.5,7.0) --node[above,text=blue!75]{\texttt{\textbackslash{}labelsep}} (9.5,7.0);
%     \draw[color=lightgray,fill=lightgray,fill opacity=0.15] (4.5,6.5) rectangle (7.5,6.85);% labelwidth
%     \draw[color=lightgray,fill=lightgray,fill opacity=0.15] (7.5,6.5) rectangle (9.5,6.85);% labelsep
%     \draw[color=lightgray,fill=lightgray,fill opacity=0.15] (9.5,6.5) rectangle (15,6.85);% item width
%     \draw[color=lightgray,fill=lightgray,fill opacity=0.15] (9.5,6.0) rectangle (15,6.35);
%     \draw[color=lightgray,fill=lightgray,fill opacity=0.15] (9.5,5.5) rectangle (15,5.85);
%     \draw[<->,color=gray] (4.5,5.35) --node[below,text=blue!75]%
%       {\begin{tabular}{c}
%          \texttt{\textbackslash{}leftmargin + \textbackslash{}itemindent} \\
%          \textcolor{keyname}{\texttt{list-indent}}
%        \end{tabular}} (9.5,5.35);
%   \end{tikzpicture}
%   \caption{Representation of horizontal lengths concept in list in \mypkg*[no-index]{enumext}.}
%   \label{fig:nine}
% \end{figure}
%
% Where the default values will look like in the figure~\ref{fig:last}.
%
% \begin{figure}[htbp]
%   \centering
%   \begin{tikzpicture}[line cap=round,line join=round,x=0.8cm,y=0.8cm,every node/.style={font=\footnotesize}]
%     \draw[<->,color=gray] (-0.5,6.175) --node[below,text=lightgray,font=\footnotesize\ttfamily]{margin page} (2,6.175);
%     \draw [line width=0.5pt,dash pattern=on 2pt off 2pt,color=gray] (2,4.85) -- (2,7.5);
%     \draw[<->,color=gray] (2.0,7.0) --node[above,text=blue!75]{\vphantom{p}\texttt{labelwidth}} (5.0,7.0);
%     \draw[<->,color=gray] (5.0,7.0) --node[above,text=blue!75]{\texttt{labelsep}} (7.0,7.0);
%     \draw[color=lightgray,fill=lightgray,fill opacity=0.15] (2.0,6.5) rectangle (5.0,6.85);% labelwidth
%     \draw[color=lightgray,fill=lightgray,fill opacity=0.15] (5.0,6.5) rectangle (7.0,6.85);% labelsep
%     \draw[color=lightgray,fill=lightgray,fill opacity=0.15] (7.0,6.5) rectangle (15,6.85);% item width
%     \draw[color=lightgray,fill=lightgray,fill opacity=0.15] (7.0,6.0) rectangle (15,6.35);
%     \draw[color=lightgray,fill=lightgray,fill opacity=0.15] (7.0,5.5) rectangle (15,5.85);
%     \draw[<->,color=gray] (2.0,5.35) --node[below,text=blue!75]%
%       {\begin{tabular}{c}
%          \texttt{\textbackslash{}leftmargin + \textbackslash{}itemindent} \\
%          \textcolor{keyname}{\texttt{list-indent}}
%        \end{tabular}}
%     (7.0,5.35);
%   \end{tikzpicture}
%   \caption{Default horizontal lengths in \mypkg*[no-index]{enumext}.}
%   \label{fig:last}
% \end{figure}
%
% \begin{macro}[int]{\@@_calc_hspace:NNNNNNN, \@@_calc_hspace:ccccccc}
%  The function \myhlc|\@@_calc_hspace:NNNNNNN| takes seven arguments to be
%  able to determine horizontal spaces for all list environment:
%
%  \begin{verbatim}
%     #1: \l__enumext_labelwidth_X_dim       #2: \l__enumext_labelsep_X_dim
%     #3: \l__enumext_listoffset_X_dim       #4: \l__enumext_leftmargin_tmp_X_dim
%     #5: \l__enumext_leftmargin_X_dim       #6: \l__enumext_itemindent_X_dim
%     #7: \l__enumext_leftmargin_tmp_X_bool
%  \end{verbatim}
%
%  And returns the \emph{\enquote{adjusted}} values of |\leftmargin|
%  and |\itemindent|.
%
%
%    \begin{macrocode}
\cs_new_protected:Nn \@@_calc_hspace:NNNNNNN
  {
    \dim_compare:nNnT {#1} < { \c_zero_dim }
      {
        \msg_warning:nnnV { enumext } { width-negative }{ labelwidth } #1
        \dim_set:Nn #1 { \dim_abs:n {#1} }
      }
    \dim_compare:nNnT {#2} < { \c_zero_dim }
      {
        \msg_warning:nnnV { enumext } { width-negative }{ labelsep }  #2
        \dim_set:Nn #2 { \dim_abs:n {#2} }
      }
%    \end{macrocode}
%   If no value has been passed to the \mykey[env]{labelwidth} and \mykey[env]{labelsep} keys, the
%   default values are set for \myhlc|\l_@@_leftmargin_tmp_X_dim|.
%    \begin{macrocode}
    \bool_if:NF #7 { \dim_set:Nn #4 { #1 + #2} }
%    \end{macrocode}
%   The cases are now analyzed and the values for |\leftmargin| and
%   |\itemindent| are set.
%    \begin{macrocode}
    \dim_compare:nNnTF {#4} < { \c_zero_dim }
      {
        \dim_set:Nn #6 { #1 + #2 - #4}
        \dim_set:Nn #5 { #1 + #2 + #3 - #6 }
      }
      {
        \dim_compare:nNnT {#4} = { #1 + #2 }
          { \dim_set:Nn #6 { \c_zero_dim } }
        \dim_compare:nNnT {#4} < { #1 + #2 }
          { \dim_set:Nn #6 { #1 + #2 - #4} }
        \dim_compare:nNnT {#4} > { #1 + #2 }
          {
            \dim_set:Nn #6 { -#1 - #2 + #4}
            \dim_set:Nn #6 { #6*-1}
          }
        \dim_set:Nn #5 { #1 + #2 + #3 - #6 }
      }
  }
\cs_generate_variant:Nn \@@_calc_hspace:NNNNNNN { ccccccc }
%    \end{macrocode}
% \end{macro}
%
% \subsubsection{Setting second argument of the lists}
%
% \begin{macro}[int]{\@@_list_arg_two_i:, \@@_list_arg_two_ii:, \@@_list_arg_two_iii:,
%                    \@@_list_arg_two_iv:, \@@_list_arg_two_v:}
%   |\leftmargini|, |\leftmarginii|, |\leftmarginiii| and
%   |\leftmarginiv| are \enquote{not set}; in this case, the parameters
%   for vertical and horizontal list spacing per level are set directly.
%    \begin{macrocode}
\cs_set_protected:Nn \@@_tmp:n
  {
    \cs_new_protected:cpn { @@_list_arg_two_ #1 : }
      {
        \@@_calc_hspace:ccccccc
          { l_@@_labelwidth_ #1 _dim } { l_@@_labelsep_ #1 _dim }
          { l_@@_listoffset_ #1 _dim } { l_@@_leftmargin_tmp_ #1 _dim }
          { l_@@_leftmargin_ #1 _dim } { l_@@_itemindent_ #1 _dim }
          { l_@@_leftmargin_tmp_ #1 _bool }
        \clist_map_inline:nn
          { labelsep, labelwidth, itemindent, leftmargin, rightmargin,
            listparindent }
          { \dim_set_eq:cc {####1} { l_@@_####1 _ #1 _dim } }
        \clist_map_inline:nn { topsep, parsep, partopsep, itemsep }
          { \skip_set_eq:cc {####1} { l_@@_####1 _ #1 _skip } }
        \clist_map_inline:nn { beginparpenalty, itempenalty, endparpenalty }
          { \int_set_eq:cc {@####1} { l_@@_####1 _ #1 _int } }
        \usecounter { enumX #1 }
        \setcounter { enumX #1 }
          { \int_eval:n { \int_use:c { l_@@_start_ #1 _int } - 1 } }
        \str_if_eq:nnTF {#1} { v }
          {
            \@@_keyans_redefine_item:
            \@@_keyans_make_label:
            \@@_keyans_ref:
            \@@_keyans_fake_item_indent:
            \bool_if:cT { l_@@_show_length_ #1 _bool }
              {
                \msg_term:nnnn { enumext } { list-lengths-not-nested } { v }
                  { keyans }
              }
          }
          {
            \@@_redefine_item:
            \@@_make_label:
            \@@_standard_ref:
            \@@_fake_item_indent:
            \bool_if:cT { l_@@_show_length_ #1 _bool }
              {
                \msg_term:nnne { enumext } { list-lengths } {#1}
                  { \int_use:N \l_@@_level_int }
              }
          }
      }
  }
\clist_map_inline:nn { i, ii, iii, iv, v } { \@@_tmp:n {#1} }
%    \end{macrocode}
% \end{macro}
%
% \begin{macro}[int]{\@@_list_arg_two_vii:, \@@_list_arg_two_viii:}
%   For the horizontal environments \myenv*[no-index]{enumext*} and \myenv*[no-index]{keyans*}
%   the implementation is similar, but, the value of |\partopsep| is always
%   \mydim{0pt}. At this point the \mykey[env]{parsep} key is modified to
%   make it take the value of the \mykey[env]{itemsep} key and later, in the
%   environment definition, \mykey{parindent} is modified to make it set
%   the value of \mykey[env]{listparindent} and \mykey{parsep} to set the value
%   of |\parskip| locally.
%
% \begin{important}*
%  For compatibility with \emph{tagged} \texttt{PDF} and \hologo{LaTeX}
%  release 2026-06-01, |\setcounter| is moved to |\begin{list}|\mymarg{arg one}|{|\mymeta{arg two}|}\setcounter|
%  in the definition of the environments |enumext*| and |keyans*| (see~\S\ref{env:starred} and~\S\ref{env:keyans-starred}).
%  This resolves the nesting problem reported by \texttt{@teepeemm} in
%  \href{https://github.com/pablgonz/enumext/issues/5}{enumext* nested in enumext starts at 2 with DocumentMetadata}
%  this was caused by the \mypkg[no-index]{latex-lab-block}\cite{latexlab} update in this release.
% \end{important}
%
%    \begin{macrocode}
\cs_set_protected:Nn \@@_tmp:n
  {
    \cs_new_protected:cpn { @@_list_arg_two_ #1 : }
      {
        \bool_set_true:c { l_@@_leftmargin_tmp_ #1 _bool }
        \dim_zero:c { l_@@_leftmargin_tmp_ #1 _dim }
        \@@_calc_hspace:ccccccc
          { l_@@_labelwidth_ #1 _dim } { l_@@_labelsep_ #1 _dim }
          { l_@@_listoffset_ #1 _dim } { l_@@_leftmargin_tmp_ #1 _dim }
          { l_@@_leftmargin_ #1 _dim } { l_@@_itemindent_ #1 _dim }
          { l_@@_leftmargin_tmp_ #1 _bool }
        \clist_map_inline:nn
          { labelsep, labelwidth, itemindent, leftmargin, rightmargin,
            listparindent }
          { \dim_set_eq:cc {####1} { l_@@_####1 _ #1 _dim } }
        \clist_map_inline:nn { topsep, parsep, partopsep, itemsep }
          { \skip_set_eq:cc {####1} { l_@@_####1 _ #1 _skip } }
        \clist_map_inline:nn { beginparpenalty, itempenalty, endparpenalty }
          { \int_set_eq:cc {@####1} { l_@@_####1 _ #1 _int } }
        \skip_set_eq:Nc \parsep  { l_@@_itemsep_ #1 _skip }
        \skip_zero:N \partopsep
        \usecounter { enumX #1 }
        % \setcounter { enumX #1 } { \int_eval:n { \int_use:c { l_@@_start_ #1 _int } - 1 } }
        \@@_starred_ref:
        \str_if_eq:nnTF {#1} { vii }
          {
            \@@_fake_item_indent_vii:
            \bool_if:cT { l_@@_show_length_vii_bool }
              { \msg_term:nnnn { enumext } { list-lengths-not-nested } { vii }
                { enumext* } }
          }
          {
            \@@_fake_item_indent_viii:
            \bool_if:cT { l_@@_show_length_ #1 _bool }
              { \msg_term:nnnn { enumext } { list-lengths-not-nested } {#1}
                { keyans* } }
          }
      }
  }
\clist_map_inline:nn { vii, viii } { \@@_tmp:n {#1} }
%    \end{macrocode}
% \end{macro}
%
% \subsection{The environment \env{enumext}}\label{env:enumext}
%
% \subsubsection{Definition of error messages for \env{enumext}}
%
% \begin{macro}[int]{list-too-deep}
% Message for the nesting depth of the environment \myenv*[no-index]{enumext}.
%    \begin{macrocode}
\msg_new:nnnn { enumext } { list-too-deep }
  {
    Too~deep~nesting~for~'enumext'~\msg_line_context:.
  }
  {
    The~maximum~level~of~nesting~is~4.
  }
%    \end{macrocode}
% \end{macro}
%
% \subsubsection{Internal functions for \env{enumext}}
%
% \begin{macro}[int]{\@@_safe_exec:}\label{fun:safe-exec-enumext}
%   The \myhlc|\@@_safe_exec:| function first calls the function \myhlc|\@@_is_not_nested:|
%   which sets \myhlc|\g_@@_standard_bool| to |true| if it
%   is NOT nested within \myenv*[no-index]{enumext*}, then calls the function
%   \myhlc|\@@_internal_mini_page:| to create the environment \myhlc|\@@_mini_page|,
%   increments \myhlc|\l_@@_level_int| to restrict nesting of the
%   environment, sets \myhlc|\l_@@_standard_bool| to |true| and
%   finally calls the function \myhlc|\@@_is_on_first_level:| which sets
%   \myhlc|\l_@@_standard_first_bool| to |true| only if the
%   environment is NOT nested and is at the \emph{\enquote{first level}}.
%    \begin{macrocode}
\cs_new_protected:Nn \@@_safe_exec:
  {
    \@@_is_not_nested:
    \@@_internal_mini_page:
    \int_incr:N \l_@@_level_int
    \int_compare:nNnT { \l_@@_level_int } > { 4 }
      { \msg_fatal:nn { enumext } { list-too-deep } }
    \bool_set_true:N \l_@@_standard_bool
    \bool_set_false:N \l_@@_starred_bool
    \@@_is_on_first_level:
  }
%    \end{macrocode}
% \end{macro}
%
% \begin{macro}[int]{\@@_parse_keys:n}
%  The \myhlc|\@@_parse_store_keys:n| function first clears the
%  variable \myhlc|\l_@@_series_name_str| used by the key \mykey[env]{series} and then
%  checks if it is at the \emph{\enquote{first level}}; if so, it
%  processes the \mymeta{keys} and then executes the function
%  \myhlc|\@@_parse_series:n| used by the key \mykey{series} and calls the
%  function \myhlc|\@@_nested_base_line_fix:| used by the key
%  \mykey[env]{base-fix}; otherwise it passes the \mymeta{keys} to the inner
%  levels of the environment, then executes the function \myhlc|\@@_store_active_keys:n|
%  and reprocesses the \mymeta{keys} to pass them to the \emph{sequence}
%  if the key \mykey[env]{save-key} is not active.
%    \begin{macrocode}
\cs_new_protected:Nn \@@_parse_keys:n
  {
    \tl_if_novalue:nF {#1}
      {
        \str_clear:N \l_@@_series_name_str
        \int_compare:nNnTF { \l_@@_level_int } = { 1 }
          {
            \keys_set:nn { enumext / level-1 } {#1}
            \bool_if:NF \l_@@_print_keyans_cmd_bool
              {
                \@@_parse_series:n {#1}
              }
            \@@_nested_base_line_fix:
          }
          {
            \exp_args:Ne \keys_set:nn
              { enumext / level-\int_use:N \l_@@_level_int } {#1}
            \bool_if:NF \l_@@_print_keyans_cmd_bool
              {
                \@@_parse_series:n {#1}
              }
          }
        \@@_store_active_keys:n {#1}
      }
  }
%    \end{macrocode}
% \end{macro}
%
% \begin{macro}[int]{\@@_start_store_level:}
%   The \myhlc|\@@_start_store_level:| function activates the \emph{\enquote{storing structure}}
%   mechanism in the \emph{sequence} for the command
%   |\anskey| and the environment \myenv*[no-index]{anskey*}.
%    \begin{macrocode}
\cs_new_protected:Nn \@@_start_store_level:
  {
    \bool_lazy_all:nT
      {
        { \bool_if_p:N \l_@@_store_active_bool }
        { \bool_not_p:n { \l_@@_keyans_env_bool } }
        { \bool_if_p:N \g_@@_standard_bool }
      }
      {
        \int_compare:nNnT { \l_@@_level_int } > { 1 }
          {
            \bool_set_true:c { l_@@_store_upper_level_ \@@_level: _bool }
            \@@_store_level_open:
          }
      }
%    \end{macrocode}
%   If \myenv*[no-index]{enumext} are nested in \myenv*[no-index]{enumext*} add \myhlc|\@@_store_level_open:|
%   to preserve the \emph{\enquote{storing structure}}.
%    \begin{macrocode}
    \bool_lazy_all:nT
      {
        { \bool_if_p:N \l_@@_store_active_bool }
        { \bool_not_p:n { \l_@@_keyans_env_bool } }
        { \int_compare_p:nNn { \l_@@_level_h_int } = { 1 } }
      }
      {
        \int_compare:nNnT { \l_@@_level_int } > { 0 }
          {
            \bool_set_true:c { l_@@_store_upper_level_ \@@_level: _bool }
            \@@_store_level_open:
          }
      }
  }
%    \end{macrocode}
% \end{macro}
%
% \begin{macro}[int]{\@@_stop_store_level:}
%   The \myhlc|\@@_stop_store_level:| function stops the \emph{\enquote{storing structure}}
%   mechanism in the \emph{sequence} for the command
%   |\anskey| and the environment \myenv*[no-index]{anskey*}.
%    \begin{macrocode}
\cs_new_protected:Nn \@@_stop_store_level:
  {
    \bool_if:cT { l_@@_store_upper_level_ \@@_level: _bool }
      {
        \@@_store_level_close:
      }
  }
%    \end{macrocode}
% \end{macro}
%
% \begin{macro}[int]{\@@_multicols_start:}
%   The function \myhlc|\@@_multicols_start:| will start the
%   \myenv[no-index]{multicols} environment according to the value passed by the
%   \mykey[env]{columns} key, then set the default value for
%   |\columnsep| when \mykey[env]{columns-sep}=\mydim{0pt} and set the
%   value of |\multicolsep| equal to zero and leave |\columnseprule|
%   equal to zero for inner levels.
%    \begin{macrocode}
\cs_new_protected:Nn \@@_multicols_start:
  {
    \int_compare:nNnT
      { \int_use:c { l_@@_columns_ \@@_level: _int } } > { 1 }
      {
        \dim_compare:nNnT
          { \dim_use:c { l_@@_columns_sep_ \@@_level: _dim } } = { \c_zero_dim }
          {
            \dim_set:cn { l_@@_columns_sep_ \@@_level: _dim }
              {
                ( \dim_use:c { l_@@_labelwidth_ \@@_level: _dim }
                  + \dim_use:c { l_@@_labelsep_ \@@_level: _dim }
                ) / \int_use:c { l_@@_columns_ \@@_level: _int }
                - \dim_use:c { l_@@_listoffset_ \@@_level: _dim }
              }
          }
        \dim_set_eq:Nc \columnsep { l_@@_columns_sep_ \@@_level: _dim  }
        \int_compare:nNnT { \l_@@_level_int } > { 1 }
          {
            \dim_zero:N \columnseprule
          }
%    \end{macrocode}
%   The \emph{vertical spacing} settings for the
%   \myenv[no-index]{multicols} environment are calculated using the function \myhlc|\@@_multi_addvspace:|,
%   the \emph{\enquote{vertical adjust spacing}} is applied, and then the
%   \myenv[no-index]{multicols} environment is started.
%    \begin{macrocode}
        \bool_if:cF { l_@@_minipage_active_ \@@_level: _bool }
          {
            \skip_zero:N \multicolsep
            \@@_multi_addvspace:
          }
        \raggedcolumns
        \begin{multicols}{ \int_use:c { l_@@_columns_ \@@_level: _int } }
      }
  }
%    \end{macrocode}
% \end{macro}
%
% \begin{macro}[int]{\@@_multicols_stop:}
%   The function \myhlc|\@@_multicols_stop:| will stop the \myenv[no-index]{multicols}
%   environment and apply the \emph{\enquote{vertical adjust}} spacing.
%   For compatibility with \emph{tagged} \texttt{PDF}, the closing of the
%   \myenv[no-index]{list} environment is executed here along with \myhlc|\@@_stop_store_level:|.
%    \begin{macrocode}
\cs_new_protected:Nn \@@_multicols_stop:
  {
    \int_compare:nNnTF
      { \int_use:c { l_@@_columns_ \@@_level: _int } } > { 1 }
      {
        \@@_stop_list:
        \@@_stop_store_level:
        \end{multicols}
        \@@_unskip_unkern:
        \@@_unskip_unkern:
        \par\addvspace{ \skip_use:c { l_@@_multicols_below_ \@@_level: _skip } }
      }
      {
        \@@_stop_list:
        \@@_stop_store_level:
      }
  }
%    \end{macrocode}
% \end{macro}
%
% \begin{macro}[int]{\@@_before_list:}
%   The function \myhlc|\@@_before_list:| first calls the function
%   \myhlc|\@@_vspace_above:| used by the keys \mykey[env]{above} and
%   \mykey[env]{above*}, then calls the function \myhlc|\@@_before_args_exec:|
%   used by the key \mykey[env]{before*} and finally execute the function
%   \myhlc|\@@_check_ans_active:| for the check answer mechanism.
%    \begin{macrocode}
\cs_new_protected:Nn \@@_before_list:
  {
    \@@_vspace_above:
    \@@_before_args_exec:
    \@@_check_ans_active:
%    \end{macrocode}
%   When the \mykey[env]{mini-env} key is active it will set the value of the
%   \myhlc|\l_@@_minipage_right_X_dim| to be the \emph{width} of the \myhlc|\@@_mini_page|
%   environment on the \emph{\enquote{right side}}, using this value together with the value
%   of the \myhlc|\l_@@_minipage_hsep_X_dim| set by the \mykey[env]{mini-sep} key,
%   the value of \myhlc|\l_@@_minipage_left_X_dim| will be set, which will be
%   the \emph{width} of \myhlc|\@@_mini_page| environment on the
%   \emph{\enquote{left side}}, always having a current |\linewidth|
%   as \emph{maximum width} between them.
%    \begin{macrocode}
    \dim_compare:nNnT
      { \dim_use:c { l_@@_minipage_right_ \@@_level: _dim } } > { \c_zero_dim }
      {
        \dim_set:cn { l_@@_minipage_left_ \@@_level: _dim }
          {
            \linewidth
            - \dim_use:c { l_@@_minipage_right_ \@@_level: _dim }
            - \dim_use:c { l_@@_minipage_hsep_ \@@_level: _dim }
          }
%    \end{macrocode}
%    The boolean variable \myhlc|\l_@@_minipage_active_X_bool| will be
%    activated and the integer variable \myhlc|\g_@@_minipage_stat_int|
%    used by the |\miniright| command will be incremented, then
%    the function \myhlc|\@@_minipage_add_space:| is called and the
%    \myhlc|\@@_mini_page| environment on the \emph{\enquote{left
%    side}} will be initialized followed by the \emph{\enquote{vertical spacing}}
%    applied to preserve the \emph{\enquote{baseline}} between the \emph{left}
%    and \emph{right} side environments.
%    After these actions, the function \myhlc|\@@_multicols_start:| is
%    called to handle the \myenv[no-index]{multicols} environment.
%    \begin{macrocode}
        \bool_set_true:c { l_@@_minipage_active_ \@@_level: _bool }
        \int_gincr:N \g_@@_minipage_stat_int
        \@@_minipage_add_space:
        \noindent
        \@@_mini_page{ \dim_use:c { l_@@_minipage_left_ \@@_level: _dim } }
      }
    \@@_multicols_start:
  }
%    \end{macrocode}
% \end{macro}
%
% \begin{macro}[int]{\@@_second_part:}
%   The function \myhlc|\@@_second_part:| first checks the state of
%   the boolean variable \myhlc|\l_@@_minipage_active_X_bool|; if it is
%   |true|, a small test is executed to check if the use of
%   |\miniright| was omitted (the \myhlc|\@@_mini_page|
%   environment has not been closed), then it closes \myhlc|\@@_mini_page|
%   and adds the \emph{adjusted vertical space} \myhlc|\l_@@_minipage_after_skip|,
%   otherwise it closes the \myenv[no-index]{multicols} environment.
%    \begin{macrocode}
\cs_new_protected:Nn \@@_second_part:
  {
    \bool_if:cTF { l_@@_minipage_active_ \@@_level: _bool }
      {
        \int_compare:nNnT { \g_@@_minipage_stat_int } = { 1 }
          {
            \msg_warning:nn { enumext } { missing-miniright }
            \miniright
          }
        \int_gzero:N \g_@@_minipage_stat_int
        \@@_unskip_unkern: % remove topsep + [partopsep]
        \end@@_mini_page
      }
      {
        \@@_multicols_stop:
      }
%    \end{macrocode}
%   Now the functions \myhlc|\@@_after_stop_list:| used by
%   the key \mykey[env,foreachkeyans]{after}, \myhlc|\@@_check_ans_key_hook:| used by the key
%   \mykey[env]{check-ans}, and \myhlc|\@@_vspace_below:| used by the keys \mykey[env]{below}
%   and \mykey[env]{below*} are executed. Finally it sets \myhlc|\l_@@_standard_bool| to false and calls
%   the function \myhlc|\@@_resume_save_counter:| used by the \mykey[env]{series},
%   \mykey[env]{resume} and \mykey[env]{resume*} keys.
%    \begin{macrocode}
    \@@_after_stop_list:
    \@@_check_ans_key_hook:
    \@@_vspace_below:
    \bool_set_false:N \l_@@_standard_bool
    \bool_if:NF \l_@@_print_keyans_cmd_bool
      {
        \@@_standard_save_counter:
      }
  }
%    \end{macrocode}
% \end{macro}
%
% \begin{macro}[int]{\@@_set_item_width:}\label{fun:item-width-enumext}
%   The function \myhlc|\@@_set_item_width:| will set the value of |\itemwidth|
%   taking into account the value established by the \mykey[env]{list-offset}
%   key for each level of the environment.
%    \begin{macrocode}
\cs_new_protected:Nn \@@_set_item_width:
  {
    \dim_set:Nn \itemwidth { \linewidth }
    \dim_compare:nT
      {
        \dim_use:c { l_@@_listoffset_ \@@_level: _dim } != \c_zero_dim
      }
      {
        \dim_sub:Nn \itemwidth
          {
            \dim_use:c { l_@@_listoffset_ \@@_level: _dim }
          }
      }
  }
%    \end{macrocode}
% \end{macro}
%
% \begin{macro}{enumext}
%   Now create the \myenv*{enumext} environment based on \myenv[no-index]{list}
%   environment by levels.
%    \begin{macrocode}
\NewDocumentEnvironment{enumext}{ O{} }
  {
    \@@_safe_exec:
    \@@_parse_keys:n {#1}
    \@@_before_list:
    \@@_start_store_level:
    \@@_start_list:nn
      { \tl_use:c { l_@@_label_ \@@_level: _tl } }
      {
        \use:c { @@_list_arg_two_ \@@_level: : }
        \@@_before_keys_exec:
      }
    \@@_set_item_width:
    \@@_after_args_exec:
  }
  {
    \@@_second_part:
  }
%    \end{macrocode}
% \end{macro}
%
% Since the \mykey{check-ans} check must not be executed by levels
% but on the complete list, it is taken out of the \myenv*[no-index]{enumext}
% environment using the \emph{\enquote{hook}} function \myhlc|\@@_after_env:nn|.
%    \begin{macrocode}
\@@_after_env:nn {enumext}
  {
    \@@_execute_after_env:
  }
%    \end{macrocode}
%
% \subsection{The environment \env{keyans}}\label{env:keyans}
%
% The environment \myenv*[no-index]{keyans} also based on lists.
% The main differences with the \myenv*[no-index]{enumext} environment are the
% \emph{nesting} and the way the \emph{answers} (choice) will be stored and checked, this
% environment is intended exclusively for \emph{\enquote{multiple choice questions}}.
%
% \subsubsection{Definition of error messages for \env{keyans}}
%
% \begin{macro}[int]{keyans-nested}
% Message used by \myenv*[no-index]{keyans} environment.
%    \begin{macrocode}
\msg_new:nnn { enumext } { keyans-nested }
  {
    The~environment~'keyans'~can't~be~nested~\msg_line_context:.
  }
%    \end{macrocode}
% \end{macro}
%
% \subsubsection{Internal functions for \env{keyans}}
%
% \begin{macro}[int]{\@@_keyans_safe_exec:}
%   The \myenv*[no-index]{keyans} environment will only be available if the
%   \mykey[env]{save-ans} key is active and can only be used at the
%   \emph{\enquote{first level}} within the \myenv*[no-index]{enumext} environment.
%   The environment must not be nested, so a maximum is set at this point.
%   If the conditions are not met, an error message is returned.
%    \begin{macrocode}
\cs_new_protected:Nn \@@_keyans_safe_exec:
  {
    \bool_if:NF \l_@@_store_active_bool
      {
        \msg_error:nnnn { enumext } { wrong-place }{ keyans }{ save-ans }
      }
    \int_incr:N \l_@@_keyans_level_int
    \bool_set_true:N \l_@@_keyans_env_bool
    \@@_keyans_name_and_start:
    % Set false for interfering with enumext nested in keyans
    % (yes, its possible and crayze)
    \bool_set_false:N \l_@@_store_active_bool
    \int_compare:nNnT { \l_@@_keyans_level_int } > { 1 }
      {
        \msg_error:nn { enumext } { keyans-nested }
      }
    \int_compare:nNnT { \l_@@_level_int } > { 1 }
      {
        \msg_error:nn { enumext } { keyans-wrong-level }
      }
  }
%    \end{macrocode}
% \end{macro}
%
% \begin{macro}[int]{\@@_keyans_parse_keys:n}
%    Parse \myoarg{key \textnormal{\textcolor{gray}{=}} val} for
%    \myenv*[no-index]{keyans} environment.
%    \begin{macrocode}
\cs_new_protected:Nn \@@_keyans_parse_keys:n
  {
    \keys_set:nn { enumext / keyans } {#1}
  }
%    \end{macrocode}
% \end{macro}
%
% \begin{macro}[int]{\@@_before_list_v:, \@@_keyans_multicols_start:,
%                    \@@_keyans_multicols_stop:, \@@_second_part_v:}
%    Same implementation as the one used in the \myenv*[no-index]{enumext} environment.
%    \begin{macrocode}
\cs_new_protected:Nn \@@_before_list_v:
  {
    \@@_vspace_above_v:
    \@@_before_args_exec_v:
    \dim_compare:nNnT { \l_@@_minipage_right_v_dim } > { \c_zero_dim }
      {
        \dim_set:Nn \l_@@_minipage_left_v_dim
          {
            \linewidth - \l_@@_minipage_right_v_dim - \l_@@_minipage_hsep_v_dim
          }
        \bool_set_true:N \l_@@_minipage_active_v_bool
        \int_gincr:N \g_@@_minipage_stat_int
        \@@_keyans_minipage_add_space:
        \@@_mini_page{ \l_@@_minipage_left_v_dim }
      }
    \@@_keyans_multicols_start:
  }
\cs_new_protected:Nn \@@_keyans_multicols_start:
  {
    \int_compare:nNnT { \l_@@_columns_v_int } > { 1 }
      {
        \dim_compare:nNnT { \l_@@_columns_sep_v_dim } = { \c_zero_dim }
          {
            \dim_set:Nn \l_@@_columns_sep_v_dim
              {
                (
                  \l_@@_labelwidth_v_dim + \l_@@_labelsep_v_dim
                ) / \l_@@_columns_v_int
                - \l_@@_listoffset_v_dim
              }
          }
        \dim_set_eq:NN \columnsep \l_@@_columns_sep_v_dim
        \dim_zero:N \columnseprule % no rule here
        \bool_if:NF \l_@@_minipage_active_v_bool
          {
            \skip_zero:N \multicolsep
            \@@_keyans_multi_addvspace:
          }
        \raggedcolumns
        \begin{multicols}{ \l_@@_columns_v_int }
      }
  }
\cs_new_protected:Nn \@@_keyans_multicols_stop:
  {
    \int_compare:nNnTF { \l_@@_columns_v_int } > { 1 }
      {
        \@@_stop_list:
        \end{multicols}
        \@@_unskip_unkern:
        \@@_unskip_unkern:
        \par\addvspace{ \l_@@_multicols_below_v_skip }
      }
      {
        \@@_stop_list:
      }
  }
\cs_new_protected:Nn \@@_second_part_v:
  {
    \bool_if:NTF \l_@@_minipage_active_v_bool
      {
        \int_compare:nNnT { \g_@@_minipage_stat_int } = { 1 }
          {
            \msg_warning:nn { enumext } { missing-miniright }
            \miniright
          }
        \int_gzero:N \g_@@_minipage_stat_int
        \@@_unskip_unkern: % remove \topsep + [\partopsep]
        \end@@_mini_page
        \par\addvspace{ \l_@@_minipage_after_skip }
      }
      {
        \@@_keyans_multicols_stop:
      }
    \bool_set_false:N \l_@@_keyans_env_bool
    \@@_after_stop_list_v:
    \@@_vspace_below_v:
  }
%    \end{macrocode}
% \end{macro}
%
% \begin{macro}[int]{\@@_keyans_set_item_width:}\label{fun:item-width-keyans}
%   The function \myhlc|\@@_keyans_set_item_width:| will set the value of |\itemwidth|
%   taking into account the value established by the \mykey[env]{list-offset}
%   key.
%    \begin{macrocode}
\cs_new_protected:Nn \@@_keyans_set_item_width:
  {
    \dim_set:Nn \itemwidth { \linewidth }
    \dim_compare:nT
      {
        \l_@@_listoffset_v_dim != \c_zero_dim
      }
      {
        \dim_sub:Nn \itemwidth { \l_@@_listoffset_v_dim }
      }
  }
%    \end{macrocode}
% \end{macro}
%
% \begin{macro}{keyans}
%   Defines the environment \myenv*{keyans}, also based on lists.
%    \begin{macrocode}
\NewDocumentEnvironment{keyans}{ O{} }
  {
    \@@_keyans_safe_exec:
    \@@_keyans_parse_keys:n {#1}
    \@@_before_list_v:
    \@@_start_list:nn
      { \tl_use:N \l_@@_label_v_tl }
      {
        \@@_list_arg_two_v:
        \@@_before_keys_exec_v:
      }
    \@@_keyans_set_item_width:
    \@@_after_args_exec_v:
  }
  {
    \@@_check_starred_cmd:n { item }
    \@@_second_part_v:
  }
%    \end{macrocode}
% \end{macro}
%
% \subsection{Tagging PDF support for non-standard list environments}\label{subsec:tag-nonstandard-list}
%
% The \hologo{LaTeX} release 2022-06-01 brings automatic support for
% \emph{tagged} \texttt{PDF} in several aspects, including the standard
% \emph{list environments} and the \myenv[no-index]{list} environment. Unfortunately
% non-standard \emph{list environments} like \myenv*[no-index]{keyanspic} or the
% horizontal list environments \myenv*[no-index]{enumext*} and \myenv*[no-index]{keyans*} are
% not structured in a nice way, i.e. the expected result in the \texttt{PDF}
% file is the expected one, but the underlying structure is not correct. In
% simple terms, for \emph{tagged} \texttt{PDF} a \myenv[no-index]{list} environment
% is a \myenv[no-index]{list} environment, no matter what it looks like in the \texttt{PDF} file.
%
% To maintain a correct \myenv[no-index]{list} structure when \myhlc|\DocumentMetadata| is
% active, it is necessary to do some things manually using
% \mypkg{tagpdf}\cite{tagpdf} and \mypkg{ltsockets}\cite{ltsockets}.
% This implementation is an adaptation of the answer given by Pablo González
% thanks to Ulrike Fischer's comments in \href{https://tex.stackexchange.com/a/725902}{How can I modify my \texttt{\textbackslash item} redefinition to be compatible with \texttt{tagging-pdf}}.
%
% \subsubsection{Socket for tagging support in \env{enumext*} and \env{keyans*}}
%
% \begin{macro}[int]{start-list-tags, stop-start-tags, stop-list-tags,
%                    \@@_start_list_tag:n, \@@_stop_start_list_tag:,
%                    \@@_stop_list_tag:n}
%   First, the necessary \texttt{sockets} and their
%   behavior are defined for \myenv*[no-index]{enumext*} and \myenv*[no-index]{keyans*}.
%    \begin{macrocode}
\socket_new:nn {tagsupport/@@/starred}{ 1 }
\socket_new_plug:nnn {tagsupport/@@/starred} {start-list-tags}
  {
    \tag_resume:n {#1}
      \tag_struct_begin:n {tag=LI}
        \tag_struct_begin:n {tag=Lbl}
          \tag_mc_begin:n {tag=Lbl}
  }
\socket_new_plug:nnn {tagsupport/@@/starred} {stop-start-tags}
  {
          \tag_mc_end:
        \tag_struct_end:n {tag=Lbl}
        \tag_struct_begin:n {tag=LBody}
          % {tag=text-unit} now "semantic-para"
          % (https://github.com/latex3/latex2e/commit/83f054afc042e82a664085b6f2fd7c8d61705fb9)
          \tag_struct_begin:n {tag=semantic-para}
            % {tag=text} now "text-block"
            \tag_struct_begin:n {tag=text-block}
  }
\socket_new_plug:nnn {tagsupport/@@/starred} {stop-list-tags}
  {
            \tag_struct_end:n {tag=text-block}
          \tag_struct_end:n {tag=semantic-para} % {tag=text-unit}
        \tag_struct_end:n {tag=LBody}
      \tag_struct_end:n {tag=LI}
    \tag_suspend:n {#1}
  }
%    \end{macrocode}
%
% And now they are wrapped so that they are only active when \myhlc|\DocumentMetadata| is present.
%    \begin{macrocode}
\cs_new_protected_nopar:Nn \@@_start_list_tag:n
  {
    \IfDocumentMetadataT
      {
        \socket_assign_plug:nn {tagsupport/@@/starred} {start-list-tags}
        \socket_use:nn {tagsupport/@@/starred} {#1}
      }
  }
\cs_new_protected_nopar:Nn \@@_stop_start_list_tag:
  {
    \IfDocumentMetadataT
      {
        \socket_assign_plug:nn {tagsupport/@@/starred} {stop-start-tags}
        \socket_use:nn {tagsupport/@@/starred} { }
      }
  }
\cs_new_protected_nopar:Nn \@@_stop_list_tag:n
  {
    \IfDocumentMetadataT
      {
        \socket_assign_plug:nn {tagsupport/@@/starred} {stop-list-tags}
        \socket_use:nn {tagsupport/@@/starred} {#1}
      }
  }
%    \end{macrocode}
% \end{macro}
%
% \subsubsection{Socket for tagging support in \env{keyanspic}}
%
% \begin{macro}[int]{start-list-tags, stop-start-tags, stop-list-tags,
%                    \@@_anspic_start_list_tag:, \@@_anspic_stop_start_list_tag:,
%                    \@@_anspic_stop_list_tag:}
%   First, the necessary \texttt{sockets} and their
%   behavior are defined for the \myenv*[no-index]{keyanspic} environment.
%    \begin{macrocode}
\socket_new:nn {tagsupport/@@/keyanspic}{ 0 }
\socket_new_plug:nnn {tagsupport/@@/keyanspic} {start-list-tags}
  {
    \tag_resume:n {keyanspic}
      \tag_struct_begin:n {tag=LI}
        \tag_struct_begin:n {tag=Lbl}
          \tag_mc_begin:n {tag=Lbl}
  }
\socket_new_plug:nnn {tagsupport/@@/keyanspic} {stop-start-tags}
  {
          \tag_mc_end:
        \tag_struct_end:n {tag=Lbl}
        \tag_struct_begin:n {tag=LBody}
          \tag_struct_begin:n {tag=semantic-para}
            \tag_struct_begin:n {tag=text-block}
              \tag_mc_begin:n {tag=text-block}
  }
\socket_new_plug:nnn {tagsupport/@@/keyanspic} {stop-list-tags}
  {
              \tag_mc_end:
            \tag_struct_end:n {tag=text-block}
          \tag_struct_end:n {tag=semantic-para}
        \tag_struct_end:n {tag=LBody}
      \tag_struct_end:n {tag=LI}
    \tag_suspend:n {keyanspic}
  }
%    \end{macrocode}
%   And now they are wrapped so that they are only active when
%   \myhlc|\DocumentMetadata| is present.
%    \begin{macrocode}
\cs_new_protected_nopar:Nn \@@_anspic_start_list_tag:
  {
    \IfDocumentMetadataT
      {
        \socket_assign_plug:nn {tagsupport/@@/keyanspic} {start-list-tags}
        \socket_use:n {tagsupport/@@/keyanspic}
      }
  }
\cs_new_protected_nopar:Nn \@@_anspic_stop_start_list_tag:
  {
    \IfDocumentMetadataT
      {
        \socket_assign_plug:nn {tagsupport/@@/keyanspic} {stop-start-tags}
        \socket_use:n {tagsupport/@@/keyanspic}
      }
  }
\cs_new_protected_nopar:Nn \@@_anspic_stop_list_tag:
  {
    \IfDocumentMetadataT
      {
        \socket_assign_plug:nn {tagsupport/@@/keyanspic} {stop-list-tags}
        \socket_use:n {tagsupport/@@/keyanspic}
      }
  }
%    \end{macrocode}
% \end{macro}
%
% \subsection{The environment \env{keyanspic} and \cs{anspic}}\label{subsec:keyanspic-anspic}
%
% The \myenv*[no-index]{keyanspic} environment is a \myenv[no-index]{list} based environment that
% uses the same configuration for \emph{\enquote{spacing}} and \mymeta{label}
% as the \myenv*[no-index]{keyans} environment, but it does not use |\item|.
% The \mymeta{contents} are passed to the environment by means of the
% |\anspic| command as replacement for |\item| command and placed
% inside \myenv[no-index]{minipage} environments, with the \mymeta{label}
% centered \emph{\enquote{above}} or \emph{\enquote{below}}, adjusting
% \emph{widths} and \emph{position} according to the options passed to
% the environment.
%
% \begin{figure}[ht]
%   \centering
%   \begin{tikzpicture}[line cap=round,line join=round,x=0.8cm,y=0.8cm,every node/.style={font=\footnotesize}]
%     % space above
%     \draw[color=white] (7.0,6.10) rectangle (17.5,6.5) node[pos=.5,text=black]%
%       { \textcolor{keyname}{\texttt{layout-top}} };
%     % space between
%     \draw[color=white] (7.0,4.5) rectangle (17.5,4.9) node[pos=.5,text=black]%
%       { \textcolor{keyname}{\texttt{layout-sep}} };
%     % space below
%     \draw[color=white] (7.0,2.9) rectangle (17.5,3.3) node[pos=.5,text=black]%
%       { \textcolor{keyname}{\texttt{layout-top}} };
%     % labelwidth
%     \draw[<->,color=gray] (2.0,7.0) --node[above,text=keyname]{\vphantom{p}\texttt{labelwidth}} (5.0,7.0);
%     \draw[<->,color=gray] (5.0,7.0) --node[above,text=keyname]{\texttt{labelsep}} (7.0,7.0);
%     \draw[color=lightgray,fill=lightgray,fill opacity=0.15] (2.0,6.5) rectangle (5.0,6.85);% labelwidth
%     \draw[color=lightgray,fill=lightgray,fill opacity=0.15] (5.0,6.5) rectangle (7.0,6.85);% labelsep
%     \draw[color=lightgray,fill=lightgray,fill opacity=0.15] (7.0,6.5) rectangle (17.5,6.85);% item width
%     % minipage one
%     \draw[color=lightgray,fill=lightgray,fill opacity=0.15] (7.0,5.25) rectangle (10.5,6.10) node[pos=.5,text=black]%
%       {\begin{tabular}{c}
%          \texttt{minipage one} \\
%          \textcolor{blue}{\emph{drawing or tabular}}
%        \end{tabular}};
%     % label one
%     \draw[color=lightgray,fill=lightgray,fill opacity=0.15] (7.0,4.9) rectangle (10.5,5.25) node[pos=.5,text=black]%
%       { \textcolor{blue}{\texttt{label}} };
%     % minipage two
%     \draw[color=lightgray,fill=lightgray,fill opacity=0.15] (10.5,5.25) rectangle (14,6.10) node[pos=.5,text=black]%
%       {\begin{tabular}{c}
%          \texttt{minipage two} \\
%          \textcolor{blue}{\emph{drawing or tabular}}
%        \end{tabular}};
%     % label two
%     \draw[color=lightgray,fill=lightgray,fill opacity=0.15] (10.5,4.9) rectangle (14,5.25) node[pos=.5,text=black]%
%       { \textcolor{blue}{\texttt{label}} };
%     % minipage three
%     \draw[color=lightgray,fill=lightgray,fill opacity=0.15] (14,5.25) rectangle (17.5,6.10) node[pos=.5,text=black]%
%      {\begin{tabular}{c}
%          \texttt{minipage three} \\
%          \textcolor{blue}{\emph{drawing or tabular}}
%        \end{tabular}};
%     % label three
%     \draw[color=lightgray,fill=lightgray,fill opacity=0.15] (14,4.9) rectangle (17.5,5.25) node[pos=.5,text=black]%
%       { \textcolor{blue}{\texttt{label}} };
%     % minipage four
%     \draw[color=lightgray,fill=lightgray,fill opacity=0.15] (7.0,3.65) rectangle (12.25,4.5) node[pos=.5,text=black]%
%      {\begin{tabular}{c}
%          \texttt{minipage four} \\
%          \textcolor{blue}{\emph{drawing or tabular}}
%        \end{tabular}};
%     % label four
%     \draw[color=lightgray,fill=lightgray,fill opacity=0.15] (7.0,3.3) rectangle (12.25,3.65) node[pos=.5,text=black]%
%       { \textcolor{blue}{\texttt{label}} };
%     % minipage five
%     \draw[color=lightgray,fill=lightgray,fill opacity=0.15] (12.25,3.65) rectangle (17.5,4.5) node[pos=.5,text=black]%
%      {\begin{tabular}{c}
%          \texttt{minipage five} \\
%          \textcolor{blue}{\emph{drawing or tabular}}
%        \end{tabular}};
%     % label five
%     \draw[color=lightgray,fill=lightgray,fill opacity=0.15] (12.25,3.3) rectangle (17.5,3.65) node[pos=.5,text=black]%
%       { \textcolor{blue}{\texttt{label}} };
%     \draw[color=lightgray,fill=lightgray,fill opacity=0.15] (2.0,2.55) rectangle (5.0,2.9);% labelwidth
%     \draw[color=lightgray,fill=lightgray,fill opacity=0.15] (5.0,2.55) rectangle (7.0,2.9);% labelsep
%     \draw[color=lightgray,fill=lightgray,fill opacity=0.15] (7.0,2.55) rectangle (17.5,2.9);% item width
%   \end{tikzpicture}
%   \caption{Representation of the \myenv*{keyanspic} spacing in \mypkg*[no-index]{enumext}.}
%   \label{fig:keyanspic}
% \end{figure}
%
% In order for the \myenv*[no-index]{keyanspic} environment and the |\anspic|
% command to work correctly, some variables must be set and exported in
% the first part of the environment definition and pass them to |\anspic|
% which is executed in the second part of the environment. This
% implementation is adapted from the answer given by Enrico Gregorio
% (\texttt{@egreg}) in \href{https://tex.stackexchange.com/a/577705/7832}{How to process the body of an environment and divide it by a \textbackslash macro?}.
%
% \subsubsection{Definition of error messages for \env{keyanspic}}
%
% \begin{macro}[int]{keyanspic-nested, keyanspic-wrong-level,
%                    keyanspic-item-cmd}
% Messages used by \myenv*[no-index]{keyanspic} environment.
%    \begin{macrocode}
\msg_new:nnn { enumext } { keyanspic-nested }
  {
    The~environment~'keyanspic'~can't~be~nested~\msg_line_context:.
  }
\msg_new:nnnn { enumext } { keyanspic-wrong-level }
  {
    Wrong~level~position~for~'keyanspic'~\msg_line_context:.
  }
  {
    The~environment~'keyanspic'~can~only~be~in~the~first~level.
  }
\msg_new:nnn { enumext } { keyanspic-item-cmd }
  {
    Can't~use~'\token_to_str:N \item'~in~'keyanspic'~\msg_line_context:.
  }
%    \end{macrocode}
% \end{macro}
%
% \subsubsection{The environment \env{keyanspic}}\label{env:keyanspic}
%
% \begin{macro}[int]{label-pos, label-sep, layout-sty, layout-sep,
%                    layout-top, mark-ans, mark-pos, mark-sep, save-sep,
%                    wrap-opt, wrap-ans*, show-ans, show-pos, unknown}
%   First the keys are defined. The position of the
%   \mymeta{label} centered \enquote{above} or \enquote{below} is handled by \mykey[env]{label-pos},
%   and its vertical separation from the \emph{drawing or tabular} by the key
%   \mykey[env]{label-sep}.
%   The \emph{\enquote{layout style}} is handled by the key
%   \mykey[env]{layout-sty}, which takes two values separated by comma \mymarg{n°
%   upper, n° lower} and determines the number of \myenv[no-index]{minipage}
%   environments in which all arguments of |\anspic| are printed at the
%   \enquote{upper} and \enquote{lower} within the environments separated
%   by the value of the key \mykey[env]{layout-sep}. The vertical space
%   \enquote{top} and \enquote{bottom} of the environment is handled
%   by the key \mykey[env]{layout-top}.
%    \begin{macrocode}
\keys_define:nn { enumext / keyanspic }
  {
    label-pos  .choices:nn =
      { above , below }
      {
        \int_case:nn { \l_keys_choice_int }
          {
            {1}
              {
                \bool_set_true:N \l_@@_anspic_label_above_bool
                \str_set:Nn \l_@@_anspic_mini_pos_str { t }
              }
            {2}
              {
                \bool_set_false:N \l_@@_anspic_label_above_bool
                \str_set:Nn \l_@@_anspic_mini_pos_str { b }
              }
          }
      },
    label-pos  .initial:n = below,
    label-pos  .value_required:n = true,
    label-pos / unknown .code:n =
      \msg_error:nneee { enumext } { unknown-choice }
        { label-pos } { above,~ below } { \exp_not:n {#1} },
    label-sep  .skip_set:N = \l_@@_anspic_label_sep_skip,
    label-sep  .value_required:n = true,
    layout-sty .tl_set:N = \l_@@_anspic_layout_style_tl,
    layout-sty .value_required:n = true,
    layout-sep .meta:nn = { enumext / keyans } { parsep = {#1} },
    layout-sep .value_required:n = true,
    layout-top .meta:nn = { enumext / keyans } { topsep = {#1} },
    layout-top .value_required:n = true,
    mark-ans   .meta:nn = { enumext / keyans } { mark-ans = {#1} },
    mark-ans   .value_required:n = true,
    mark-pos   .meta:nn = { enumext / keyans } { mark-pos = {#1} },
    mark-pos   .value_required:n = true,
    mark-sep   .meta:nn = { enumext / keyans } { mark-sep = {#1} },
    mark-sep   .value_required:n = true,
    save-sep   .meta:nn = { enumext / keyans } { save-sep = {#1} },
    save-sep   .value_required:n = true,
    wrap-opt   .meta:nn = { enumext / keyans } { wrap-opt = {#1} },
    wrap-opt   .value_required:n = true,
    wrap-ans*  .meta:nn = { enumext / keyans } { wrap-ans* = {#1} },
    wrap-ans*  .value_required:n = true,
    show-ans   .meta:nn = { enumext / keyans } { show-ans = {#1} },
    show-ans   .value_required:n = true,
    show-pos   .meta:nn = { enumext / keyans } { show-pos = {#1} },
    show-pos   .value_required:n = true,
    unknown    .code:n =
      {
        \tl_set:Ne \l_@@_envir_name_tl { keyanspic }
        \@@_keyans_unknown_keys:n {#1}
      },
  }
%    \end{macrocode}
% \end{macro}
%
% \begin{macro}[int]{\@@_keyans_pic_safe_exec:, \@@_keyans_pic_parse_keys:n,
%                    \@@_keyans_pic_skip_abs:N, \@@_keyans_pic_arg_two:}
%   The function \myhlc|\@@_keyans_pic_safe_exec:| check the nested level
%   position inside the \myenv*[no-index]{enumext} environment.
%    \begin{macrocode}
\cs_new_protected:Nn \@@_keyans_pic_safe_exec:
  {
    \int_incr:N \l_@@_keyans_pic_level_int
    \int_compare:nNnT { \l_@@_keyans_pic_level_int } > { 1 }
      {
        \msg_error:nn { enumext } { keyanspic-nested }
      }
    \@@_keyans_name_and_start:
  }
%    \end{macrocode}
%   Parse \myoarg{key \textnormal{\textcolor{gray}{=}} val} for
%   \myenv*[no-index]{keyanspic} environment.
%    \begin{macrocode}
\cs_new_protected:Nn \@@_keyans_pic_parse_keys:n
  {
    \tl_if_novalue:nF {#1}
      {
        \keys_set:nn { enumext / keyanspic } {#1}
      }
  }
%    \end{macrocode}
%   The function \myhlc|\@@_keyans_pic_skip_abs:N| will return a positive
%   value |\parsep| from \myenv*[no-index]{keyans} environment.
%    \begin{macrocode}
\cs_new_protected:Nn \@@_keyans_pic_skip_abs:N
  {
    \dim_compare:nNnT {#1} < { \c_zero_dim }
      {
        \skip_set:Nn #1 { -#1 }
      }
  }
%    \end{macrocode}
%   The \myhlc|\@@_keyans_pic_arg_two:| function is used in the
%   \emph{second argument} of the \myenv[no-index]{list} environment that defines
%   the \myenv*[no-index]{keyanspic} environment, taking the
%   configuration of the \emph{\enquote{spaces}} and the keys
%   \mykey[env]{label}, \mykey[env]{wrap-label}, \mykey[env]{parsep} and \mykey[env]{topsep}
%   from the \myenv*[no-index]{keyans} environment.
%   The first step is to set the boolean variable
%   \myhlc|\l_@@_leftmargin_tmp_v_bool| handled by the \mykey[env]{list-indent} key
%   to |false|, then copy the definition of the second list argument
%   from the \myenv*[no-index]{keyans} environment definition and make sure that
%   |\parsep| does not have a negative value.
%    \begin{macrocode}
\cs_new_protected:Nn \@@_keyans_pic_arg_two:
  {
    \bool_set_false:N \l_@@_leftmargin_tmp_v_bool
    \@@_list_arg_two_v:
    \@@_keyans_pic_skip_abs:N \parsep
%    \end{macrocode}
%   Now the counter |enumXv| of the \myenv*[no-index]{keyans} environment is incremented
%   and the \emph{total height} of the \mymeta{label} is saved in
%   \myhlc|\l_@@_anspic_label_htdp_dim| used by |\anspic|, and the
%   values of |\parsep| are adjusted only if the key \mykey[env]{label-pos} is set to \emph{below}.
%    \begin{macrocode}
    \bool_if:NF \l_@@_anspic_label_above_bool
      {
        \stepcounter { enumXv }
        \hbox_set:Nn \l_@@_anspic_label_box { \l_@@_label_v_tl }
        \dim_set:Nn \l_@@_anspic_label_htdp_dim
          {
            \box_ht_plus_dp:N \l_@@_anspic_label_box
          }
        \skip_add:Nn \parsep
          {
            \l_@@_anspic_label_htdp_dim
            + \box_dp:N \strutbox
            + \l_@@_anspic_label_sep_skip
          }
      }
%    \end{macrocode}
%   Finally the value of |\leftmargin| and |\topsep| is \emph{adjusted},
%   then |\listparindent|, |\partopsep| and |\itemsep| are set to zero so that
%   the \emph{horizontal} and \emph{vertical} space is not affected.
%    \begin{macrocode}
    \dim_add:Nn \leftmargin { -\l_@@_labelwidth_v_dim - \l_@@_labelsep_v_dim }
    \ignorespaces
    \skip_add:Nn \topsep { 0.5\box_dp:N \strutbox }
    \dim_zero:N  \listparindent
    \skip_zero:N \partopsep
    \skip_zero:N \itemsep
  }
%    \end{macrocode}
% \end{macro}
%
% \begin{macro}{keyanspic}
%   Defines the environment \myenv*{keyanspic}. For compatibility
%   with \emph{tagged} \texttt{PDF} the |\begin{|\myenv[no-index]{list}|}|
%   form must be used, along with a lot of conditional code using \myhlc|\IfDocumentMetadataTF|.
%   First, the code for automatic \emph{tagged} \texttt{PDF}
%   for \myenv[no-index]{list} environments is stopped, |\item| is redefined so that it cannot be
%   used, and the code for automatic \emph{tagged} \texttt{PDF} is stopped for
%   the \myenv*[no-index]{keyanspic} environment.
%    \begin{macrocode}
\NewDocumentEnvironment{keyanspic}{ o }
  {
    \@@_keyans_pic_safe_exec:
    \@@_keyans_pic_parse_keys:n {#1}
    \begin{list} { } { \@@_keyans_pic_arg_two: }
    \IfDocumentMetadataT
      {
        \tag_suspend:n {list}
      }
    \item[] \scan_stop:
    \RenewDocumentCommand \item {}
      {
        \msg_error:nn { enumext } { keyanspic-item-cmd }
      }
    \IfDocumentMetadataT
      {
        \tag_resume:n {keyanspic}
        \tag_tool:n {para/tagging=false} % no socket for this now
        \tag_suspend:n {keyanspic}
      }
  }
%    \end{macrocode}
%   In the second part of the environment definition, the code for
%   \emph{tagged} \texttt{PDF} is placed manually and the command is executed:
%   |\anspic| is run using the \myhlc|\@@_anspic_exec:| function.
%    \begin{macrocode}
  {
    \IfDocumentMetadataT
      {
        \tag_resume:n {keyanspic}
      }
    \@@_anspic_exec:
    \IfDocumentMetadataT
      {
        \tag_suspend:n {keyanspic}
      }
    \end{list}
    \IfDocumentMetadataT
      {
        \tag_struct_end:n {tag=L}
      }
%    \end{macrocode}
%   Finally it checks if |\anspic*| has been used, sets the counter |enumXvi|
%   to zero and applies the \enquote{adjusted} vertical space at the bottom.
%    \begin{macrocode}
    \@@_check_starred_cmd:n { anspic }
    \setcounter { enumXvi } { 0 }
    \bool_if:NTF \l_@@_anspic_label_above_bool
      {
        \par\addvspace{ 0.5\box_dp:N \strutbox }
      }
      {
        \par
        \addvspace
          {
            \dim_eval:n
              {
                \l_@@_anspic_label_htdp_dim + \box_ht_plus_dp:N \strutbox
                + \l_@@_anspic_label_sep_skip + \l_@@_topsep_v_skip
              }
          }
      }
  }
%    \end{macrocode}
% \end{macro}
%
% \subsubsection{The command \cs{anspic}}\label{cmd:anspic}
% The |\anspic| command takes three arguments, the \emph{starred versions} |\anspic*|\myoarg{content}
% \emph{store} the current \mymeta{label} next to the \emph{optional argument}
% \myoarg{content} in the \emph{sequence} and \emph{prop list}
% defined by \mykey{save-ans} key. The third \emph{mandatory argument}
% \mymarg{drawing or tabular} is NOT \emph{stored} in the \emph{sequence} or
% \emph{prop list}.
%
% \begin{important}
% One of the complications here to make the \myenv*[no-index]{keyanspic}
% environment compatible with \emph{tagged} PDF is the position of
% \mymeta{label}, the |\anspic| command processes the arguments in
% order, where `|#1|' and `|#2|' correspond to \mymeta{label} and `|#3|' to the
% mandatory argument and puts all this inside a \myenv[no-index]{minipage}
% environment. If `|#1|' and `|#2|', that is \mymeta{label}, is above `|#3|'
% there are no problems with \emph{tagged} PDF, but if `|#3|' comes first
% the list created with \emph{tagged} PDF will not be correct.
% \end{important}
%
% \begin{macro}[int]{\@@_anspic_body_dim:n, \@@_anspic_label:nn,
%                    \@@_anspic_label_pos:nnn, \@@_anspic_args:nnn,
%                    \@@_anspic_print:n, \@@_anspic_print:e,
%                    \@@_anspic_print:V, \@@_anspic_row:n, \@@_anspic_exec:}
%   The \myhlc|\@@_anspic_body_dim:n| function will set the value of
%   \myhlc|\l_@@_anspic_body_htdp_dim| equal to the \enquote{height plus depth}
%   of the \emph{mandatory argument} if the key \mykey[env]{label-pos} is set
%   \enquote{below}.
%    \begin{macrocode}
\cs_new_protected:Nn \@@_anspic_body_dim:n
  {
    \bool_if:NF \l_@@_anspic_label_above_bool
      {
        \IfDocumentMetadataT
          {
            \tag_suspend:n {keyanspic}
          }
        \vbox_set:Nn \l_@@_anspic_body_box {#1}
        \dim_set:Nn \l_@@_anspic_body_htdp_dim
          {
            \box_ht_plus_dp:N \l_@@_anspic_body_box
          }
        \IfDocumentMetadataT
          {
            \tag_resume:n {keyanspic}
          }
      }
  }
%    \end{macrocode}
%   The \myhlc|\@@_anspic_label:nn| function will process inside |\makebox| the
%   \emph{starred argument} `|*|' and \emph{optional argument} passed
%   to the command. Here the \mymeta{label} and
%   \emph{optional argument} are stored in \emph{prop list} and \emph{sequence}
%   and the \mykey[env]{show-ans}, \mykey[env]{show-pos}, \mykey[env]{font},
%   \mykey[env]{wrap-label}, \mykey[env]{wrap-ans*} and \mykey[env]{wrap-opt} keys are executed.
%    \begin{macrocode}
\cs_new_protected:Nn \@@_anspic_label:nn
  {
    \makebox[ \l_@@_anspic_mini_width_dim ][ c ]
      {
        \bool_if:nTF {#1}
          {
            \bool_set_true:N \l_@@_item_wrap_key_bool
            \bool_set_true:N \l_@@_wrap_label_v_bool
            \@@_keyans_save_item_opt:n {#2}
            \@@_keyans_addto_prop:n {#2}
            \@@_keyans_store_ref:
            \@@_keyans_addto_seq:n {#2}
            \int_gincr:N \g_@@_check_starred_cmd_int
            \@@_keyans_show_ans:
            \@@_keyans_show_pos:
            \makebox[ \l_@@_labelwidth_v_dim ][c]
              {
                \tl_use:N \l_@@_label_font_style_v_tl
                \@@_keyans_wrapper_label:n { \l_@@_label_vi_tl }
              }
            \skip_horizontal:n { \l_@@_labelsep_v_dim }
            \@@_keyans_show_item_opt:
          }
          {
            \bool_set_false:N \l_@@_item_wrap_key_bool
            \tl_use:N \l_@@_label_font_style_v_tl
            \@@_wrapper_label_v:n { \l_@@_label_vi_tl }
          }
      }
  }
%    \end{macrocode}
%   The function \myhlc|\@@_anspic_label_pos:nnn| will be in charge of
%   handling the \emph{\enquote{counter}} and the position of the \mymeta{label},
%   set by \mykey{label-pos} key which will have the same configuration
%   as the \myenv*[no-index]{keyans} environment.
%    \begin{macrocode}
\cs_new_protected:Nn \@@_anspic_label_pos:nnn
  {
    \stepcounter { enumXvi }
    \@@_anspic_body_dim:n {#3}
    \bool_if:NTF \l_@@_anspic_label_above_bool
      {
        \@@_anspic_label:nn {#1} {#2}
      }
      {
        \raisebox
          {
            -\dim_eval:n
              {
                \l_@@_anspic_label_htdp_dim
                + \l_@@_anspic_body_htdp_dim
                + \box_dp:N \strutbox
                + \l_@@_anspic_label_sep_skip
              }
          }
          [ 0pt ] [ 0pt ]
          {
            \@@_anspic_label:nn {#1} {#2}
          }
      }
  }
%
%    \end{macrocode}
%   The \myhlc|\@@_anspic_args:nnn| function will be responsible for placing
%   the code compatible with \emph{tagged} \texttt{PDF} and the arguments
%   within the \myhlc|\l_@@_anspic_args_seq| sequence which will be processed
%   by the \myhlc|\@@_anspic_print:n| function in the second part of the
%   definition of the \myenv*[no-index]{keyanspic} environment.
%    \begin{macrocode}
\cs_new_protected:Nn \@@_anspic_args:nnn
  {
    \@@_anspic_start_list_tag:
    \@@_anspic_label_pos:nnn {#1} {#2} {#3}
    \@@_anspic_stop_start_list_tag:
    \bool_if:NTF \l_@@_anspic_label_above_bool
      {
        \\[\l_@@_anspic_label_sep_skip] #3
      }
      {
        \\ #3
      }
    \@@_anspic_stop_list_tag:
  }
%    \end{macrocode}
%   The value \mymarg{n° upper, n° lower} passed to the \mykey[env]{layout-sty}
%   key is split by comma and is handled directly by the function
%   \myhlc|\@@_anspic_print:n| and passed to the function \myhlc|\@@_anspic_row:n|.
%    \begin{macrocode}
\cs_new_protected:Nn \@@_anspic_print:n
  {
    \clist_map_function:nN {#1} \@@_anspic_row:n
  }
\cs_generate_variant:Nn \@@_anspic_print:n { e, V }
%    \end{macrocode}
%   The function \myhlc|\@@_anspic_row:n| will set the \emph{widths} for the
%   \myenv[no-index]{minipage} environments and place \emph{all arguments} passed
%   to |\anspic| \emph{saved} in the \myhlc|\l_@@_anspic_args_seq| sequence
%   inside them.
%    \begin{macrocode}
\cs_new_protected:Nn \@@_anspic_row:n
  {
    \dim_set:Nn \l_@@_anspic_mini_width_dim { \linewidth / #1 }
    \int_set:Nn \l_@@_anspic_above_int { \l_@@_anspic_below_int }
    \int_set:Nn \l_@@_anspic_below_int { \l_@@_anspic_above_int + #1 }
    \int_step_inline:nnn
      { \l_@@_anspic_above_int + 1 }
      { \l_@@_anspic_below_int }
      {
        \IfDocumentMetadataT
          {
            \tag_suspend:n {minipage}
          }
        \begin{minipage}[ \l_@@_anspic_mini_pos_str ]{
          \l_@@_anspic_mini_width_dim }
          \centering
          \seq_item:Nn \l_@@_anspic_args_seq {##1}
        \end{minipage}
        \IfDocumentMetadataT
          {
            \tag_resume:n {minipage}
          }
      }
    \par
  }
%    \end{macrocode}
%   The \myhlc|\@@_anspic_exec:| function will execute all the code in the
%   |\anspic| command in the second argument of the |keyanspic| environment
%   definition. If the key \mykey{layout-sty} is not set, everything will
%   be printed on a \emph{single line}.
%    \begin{macrocode}
\cs_new_protected:Nn \@@_anspic_exec:
  {
    \tl_if_empty:NTF \l_@@_anspic_layout_style_tl
      {
        \@@_anspic_print:e { \seq_count:N \l_@@_anspic_args_seq }
      }
      {
        \@@_anspic_print:V \l_@@_anspic_layout_style_tl
      }
  }
%    \end{macrocode}
% \end{macro}
%
% \begin{macro}{\anspic}
%   The command is active in the \myenv*[no-index]{keyanspic}
%   environment only if the \mykey[env]{save-ans} key is present, otherwise
%   an error is returned.
%   The three arguments are handled by the function \myhlc|\@@_anspic_args:nnn|
%   and stored in the sequence \myhlc|\l_@@_anspic_args_seq| which is processed
%   by the \myenv*[no-index]{keyanspic} environment.
%    \begin{macrocode}
\NewDocumentCommand \anspic { s o +m }
  {
    \bool_if:NF \l_@@_store_active_bool
      {
        \msg_error:nnnn { enumext } { wrong-place }{ keyanspic }{ save-ans }
      }
    \int_compare:nNnT { \l_@@_level_int } > { 1 }
      {
        \msg_error:nn { enumext } { keyanspic-wrong-level }
      }
    \int_compare:nNnT { \l_@@_keyans_level_int } = { 1 }
      {
        \msg_error:nnnn { enumext } { command-wrong-place }{ anspic }{ keyans }
      }
    \seq_put_right:Nn \l_@@_anspic_args_seq
      {
        \@@_anspic_args:nnn {#1} {#2} {#3}
      }
  }
%    \end{macrocode}
% \end{macro}
%
% \subsection{The horizontal environments}\label{subsec:horizontal-envs}
%
% Generating \emph{horizontal list environments} is NOT as simple as standard
% \hologo{LaTeX} list environments. The fundamental part of the code is adapted
% from the \mypkg{shortlst} package to a more modern version using \mypkg{expl3}\cite{expl3}.
% It is not possible to redefine |\item| and |\makelabel| using |\RenewDocumentCommand| as in
% the vertical \emph{non starred} versions.
%
% To achieve the \emph{horizontal list environments}, the |\item|
% command and its \mymeta{content} are captured in a \emph{horizontal box} using |\makebox| for
% the \mykey{label} and a \myenv[no-index]{minipage} environment for the \mymeta{content} passed to
% |\item|; the \emph{optional argument} \myparg{number} is also added to |\item|
% to be able to \emph{join columns} horizontally. In simple terms, the goal is for
% |\item| to behave in the same way as in the \myenv*[no-index]{enumext} environment but
% with an additional \emph{first optional argument} \myparg{number}.
%
% A side effect is the limitation of using |\item| in this way \emph{without}
% using |\RenewDocumentCommand|, which loses the original definition and
% affects the \emph{standard list environments} provided by \hologo{LaTeX}
% and any environment defined using base \myenv[no-index]{list} environment, including:
% \myenv[no-index]{itemize}, \myenv[no-index]{enumerate}, \myenv[no-index]{description}, \myenv[no-index]{quote},
% \myenv[no-index]{quotation}, \myenv[no-index]{verse}, \myenv[no-index]{center}, \myenv[no-index]{flushleft},
% \myenv[no-index]{flushright}, \myenv[no-index]{verbatim}, \myenv[no-index]{tabbing}, \myenv[no-index]{trivlist},
% \myenv[no-index]{list}  and all environments created with |\newtheorem|.
%
% \begin{important}*
% One way to get around this is to use something like:
%
% \begin{Verbatim}
%      \AddToHook{env/enumerate/before}{recover original \item definition}
% \end{Verbatim}
%
% inside \myenv[no-index]{minipage}, but in partial tests this does not have the
% desired effect and the vertical and horizontal spacing is distorted.
% For now this remains a limitation, and its implementation may be
% reconsidered in the future.
% \end{important}
%
% \begin{important}
%   For compatibility with the \emph{tagged} \texttt{PDF}, the environments are closed
%   according to the presence or not of the \mykey{mini-env} key.
% \end{important}
%
% \subsubsection{Definition of error messages}
%
% \begin{macro}[int]{nested-horizontal, item-joined, item-joined-columns}
% Messages used by \myenv*[no-index]{enumext*} and \myenv*[no-index]{keyans*} environments.
%    \begin{macrocode}
\msg_new:nnn { enumext } { nested-horizontal }
  {
    The~environment~\l_@@_envir_name_tl \c_space_tl can't~be~nested~in~'#1'~
    \msg_line_context:.
  }
\msg_new:nnn { enumext } { item-joined }
  {
    Items~joined~(#1)~>~#2~columns~\msg_line_context:.
  }
\msg_new:nnn { enumext } { item-joined-columns }
  {
    Not~enough~space~to~join~items~(#1)~>~#2~\msg_line_context:.
  }
%    \end{macrocode}
% \end{macro}
%
% \subsubsection{Functions for item box width}
%
% \begin{macro}[int]{\@@_starred_columns_set_vii:, \@@_starred_columns_set_viii:}
%   Sets the default value for the \emph{width of the box} containing the
%   \mymeta{content} of the items for \myenv*[no-index]{enumext*} environment.
%    \begin{macrocode}
\cs_new_protected:Nn \@@_starred_columns_set_vii:
  {
    \dim_compare:nNnT { \l_@@_columns_sep_vii_dim } = { \c_zero_dim }
      {
        \dim_set:Nn \l_@@_columns_sep_vii_dim
          {
            ( \l_@@_labelwidth_vii_dim + \l_@@_labelsep_vii_dim )
            / \l_@@_columns_vii_int
          }
      }
    \int_set:Nn \l_@@_tmpa_vii_int { \l_@@_columns_vii_int - 1 }
    \dim_set:Nn \l_@@_item_width_vii_dim
      {
        ( \linewidth - \l_@@_columns_sep_vii_dim * \l_@@_tmpa_vii_int )
        / \l_@@_columns_vii_int
        - \l_@@_labelwidth_vii_dim
        - \l_@@_labelsep_vii_dim
      }
%    \end{macrocode}
%  When the key \mykey[env]{rightmargin} is active the values must be adjusted.
%    \begin{macrocode}
    \dim_compare:nNnT { \l_@@_rightmargin_vii_dim } > { \c_zero_dim }
      {
        \dim_sub:Nn \l_@@_item_width_vii_dim
          {
            ( \l_@@_rightmargin_vii_dim * \l_@@_tmpa_vii_int )
            / \l_@@_columns_vii_int
          }
        \dim_add:Nn \l_@@_columns_sep_vii_dim
          {
            \l_@@_rightmargin_vii_dim
          }
      }
  }
%    \end{macrocode}
%  Same implementation for the \myenv*[no-index]{keyans*} environment.
%    \begin{macrocode}
\cs_new_protected:Nn \@@_starred_columns_set_viii:
  {
    \dim_compare:nNnT { \l_@@_columns_sep_viii_dim } = { \c_zero_dim }
      {
        \dim_set:Nn \l_@@_columns_sep_viii_dim
          {
            ( \l_@@_labelwidth_viii_dim + \l_@@_labelsep_viii_dim )
            / \l_@@_columns_viii_int
          }
      }
    \int_set:Nn \l_@@_tmpa_viii_int { \l_@@_columns_viii_int - 1 }
    \dim_set:Nn \l_@@_item_width_viii_dim
      {
        ( \linewidth - \l_@@_columns_sep_viii_dim * \l_@@_tmpa_viii_int )
        / \l_@@_columns_viii_int
        - \l_@@_labelwidth_viii_dim
        - \l_@@_labelsep_viii_dim
      }
    \dim_compare:nNnT { \l_@@_rightmargin_viii_dim } > { \c_zero_dim }
      {
        \dim_sub:Nn \l_@@_item_width_viii_dim
          {
            ( \l_@@_rightmargin_viii_dim * \l_@@_tmpa_viii_int )
            / \l_@@_columns_viii_int
          }
        \dim_add:Nn \l_@@_columns_sep_viii_dim
          {
            \l_@@_rightmargin_viii_dim
          }
      }
  }
%    \end{macrocode}
% \end{macro}
%
% \subsubsection{Functions for join item columns}
%
% \begin{macro}[int]{\@@_starred_joined_item_vii:n, \@@_starred_joined_item_viii:n}
%   The functions \myhlc|\@@_starred_joined_item_vii:n| and \myhlc|\@@_starred_joined_item_viii:n|
%   will set the \emph{width} of the box in which the \mymeta{content}
%   passed to |\item|\myparg{columns} will be stored together with
%   the value of |\itemwidth| for the \myenv*[no-index]{enumext*} environment.
%    \begin{macrocode}
\cs_new_protected:Nn \@@_starred_joined_item_vii:n
  {
    \int_set:Nn \l_@@_joined_item_vii_int {#1}
    \int_compare:nNnT { \l_@@_joined_item_vii_int } > { \l_@@_columns_vii_int }
      {
        \msg_warning:nnee { enumext } { item-joined }
          { \int_use:N \l_@@_joined_item_vii_int }
          { \int_use:N \l_@@_columns_vii_int }
        \int_set:Nn \l_@@_joined_item_vii_int
          {
            \l_@@_columns_vii_int - \l_@@_item_column_pos_vii_int + 1
          }
      }
    \int_compare:nNnT
      { \l_@@_joined_item_vii_int }
        >
      { \l_@@_columns_vii_int - \l_@@_item_column_pos_vii_int + 1 }
      {
        \msg_warning:nnee { enumext } { item-joined-columns }
          { \int_use:N \l_@@_joined_item_vii_int }
          {
            \int_eval:n
              { \l_@@_columns_vii_int - \l_@@_item_column_pos_vii_int + 1 }
          }
        \int_set:Nn \l_@@_joined_item_vii_int
          {
            \l_@@_columns_vii_int - \l_@@_item_column_pos_vii_int + 1
          }
      }
    \int_compare:nNnTF { \l_@@_joined_item_vii_int } > { 1 }
      {
        \int_set_eq:NN \l_@@_joined_item_aux_vii_int \l_@@_joined_item_vii_int
        \int_decr:N \l_@@_joined_item_aux_vii_int
        \int_add:Nn \l_@@_item_column_pos_vii_int
          { \l_@@_joined_item_aux_vii_int }
        \int_gadd:Nn \g_@@_item_count_all_vii_int
          { \l_@@_joined_item_aux_vii_int }
        \dim_set:Nn \l_@@_joined_width_vii_dim
          {
            \l_@@_item_width_vii_dim * \l_@@_joined_item_vii_int
            + (
                \l_@@_labelwidth_vii_dim + \l_@@_labelsep_vii_dim
                + \l_@@_columns_sep_vii_dim
              ) * \l_@@_joined_item_aux_vii_int
          }
        \dim_set_eq:NN \itemwidth \l_@@_joined_width_vii_dim
      }
      {
        \dim_set_eq:NN \l_@@_joined_width_vii_dim \l_@@_item_width_vii_dim
        \dim_set_eq:NN \itemwidth \l_@@_item_width_vii_dim
      }
  }
%    \end{macrocode}
%   Same implementation for the \myenv*[no-index]{keyans*} environment.
%    \begin{macrocode}
\cs_new_protected:Nn \@@_starred_joined_item_viii:n
  {
    \int_set:Nn \l_@@_joined_item_viii_int {#1}
    \int_compare:nNnT { \l_@@_joined_item_viii_int } >
      { \l_@@_columns_viii_int }
      {
        \msg_warning:nnee { enumext } { item-joined }
          { \int_use:N \l_@@_joined_item_viii_int }
          { \int_use:N \l_@@_columns_viii_int }
        \int_set:Nn \l_@@_joined_item_viii_int
          {
            \l_@@_columns_viii_int - \l_@@_item_column_pos_viii_int + 1
          }
      }
    \int_compare:nNnT
      { \l_@@_joined_item_viii_int }
        >
      { \l_@@_columns_viii_int - \l_@@_item_column_pos_viii_int + 1 }
      {
        \msg_warning:nnee { enumext } { item-joined-columns }
          { \int_use:N \l_@@_joined_item_viii_int }
          {
            \int_eval:n
              { \l_@@_columns_viii_int - \l_@@_item_column_pos_viii_int + 1 }
          }
        \int_set:Nn \l_@@_joined_item_viii_int
          {
            \l_@@_columns_viii_int - \l_@@_item_column_pos_viii_int + 1
          }
      }
    \int_compare:nNnTF { \l_@@_joined_item_viii_int } > { 1 }
      {
        \int_set_eq:NN \l_@@_joined_item_aux_viii_int \l_@@_joined_item_viii_int
        \int_decr:N \l_@@_joined_item_aux_viii_int
        \int_add:Nn \l_@@_item_column_pos_viii_int
          { \l_@@_joined_item_aux_viii_int }
        \int_gadd:Nn \g_@@_item_count_all_viii_int
          { \l_@@_joined_item_aux_viii_int }
        \dim_set:Nn \l_@@_joined_width_viii_dim
          {
            \l_@@_item_width_viii_dim * \l_@@_joined_item_viii_int
            + ( \l_@@_labelwidth_viii_dim + \l_@@_labelsep_viii_dim
                + \l_@@_columns_sep_viii_dim
              )*\l_@@_joined_item_aux_viii_int
          }
        \dim_set_eq:NN \itemwidth \l_@@_joined_width_viii_dim
      }
      {
        \dim_set_eq:NN \l_@@_joined_width_viii_dim \l_@@_item_width_viii_dim
        \dim_set_eq:NN \itemwidth \l_@@_item_width_viii_dim
      }
  }
%    \end{macrocode}
% \end{macro}
%
% \subsubsection{Functions for \texttt{mini-env}, \texttt{mini-right} and \texttt{mini-right*} keys}\label{subsubsec:mini-env}
%
% \begin{macro}[int]{\@@_start_mini_vii:, \@@_stop_mini_vii:}
%   The implementation of the \mykey[env]{mini-env} key support is almost
%   identical to the one used in the \myenv*[no-index]{enumext} and \myenv*[no-index]{keyans}
%   environments, the difference is that the \myhlc|\@@_mini_page|
%   environment on the \emph{\enquote{right side}} is executed
%   \emph{\enquote{after}} closing the environment, so it is necessary to
%   make a global copy of the variable \myhlc|\l_@@_minipage_right_vii_dim| in
%   the variable \myhlc|\g_@@_minipage_right_vii_dim|.%
%    \begin{macrocode}
\cs_new_protected:Nn \@@_start_mini_vii:
  {
    \dim_compare:nNnT { \l_@@_minipage_right_vii_dim } > { \c_zero_dim }
      {
        \dim_set:Nn \l_@@_minipage_left_vii_dim
          {
            \linewidth
            - \l_@@_minipage_right_vii_dim
            - \l_@@_minipage_hsep_vii_dim
          }
        \bool_set_true:N \l_@@_minipage_active_vii_bool
        \dim_gset_eq:NN
          \g_@@_minipage_right_vii_dim
          \l_@@_minipage_right_vii_dim
        \@@_mini_addvspace_vii:
        \nointerlineskip\noindent
        \@@_mini_page{ \l_@@_minipage_left_vii_dim }
      }
  }
%    \end{macrocode}
%   The function \myhlc|\@@_stop_mini_vii:| closes the \myhlc|\@@_mini_page|
%   environment on the \emph{\enquote{left side}}, applies |\hfill| and
%   sets the variable \myhlc|\g_@@_minipage_active_vii_bool| to |true|
%   which will be used in the function \myhlc|\@@_after_env:nn| to execute
%   the \myenv[no-index]{minipage} on the \emph{\enquote{right side}}. At this point
%   the \myhlc|\@@_stop_list:| and \myhlc|\@@_stop_store_level_vii:|
%   functions are executed, stopping the \myenv[no-index]{list} environment and the level saving
%   mechanism for storage in \emph{sequence} of the |\anskey| command
%   and |anskey*| environment. This function is passed to the
%   \myhlc|\@@_after_list_vii:| function in the second part of the \myenv*[no-index]{enumext*}
%   environment definition (\S\ref{env:starred}).
%    \begin{macrocode}
\cs_new_protected:Nn \@@_stop_mini_vii:
  {
    \bool_if:NTF \l_@@_minipage_active_vii_bool
      {
        \@@_stop_list:
        \@@_stop_store_level_vii:
        \IfDocumentMetadataT { \tag_resume:n {enumext*} }
        \end@@_mini_page
        \hfill
        \bool_gset_true:N \g_@@_minipage_active_vii_bool
      }
      {
        \@@_stop_list:
        \@@_stop_store_level_vii:
      }
  }
%    \end{macrocode}
% \end{macro}
%
%  Finally the \mymarg{code} passed to the \mykey{mini-right}
%  or \mykey{mini-right*} keys stored in the variable \myhlc|\g_@@_miniright_code_vii_tl|
%  in the \myenv[no-index]{minipage} environment on the \emph{\enquote{right
%  side}} is executed. For compatibility with the \mypkg{caption} package and possibly
%  other \mymarg{code} passed to this key, it is passed to a box and
%  then printed.
%    \begin{macrocode}
\@@_after_env:nn {enumext*}
  {
    \bool_if:NT \g_@@_minipage_active_vii_bool
      {
        \@@_minipage:w [ t ] { \g_@@_minipage_right_vii_dim }
          \legacy_if_gset_false:n { @minipage }
          \skip_vertical:N \c_zero_skip
          \par\addvspace { \g_@@_minipage_right_skip }
          \bool_if:NF \g_@@_minipage_center_vii_bool
            {
              \tl_gput_left:Nn \g_@@_miniright_code_vii_tl
                {
                  \centering
                }
            }
          \vbox_set_top:Nn \l_@@_miniright_code_vii_box
            {
              \tl_use:N \g_@@_miniright_code_vii_tl
            }
          \box_use_drop:N \l_@@_miniright_code_vii_box
          \skip_vertical:N \c_zero_skip
        \@@_endminipage:
        \par\addvspace{ \g_@@_minipage_after_skip }
      }
    \bool_gset_false:N \g_@@_minipage_active_vii_bool
    \bool_gset_true:N \g_@@_minipage_center_vii_bool
    \tl_gclear:N \g_@@_miniright_code_vii_tl
    \dim_gzero:N \g_@@_minipage_right_vii_dim
    \bool_gset_false:N \g_@@_starred_bool
  }
%    \end{macrocode}
%
% \begin{macro}[int]{\@@_start_mini_viii:, \@@_stop_mini_viii:}
%   The implementation of the \mykey[env]{mini-env}, \mykey[env]{mini-right} and
%   \mykey[env]{mini-right*} keys is identical to the one used in the
%   \myenv*[no-index]{enumext*} environment.
%    \begin{macrocode}
\cs_new_protected:Nn \@@_start_mini_viii:
  {
    \dim_compare:nNnT { \l_@@_minipage_right_viii_dim } > { \c_zero_dim }
      {
        \dim_set:Nn \l_@@_minipage_left_viii_dim
          {
            \linewidth
            - \l_@@_minipage_right_viii_dim
            - \l_@@_minipage_hsep_viii_dim
          }
        \bool_set_true:N \l_@@_minipage_active_viii_bool
        \dim_gset_eq:NN
          \g_@@_minipage_right_viii_dim
          \l_@@_minipage_right_viii_dim
        \@@_mini_addvspace_viii:
        \nointerlineskip\noindent
        \@@_mini_page{ \l_@@_minipage_left_viii_dim }
      }
  }
\cs_new_protected:Nn \@@_stop_mini_viii:
  {
    \bool_if:NTF \l_@@_minipage_active_viii_bool
      {
        \@@_stop_list:
        \IfDocumentMetadataT { \tag_resume:n {keyans*} }
        \end@@_mini_page
        \hfill
        \bool_gset_true:N \g_@@_minipage_active_viii_bool
      }
      {
        \@@_stop_list:
      }
  }
\@@_after_env:nn {keyans*}
  {
    \bool_if:NT \g_@@_minipage_active_viii_bool
      {
        \@@_mini_page{ \g_@@_minipage_right_viii_dim }
          \par\addvspace { \g_@@_minipage_right_skip }
          \bool_if:NF \g_@@_minipage_center_viii_bool
            {
              \tl_gput_left:Nn \g_@@_miniright_code_viii_tl
                {
                  \centering
                }
            }
          \vbox_set_top:Nn \l_@@_miniright_code_viii_box
            {
              \tl_use:N \g_@@_miniright_code_viii_tl
            }
          \box_use_drop:N \l_@@_miniright_code_viii_box
        \end@@_mini_page
        \par\addvspace{ \g_@@_minipage_after_skip }
      }
    \bool_gset_false:N \g_@@_minipage_active_viii_bool
    \bool_gset_true:N \g_@@_minipage_center_viii_bool
    \tl_gclear:N \g_@@_miniright_code_viii_tl
    \dim_gzero:N \g_@@_minipage_right_viii_dim
  }
%    \end{macrocode}
% \end{macro}
%
% \subsection{The environment \env{enumext*}}\label{env:starred}
%
% \begin{macro}{enumext*}
%    First the environment is generated and a
%    temporary definition is given to \myhlc|\@@_stop_item_tmp_vii:| equal to
%    \myhlc|\@@_first_item_tmp_vii:| and next to |\item| equal to \myhlc|\@@_start_item_tmp_vii:|
%    which are redefined later. Unlike the implementation used by
%    the \mypkg[no-index]{shortlst} package, the values of
%    |\rightskip| and |\@rightskip| are not set equal to |\@flushglue|, whose value is
%    \mydim{0.0pt plus 1.0 fil}; in the tests performed this fails in
%    some circumstances and different results are obtained when using
%    \hologo{pdfTeX} and \hologo{LuaTeX}.
%    \begin{macrocode}
\NewDocumentEnvironment{enumext*}{ o }
  {
    \@@_safe_exec_vii:
    \@@_parse_keys_vii:n {#1}
    \@@_before_list_vii:
    \@@_start_store_level_vii:
    \@@_start_list:nn { }
      {
        \@@_list_arg_two_vii:
        \@@_before_keys_exec_vii:
      }
    \IfDocumentMetadataT { \tag_suspend:n {enumext*} }
    \@@_starred_columns_set_vii:
    \item[] \scan_stop:
%    \end{macrocode}
%   For compatibility with \emph{tagged} \texttt{PDF} in \hologo{LaTeX}
%   release 2026-06-01, |\setcounter| is set at this point in the
%   environment definition, thereby avoiding the problem reported by
%   \texttt{@teepeemm} in \href{https://github.com/pablgonz/enumext/issues/5}{enumext* nested in enumext starts at 2 with DocumentMetadata}
%   this was caused by the \mypkg[no-index]{latex-lab-block}\cite{latexlab} update in this release.
%    \begin{macrocode}
    \setcounter { enumXvii }
      { \int_eval:n { \int_use:c { l_@@_start_vii_int } - 1 } }
    \cs_set_eq:NN \@@_stop_item_tmp_vii: \@@_first_item_tmp_vii:
    \cs_set_eq:NN \item \@@_start_item_tmp_vii:
    \ignorespaces
  }
  {
    % {tag=text-unit}
    \IfDocumentMetadataT { \tag_struct_end:n {tag=semantic-para} }
    \@@_stop_item_tmp_vii:
    \@@_remove_extra_parsep_vii:
    \@@_after_list_vii:
  }
%    \end{macrocode}
% \end{macro}
%
% \begin{macro}[int]{\@@_safe_exec_vii:}\label{fun:safe-exec-starred}
%   The function \myhlc|\@@_is_not_nested:| is called first, which sets
%   \myhlc|\g_@@_starred_bool| to true if it is NOT nested within \myenv*[no-index]{enumext};
%   then the function \myhlc|\@@_internal_mini_page:| is called to create the
%   environment \myhlc|\@@_mini_page|, \myhlc|\l_@@_level_h_int| is incremented
%   to restrict nesting of the environment, \myhlc|\l_@@_starred_bool| is set
%   to true, and finally the function \myhlc|\@@_is_on_first_level:| is called, which
%   sets \myhlc|\l_@@_starred_first_bool| to
%   true if it is not nested, allowing the \emph{\enquote{storage system}}
%   to be used.
%    \begin{macrocode}
\cs_new_protected:Nn \@@_safe_exec_vii:
  {
    \@@_is_not_nested:
    \@@_internal_mini_page:
    \int_incr:N \l_@@_level_h_int
    \int_compare:nNnT { \l_@@_level_h_int } > { 1 }
      {
        \msg_error:nn { enumext } { nested }
      }
    \int_compare:nNnT { \l_@@_keyans_level_h_int } = { 1 }
      {
        \msg_error:nnn { enumext } { nested-horizontal } { keyans*}
      }
    \bool_set_true:N \l_@@_starred_bool
    \bool_set_false:N \l_@@_standard_bool
    \@@_is_on_first_level:
  }
%    \end{macrocode}
% \end{macro}
%
% \begin{macro}[int]{\@@_parse_keys_vii:n}
%   It first checks the state of the variable \myhlc|\l_@@_resume_count_vii_bool|
%   set by the key \mykey[env]{resume} and calls the function \myhlc|\@@_resume_last_counter:|
%   if it is |true|, then clears the variable \myhlc|\l_@@_series_name_str|
%   used by the key \mykey[env]{series}, processes the environment \myoarg{key \textnormal{\textcolor{gray}{=}} val}
%   and executes the function \myhlc|\@@_parse_series:n| used by the key
%   \mykey{series}, then executes the function \myhlc|\@@_store_active_keys_vii:n|
%   and reprocesses the \mymeta{keys} to pass them to the storage \emph{sequence}
%   if the key \mykey[env]{save-key} is not active.
%
% \begin{important}*
%  Here it is necessary to check the status of \myhlc|\l_@@_resume_count_vii_bool|
%  in case the key \mykey{resume} is set using |\setenumext{enumext*}{resume}|.
% \end{important}
%
%    \begin{macrocode}
\cs_new_protected:Nn \@@_parse_keys_vii:n
  {
    \bool_if:NT \l_@@_resume_count_vii_bool
      {
        \@@_resume_last_counter:
      }
    \tl_if_novalue:nF {#1}
      {
        \str_clear:N \l_@@_series_name_str
        \keys_set:nn { enumext / enumext* } {#1}
        \bool_if:NF \l_@@_print_keyans_cmd_bool
          {
            \@@_parse_series:n {#1}
          }
        \@@_store_active_keys_vii:n {#1}
      }
  }
%    \end{macrocode}
% \end{macro}
%
% \begin{macro}[int]{\@@_before_list_vii:}
%   The function \myhlc|\@@_before_list_vii:| first calls the function
%   \myhlc|\@@_vspace_above_vii:| used by the keys \mykey[env]{above} and
%   \mykey[env]{above*}, then calls the function \myhlc|\@@_check_ans_active:| for
%   the check answer mechanism and finally calls the functions
%   \myhlc|\@@_before_args_exec_vii:| and \myhlc|\@@_start_mini_vii:| used by the
%   keys \mykey[env]{before*}, \mykey[env]{mini-env}, \mykey[env]{mini-right} and
%   \mykey[env]{mini-right*}.
%    \begin{macrocode}
\cs_new_protected:Nn \@@_before_list_vii:
  {
    \@@_vspace_above_vii:
    \@@_check_ans_active:
    \@@_before_args_exec_vii:
    \@@_start_mini_vii:
  }
%    \end{macrocode}
% \end{macro}
%
% \begin{macro}[int]{\@@_after_list_vii:}
%   The function \myhlc|\@@_after_list_vii:| first calls the function
%   \myhlc|\@@_stop_mini_vii:| which internally calls \myhlc|\@@_stop_list:| and
%   \myhlc|\@@_stop_store_level_vii:| (\S\ref{subsubsec:mini-env}) used
%   by the keys \mykey[env]{mini-env}, \mykey[env]{mini-right} and \mykey[env]{mini-right*},
%   then to the functions \myhlc|\@@_after_stop_list_vii:| used by the key
%   \mykey[env,foreachkeyans]{after}, \myhlc|\@@_check_ans_key_hook:| used by the key \mykey[env]{check-ans},
%   \myhlc|\@@_vspace_below_vii:| used by the keys \mykey[env]{below} and
%   \mykey[env]{below*}. Finally set \myhlc|\l_@@_starred_bool| to false and call
%   the \myhlc|\@@_resume_save_counter:| function used by the \mykey[env]{series},
%   \mykey[env]{resume} and \mykey[env]{resume*} keys.
%    \begin{macrocode}
\cs_new_protected:Nn \@@_after_list_vii:
  {
    \@@_stop_mini_vii:
    \@@_after_stop_list_vii:
    \@@_check_ans_key_hook:
    \@@_vspace_below_vii:
    \bool_set_false:N \l_@@_starred_bool
    \bool_if:NF \l_@@_print_keyans_cmd_bool
      {
        \@@_starred_save_counter:
      }
  }
%    \end{macrocode}
% \end{macro}
%
% \begin{macro}[int]{\@@_start_store_level_vii:, \@@_stop_store_level_vii:}
%   The \myhlc|\@@_start_store_level_vii:| and \myhlc|\@@_stop_store_level_vii:|
%   functions activate the \emph{\enquote{storing structure}} mechanism
%   in \emph{sequence} for |\anskey| command and \myenv*[no-index]{anskey*}
%   environment if \myenv*[no-index]{enumext*} are nested in \myenv*[no-index]{enumext}.
%    \begin{macrocode}
\cs_new_protected:Nn \@@_start_store_level_vii:
  {
    \bool_if:NT \l_@@_store_active_bool
      {
        \int_compare:nNnT { \l_@@_level_int } > { 0 }
          {
            \@@_store_level_open_vii:
          }
      }
  }
\cs_new_protected:Nn \@@_stop_store_level_vii:
  {
    \bool_if:NT \l_@@_store_active_bool
      {
        \int_compare:nNnT { \l_@@_level_int } > { 0 }
          {
            \@@_store_level_close_vii:
          }
      }
  }
%    \end{macrocode}
% \end{macro}
%
% \subsubsection{The command \cs{item} in \env{enumext*}}\label{cmd:item-enumext-starred}
%
% \begin{macro}[int]{\@@_first_item_tmp_vii:}
%
% The \myhlc|\@@_first_item_tmp_vii:| function will remove horizontal space
% equal to |\labelwidth| plus |\labelsep| to the left of the \emph{\enquote{first}}
% |\item| in the environment at the point of execution of this function,
% where it is equal to the \myhlc|\@@_stop_item_tmp_vii:| function inside the
% environment body definition.
%    \begin{macrocode}
\cs_new_protected_nopar:Nn \@@_first_item_tmp_vii:
  {
    \skip_horizontal:n
      {
        -\l_@@_labelwidth_vii_dim - \l_@@_labelsep_vii_dim
      }
    \ignorespaces
  }
%    \end{macrocode}
% \end{macro}
%
% \begin{macro}[int]{\@@_start_item_tmp_vii:, \@@_item_peek_args_vii:,
%                    \@@_joined_item_vii:w, \@@_standard_item_vii:w,
%                    \@@_starred_item_vii:w, \@@_starred_item_vii_aux_i:w,
%                    \@@_starred_item_vii_aux_ii:w, \@@_starred_item_vii_aux_iii:w}
%   First the function \myhlc|\@@_stop_item_tmp_vii:| is called, which is
%   redefined later, then the value of
%   \myhlc|\l_@@_item_column_pos_vii_int|, which counts the items by rows, and
%   the value of \myhlc|\g_@@_item_count_all_vii_int|, which counts the
%   total of items in the environment, are incremented. After that the
%   function \myhlc|\@@_item_peek_args_vii:| is called, which handles the
%   arguments passed to |\item|.
%    \begin{macrocode}
\cs_new_protected_nopar:Nn \@@_start_item_tmp_vii:
  {
    \@@_stop_item_tmp_vii:
    \int_incr:N \l_@@_item_column_pos_vii_int
    \int_gincr:N \g_@@_item_count_all_vii_int
    \@@_item_peek_args_vii:
  }
%    \end{macrocode}
%   The function \myhlc|\@@_item_peek_args_vii:| handles the |\item|\myparg{number}.
%   It looks for the argument \enquote{\textcolor{star}{\ttfamily (}};
%   if it is present, it calls the function \myhlc|\@@_joined_item_vii:w| \myparg{number},
%   which is in charge of joining the items in the same row; if it is
%   not present, it sets the default value \textcolor{gray}{\ttfamily(\mydim{1})}.
%    \begin{macrocode}
\cs_new_protected:Nn \@@_item_peek_args_vii:
  {
    \peek_meaning:NTF (
      { \@@_joined_item_vii:w }
      { \@@_joined_item_vii:w (1) }
  }
%    \end{macrocode}
%    The function \myhlc|\@@_joined_item_vii:w| will first call the
%    function \myhlc|\@@_starred_joined_item_vii:n| in charge of setting the
%    \emph{width} of the box that stores the content passed to |\item|. Then it
%    looks for the argument \enquote{\textcolor{star}{\ttfamily *}}; if it is present, it calls the
%    function \myhlc|\@@_starred_item_vii:w|, otherwise it calls the
%    function \myhlc|\@@_standard_item_vii:w|.
%    \begin{macrocode}
\cs_new_protected:Npn \@@_joined_item_vii:w (#1)
  {
    \@@_starred_joined_item_vii:n {#1}
    \peek_meaning_remove:NTF *
      { \@@_starred_item_vii:w  }
      { \@@_standard_item_vii:w }
  }
%    \end{macrocode}
%   The function \myhlc|\@@_standard_item_vii:w| will first look for the
%   argument `|[|', if present it will set the state of the variable
%   \myhlc|\l_@@_wrap_label_opt_vii_bool| equal to the state of the variable
%   \myhlc|\l_@@_wrap_label_opt_vii_bool| handled by the key
%   \mykey[env]{wrap-label*} and finally execute the \emph{non-enumerated}
%   version |\item|\myoarg{custom} by means of the function
%   \myhlc|\@@_start_item_vii:w|; otherwise it sets the value of the
%   variable \myhlc|\l_@@_wrap_label_vii_bool| handled by the
%   \mykey[env]{wrap-label} key to true and set the switch |\if@noitemarg| to
%   true to execute the enumerated version of |\item| by means of the
%   function \myhlc|\@@_start_item_vii:w| |[| \myhlc|\l_@@_label_vii_tl| |]|.
%    \begin{macrocode}
\cs_new_protected:Npn \@@_standard_item_vii:w
  {
    \bool_set_false:N \l_@@_item_starred_vii_bool
    \peek_meaning:NTF [
      {
        \bool_set_eq:NN \l_@@_wrap_label_vii_bool \l_@@_wrap_label_opt_vii_bool
        \@@_start_item_vii:w
      }
      {
        \bool_set_true:N \l_@@_wrap_label_vii_bool
        \legacy_if_set_true:n { @noitemarg }
        \@@_start_item_vii:w [ \l_@@_label_vii_tl ] \ignorespaces
      }
  }
%    \end{macrocode}
%   The function \myhlc|\@@_starred_item_vii:w| together with the specified
%   auxiliary functions \myhlc|aux_i:w|, \myhlc|aux_ii:w|, and \myhlc|aux_iii:w| execute |\item*|,
%   |\item*|\myoarg{symbol} and |\item*|\myoarg{symbol}\myoarg{offset}.
%    \begin{macrocode}
\cs_new_protected:Npn \@@_starred_item_vii:w
  {
    \bool_set_true:N \l_@@_item_starred_vii_bool
    \bool_set_true:N \l_@@_wrap_label_vii_bool
    \peek_meaning:NTF [
      { \@@_starred_item_vii_aux_i:w }
      { \@@_starred_item_vii_aux_ii:w }
  }
\cs_new_protected:Npn \@@_starred_item_vii_aux_i:w [#1]
  {
    \tl_gset:Nn \g_@@_item_symbol_aux_vii_tl {#1}
    \@@_starred_item_vii_aux_ii:w
  }
\cs_new_protected:Npn \@@_starred_item_vii_aux_ii:w
  {
    \peek_meaning:NTF [
      { \@@_starred_item_vii_aux_iii:w }
      {
        \dim_set_eq:NN \l_@@_item_symbol_sep_vii_dim \l_@@_labelsep_vii_dim
        \legacy_if_set_true:n { @noitemarg }
        \@@_start_item_vii:w [ \l_@@_label_vii_tl ] \ignorespaces
      }
  }
\cs_new_protected:Npn \@@_starred_item_vii_aux_iii:w [#1]
  {
    \dim_set:Nn \l_@@_item_symbol_sep_vii_dim {#1}
    \legacy_if_set_true:n { @noitemarg }
    \@@_start_item_vii:w [ \l_@@_label_vii_tl ] \ignorespaces
  }
%    \end{macrocode}
% \end{macro}
%
% \begin{macro}[int]{\@@_fake_make_label_vii:n}
%   The \myhlc|\@@_fake_make_label_vii:n| function will be in charge of
%   handling the definition of |\item|. First it increments the counter
%   \icounter[no-index]{enumXvii} for the enumerated items and activates support for
%   the \emph{check answers} mechanism, followed by support for
%   |\item*|\myoarg{symbol}\myoarg{offset} if present, then the
%   \mykey[env]{wrap-label} and \mykey[env]{wrap-label*} keys, executed using
%   |\makebox| whose width is given by the \mykey[env]{labelwidth} key and
%   position by the \mykey[env]{align} key; inside its argument it
%   executes the \mykey[env]{font} key together with the function defined by the
%   \mykey{wrap-label} or \mykey{wrap-label*} keys. Finally it executes the
%   \mykey[env]{labelsep} key applying a \myhlc|\skip_horizontal:N| and |\ignorespaces|.
%
% \begin{important}
%   For compatibility with \emph{tagged} \texttt{PDF} and \mypkg[no-index]{hyperref}
%   when an environment \myenv*[no-index]{enumext} is nested in \myenv*[no-index]{enumext*}
%   and the key \mykey[env]{save-ans} is not active need setting the |\if@hyper@item|
%   switch to |true|. The explanation for this is given
%   by the master Heiko Oberdiek on \href{https://tex.stackexchange.com/a/404911}{\textbackslash
%   refstepcounter\{enumi\} twice (or more) creates destination with the
%   same identifier}.
%   This patch is only needed if you are running |pdflatex| and not if
%   you are running |lualatex|
% \end{important}
%
%    \begin{macrocode}
\cs_new_protected_nopar:Nn \@@_fake_make_label_vii:n
  {
    \legacy_if:nT { @noitemarg }
      {
        \legacy_if_set_false:n { @noitemarg }
        \legacy_if:nT { @nmbrlist }
          {
            \IfDocumentMetadataT
              {
                \bool_if:NT \l_@@_hyperref_bool
                  {
                    \legacy_if_set_true:n { @hyper@item }
                  }
              }
            \refstepcounter{enumXvii}
            \bool_if:NT \l_@@_check_answers_bool
              {
                \int_gincr:N \g_@@_item_number_int
                \bool_set_true:N \l_@@_item_number_bool
              }
          }
      }
    \bool_if:NT \l_@@_item_starred_vii_bool
      {
        \tl_if_blank:VT \g_@@_item_symbol_aux_vii_tl
          {
            \tl_gset_eq:NN
              \g_@@_item_symbol_aux_vii_tl \l_@@_item_symbol_vii_tl
          }
        \mode_leave_vertical:
        \skip_horizontal:n { -\l_@@_item_symbol_sep_vii_dim }
        \hbox_overlap_left:n { \g_@@_item_symbol_aux_vii_tl }
        \skip_horizontal:N \l_@@_item_symbol_sep_vii_dim
        \tl_gclear:N \g_@@_item_symbol_aux_vii_tl
      }
    \makebox[ \l_@@_labelwidth_vii_dim ][ \l_@@_align_label_vii_str ]
      {
        \tl_use:N \l_@@_label_font_style_vii_tl
        \bool_if:NTF \l_@@_wrap_label_vii_bool
          {
            \@@_wrapper_label_vii:n {#1}
          }
          {#1}
      }
    \skip_horizontal:N \l_@@_labelsep_vii_dim \ignorespaces
  }
%    \end{macrocode}
% \end{macro}
%
% \subsubsection{Real definition of \cs{item} in \env{enumext*}}
%
% The functions \myhlc|\@@_start_item_vii:w| and \myhlc|\@@_stop_item_vii:| executing
% the true definition of |\item| inside the \myenv*[no-index]{enumext*} environment, unlike
% the implementation in \mypkg[no-index]{shortlst}, an extra group is NOT used
% and the plain form of the \myenv[no-index]{lrbox} environment.
%
% \begin{macro}[int]{\@@_start_item_vii:w, \@@_stop_item_vii:}
%   First the value of \myhlc|\@@_stop_item_tmp_vii:| is set equal to
%   \myhlc|\@@_stop_item_vii:|, which is defined later; after that it
%   starts capturing |\item| and \emph{\enquote{item content}} in a \emph{horizontal box}
%   where the width will be |\itemwidth| plus |\labelwidth| plus |\labelsep|.
%    \begin{macrocode}
\cs_new_protected_nopar:Npn \@@_start_item_vii:w [#1]
  {
    \cs_set_eq:NN \@@_stop_item_tmp_vii: \@@_stop_item_vii:
    \hbox_set_to_wd:Nnw \l_@@_item_text_vii_box
      {
        \l_@@_joined_width_vii_dim
        + \l_@@_labelwidth_vii_dim
        + \l_@@_labelsep_vii_dim
      }
%    \end{macrocode}
%   Redefine the |\footnote| command.
%    \begin{macrocode}
      \@@_renew_footnote_starred:
%    \end{macrocode}
%   Now the \emph{sockets} for \emph{tagging} \texttt{PDF}
%   support are inserted and |\item| is run.
%    \begin{macrocode}
      \@@_start_list_tag:n {enumext*}
      \@@_fake_make_label_vii:n {#1}
      \@@_stop_start_list_tag:
%    \end{macrocode}
%   Finally the \myenv[no-index]{minipage} environment is opened, the
%   \emph{\enquote{item content}} is captured, |\parindent| takes the value of the
%   key \mykey[env]{listparindent} and |\parskip| takes the value of the key
%   \mykey[env]{parsep}, then the keys \mykey[env]{itemindent} and
%   \mykey[env]{first} are executed.
%
% \begin{important}*
%  Here the use of |\unskip| and \myhlc|\skip_horizontal:n| with the value of
%  \mykey{listparindent} is necessary, otherwise an unwanted space is
%  created when using |\item|\myoarg{opt} and the value passed to the key
%  \mykey{itemindent} is incremented.
% \end{important}
%
%    \begin{macrocode}
      \@@_minipage:w [ t ]{ \l_@@_joined_width_vii_dim }
        \dim_set_eq:NN \parindent \l_@@_listparindent_vii_dim
        \skip_set_eq:NN \parskip \l_@@_parsep_vii_skip
        \@@_unskip_unkern:
        \@@_unskip_unkern:
        \skip_horizontal:n { -\l_@@_listparindent_vii_dim } \ignorespaces
        \tl_use:N \l_@@_fake_item_indent_vii_tl
        \tl_use:N \l_@@_after_list_args_vii_tl
  }
%    \end{macrocode}
%   The \myhlc|\@@_stop_item_vii:| function will finish fetching |\item|
%   and \emph{\enquote{item content}} by closing the \myenv[no-index]{minipage}
%   environment, the \emph{sockets} for \emph{tagging} \texttt{PDF} and the
%   \emph{horizontal box}.
%    \begin{macrocode}
\cs_new_protected_nopar:Nn \@@_stop_item_vii:
  {
      \@@_endminipage:
    \@@_stop_list_tag:n {enumext*}
    \hbox_set_end:
%    \end{macrocode}
%   Here the \emph{warnings} are reduced a bit by setting the value
%   of |\hbadness| to \mydim{10000}, print |\item| and \emph{\enquote{item content}}
%   from the \emph{horizontal box}.
%    \begin{macrocode}
    \int_set:Nn \hbadness { 10000 }
    \box_use_drop:N \l_@@_item_text_vii_box
%    \end{macrocode}
%   Finally apply the \emph{vertical space} between rows set by \mykey[env]{itemsep}
%   key passed to |\parsep| using |\par\noindent| and \emph{horizontal space}
%   between columns set by \mykey[env]{columns-sep} key using \myhlc|\skip_horizontal:N|.
%    \begin{macrocode}
    \int_compare:nNnTF
      { \l_@@_item_column_pos_vii_int } = { \l_@@_columns_vii_int }
      {
        \par\noindent
        \int_zero:N \l_@@_item_column_pos_vii_int
      }
      {
        \skip_horizontal:N \l_@@_columns_sep_vii_dim
      }
  }
%    \end{macrocode}
% \end{macro}
%
% \begin{macro}[int]{\@@_remove_extra_parsep_vii:}
%   Remove the extra \emph{vertical space} equal to |\parsep=\itemsep| when the
%   total number of |\item| is divisible by the number of |\item| in the last
%   row of the environment. Here the use of |\unskip| or |\removelastskip|
%   fails and does not obtain the expected result; using |\vspace| is
%   the option and in this case a simplified version can be used since
%   it is always in \mymeta{vertical mode}.
%    \begin{macrocode}
\cs_new_protected:Nn \@@_remove_extra_parsep_vii:
  {
    \int_compare:nNnT
      {
        \int_mod:nn
          {  \g_@@_item_count_all_vii_int } { \l_@@_columns_vii_int }
      }
      =
      { 0 }
      {
        \para_end:
        \skip_vertical:n { -\l_@@_itemsep_vii_skip }
        \skip_vertical:N \c_zero_skip
        \int_gzero:N \g_@@_item_count_all_vii_int
      }
  }
%    \end{macrocode}
% \end{macro}
%
% Since the \mykey{check-ans} check must not be executed by levels
% but on the complete list, it is taken out of the \myenv*[no-index]{enumext*}
% environment using the \emph{\enquote{hook}} function \myhlc|\@@_after_env:nn|.
%    \begin{macrocode}
\@@_after_env:nn {enumext*}
  {
    \@@_execute_after_env:
  }
%    \end{macrocode}
%
% \subsection{The environment \env{keyans*}}\label{env:keyans-starred}
%
% \begin{macro}{keyans*}
%   The implementation of \myenv*{keyans*} environment is the
%   similar as that used by the \myenv*[no-index]{enumext*} environment except for the
%   \myhlc|\@@_check_starred_cmd:n| function added in the second part.
%    \begin{macrocode}
\NewDocumentEnvironment{keyans*}{ o }
  {
    \@@_safe_exec_viii:
    \@@_parse_keys_viii:n {#1}
    \@@_before_list_viii:
    \@@_start_list:nn { }
      {
        \@@_list_arg_two_viii:
        \@@_before_keys_exec_viii:
      }
    \IfDocumentMetadataT { \tag_suspend:n {keyans*} }
    \@@_starred_columns_set_viii:
    \item[] \scan_stop:
    % Set start counter
    \setcounter { enumXviii }
      { \int_eval:n { \int_use:c { l_@@_start_viii_int } - 1 } }
    \cs_set_eq:NN \@@_stop_item_tmp_viii: \@@_first_item_tmp_viii:
    \cs_set_eq:NN \item \@@_start_item_tmp_viii:
    \ignorespaces
  }
  {
    % {tag=text-unit}
    \IfDocumentMetadataT { \tag_struct_end:n {tag=semantic-para} }
    \@@_stop_item_tmp_viii:
    \@@_remove_extra_parsep_viii:
    \@@_check_starred_cmd:n { item }
    \@@_after_list_viii:
  }
%    \end{macrocode}
% \end{macro}
%
% \begin{macro}[int]{\@@_safe_exec_viii:}
%   The \myhlc|\@@_safe_exec_viii:| function first checks if the
%   \mykey[env]{save-ans} key is active; only when this is true is the
%   environment available. It increments the value of
%   \myhlc|\l_@@_keyans_level_h_int| and returns an error message when the
%   environment is nested, then calls the
%   \myhlc|\@@_keyans_name_and_start:| function in charge of saving the name of
%   the environment and the line it is running on, then checks if
%   \myenv*[no-index]{keyans*} is being nested in \myenv*[no-index]{enumext*}, returning an
%   error, and sets \myhlc|\l_@@_starred_bool| to true; finally it
%   checks if it is at the appropriate level within the
%   \myenv*[no-index]{enumext} environment.
%    \begin{macrocode}
\cs_new_protected:Nn \@@_safe_exec_viii:
  {
    \bool_if:NF \l_@@_store_active_bool
      {
        \msg_error:nnnn { enumext } { wrong-place }{ keyans* }{ save-ans }
      }
    \int_incr:N \l_@@_keyans_level_h_int
    \int_compare:nNnT { \l_@@_keyans_level_h_int } > { 1 }
      {
        \msg_error:nn { enumext } { nested }
      }
    \@@_keyans_name_and_start:
    \bool_if:NT \l_@@_starred_bool
      {
        \msg_error:nnn { enumext } { nested-horizontal } { enumext* }
      }
    \bool_set_true:N \l_@@_starred_bool
    % Set false for interfering with enumext nested in keyans*
    % (yes, its possible and crayze)
    \bool_set_false:N \l_@@_store_active_bool
    \int_compare:nNnT { \l_@@_level_int } > { 1 }
      {
        \msg_error:nn { enumext } { keyans-wrong-level }
      }
  }
%    \end{macrocode}
% \end{macro}
%
% \begin{macro}[int]{\@@_parse_keys_viii:n}
%    Parse \myoarg{key \textnormal{\textcolor{gray}{=}} val}
%    for \myenv*[no-index]{keyans*}.
%    \begin{macrocode}
\cs_new_protected:Nn \@@_parse_keys_viii:n
  {
    \tl_if_novalue:nF {#1}
      {
        \keys_set:nn { enumext / keyans* } {#1}
      }
  }
%    \end{macrocode}
% \end{macro}
%
% \begin{macro}[int]{\@@_before_list_viii:}
%    The function \myhlc|\@@_before_list_viii:| will add the vertical
%    spacing on the environment if the \mykey[env]{above} key is active next to
%    the \mymarg{code} defined by the \mykey[env]{before*} key if
%    it is active (function \myhlc|\@@_before_args_exec_viii:|), then call the
%    function \myhlc|\@@_start_mini_viii:| handled by
%    \mykey[env]{mini-env}.
%    \begin{macrocode}
\cs_new_protected:Nn \@@_before_list_viii:
  {
    \@@_vspace_above_viii:
    \@@_before_args_exec_viii:
    \@@_start_mini_viii:
  }
%    \end{macrocode}
% \end{macro}
%
% \begin{macro}[int]{\@@_after_list_viii:}
%    The function \myhlc|\@@_after_list_viii:| first calls the function \myhlc|\@@_stop_mini_viii:|,
%    then apply the \mymarg{code} handled by the \mykey[env,foreachkeyans]{after} key
%    together with the \emph{vertical space} handled by the \mykey[env]{below} key if
%    they are present.
%    \begin{macrocode}
\cs_new_protected:Nn \@@_after_list_viii:
  {
    \@@_stop_mini_viii:
    \@@_after_stop_list_viii:
    \@@_vspace_below_viii:
  }
%    \end{macrocode}
% \end{macro}
%
% \subsubsection{The command \cs{item} in \env{keyans*}}\label{cmd:item-keyans-starred}
%
% The idea here is to make the |\item| command behave in the same way as
% in the \myenv*[no-index]{keyans} environment with the difference of the \emph{optional argument}
% \myparg{number} which works in the same way as in the \myenv*[no-index]{enumext*}
% environment.
%
% In simple terms, the goal is to store the \mymeta{label} next to the \myoarg{content} if it is present
% in the \emph{sequence} and \emph{prop list} defined by \mykey{save-ans} key for
% |\item*|, |\item*|\myoarg{content}, |\item|\myparg{number}|*| and
% |\item|\myparg{number}|*|\myoarg{content} commands.
%
% \begin{macro}[int]{\@@_first_item_tmp_viii:}
%   The \myhlc|\@@_first_item_tmp_viii:| function will remove horizontal space
%   equal to |\labelwidth| plus |\labelsep| to the left of the \emph{\enquote{first}}
%   |\item| in the environment at the point of execution of this function,
%   where it is equal to the \myhlc|\@@_stop_item_tmp_viii:| function inside the
%   environment body definition.
%    \begin{macrocode}
\cs_new_protected_nopar:Nn \@@_first_item_tmp_viii:
  {
    \skip_horizontal:n
      {
        -\l_@@_labelwidth_viii_dim - \l_@@_labelsep_viii_dim
      }
    \ignorespaces
  }
%    \end{macrocode}
% \end{macro}
%
% \begin{macro}[int]{\@@_start_item_tmp_viii:, \@@_item_peek_args_viii:,
%                    \@@_joined_item_viii:w, \@@_standard_item_viii:w}
%   First the function \myhlc|\@@_stop_item_tmp_viii:| is called, which is
%   redefined later, then the value of
%   \myhlc|\l_@@_item_column_pos_viii_int|, which counts the items by rows, and
%   the value of \myhlc|\g_@@_item_count_all_viii_int|, which counts the
%   total of items in the environment, are incremented. After that the
%   function \myhlc|\@@_item_peek_args_viii:| is called, which handles the
%   arguments passed to |\item|.
%    \begin{macrocode}
\cs_new_protected_nopar:Nn \@@_start_item_tmp_viii:
  {
    \@@_stop_item_tmp_viii:
    \int_incr:N \l_@@_item_column_pos_viii_int
    \int_gincr:N \g_@@_item_count_all_viii_int
    \@@_item_peek_args_viii:
  }
%    \end{macrocode}
%   The function \myhlc|\@@_item_peek_args_viii:| handles the |\item|\myparg{number}.
%   It looks for the argument \enquote{\textcolor{star}{\ttfamily (}}; if it is
%   present, it calls the function \myhlc|\@@_joined_item_viii:w| \myparg{number},
%   which is in charge of joining the items in the same row; if it is
%   not present, it sets the default value \textcolor{gray}{\ttfamily(\mydim{1})}.
%    \begin{macrocode}
\cs_new_protected:Nn \@@_item_peek_args_viii:
  {
    \peek_meaning:NTF (
      { \@@_joined_item_viii:w }
      { \@@_joined_item_viii:w (1) }
  }
%    \end{macrocode}
%   The function \myhlc|\@@_joined_item_viii:w| will first call the
%   function \myhlc|\@@_starred_joined_item_viii:n| in charge of setting the
%   \emph{width} of the box that stores the content passed to |\item|. Then it
%   looks for the argument \enquote{\textcolor{star}{\ttfamily *}}; if it is present, it calls the
%   function \myhlc|\@@_starred_item_viii:w|, otherwise it calls the
%   function \myhlc|\@@_standard_item_viii:w|.
%    \begin{macrocode}
\cs_new_protected:Npn \@@_joined_item_viii:w (#1)
  {
    \@@_starred_joined_item_viii:n {#1}
    \peek_meaning_remove:NTF *
      { \@@_starred_item_viii:w }
      { \@@_standard_item_viii:w }
  }
%    \end{macrocode}
%   The function \myhlc|\@@_standard_item_viii:w| will first look for the
%   argument `|[|', if present it will set the state of the variable
%   \myhlc|\l_@@_wrap_label_opt_viii_bool| equal to the state of the variable
%   \myhlc|\l_@@_wrap_label_opt_viii_bool| handled by the key
%   \mykey[env]{wrap-label*} and finally execute the \emph{non-enumerated}
%   version |\item|\myoarg{custom} by means of the function
%   \myhlc|\@@_start_item_viii:w|; otherwise it sets the value of the
%   variable \myhlc|\l_@@_wrap_label_viii_bool| handled by the
%   \mykey[env]{wrap-label} key to true and set the switch |\if@noitemarg| to
%   true to execute the enumerated version of |\item| by means of the
%   function \myhlc|\@@_start_item_viii:w| |[| \myhlc|\l_@@_label_viii_tl| |]|.
%    \begin{macrocode}
\cs_new_protected:Npn \@@_standard_item_viii:w
  {
    \bool_set_false:N \l_@@_item_starred_viii_bool
    \bool_set_false:N \l_@@_item_wrap_key_bool
    \peek_meaning:NTF [
      {
        \bool_set_eq:NN \l_@@_wrap_label_viii_bool
          \l_@@_wrap_label_opt_viii_bool
        \@@_start_item_viii:w
      }
      {
        \bool_set_true:N \l_@@_wrap_label_viii_bool
        \legacy_if_set_true:n { @noitemarg }
        \@@_start_item_viii:w [ \l_@@_label_viii_tl ] \ignorespaces
      }
  }
%    \end{macrocode}
% \end{macro}
%
% \begin{macro}[int]{\@@_starred_item_viii:w, \@@_starred_item_viii_aux_i:w,
%                    \@@_starred_item_viii_aux_ii:w, \@@_keyans_starred_item_star:}
%   The function \myhlc|\@@_starred_item_viii:w| together with the specified
%   auxiliary functions \myhlc|aux_i:w| and \myhlc|aux_ii:w| execute |\item*|
%   and |\item*|\myoarg{content}.
%    \begin{macrocode}
\cs_new_protected:Npn \@@_starred_item_viii:w
  {
    \bool_set_true:N \l_@@_item_starred_viii_bool
    \bool_set_true:N \l_@@_item_wrap_key_bool
    \bool_set_true:N \l_@@_wrap_label_viii_bool
    \peek_meaning:NTF [
      { \@@_starred_item_viii_aux_i:w }
      { \@@_starred_item_viii_aux_ii:w }
  }
%    \end{macrocode}
%   The function \myhlc|\@@_starred_item_viii_aux_i:w| will save the
%   \emph{optional argument} to |\item*| in \myhlc|\l_@@_store_current_opt_arg_tl|
%   and will save this argument along with the spacing set by the key \mykey[env]{save-sep} in
%   variable \myhlc|\l_@@_store_current_label_tl| if present, then call the
%   function \myhlc|\@@_starred_item_viii_aux_ii:w|.
%    \begin{macrocode}
\cs_new_protected:Npn \@@_starred_item_viii_aux_i:w [#1]
  {
    \tl_clear:N \l_@@_store_current_label_tl
    \tl_if_novalue:nF {#1}
      {
        \tl_if_empty:NF \l_@@_store_keyans_item_opt_sep_viii_tl
          {
            \tl_put_right:NV \l_@@_store_current_label_tl
              \l_@@_store_keyans_item_opt_sep_viii_tl
            \tl_put_right:Nn \l_@@_store_current_label_tl {#1}
          }
        \tl_set:Nn \l_@@_store_current_opt_arg_tl {#1}
      }
    \@@_starred_item_viii_aux_ii:w
  }
\cs_new_protected:Npn \@@_starred_item_viii_aux_ii:w
  {
    \legacy_if_set_true:n { @noitemarg }
    \@@_start_item_viii:w [ \l_@@_label_viii_tl ] \ignorespaces
  }
%    \end{macrocode}
%   The function \myhlc|\@@_keyans_starred_item_star:| will be in charge of storing
%   the current \mymeta{label} for |\item*| followed by the
%   \myoarg{content} for |\item*|\myoarg{content} if
%   present in the \emph{sequence} and \emph{prop list} set by the
%   \mykey[env]{save-ans} key. In this same function the keys \mykey[env]{show-ans},
%   \mykey[env]{show-pos}, \mykey[env]{mark-sep} and \mykey[env]{save-ref} are implemented.
%    \begin{macrocode}
\cs_new_protected:Nn \@@_keyans_starred_item_star:
  {
    \tl_put_left:Ne \l_@@_store_current_label_tl { \l_@@_label_viii_tl }
    \@@_store_addto_prop:V \l_@@_store_current_label_tl
    \@@_keyans_store_ref:
    \tl_put_left:Nn \l_@@_store_current_label_tl { \item }
    \@@_keyans_addto_seq_link:
    \int_gincr:N \g_@@_check_starred_cmd_int
    \dim_compare:nNnT { \l_@@_mark_sym_sep_viii_dim } = { \c_zero_dim }
      {
        \dim_set:Nn \l_@@_mark_sym_sep_viii_dim { \l_@@_labelsep_viii_dim }
      }
    \bool_if:NT \l_@@_show_answer_bool
      {
        \tl_set_eq:NN \l_@@_mark_answer_sym_tl \l_@@_mark_answer_sym_viii_tl
        \str_set_eq:NN \l_@@_mark_position_str \l_@@_mark_position_viii_str
        \@@_print_keyans_box:NN
          \l_@@_labelwidth_viii_dim \l_@@_mark_sym_sep_viii_dim
      }
    \bool_if:NT \l_@@_show_position_bool
      {
        \tl_set:Ne \l_@@_mark_answer_sym_tl
          {
            \group_begin:
              \exp_not:N \normalfont
              \exp_not:N \footnotesize [ \int_eval:n
                {
                  \prop_count:c { g_@@_ \l_@@_store_name_tl _prop }
                }
              ]
            \group_end:
          }
        \str_set_eq:NN \l_@@_mark_position_str \l_@@_mark_position_viii_str
        \@@_print_keyans_box:NN
          \l_@@_labelwidth_viii_dim \l_@@_mark_sym_sep_viii_dim
      }
  }
%    \end{macrocode}
% \end{macro}
%
% \begin{macro}[int]{\@@_keyans_wrapper_label_viii:n, \@@_fake_make_label_viii:n,
%                    \@@_wrapper_label_viii:n}
%    The implementation at this point is very similar to that of the
%    \myenv*[no-index]{enumext*} environment.
%    \begin{macrocode}
\cs_new_protected:Nn \@@_keyans_wrapper_label_viii:n
  {
    \bool_lazy_all:nT
      {
        { \bool_if_p:N \l_@@_wrap_label_viii_bool         }
        { \bool_if_p:N \l_@@_show_answer_bool             }
        { \bool_if_p:N \l_@@_item_wrap_key_bool           }
        { \cs_if_exist_p:N \@@_keyans_wrapper_item_viii:n }
      }
      {
        \cs_set_eq:NN
          \@@_wrapper_label_viii:n \@@_keyans_wrapper_item_viii:n
      }
    \bool_if:NTF \l_@@_wrap_label_viii_bool
      {
        \@@_wrapper_label_viii:n {#1}
      }
      {#1}
  }
\cs_new_protected_nopar:Nn \@@_fake_make_label_viii:n
  {
    \legacy_if:nT { @noitemarg }
      {
        \legacy_if_set_false:n { @noitemarg }
        \legacy_if:nT { @nmbrlist }
          {
            \refstepcounter{enumXviii}
          }
      }
    \bool_if:NT \l_@@_item_starred_viii_bool
      {
        \@@_keyans_starred_item_star:
      }
    \makebox[ \l_@@_labelwidth_viii_dim ][ \l_@@_align_label_viii_str ]
      {
        \tl_use:N \l_@@_label_font_style_viii_tl
        \@@_keyans_wrapper_label_viii:n {#1}
      }
    \skip_horizontal:N \l_@@_labelsep_viii_dim \ignorespaces
  }
%    \end{macrocode}
% \end{macro}
%
% \subsubsection{Real definition of \cs{item} in \env{keyans*}}
%
% \begin{macro}[int]{\@@_start_item_viii:w, \@@_stop_item_viii:}
%    The implementation at this point is very similar to that of the
%    \myenv*[no-index]{enumext*} environment.
%    \begin{macrocode}
\cs_new_protected_nopar:Npn \@@_start_item_viii:w [#1]
  {
    \cs_set_eq:NN \@@_stop_item_tmp_viii: \@@_stop_item_viii:
    \hbox_set_to_wd:Nnw \l_@@_item_text_viii_box
      {
        \l_@@_joined_width_viii_dim
        + \l_@@_labelwidth_viii_dim
        + \l_@@_labelsep_viii_dim
      }
      \@@_renew_footnote_starred:
      \@@_start_list_tag:n {keyans*}
      \@@_fake_make_label_viii:n {#1}
      \@@_stop_start_list_tag:
      \@@_minipage:w [ t ]{ \l_@@_joined_width_viii_dim  }
        \dim_set_eq:NN \parindent \l_@@_listparindent_viii_dim
        \skip_set_eq:NN \parskip \l_@@_parsep_viii_skip
        \@@_unskip_unkern:
        \@@_unskip_unkern:
        \skip_horizontal:n { -\l_@@_listparindent_viii_dim } \ignorespaces
        \tl_use:N \l_@@_fake_item_indent_viii_tl
        \bool_if:NT \l_@@_item_starred_viii_bool
          {
            \@@_keyans_show_item_opt_viii:
          }
        \tl_use:N \l_@@_after_list_args_viii_tl
  }
\cs_new_protected_nopar:Nn \@@_stop_item_viii:
  {
      \@@_endminipage:
    \@@_stop_list_tag:n {keyans*}
    \hbox_set_end:
    \int_set:Nn \hbadness { 10000 }
    \box_use_drop:N \l_@@_item_text_viii_box
    \int_compare:nNnTF
      { \l_@@_item_column_pos_viii_int } = { \l_@@_columns_viii_int }
      {
        \par\noindent
        \int_zero:N \l_@@_item_column_pos_viii_int
      }
      {
        \skip_horizontal:N \l_@@_columns_sep_viii_dim
      }
  }
%    \end{macrocode}
% \end{macro}
%
% \begin{macro}[int]{\@@_remove_extra_parsep_viii:}
%    The implementation at this point is very similar to that of the
%    \myenv*[no-index]{enumext*} environment.
%    \begin{macrocode}
\cs_new_protected:Nn \@@_remove_extra_parsep_viii:
  {
    \int_compare:nNnT
      {
        \int_mod:nn
          { \g_@@_item_count_all_viii_int }
          { \l_@@_columns_viii_int }
      }
      =
      { 0 }
      {
        \para_end:
        \skip_vertical:n { -\l_@@_itemsep_viii_skip }
        \skip_vertical:N \c_zero_skip
        \int_gzero:N \g_@@_item_count_all_viii_int
      }
  }
%    \end{macrocode}
% \end{macro}
%
% \subsection{The command \cs{getkeyans}}\label{cmd:getkeyans}
%
% \subsubsection{Definition of error messages for \cs[no-index]{getkeyans}}
%
% \begin{macro}[int]{missing-colon}
% Message used by |\getkeyans| command.
%    \begin{macrocode}
\msg_new:nnnn { enumext } { missing-colon }
  {
    The~argument~'#1'~has~no~colon~\msg_line_context:.
  }
  {
    The~argument~of~'\token_to_str:N \getkeyans'~must~have~the~form~
    'store~name:position'.
  }
%    \end{macrocode}
% \end{macro}
%
% \subsubsection{Internal functions for \cs[no-index]{getkeyans}}
%
% \begin{macro}[int]{\@@_getkeyans_aux:n, \@@_getkeyans:nn}
%   The internal function \myhlc|\@@_getkeyans_aux:n| is in charge of
%   \emph{splitting} the \emph{mandatory argument} using \enquote{\textcolor{red}{:}}.
%   If \enquote{\textcolor{red}{:}} is omitted it will return an error.
%    \begin{macrocode}
\cs_new_protected:Nn \@@_getkeyans_aux:n
  {
    \str_if_in:nnTF {#1} { : }
      {
        \use:e
          {
            \cs_set:Npn \exp_not:N \@@_tmp:w ##1 \c_colon_str ##2 \scan_stop:
              { {##1} {##2} }
          }
        \exp_after:w \@@_getkeyans:nn \@@_tmp:w #1 \scan_stop:
      }
      { \msg_error:nnn { enumext } { missing-colon } {#1} }
  }
%    \end{macrocode}
%   The internal function \myhlc|\@@_getkeyans:nn| will check for the existence
%   of the \emph{prop list}, if it does not exist it will return an error
%   message, then it will fetch the content specified by the \emph{second argument}
%   from \emph{prop list}.
%    \begin{macrocode}
\cs_new_protected:Nn \@@_getkeyans:nn
  {
    \prop_if_exist:cTF { g_@@_ #1 _prop }
      {
        \prop_item:cn { g_@@_ #1 _prop }{#2}
      }
      {
        \msg_error:nnn { enumext } { undefined-storage-anskey } {#1}
      }
  }
%    \end{macrocode}
% \end{macro}
%
% \begin{macro}{\getkeyans}
%   The \ics*{getkeyans} command takes a \emph{mandatory argument} of the
%   form \mymarg{store name \textcolor{red}{:} position}.
%   Retrieve a \emph{\enquote{single content}}  stored by |\anskey|, |\anspic*|
%   and |\item*| and \myenv*[no-index]{anskey*} from \emph{prop list} defined by \mykey[env]{save-ans} key.
%    \begin{macrocode}
\NewDocumentCommand \getkeyans { m }
  {
    \exp_args:Ne \@@_getkeyans_aux:n
      { \tl_to_str:e { \text_expand:n {#1} } }
  }
%    \end{macrocode}
% \end{macro}
%
% \subsection{The command \cs{printkeyans}}\label{cmd:printkeyans}
%
% The |\printkeyans| command prints \emph{\enquote{all stored
% content}} in the \emph{sequence} defined by the \mykey{save-ans} key.
%
% \subsubsection{Definition of variables for \cs[no-index]{printkeyans}}
%
% \begin{variable}[int]{\c_@@_print_envs_clist, \c_@@_default_vars_clist,
%                       \g_@@_default_NAME_X_TYPE, \c_@@_default_label_X_tl}
%   The command \ics*{printkeyans} must not depend on the values that
%   \ics*{setenumext} sets for the environments, so the default values of the
%   package are saved when it is loaded. The constant
%   \myhlc|\c_@@_print_envs_clist| lists the environments that
%   \ics*{printkeyans} can open and the constant
%   \myhlc|\c_@@_default_vars_clist| lists, as pairs \mymeta{name}\mymeta{type},
%   the variables of each environment whose default value is saved in
%   \myhlc|\g_@@_default_NAME_X_TYPE|. The default \mymeta{label} is saved by the
%   \mykey[env]{label} key in \myhlc|\c_@@_default_label_X_tl|.
%    \begin{macrocode}
\clist_const:Nn \c_@@_print_envs_clist
  {
    {level-1}{i}, {level-2}{ii}, {level-3}{iii}, {level-4}{iv}, {enumext*}{vii}
  }
\clist_const:Nn \c_@@_default_vars_clist
  {
    {labelsep}{dim}, {labelwidth}{dim}, {fake_item_indent}{dim},
    {listoffset}{dim}, {rightmargin}{dim}, {listparindent}{dim},
    {leftmargin_tmp}{dim}, {leftmargin_tmp}{bool}, {columns_sep}{dim},
    {minipage_hsep}{dim}, {topsep}{skip}, {partopsep}{skip}, {parsep}{skip},
    {itemsep}{skip}, {columns}{int}, {start}{int}, {beginparpenalty}{int},
    {itempenalty}{int}, {endparpenalty}{int}, {label_fill_left}{tl},
    {label_fill_right}{tl}, {align_label_pos}{str}, {label_font_style}{tl},
    {after_list_args}{tl}, {wrap_label_opt}{bool}
  }
%    \end{macrocode}
% \end{variable}
%
% \subsubsection{Definition of error messages for \cs[no-index]{printkeyans}}
%
% \begin{macro}[int]{print-starred}
%  Messages used by |\printkeyans*| command.
%    \begin{macrocode}
\msg_new:nnn { enumext } { print-starred }
  {
    '\token_to_str:N \printkeyans*':~The~sequence~'#1'~already~contains~
    #2~environment~\msg_line_context:.
  }
%    \end{macrocode}
% \end{macro}
%
% \subsubsection{Definition of keys for \cs[no-index]{printkeyans}}
%
% \begin{macro}[int]{print*, print-*, print-1, print-2, print-3, print-4}
% First a set of \mymeta{filtered keys} is defined, with
% which the options of the different nesting levels are controlled for
% the environments \myenv*[no-index]{enumext} and \myenv*[no-index]{enumext*} by storing their
% values in the list of tokens \myhlc|\l_@@_print_keyans_X_tl|.
%
% The default values only contain what the printed content changes with respect
% to the default values of the package: \mykey[env]{nosep}, \mykey[env]{first} and
% \mykey[env]{font}, plus \mykey[env]{columns} for the \emph{\enquote{first level}}.
% The remaining values are restored by the function \myhlc|\@@_default_restore:|
% before these \mymeta{keys} are run (\S\ref{cmd:printkeyans}).
%
% The variable \myhlc|\l_@@_print_keyans_starred_tl| will have the default \mymeta{keys}
% for \ics*{printkeyans*} and will be set by |\setenumext|\myoarg{print\textnormal{\textcolor{star}{\texttt{*}}}}
% and the variable \myhlc|\l_@@_print_keyans_vii_tl| will have the default
% keys for the environment \myenv*[no-index]{enumext*} nested within the \emph{sequence}
% and will be set by |\setenumext|\myoarg{print \textnormal{\textcolor{gray}{,}}\textnormal{\textcolor{star}{\texttt{*}}}},
% the rest of the variables will be for the environment \myenv*[no-index]{enumext} and will be set
% by |\setenumext|\myoarg{print \textnormal{\textcolor{gray}{,}} level}.
%
%    \begin{macrocode}
\keys_define:nn  { enumext / print }
  {
    print*  .code:n =
      {
        \keys_precompile:neN { enumext / enumext* }
          { \@@_filter_save_key:n {#1} }
          \l_@@_print_keyans_starred_tl
      }, % starred cmd
    print*  .initial:n =
      {
        nosep,
        columns = 2,
        first   = \small,
        font    = \small
      },
    print-1 .code:n =
      {
        \keys_precompile:neN { enumext / level-1 }
          { \@@_filter_save_key:n {#1} }
          \l_@@_print_keyans_i_tl
      },
    print-1 .initial:n =
      {
        nosep,
        columns = 2,
        first   = \small,
        font    = \small
      },
    print-2 .code:n =
      {
        \keys_precompile:neN { enumext / level-2 }
          { \@@_filter_save_key:n {#1} }
          \l_@@_print_keyans_ii_tl
      },
    print-2 .initial:n =
      {
        nosep,
        first = \small,
        font  = \small
      },
    print-3 .code:n =
      {
        \keys_precompile:neN { enumext / level-3 }
          { \@@_filter_save_key:n {#1} }
          \l_@@_print_keyans_iii_tl
      },
    print-3 .initial:n =
      {
        nosep,
        first = \small,
        font  = \small
      },
    print-4 .code:n =
      {
        \keys_precompile:neN { enumext / level-4 }
          { \@@_filter_save_key:n {#1} }
          \l_@@_print_keyans_iv_tl
      },
    print-4 .initial:n =
      {
        nosep,
        first = \small,
        font  = \small
      },
    print-* .code:n =
      {
        \keys_precompile:neN { enumext / enumext* }
          { \@@_filter_save_key:n {#1} }
          \l_@@_print_keyans_vii_tl
      }, % starred nested
    print-* .initial:n =
      {
        nosep,
        first = \small,
        font  = \small
      },
  }
%    \end{macrocode}
% \end{macro}
%
% \subsubsection{Internal functions for \cs[no-index]{printkeyans}}
%
% \begin{macro}[int]{\@@_default_save:nnn, \@@_default_restore:nnn,
%                    \@@_default_save_env:nn, \@@_default_restore_env:nn,
%                    \@@_default_restore:}
%   The function \myhlc|\@@_default_save:nnn| saves in
%   \myhlc|\g_@@_default_NAME_X_TYPE| the value of the variable
%   \myhlc|\l_@@_NAME_X_TYPE| if it exists and the function
%   \myhlc|\@@_default_restore:nnn| copies it back.
%    \begin{macrocode}
\cs_new_protected:Nn \@@_default_save:nnn
  {
    \cs_if_exist:cT { l_@@_ #1 _ #3 _ #2 }
      {
        \use:c { #2 _new:c } { g_@@_default_ #1 _ #3 _ #2 }
        \use:c { #2 _gset_eq:cc }
          { g_@@_default_ #1 _ #3 _ #2 } { l_@@_ #1 _ #3 _ #2 }
      }
  }
\cs_new_protected:Nn \@@_default_restore:nnn
  {
    \cs_if_exist:cT { g_@@_default_ #1 _ #3 _ #2 }
      {
        \use:c { #2 _set_eq:cc }
          { l_@@_ #1 _ #3 _ #2 } { g_@@_default_ #1 _ #3 _ #2 }
      }
  }
%    \end{macrocode}
%   The function \myhlc|\@@_default_save_env:nn| saves all the variables of
%   \myhlc|\c_@@_default_vars_clist| for one environment, together with the
%   function of the \mykey[env]{wrap-label} key. The function
%   \myhlc|\@@_default_restore_env:nn| restores them and runs the
%   \mykey[env]{label} key with the saved \mymeta{label}, since this key also
%   defines the counter style and the width of the label.
%    \begin{macrocode}
\cs_new_protected:Nn \@@_default_save_env:nn
  {
    \clist_map_inline:Nn \c_@@_default_vars_clist
      { \@@_default_save:nnn ##1 {#2} }
    \cs_gset_eq:cc { @@_default_wrapper_label_ #2 :n }
      { @@_wrapper_label_ #2 :n }
  }
\cs_new_protected:Nn \@@_default_restore_env:nn
  {
    \clist_map_inline:Nn \c_@@_default_vars_clist
      { \@@_default_restore:nnn ##1 {#2} }
    \cs_set_eq:cc { @@_wrapper_label_ #2 :n }
      { @@_default_wrapper_label_ #2 :n }
    \keys_set:ne { enumext / #1 }
      { label = { \exp_not:v { c_@@_default_label_ #2 _tl } } }
  }
%    \end{macrocode}
%   Finally the function \myhlc|\@@_default_restore:| restores the default
%   values of all the environments and the values are saved when the package
%   is loaded.
%    \begin{macrocode}
\cs_new_protected:Nn \@@_default_restore:
  {
    \clist_map_inline:Nn \c_@@_print_envs_clist
      { \@@_default_restore_env:nn ##1 }
  }
\clist_map_inline:Nn \c_@@_print_envs_clist
  { \@@_default_save_env:nn #1 }
%    \end{macrocode}
% \end{macro}
%
% \begin{macro}[int]{\@@_printkeyans:nnn}
%   The internal function \myhlc|\@@_printkeyans:nnn| will check
%   for the existence of the \emph{sequence}, if it does not exist
%   it will return an error message, then it will check if not empty.
%    \begin{macrocode}
\cs_new_protected:Nn \@@_printkeyans:nnn
  {
    \seq_if_exist:cTF { g_@@_ #3 _seq }
      {
        \seq_if_empty:cF { g_@@_ #3 _seq }
          {
%    \end{macrocode}
%   If the \emph{starred argument} `|*|' is present, it checks that
%   the environment \myenv*[no-index]{enumext*} is not saved in the \emph{sequence},
%   then executes the variable \myhlc|\l_@@_print_keyans_starred_tl| that
%   contains the default \mymeta{keys} for the environment \myenv*[no-index]{enumext*},
%   sets \myhlc|\l_@@_base_line_fix_bool| and \myhlc|\l_@@_print_keyans_star_bool|
%   to true for \emph{baseline correction}, opens the \myenv*[no-index]{enumext*}
%   environment passing the \emph{optional argument} and maps the \emph{sequence},
%   then sets \myhlc|\l_@@_base_line_fix_bool| and \myhlc|\l_@@_print_keyans_star_bool|
%   to false.
%    \begin{macrocode}
            \bool_if:nTF {#1}
              {
                \seq_if_in:cnTF { g_@@_ #3 _seq } { \end{enumext*} }
                  {
                    \msg_error:nnnn { enumext } { print-starred } {#3}
                      { enumext* }
                  }
                  {
                    \tl_use:N \l_@@_print_keyans_starred_tl
                    \bool_set_true:N \l_@@_base_line_fix_bool
                    \bool_set_true:N \l_@@_print_keyans_star_bool
                    \begin{enumext*}[#2]
                      \seq_map_inline:cn { g_@@_ #3 _seq } {##1}
                    \end{enumext*}
                    \bool_set_false:N \l_@@_base_line_fix_bool
                    \bool_set_false:N \l_@@_print_keyans_star_bool
                  }
              }
%    \end{macrocode}
%   Otherwise it will open the environment \myenv*[no-index]{enumext} passing
%   the \emph{optional argument} to the \emph{\enquote{first level}}
%   then map the \emph{sequence}.
%    \begin{macrocode}
              {
                \begin{enumext}[#2]
                  \seq_map_inline:cn { g_@@_ #3 _seq } {##1}
                \end{enumext}
              }
          }
      }
      {
        \msg_error:nnn { enumext } { undefined-storage-anskey } {#3}
      }
  }
%    \end{macrocode}
% \end{macro}
%
% \begin{macro}{\printkeyans}
%   Create a user command to print \emph{\enquote{all stored content}}
%   in \emph{sequence} for |\anskey|, \myenv*[no-index]{anskey*}, |\item*| and |\anspic*|.
%   Within a group the default values of the package are restored with
%   \myhlc|\@@_default_restore:|, then the \emph{\enquote{precompiled keys}} are
%   run and finally the internal function \myhlc|\@@_printkeyans:nnn| is called.
%    \begin{macrocode}
\NewDocumentCommand \printkeyans { s O{} m }
  {
    \group_begin:
      \@@_default_restore:
      \bool_set_true:N \l_@@_print_keyans_cmd_bool
      \tl_use:N \l_@@_print_keyans_i_tl
      \tl_use:N \l_@@_print_keyans_ii_tl
      \tl_use:N \l_@@_print_keyans_iii_tl
      \tl_use:N \l_@@_print_keyans_iv_tl
      \tl_use:N \l_@@_print_keyans_vii_tl
      \@@_printkeyans:nnn {#1} {#2} {#3}
      \bool_set_false:N \l_@@_print_keyans_cmd_bool
    \group_end:
  }
%    \end{macrocode}
% \end{macro}
%
% \begin{important}*
%  The reason for storing \mymeta{keys} in token lists using
%  \cs{keys_precompile:neN} is because the keys are set via
%  |\setenumext| but are later executed by running the command
%  |\printkeyans| and they are not handled directly by its
%  \emph{optional argument}, except those related to the \emph{first}
%  opening level.
% \end{important}
%
% \subsection{The command \cs{setenumext}}\label{cmd:setenumext}
%
% The command |\setenumext| is in charge of managing the
% \mymeta{keys} passed to all environments and to the |\printkeyans|
% command. Precautions must be taken with the \myenv*[no-index]{enumext*} and \myenv*[no-index]{enumext}
% environments so as not to capture \mymeta{keys} that cause problems.
%
% \subsubsection{Definition of variables for \cs[no-index]{setenumext}}
%
% \begin{variable}[int]{\c_@@_all_families_seq}
%    They are stored in the constant sequence \myhlc|\c_@@_all_families_seq|
%    separated by commas.
%    \begin{macrocode}
\seq_const_from_clist:Nn \c_@@_all_families_seq
  {
    enumext-1, enumext-2, enumext-3, enumext-4, keyans, enumext*,
    keyans*, print-1, print-2, print-3, print-4, print-*, print*,
  }
%    \end{macrocode}
% \end{variable}
%
% \subsubsection{Definition of error messages for \cs[no-index]{setenumext}}
%
% \begin{macro}[int]{invalid-key, unknown-key-family}
% Message used by \myoarg{key \textnormal{\textcolor{gray}{=}} val}
% system and |\setenumext| command.
%    \begin{macrocode}
\msg_new:nnn { enumext } { invalid-key }
  {
    The~key~'#1'~is~not~known~at~level~#2~\msg_line_context:.
  }
\msg_new:nnn { enumext } { unknown-key-family }
  {
    Unknown~key~family~'\l_keys_key_str'~for~enumext~\msg_line_context:.
  }
%    \end{macrocode}
% \end{macro}
%
% \subsubsection{Definition of keys for \cs[no-index]{setenumext}}
%
% \begin{macro}[int]{enumext-1, enumext-2, enumext-3, enumext-4, enumext*,
%                    keyans, keyans*, print*, print-1, print-2, print-3,
%                    print-4, print-*, unknown}
% Now define a \emph{\enquote{meta families}} of \mymeta{keys}
% to access from |\setenumext|.
%
%    \begin{macrocode}
\keys_define:nn { enumext / meta-families }
  {
    enumext-1 .code:n =
      {
        \keys_set:ne { enumext / level-1 }
          {
            \@@_filter_level:n {#1}
          }
      },
    enumext-2 .code:n =
      {
        \keys_set:ne { enumext / level-2 }
          {
            \@@_filter_level:n {#1}
          }
      },
    enumext-3 .code:n =
      {
        \keys_set:ne { enumext / level-3 }
          {
            \@@_filter_level:n {#1}
          }
      },
    enumext-4 .code:n =
      {
        \keys_set:ne { enumext / level-4 }
          {
            \@@_filter_level:n {#1}
          }
      },
    enumext*  .code:n =
      {
        \keys_set:ne { enumext / enumext* }
          {
            \@@_filter_level:n {#1}
          }
      },
    keyans    .meta:nn = { enumext / keyans } {#1},
    keyans*   .meta:nn = { enumext / keyans* } {#1},
    print*    .meta:nn = { enumext / print } { print* = {#1} },
    print-1   .meta:nn = { enumext / print } { print-1 = {#1} },
    print-2   .meta:nn = { enumext / print } { print-2 = {#1} },
    print-3   .meta:nn = { enumext / print } { print-3 = {#1} },
    print-4   .meta:nn = { enumext / print } { print-4 = {#1} },
    print-*   .meta:nn = { enumext / print } { print-* = {#1} },
    unknown   .code:n = { \msg_error:nn { enumext } { unknown-key-family } },
  }
%    \end{macrocode}
% \end{macro}
%
% \subsubsection{Internal functions for \cs[no-index]{setenumext}}
%
% \begin{macro}[int]{\@@_filter_level:n, \@@_filter_level_key:n,
%                    \@@_filter_level_pair:nn}
%   The function \myhlc|\@@_filter_level:n| will be in charge of
%   filtering the \mymeta{keys} passed to the \myenv*[no-index]{enumext}
%   and \myenv*[no-index]{enumext*} environments.
%    \begin{macrocode}
\cs_new:Nn \@@_filter_level:n
  {
    \use:e
      {
        \keyval_parse:NNn
          \@@_filter_level_key:n
          \@@_filter_level_pair:nn {#1}
      }
  }
%    \end{macrocode}
%   The function \myhlc|\@@_filter_level_key:n| will be responsible for
%   filtering the \mymeta{keys} that are passed \emph{\enquote{without value}} by
%   excluding the keys \mykey[env]{resume*}, \mykey[env]{reset} and \mykey[env]{reset*}
%   passed to the \myenv*[no-index]{enumext} and \myenv*[no-index]{enumext*} environments.
%    \begin{macrocode}
\cs_new:Nn \@@_filter_level_key:n
  {
    \str_case:nnF {#1}
      {
        { resume* } { } { reset } { } { reset* } { }
      }
      { , { \exp_not:n {#1} } }
  }
%    \end{macrocode}
%   The function \myhlc|\@@_filter_level_pair:nn| will be responsible for
%   filtering the \mymeta{keys} that are passed \emph{\enquote{with value}} by
%   excluding the \mykey[env]{series}, \mykey[env]{resume} and \mykey[env]{save-ans} keys
%   passed to the \myenv*[no-index]{enumext} and \myenv*[no-index]{enumext*} environments.
%    \begin{macrocode}
\cs_new:Nn \@@_filter_level_pair:nn
  {
    \str_case:nnF {#1}
      {
        { series } { }  { save-ans } { } { resume } { }
      }
      { , { \exp_not:n {#1} } = { \exp_not:n {#2} } }
  }
%    \end{macrocode}
% \end{macro}
%
% \begin{macro}[int]{\@@_set_parse:n, \@@_set_error:nn}
%   Internal functions used by the |\setenumext| command.
%    \begin{macrocode}
\cs_new_protected:Nn \@@_set_parse:n
  {
    \tl_set:Ne \l_@@_setkey_tmpb_tl { \tl_trim_spaces:n {#1} }
    \clist_map_inline:nn { 0, 1, 2, 3, 4, * } % <- max level
      { \tl_remove_all:Nn \l_@@_setkey_tmpb_tl {##1} }
    \tl_if_empty:NTF \l_@@_setkey_tmpb_tl
      {
        \seq_put_right:Ne \l_@@_setkey_tmpa_seq
          { \tl_trim_spaces:n {#1} }
      }
      { \@@_set_error:nn {#1} { } }
  }
\cs_new_protected:Nn \@@_set_error:nn
  { \msg_error:nnnn { enumext } { invalid-key } {#1} {#2} }
%    \end{macrocode}
% \end{macro}
%
% \begin{macro}{\setenumext}
%    Defines the user command \ics*{setenumext}.
%    \begin{macrocode}
\NewDocumentCommand \setenumext { O{enumext,1} +m }
  {
    \seq_clear:N \l_@@_setkey_tmpa_seq
    \seq_set_from_clist:Nn \l_@@_setkey_tmpb_seq {#1}
    \int_set:Nn \l_@@_setkey_tmpa_int
      {
        \seq_count:N \l_@@_setkey_tmpb_seq
      }
    \int_compare:nNnTF { \l_@@_setkey_tmpa_int } > { 1 }
      {
        \seq_pop_left:NN \l_@@_setkey_tmpb_seq \l_@@_setkey_tmpa_tl
        \seq_map_function:NN \l_@@_setkey_tmpb_seq \@@_set_parse:n
        \seq_set_map_e:NNn \l_@@_setkey_tmpa_seq \l_@@_setkey_tmpa_seq
          {
            \tl_use:N \l_@@_setkey_tmpa_tl - ##1
          }
      }
      {
        \seq_put_right:Ne \l_@@_setkey_tmpa_seq { \tl_trim_spaces:n {#1} }
      }
    \seq_if_empty:NTF \l_@@_setkey_tmpa_seq
      { \seq_map_inline:Nn \c_@@_all_families_seq }
      { \seq_map_inline:Nn \l_@@_setkey_tmpa_seq }
      {
        \keys_set:nn { enumext / meta-families } { ##1 = {#2} }
      }
  }
%    \end{macrocode}
% \end{macro}
%
% \subsection{The command \cs{setenumextmeta}}\label{cmd:setenumextmeta}
%
% The command \ics*{setenumextmeta} will be responsible for adding new
% \emph{\enquote{meta-keys}} for the \myenv*[no-index]{enumext} and \myenv*[no-index]{enumext*} environments.
% The implementation code was given by Jonathan P. Spratte (\texttt{@Skillmon})
% answer in \href{https://tex.stackexchange.com/a/747159}{Simplify syntax for command that adds \texttt{.meta} key to existing keys (l3keys)}.
%
% \subsubsection{Definition of error messages for \cs[no-index]{setenumextmeta}}
%
% \begin{macro}[int]{already-defined, prohibited-unknown}
% Messages used by |\setenumextmeta| command.
%    \begin{macrocode}
\msg_new:nnn { enumext } { already-defined }
  {
    The~key~'#1'~is~already~defined~\msg_line_context:.
  }
\msg_new:nnn { enumext } { prohibited-unknown }
  {
    The~name~'unknown'~can't~be~chosen~for~a~meta~key~\msg_line_context:.
  }
%    \end{macrocode}
% \end{macro}
%
% \subsubsection{Definition of keys for \cs[no-index]{setenumextmeta}}
%
% \begin{macro}[int]{1, 2, 3, 4, enumext, enumext*}
%    First, \mymeta{keys} of type \myhlc|.code:n| are created for \emph{\enquote{all levels}}
%    of the \myenv*[no-index]{enumext} environment.
%    \begin{macrocode}
\int_step_inline:nn { 4 }
  {
    \keys_define:nn { enumext }
      {
        #1 .code:n = { \str_set:Nn \l_@@_meta_path_str { level-#1 } },
        #1 .value_forbidden:n = true,
      }
  }
%    \end{macrocode}
%   And now the \mymeta{keys} are defined for the environments using
%   \myhlc|.code:n| for the \myenv*[no-index]{enumext} environment and \myhlc|.meta:n| for the
%   \myenv*[no-index]{enumext*} environment.
%    \begin{macrocode}
\clist_map_inline:nn { enumext }
  {
    \keys_define:nn { enumext }
      {
        % ignored for now, might do something useful in the future
        #1  .code:n = { },
        #1  .value_forbidden:n = true,
        #1* .code:n = { \str_set:Nn \l_@@_meta_path_str { #1* } },
        #1* .value_forbidden:n = true,
      }
  }
\keys_define:nn { enumext }
  {
    * .meta:n = enumext*,
    * .value_forbidden:n = true,
  }
%    \end{macrocode}
% \end{macro}
%
% \subsubsection{Internal functions for \cs[no-index]{setenumextmeta}}
%
% \begin{macro}[int]{\@@_key_set_meta:nnn, \@@_key_def_meta:nnn,
%                    \@@_key_def_meta:Vnn}
%   The internal functions \myhlc|\@@_key_set_meta:nnn| and \myhlc|\@@_key_def_meta:nnn|
%   will check the \emph{optional argument} and create the \emph{\enquote{meta-key}}.
%    \begin{macrocode}
\cs_new_protected:Npn \@@_key_set_meta:nnn #1
  {
    \keys_set:nn { enumext } {#1}
    \@@_key_def_meta:Vnn \l_@@_meta_path_str
  }
\cs_new_protected:Nn \@@_key_def_meta:nnn
  {
    \bool_lazy_or:nnTF
      { \keys_if_exist_p:nn { enumext / #1} {#2} }
      { \keys_if_exist_p:nn { enumext / enumext* } {#2} }
      { \msg_error:nnn { enumext } { already-defined } {#2} }
      {
        \keys_define:nn { enumext / #1 }
          {
            #2 .meta:n = {#3},
            #2 .value_forbidden:n = true
          }
      }
  }
\cs_generate_variant:Nn \@@_key_def_meta:nnn { V }
%    \end{macrocode}
% \end{macro}
%
% \begin{macro}{\setenumextmeta}
%   Now the user command is created, taking care that \texttt{unknown}
%   cannot be passed as an argument.
%    \begin{macrocode}
\NewDocumentCommand \setenumextmeta { s O{enumext,1} m +m }
  {
    \str_if_eq:eeTF { \tl_trim_spaces:n {#3} } { unknown }
      { \msg_error:nn { enumext } { prohibited-unknown } }
      {
        \bool_if:nTF {#1}
          {
            \int_step_inline:nn { 4 }
              { \@@_key_set_meta:nnn { enumext, ##1 } {#3} {#4} }
            \@@_key_set_meta:nnn { enumext* } {#3} {#4}
          }
          { \@@_key_set_meta:nnn {#2} {#3} {#4} }
      }
  }
%    \end{macrocode}
% \end{macro}
%
% \subsection{The command \cs{foreachkeyans}}\label{cmd:foreachkeyans}
%
% The command \ics*{foreachkeyans} will execute a \emph{loop} over the \emph{prop list} and return its contents.
% The implementation code is adapted from the answer provided by Enrico Gregorio (\texttt{@egreg}) in
% \href{https://tex.stackexchange.com/a/721130}{Expand a \texttt{.cs} defined by key inside the function}.
%
% \subsubsection{Definition of error messages for \cs[no-index]{foreachkeyans}}
%
% \begin{macro}[int]{for-key-unknown, for-key-value-unknown}
% Message used by unknown \mymeta{keys} in |\foreachkeyans|.
%    \begin{macrocode}
\msg_new:nnnn { enumext } { for-key-unknown }
  {
    The~key~'#1'~is~unknown~by~'\token_to_str:N \foreachkeyans'~and~is~being~
    ignored~\msg_line_context:.
  }
  {
    The~command~'\token_to_str:N \foreachkeyans'~does~not~have~a~key~called~
    '#1'.~Check~that~you~have~spelled~the~key~name~correctly.
  }
\msg_new:nnnn { enumext } { for-key-value-unknown }
  {
    The~key~'#1=#2'~is~unknown~by~'\token_to_str:N \foreachkeyans'~and~is~
    being~ignored~\msg_line_context:.
  }
  {
    The~command~'\token_to_str:N \foreachkeyans'~does~not~have~a~key~called~
    '#1'.~Check~that~you~have~spelled~the~key~name~correctly.
  }
%    \end{macrocode}
% \end{macro}
%
% \subsubsection{Definition of keys for \cs[no-index]{foreachkeyans}}
%
% \begin{macro}[int]{before, after, sep, start, step, stop, wrapper, unknown}
%   A set of \mymeta{keys} is defined for the command and the
%   default values of these are saved in \myhlc|\g_@@_foreach_default_keys_tl| to avoid
%   the use of a group.
%    \begin{macrocode}
\keys_define:nn { enumext / foreach }
  {
    before  .tl_set:N = \l_@@_foreach_before_tl,
    before  .value_required:n = true,
    after   .tl_set:N = \l_@@_foreach_after_tl,
    after   .value_required:n = true,
    start   .int_set:N = \l_@@_foreach_start_int,
    start   .value_required:n = true,
    stop    .int_set:N = \l_@@_foreach_stop_int,
    stop    .value_required:n = true,
    step    .int_set:N = \l_@@_foreach_step_int,
    step    .value_required:n = true,
    wrapper .cs_set_protected:Np = \@@_foreach_wrapper:n #1,
    wrapper .value_required:n = true,
    sep     .tl_set:N = \l_@@_foreach_sep_tl,
    sep     .value_required:n = true,
    unknown .code:n = { \@@_parse_foreach_keys:n {#1} }
  }
\keys_precompile:nnN { enumext / foreach }
  {
    before={},after={},start=1,step=1,stop=0,wrapper=#1,sep={; }
  }
  \l_@@_foreach_default_keys_tl
%    \end{macrocode}
% \end{macro}
%
% \subsubsection{Internal functions for \cs[no-index]{foreachkeyans}}
%
% \begin{macro}[int]{\@@_parse_foreach_keys:nn, \@@_parse_foreach_keys:n}
%   Functions for handling unknown \mymeta{keys}.
%    \begin{macrocode}
\cs_new_protected:Nn \@@_parse_foreach_keys:nn
  {
    \tl_if_blank:nTF {#2}
      {
        \msg_error:nnn { enumext } { for-key-unknown } {#1}
      }
      {
        \msg_error:nnnn { enumext } { for-key-value-unknown } {#1} {#2}
      }
  }
\cs_new_protected:Nn \@@_parse_foreach_keys:n
  {
    \exp_args:NV \@@_parse_foreach_keys:nn \l_keys_key_str {#1}
  }
%    \end{macrocode}
% \end{macro}
%
% \begin{macro}[int]{\@@_foreach_keyans:nn, \@@_foreach_add_body:n}
%   The internal functions \myhlc|\@@_foreach_keyans:nn| and
%   \myhlc|\@@_foreach_add_body:n| will loop through the prop list and print
%   the contents.
%    \begin{macrocode}
\cs_new_protected:Nn \@@_foreach_keyans:nn
  {
    \tl_use:N \l_@@_foreach_default_keys_tl
    \keys_set:nn { enumext / foreach } {#1}
    \tl_set:Nn \l_@@_foreach_name_prop_tl {#2}
    \prop_if_exist:cF { g_@@_ #2 _prop }
      {
        \msg_error:nnn { enumext } { undefined-storage-anskey } {#2}
      }
    \int_compare:nNnT { \l_@@_foreach_stop_int } = { 0 }
      {
        \int_set:Nn \l_@@_foreach_stop_int
          { \prop_count:c { g_@@_ #2 _prop } }
      }
    \seq_clear:N \l_@@_foreach_print_seq
    \int_step_function:nnnN
      { \l_@@_foreach_start_int }
      { \l_@@_foreach_step_int }
      { \l_@@_foreach_stop_int }
      \@@_foreach_add_body:n
      \seq_use:NV \l_@@_foreach_print_seq \l_@@_foreach_sep_tl
  }
\cs_new_protected:Nn \@@_foreach_add_body:n
  {
    \seq_put_right:Ne \l_@@_foreach_print_seq
      {
        \exp_not:V \l_@@_foreach_before_tl
        \@@_foreach_wrapper:n
          {
            \prop_item:cn { g_@@_ \l_@@_foreach_name_prop_tl _prop }{#1}
          }
        \exp_not:V \l_@@_foreach_after_tl
      }
  }
%    \end{macrocode}
% \end{macro}
%
% \begin{macro}{\foreachkeyans}
%   The command is created.
%    \begin{macrocode}
\NewDocumentCommand \foreachkeyans { +O{} m }
  {
    \@@_foreach_keyans:nn {#1} {#2}
  }
%    \end{macrocode}
% \end{macro}
%
% \subsection{General messages}\label{subsec:messages}
%
% Messages shared by several sections of the implementation; the messages used by only
% one command, key or environment are defined in its own section.
%
% \begin{macro}[int]{package-load, package-not-load}
% Message used by package-load for \mypkg[no-index]{multicol} and \mypkg[no-index]{hyperref}
% packages.
%    \begin{macrocode}
\msg_new:nnn { enumext } { package-load }
  {
    The~'#1'~package~is~already~loaded.
  }
\msg_new:nnn { enumext } { package-not-load }
  {
    The~'#1'~package~will~be~loaded~as~a~dependency.
  }
%    \end{macrocode}
% \end{macro}
%
% \begin{macro}[int]{unknown-choice}
% Message used by \mykey[env]{align}, \mykey[env]{mark-pos} and \mykey[env]{label-pos} keys.
%    \begin{macrocode}
\msg_new:nnn { enumext } { unknown-choice }
  {
    Invalid~value~'#3'~for~key~'#1';~use~one~of:~#2~\msg_line_context:.
  }
%    \end{macrocode}
% \end{macro}
%
% \begin{macro}[int]{command-wrong-place}
% Message used by |\anskey| and |\anspic| commands when they are used
% in the wrong place.
%    \begin{macrocode}
\msg_new:nnnn { enumext } { command-wrong-place }
  {
    Wrong~place~for~command~'\c_backslash_str #1'~\msg_line_context:.
  }
  {
    '\c_backslash_str #1'~works~outside~the~environment~'#2'.
  }
%    \end{macrocode}
% \end{macro}
%
% \begin{macro}[int]{undefined-storage-anskey}
% Message used by |\getkeyans|, |\printkeyans| and |\foreachkeyans| commands
% when the storage is not defined.
%    \begin{macrocode}
\msg_new:nnn { enumext } { undefined-storage-anskey }
  {
    Storage~named~'#1'~is~not~defined~\msg_line_context:.
  }
%    \end{macrocode}
% \end{macro}
%
% \begin{macro}[int]{missing-miniright}
% Message used by \myenv*[no-index]{enumext} and \myenv*[no-index]{keyans}
% when the key \mykey[env]{mini-env} needs |\miniright|.
%    \begin{macrocode}
\msg_new:nnn { enumext } { missing-miniright }
  {
    Missing~'\token_to_str:N \miniright'~\msg_line_context:.\\
    The~key~'mini-env'~needs~'\token_to_str:N \miniright'.
  }
%    \end{macrocode}
% \end{macro}
%
% \begin{macro}[int]{nested}
% Message used when an environment is nested in itself.
%    \begin{macrocode}
\msg_new:nnn { enumext } { nested }
  {
    The~environment~\l_@@_envir_name_tl \c_space_tl can't~be~nested~
    \msg_line_context:.
  }
%    \end{macrocode}
% \end{macro}
%
% \subsection{Finish package}\label{subsec:finish}
%
% Finish package implementation.
%
%    \begin{macrocode}
\file_input_stop:
%</package>
%    \end{macrocode}
% \Finale
